\chapter{引力波}\label{chpt:gw}
\begingroup\emergencystretch=2em
把应变写成 $h_+(t),h_\times(t)$，还没有说明观者会收到怎样的信号。一个偏心双星每过近心点就发出一串短波，缓慢进动的星体可在几条紧邻的谱线上辐射。两者即使在地球处有相同的最大应变，在同一台仪器上也可能一个清楚可见，另一个淹没在噪声里。以下从源的多极矩求波形，再沿光路或机械振子把它换成可测的量。对于许多无法逐一分辨的源，可测的量又变为两组读数的相关。

线性辐射场及其应力张量沿用第一章\hyperlink{subsection.1.2.4}{「引力辐射」}，特别是式 \eqref{eq:gw-effective-stress}。双星近区的后 Newton 运动见 \S\ref{sec:exp-pn}，由保守能量与辐射通量定相位见 \S\ref{sec:pn-wave-energy}。用脉冲星轨道衰减检验这些通量的计时公式、银河系加速度修正及质量平面已在第三章\hyperlink{section.3.13}{「强场、相对论性引力的动力学检验」}给出。
逐项讨论波如何产生、传播和被测量时，按需要写出 $G,c$；用几何单位的地方另行说明。

\section{波长、曲率半径与传播}\label{sec:gw-a-optics}

在平直背景上，探测器所测潮汐和远区的平均能流已由第一章给出，见式 \eqref{eq:gw-effective-stress}。若让波从一颗星传到另一颗星，两地各可选局域惯性系，但不能以同一 Minkowski 度规代替中间整段时空。沿途哪些量可认作同一列波的性质？若波本身的能量也使时空弯曲，背景又应怎样选取？

\subsection{两种长度和两套场方程}

取背景几何变化的尺度为 $L$，在一小段光路上用正交标架记波的典型振幅 $a=|h_{\hat\mu\hat\nu}|$，约化波长记为 $\lambda$。要求 $a\ll1$，且 $\epsilon=\lambda/L\ll1$。后一条件使我们可以选一个尺度 $d$，在 $\lambda\ll d\ll L$ 的时空邻域内平均。比较不同点的张量分量时先沿短测地线平行移动。这样定义的 $\bar g$ 包含缓变部分，$g=\bar g+h$，并取 $\langle h\rangle=0$。在不同大小的平均窗上分得的 $\bar g$ 可略有不同。计算平均量时分部积分，把平均窗边缘的修正按 $\lambda/d$ 略去，把移动标架的变化按 $d/L$ 略去。对波包的包络另要求它在 $d$ 内缓变。以下若无说明用几何单位 $G=c=1$。

振幅小还不足以决定哪一个曲率项大。以 $\bar\nabla$ 为背景导数，曲率中有 $\bar R\sim L^{-2}$、$R^{(1)}\sim a/\lambda^2$、$R^{(2)}\sim a^2/\lambda^2$。第一阶项多为快振荡，平均后消去，第二阶中两个快量的乘积却有慢分量。因此哪怕 $a\ll1$，只要 $a\sim\lambda/L$，$\langle R^{(2)}\rangle$ 就可以与 $\bar R$ 同阶。若波的能量只是背景中很小的一部分，则还应有 $a\ll\lambda/L$。把 Einstein 张量以 $h$ 展开，在无快变物质源的区域，低频部分和高频部分分别满足
\begin{align}
 \bar G_{\mu\nu}
     &=8\pi\big(\langle T_{\mu\nu}\rangle+t^{\rm GW}_{\mu\nu}\big),
 &t^{\rm GW}_{\mu\nu}&=-\frac{1}{8\pi}\langle G^{(2)}_{\mu\nu}[h,h]\rangle ,
 \label{eq:gw-a-average}\\
 G^{(1)}_{\mu\nu}[h]&=0 &&\text{（快变场的领头阶）}. \notag
\end{align}
慢的二阶度规修正已收进 $\bar g$。第二行确定波怎样在给定的背景上传播，第一行把求得的波代回去修正背景。这种区分依赖于有可用的平均窗。例如两个相近频率叠加时，差频若与背景频率相当，便须保留慢变的拍频项。

为看清协变导数在其中的作用，记联络差的线性项为 $C^\rho{}_{\mu\nu}=(\bar\nabla_\mu h^\rho{}_\nu+\bar\nabla_\nu h^\rho{}_\mu-\bar\nabla^\rho h_{\mu\nu})/2$。联络差是张量，因此
\begin{align*}
 \delta R^\rho{}_{\sigma\mu\nu}
 &=\bar\nabla_\mu C^\rho{}_{\nu\sigma}-\bar\nabla_\nu C^\rho{}_{\mu\sigma},\\
 2\delta R_{\mu\nu}
 &=\bar\nabla_\rho\bar\nabla_\mu h^\rho{}_\nu
   +\bar\nabla_\rho\bar\nabla_\nu h^\rho{}_\mu
   -\bar\square h_{\mu\nu}-\bar\nabla_\nu\bar\nabla_\mu h .
\end{align*}
不能先在最后一式中交换导数再把普通导数加横线。取 $\gamma_{\mu\nu}=h_{\mu\nu}-\bar g_{\mu\nu}h/2$，令 $\bar\nabla^\mu\gamma_{\mu\nu}=0$。真空背景有 $\bar R_{\mu\nu}=0$，上式换序时两项 Ricci 曲率消失，剩下的 Riemann 曲率给出
\begin{equation}
 \bar\square\gamma_{\mu\nu}+2\bar R_{\mu\alpha\nu\beta}\gamma^{\alpha\beta}=0 .
 \label{eq:gw-a-curved-wave}
\end{equation}
在物质背景上一般还须线性化物质方程。式 \eqref{eq:gw-a-curved-wave} 的二阶导主部仍支配短波极限，不能据此把它整个照用于任意有物质的背景。

另一个细节是规范变换。在 $h\mapsto h-\pounds_\xi\bar g$ 下，线性曲率随之改变 $-\pounds_\xi\bar R^\rho{}_{\sigma\mu\nu}$。平直背景上此项为零；弯曲背景上即使在某一点令全部联络消失，它也未必为零。保留 $h$ 的小量计数要求快变的 $\xi$ 满足 $\xi=O(a\lambda)$。于是 $\bar R\bar\nabla\xi=O(a/L^2)$，比快变曲率 $a/\lambda^2$ 小两阶。这说明何以在广义坐标变换下有良好意义的是领头短波曲率，以及背景和波共同给出的总曲率。

二阶平均还可直接算。联络的二阶项是 $-h^\rho{}_\alpha C^\alpha{}_{\mu\nu}$，它在 Ricci 张量中以散度出现，快变部分积分后归入边界。在一点取背景惯性标架，并在快变场上取 TT 条件。$C^\rho{}_{\rho\lambda}=\bar\nabla_\lambda h/2=0$，故仅余 $-\langle C^\rho{}_{\nu\lambda}C^\lambda{}_{\rho\mu}\rangle$。展开乘积后，含散度的项分部积分给 $\bar\nabla^\mu h_{\mu\nu}=0$，两导数相缩并的项给 $\bar\square h=O(h/L^2)$（这里 $h$ 泛指张量分量）。因而
\[
 \langle R^{(2)}_{\mu\nu}\rangle
 =-\tfrac14\big\langle\bar\nabla_\mu h^{\rm TT}_{\alpha\beta}
                        \bar\nabla_\nu h_{\rm TT}^{\alpha\beta}\big\rangle
 \quad\text{（短波领头阶）}.
\]
同理，最后一式的迹是小量。代入式 \eqref{eq:gw-a-average} 就得到式 \eqref{eq:gw-effective-stress} 的局域协变形式。若对来自许多不同方向的非相干波再作各向同性平均，因 $\langle n_i n_j\rangle_{\bm n}=\delta_{ij}/3$，压力即为能量密度的 $1/3$。单束波只在传播方向有压强，各向同性是波束分布的条件。

\subsection{射线、偏振和波前面积}

以复场的实部表示快变扰动，写 $\gamma_{\mu\nu}=A_{\mu\nu}\exp(\mathrm i\theta/\epsilon)$，$k_\mu=\bar\nabla_\mu\theta$。令振幅随背景尺度缓变，逐次求导得到
\[
 \bar\square(A_{\mu\nu}\mathrm e^{\mathrm i\theta/\epsilon})
 =\mathrm e^{\mathrm i\theta/\epsilon}\!
 \left[-\epsilon^{-2}k^2A_{\mu\nu}
  +\mathrm i\epsilon^{-1}\{2k^\rho\bar\nabla_\rho A_{\mu\nu}
                  +(\bar\nabla_\rho k^\rho)A_{\mu\nu}\}
  +\bar\square A_{\mu\nu}\right].
\]
于是 $k^2=0$，且 $k^\rho\bar\nabla_\rho k_\mu=k^\rho\bar\nabla_\mu k_\rho=0$。射线既是波面法线也是仿射参数化的零测地线。谐和条件在最高阶给 $k^\mu A_{\mu\nu}=0$。选定屏幕上的两个物理偏振，取 $A_{\mu\nu}=A e_{\mu\nu}$、$e_{\mu\nu}e^{*\mu\nu}=1$，把共通的慢变相位吸收到 $e_{\mu\nu}$ 中。一阶输运式及其共轭分别与 $A^{*\mu\nu}$、$A^{\mu\nu}$ 缩并，相加得到
\begin{equation}
 \bar\nabla_\mu(|A|^2k^\mu)=0,\qquad
 k^\rho\bar\nabla_\rho e_{\mu\nu}=0 .
 \label{eq:gw-a-transport}
\end{equation}
偏振相对于平行移动的屏幕基保持不变。一个沿途另行选取的望远镜坐标架可以相对于它旋转，不宜把基底旋转解释成两种本征模式有不同传播速度。

令相邻射线所围的正交截面积为 $\mathcal A(s)$，$s$ 为仿射参数。截面内两个联络矢的变化给 $\mathrm d\ln\mathcal A/\mathrm ds=\bar\nabla_\mu k^\mu$，于是 $|A|^2\mathcal A$ 守恒。在平直空间，自一点发出的射束截面随 $r^2$ 增长，便得 $A\propto1/r$。若静态前景使截面积缩小一半，在源的相同发射立体角内，观者收到振幅增大 $\sqrt2$ 倍的波。沿路没有耗散，这一变化全由相邻测地线的会聚造成。焦散处 $\mathcal A\to0$，形式上的发散意味着需重新把相邻光路相干相加。

对频率还须指定观者。4-速度 $U^\mu$ 的钟读 $\omega=-k\cdot U/\epsilon$。若源与接收者之间红移为 $z$，相邻波峰的到达间隔拉长 $1+z$，每个量子的能量又少一因子 $1+z$。定义 $D_G^2=\mathcal A_o/\mathrm d\Omega_s$、$D_A^2=\mathcal A_s/\mathrm d\Omega_o$，下标 $s,o$ 分别是源与观者。无吸收时观测通量于是为 $F_o=L_s/[4\pi D_G^2(1+z)^2]$，从而 $D_L=(1+z)D_G$。光路互易使 $D_G=(1+z)D_A$。用于振幅输运的是 $D_G$，用源的光度同观测通量比较才出现 $D_L$。例如四极波在发射方向的系数含 $\mathcal M^{5/3} f_s^{2/3}$，用 $f_s=(1+z)f_o$ 及 $D_G=D_L/(1+z)$，即化为 $[(1+z)\mathcal M]^{5/3} f_o^{2/3}/D_L$。这给出第三章红移质量替换的传播依据。有限频带内如有吸收或色散，就须分别输运各频率，不能只更换距离。

\subsection{散射、吸收与测量标架}

引力波路径弯曲与能量真正转给介质是两件可分别估算的事。静态介质中若散射体数密度为 $n$、移出所选波束的截面为 $\sigma$，沿长度 $\mathrm dl$ 的损失概率为 $n\sigma\,\mathrm dl$。在互不相干近似下，式 \eqref{eq:gw-a-transport} 的强度解须乘 $\exp(-\int n\sigma\,\mathrm dl)$，振幅则取该因子的平方根。散射极强地集中于前向时应明确观测孔径。若关心方向的随机化，可用 $\sigma_{\rm tr}=\int(1-\cos\vartheta)\,\mathrm d\sigma$，它给极小偏转以很小的权重。由此可见，对有效偏折求平均与对被吸收的能量求和通常得不到同一截面。

普通物质中的微观散射很弱，宏观天体的相干响应却可有一个窄带共振。另一个极限是黑洞的直接吸收。几何光学下考察 Schwarzschild 黑洞的零测地线，径向方程可写为 $\dot r^2/E^2=1-b^2(1-2M/r)/r^2$，$b$ 为冲量参数。转向点刚好合并时，$(1-2M/r)/r^2$ 取极大值，微分给 $r=3M$，最大值为 $1/(27M^2)$。故小于 $b_c=3\sqrt3 M$ 的射线落进视界，截面 $\sigma_{\rm cap}=27\pi M^2$。截面大于视界的投影面积，差别来自外部曲率对入射束的会聚。只有 $M\omega\gg1$ 才能忽略绕射；若 $M\omega\lesssim1$，须解各角动量模的势垒穿透问题。

接收端尚有一个尺度 $d_{\rm det}$。可在中心世界线上沿钟摆放 Fermi 标架，坐标离中心为 $X^i$。对自由落体且无旋中心，展开到二阶有
\begin{align*}
 g_{00}&=-1-R_{0i0j}X^iX^j+O(X^3),\qquad
 g_{0i}=-\tfrac23R_{0jik}X^jX^k+O(X^3),\\
 g_{ij}&=\delta_{ij}-\tfrac13R_{ikjl}X^kX^l+O(X^3).
\end{align*}
这里曲率分量均沿中心标架求值。对短波 $R_{0i0j}=O(h/\lambda^2)$，下一项中的 $\nabla R$ 又多一个 $1/\lambda$。局域潮汐力近似所需的是 $d_{\rm det}\ll\lambda$。量得的相对位移仍可达到 $h d_{\rm det}$，因为加速度要积过约一个周期。若等臂干涉仪的臂长与波长相当，把静止镜子的 TT 坐标代入光的零测地线仍然可行，但要沿去程、回程逐点取扰动时刻，单独对臂端做一次测地偏离已不够。若中心靠支承力保持静止或标架有转动，$g_{00}$ 还含 $-2a_iX^i$ 等惯性项，应和上述二阶曲率项分别保留。

\subsection{平面波的二阶曲率与有限射束}

平均怎样从曲率中保留下一个正量，可用沿 $+z$ 行进的平面波独立检验。取 $u=t-z,\ v=(t+z)/2$，横向指标为 $A,B=1,2$，考虑 Rosen 坐标形式 $\mathrm ds^2=-2\,\mathrm du\,\mathrm dv+g_{AB}(u)\mathrm dx^A\mathrm dx^B$。在同一个 $u$ 面上横向度规为常数。与 $u$ 导数有关的联络为 $\Gamma^A{}_{uB}=g^{AC}g'_{CB}/2$、$\Gamma^v{}_{AB}=g'_{AB}/2$，撇表示 $\mathrm d/\mathrm du$。唯一需要的 Ricci 分量是
\[
 R_{uu}=-\tfrac12{\rm tr}(g^{-1}g'')+\tfrac14{\rm tr}(g^{-1}g'g^{-1}g') .
\]
若先把 $g_{AB}$ 取为 $\delta_{AB}+h_{AB}$、$\mathrm{tr}\,h=0$，则一阶 Ricci 张量为零而快振荡曲率一般非零。二阶得 $R_{uu}^{(2)}=\mathrm{tr}(hh'')/2+\mathrm{tr}(h'^2)/4$。用 $\mathrm{tr}(hh'')=[\mathrm{tr}(hh')]'-\mathrm{tr}(h'^2)$ 平均，得 $\langle R_{uu}^{(2)}\rangle=-\langle\mathrm{tr}(h'^2)\rangle/4$。这与前面从一般联络差所求相合。只有 $uu$ 分量出现，又说明这束波的平均源沿零方向，无法由一团静止尘埃代替。

令 $g_{AB}=\delta_{AB}+h_{AB}+j_{AB}$ 且 $j_{AB}=O(h^2)$，真空场方程便要求 $\langle \mathrm{tr}\,j''\rangle=-\langle\mathrm{tr}(h'^2)\rangle/2$。更清楚的几何量是 $u$ 作仿射参数、$x^A$ 作射线标号的一束与波迎面相遇的零线的横截面积。写 $B=g^{-1}g'/2$、$\vartheta=\mathrm{tr}\,B=\mathrm d\ln\sqrt{\det g}/\mathrm du$，直接代换可将精确方程化为
\[
 R_{uu}=-\vartheta'-{\rm tr}\,B^2,\qquad
 \vartheta'=-\tfrac12\vartheta^2-{\rm tr}\,\sigma^2,\quad
 \sigma=B-\tfrac12\vartheta\,\mathbf1 .
\]
初始平行的小束有 $\vartheta=0$，一阶波主要产生剪切 $\sigma=h'/2$，其平方却会使 $\vartheta$ 逐渐变负。单列波在真空中仍满足 $R_{\mu\nu}=0$；改变邻线会聚的是精确几何中的剪切，若改用抹去快变结构的平均几何来描述，同一效应就由式 \eqref{eq:gw-a-average} 的平均 Ricci 源提供。

求传播时所需的面积还可以由两条线性偏离方程计算。令 $E_A^\mu$ 沿一条中心射线平行移动，$k\cdot E_A=0$，偏离矢的屏幕分量写成 $\xi_A=\mathcal D_{AB}\theta_B$。初始小角 $\theta_B$ 乘上 $\mathcal D$ 给沿途的物理分离。令 $\mathcal R_{AB}=-R_{A k B k}$，则
\[
 \mathcal D''=\mathcal R\,\mathcal D,\qquad
 \mathcal D(s_0)=0,\qquad \mathcal D'(s_0)=\mathbf1 .
\]
此处把 $k$ 在起点处归一为单位本地频率。$\mathcal R$ 为对称矩阵，一半的迹给 Ricci 聚焦，去迹部分给潮汐剪切。迭代一次积分，若以 $s_0=0$ 记起点，
\[
 \mathcal D(s)=s\mathbf1+\int_0^s(s-s')\mathcal R(s')\,\mathcal D(s')\,\mathrm ds'
   \simeq s\mathbf1+\int_0^s(s-s')s'\mathcal R(s')\,\mathrm ds' .
\]
平直结果是 $s\mathbf1$。若薄潮汐层在 $s_l$ 处很窄，造成的偏离成比例于 $s_l(s-s_l)$；这给出了同一横向曲率放在源、观者附近或传播路径中央为何作用不同。$\det\mathcal D$ 改变波束面积，两个特征值之比改变波束的椭率。纯 Weyl 扰动在一次迭代中迹为零，所以可先使圆像变扁，而面积在更高阶才改变。

对任意两套矩阵解 $D_1,D_2$，Wronskian $D_1^{\mathsf T}D_2'-D_1'{}^{\mathsf T}D_2$ 对 $s$ 的导数为零。取 $D_1$ 在源点为零，$D_2$ 在观者点为零，分别按同一仿射参数取单位初导数。由两端的 Wronskian 相等可知正向传输矩阵等于反向矩阵的负转置。因此其面积行列式相等。若改以各端的本地小角来标记光线，只须把初导数乘以各自测到的 $|k\cdot U|$，就得到 $D_G/D_A=|k\cdot U_s|/|k\cdot U_o|=1+z$。这给出了前面用到的距离互易关系，源或观者的本动包含在各自的 $U$ 内。

若有两束射线均从源到达同一接收器，只输运强度还会丢失相位。设两路振幅放大率为 $a_1,a_2$、到时差为 $\Delta t$，并把两路偏振移到同一屏幕基。对相同的偏振分量，接收的频谱是 $\widetilde h_s(f)[a_1+a_2\exp(-2\pi\mathrm i f\Delta t+\mathrm i\delta)]$，其中 $\delta$ 可含穿过焦散的常数相位。可分辨的干涉条纹间隔为 $1/|\Delta t|$；若观测频箱比此宽，干涉项在箱内平均消去，才恢复两个强度相加。若暂态的持续时间远短于 $|\Delta t|$，则时域中会依次看到两次信号。即使在每一束上几何光学均适用，怎样合成两束仍取决于实验的时间、频率分辨能力。

\subsection{由 TT 坐标到自由探测器的局域坐标}

前面用 Fermi 展开给了探测器处的度规。对正入射平面波也可从显式坐标变换看到同一潮汐。取 $u=t-z$、$\mathrm ds^2=-\mathrm dt^2+\mathrm dz^2+(\delta_{ab}+h_{ab}(u))\mathrm dx^a\mathrm dx^b$，$a,b=x,y$。暂取 $c=1$。在中心线 $x^a=z=0$ 邻域内先令 $X^a=x^a+h^a{}_b(u)x^b/2$，即按瞬时正交基变换横向坐标。只做到这一步会出现 $-\dot h_{ab}X^b\mathrm dX^a\mathrm du$ 的交叉项。改取 $V=v+\dot h_{ab}X^aX^b/4$，利用 $-2\,\mathrm du\,\mathrm dv$ 中的补偿即可将其消去，得到
\[
 \mathrm ds^2=-2\,\mathrm du\,\mathrm dV+\delta_{ab}\mathrm dX^a\mathrm dX^b
          +\tfrac12\ddot h_{ab}(u)X^aX^b\,\mathrm du^2+O(h^2).
\]
这里 $v=(t+z)/2$，点是对 $u$ 求导。曲率由 $\Gamma^a{}_{uu}=-\ddot h^a{}_bX^b/2$ 给出。令处在同一波面横截面的两颗质点起初静止，则它们的 $X$ 距离满足 $\mathrm d^2X^a/\mathrm du^2=\ddot h^a{}_bX^b/2$。在 TT 系中这两颗质点可有固定的 $x^a$，度规随时刻改变；局域坐标中尺长是欧氏的，坐标位置随时刻改变。两次计算用到的是同一对质点及同一时刻的本地标尺。

若波沿 $+z$，一个横向固定坐标距离 $L$ 的两端在 TT 系同一 $t$ 上有固有空间距离 $L[1+h_{pp}(t)/2]$；这里取两端位于 $z=0$ 并令 $\bm p$ 沿臂。在同一时刻从发射端发出光脉冲，光沿臂到达另一端所需时间却是 $L+\int_0^L h_{pp}(t+l)\,\mathrm dl/2$。两者之差从 $L^2\dot h_{pp}/4$ 开始，正比于一次光行时中度规的变化。前者是指定空间切片上的尺量，后者才是以发、收事件定义的传播实验。若两端都由独立的钟标记，尚须说明两钟的校准；对同端往返用同一只钟，就能直接比较读数。后文的干涉仪响应会完整保留这段相位积分。

\section{螺旋度与有质量的张量场}\label{sec:gw-a-field}

无质量自由场的 Fierz--Pauli 作用量、螺旋度的计数及低能有效理论，见 \S\ref{sec:quantum-gravity}。从这同一个二次作用量还可问两个问题。有限波束把多少角动量传过远处的球面？若在线性场方程中加入质量项，原来由规范自由消去的分量怎样受到约束？

\subsection{旋转流中的两项}

在平直时空的辐射区取 TT 规范，二次作用量化为 $S_{\rm TT}=-(64\pi)^{-1}\int \partial_\mu h_{ij}\partial^\mu h_{ij}\,\mathrm d^4x$，本小节 $h$ 都指物理的 TT 场。正则共轭动量是 $\pi_{ij}=\dot h_{ij}/(32\pi)$。考虑绕 $z$ 轴作无穷小主动转动 $\delta\alpha$。场在固定坐标点的变化分为函数自变量的变化与张量两个指标的变化，
\[
 \frac{\delta h_{ij}}{\delta\alpha}
 =-(x\partial_y-y\partial_x)h_{ij}
     +\epsilon_{i3a}h_{aj}+\epsilon_{j3a}h_{ia}.
\]
把它与 $\pi_{ij}$ 缩并，得 Noether 荷
\begin{equation}
 J_z=\frac{1}{32\pi}\int\mathrm d^3x
 \left[-\dot h_{ij}\partial_\phi h_{ij}
        +2\dot h_{ij}\epsilon_{i3a}h_{aj}\right].
 \label{eq:gw-a-noether}
\end{equation}
第一项中 $\partial_\phi$ 对固定 Descartes 基底上的分量求导。若改用随方位转动的基底，导数还会作用到基矢上；这时不能把后一项再额外算一次。加在 Noether 流上的全散度也能在局部重新分配这两项，局域波包的总荷则不变。对无限平面波宜先取宽而有限的波包，再取单位体积的平均。

以沿 $z$ 传播的近单色波为例。在束中央可略去横向包络导数，置 $h_{xx}=-h_{yy}=A\cos(\omega t-kz)$、$h_{xy}=h_{yx}=A\sin(\omega t-kz)$。绕 $z$ 轴的指标转动使 $(h_+,h_\times)$ 变为 $(-2h_\times,2h_+)\delta\alpha$，因此自旋密度是 $(h_+\dot h_\times-h_\times\dot h_+)/(8\pi)=A^2\omega/(8\pi)$。由式 \eqref{eq:gw-plane-wave-stress} 的能量密度 $A^2\omega^2/(16\pi)$，相除即得 $2/\omega$。把正弦的符号反转，角动量反向，能量不变。乘单量子能量 $\hbar\omega$ 得 $\pm2\hbar$，与第一章由偏振态变换所得相合。线偏振是等幅的两种螺旋度相加，时间平均自旋流为零，能量仍为正。

对远方角向分布不均匀的辐射，仅以指向径向的有效能流乘力臂，$\bm r\times\bm n=0$，似乎永远得不到角动量。问题出在把 $t^{ri}$ 中小一阶的横向项提前舍掉了。球面积分的因子是 $r^2$，力臂还带一个 $r$，横向动量流的 $r^{-3}$ 项正好留下。亦可用式 \eqref{eq:gw-a-noether} 的旋转生成元直接在辐射场上作用，无须先算到次领头的局域动量流。若按轴向总角动量本征值 $m$ 展开，则一束单模有 $\dot J_z/\dot E=m/\omega$。此处 $m$ 计入两个指标的转动，未必要等于 $\pm2$。例如在球面上一个有 $m=0$ 的轴对称脉冲可发出有限能量，却不会减小源的轴向角动量。

\subsection{软引力子的规范条件与源的自洽性}

式 \eqref{eq:soft-graviton-radiated-energy} 从经典外线求出了碰撞的软辐射。把同一极限用于散射振幅，还可反问引力子能否以不同常数耦合到不同种粒子。暂取 $\hbar=c=1$，设硬过程有振幅 $\mathcal A$，又放出 4-动量 $q\to0$ 的一个无质量自旋 2 粒子。第 $A$ 条外线上附加软量子后，相邻传播子的分母化为 $(p_A+\eta_A q)^2+m_A^2=2\eta_A p_A\cdot q$，$\eta_A$ 按出入射取正负。领头顶角只需保留 $p_A^\mu p_A^\nu$；含 $q$ 的自旋与短程结构修正不会给出同一软极点。允许不同粒子有耦合 $g_A$，便有
\[
 \mathcal A(q,e)
  =\mathcal A \sum_A\eta_A g_A
      \frac{e_{\mu\nu}^*p_A^\mu p_A^\nu}{p_A\cdot q}+O(q^0).
\]
分母来自近乎在壳的外腿，局域硬顶角上另发软粒子的图在一般外部运动学下没有这个极点。将偏振张量换为纯规范 $q_\mu\xi_\nu+q_\nu\xi_\mu$，极点恰抵消，Ward 条件剩下 $\sum_A\eta_Ag_A p_A^\nu=0$。无权重的关系 $\sum_A\eta_Ap_A^\nu=0$ 是总动量守恒。对能够相互传递任意动量的两个物种做一次弹性散射，若转移动量为 $\Delta p^\nu$，便要求 $(g_1-g_2)\Delta p^\nu=0$；故这两个耦合必须相同。同理，只要引力子本身也与这些粒子交换动量，它也须按同一个耦合常数贡献到源。这里约束的是长程软极限，不决定较高导数局域顶角的一切系数。

这个结论在经典场方程中有对应的迭代。自由场的线性算符 $G_{\mu\nu}^{(1)}$ 散度恒零，故若写 $G_{\mu\nu}^{(1)}=8\pi G T_{\mu\nu}$，右边的普通散度也必须为零。但一旦让同一引力场影响物质，精确式是 $\partial_\mu T^{\mu\nu}=-\Gamma^\mu{}_{\mu\alpha}T^{\alpha\nu}-\Gamma^\nu{}_{\mu\alpha}T^{\mu\alpha}$。右边的量级为 $\partial h\,T$。在弱场方程中用 $T\sim\partial^2h/G$ 消去 $T$，即知缺的流是 $(\partial h)^2/G$ 或 $h\partial^2h/G$ 型；它正由引力能动张量补足。在作用量中需要相应的三次项 $h(\partial h)^2/G$。若简单把自由场的 $t^{\mu\nu}$ 乘上 $h_{\mu\nu}$，变分还会作用到 $t$ 中的两个 $h$，不能把这个乘积的泛函导数径直认作 $t^{\mu\nu}$。应在同阶同时修正作用量和规范变换。

把这个条件写得更明确些。令 $S=S_2+\kappa_g S_3+\cdots$、$\delta=\delta_0+\kappa_g\delta_1+\cdots$，$\delta_0\gamma_{\mu\nu}=\partial_\mu\xi_\nu+\partial_\nu\xi_\mu$。一次阶须有 $\delta_0 S_3=-\delta_1 S_2$。取 $g_{\mu\nu}=\eta_{\mu\nu}+\kappa_g\gamma_{\mu\nu}$ 并展开 $\pounds_\xi g$，非线性变换给 $\delta_1\gamma_{\mu\nu}=\xi^\rho\partial_\rho\gamma_{\mu\nu}+\gamma_{\rho\nu}\partial_\mu\xi^\rho+\gamma_{\mu\rho}\partial_\nu\xi^\rho$。这恰是三次相互作用不能独立保持线性规范对称性时所需的补偿。引力子的自由作用量、由 Einstein--Hilbert 作用量生成的各阶顶角和有效理论的导数计数均见 \S\ref{sec:quantum-gravity}。在有物质和引力交换的源计算中，这套自洽性具体要求用总有效源守恒；后面的高阶辐射矩和近区反作用都要用到它。

\subsection{从质量项得到五个自由度}

以 $\gamma_{\mu\nu}$ 表第一章规范化的自由自旋 2 场（此处不取迹反转），迹记为 $\gamma$。不含导数的 Lorentz 标量只有 $\gamma_{\mu\nu}\gamma^{\mu\nu}$ 和 $\gamma^2$。加到无质量作用量的项可取
\[
 \Delta\mathscr L=-\tfrac12 m_g^2
        (\gamma_{\mu\nu}\gamma^{\mu\nu}-b\gamma^2)
        +\tfrac{\kappa}{2}\gamma_{\mu\nu}T^{\mu\nu},
\]
暂取 $\hbar=1$。令 $K_{\mu\nu}=-2G_{\mu\nu}^{(1)}[\gamma]$，其散度恒为零，迹为 $2(\partial_\rho\partial_\sigma\gamma^{\rho\sigma}-\square\gamma)$。当 $T_{\mu\nu}$ 守恒时，对场方程取散度和迹分别得
\begin{align*}
 K_{\mu\nu}&=m_g^2(\gamma_{\mu\nu}-b\eta_{\mu\nu}\gamma)-\tfrac\kappa2 T_{\mu\nu},\\
 \partial^\mu\gamma_{\mu\nu}&=b\partial_\nu\gamma,\qquad
 2(b-1)\square\gamma=m_g^2(1-4b)\gamma-\tfrac\kappa2 T .
\end{align*}
若 $b\ne1$，最后一式是迹的传播方程，除五个自旋 2 分量外还多出一个标量。自由动能中这个分量的符号有问题。可在质量项上作 $\gamma_{\mu\nu}\mapsto\gamma_{\mu\nu}+\partial_\mu\partial_\nu\pi$ 来追踪纯纵向部分；原无质量动能对此替换不变。分部积分给 $\int(\partial_\mu\partial_\nu\pi)^2=\int(\square\pi)^2$，故纯标量留下 $-m_g^2(1-b)(\square\pi)^2/2$。这个四阶导项一般给出两个留数相反的极点。$b=1$ 恰能消掉它，这就是线性 Fierz--Pauli 质量项的调谐。

此时 $\gamma=-\kappa T/(6m_g^2)$。在源外有 $\gamma=0$、$\partial^\mu\gamma_{\mu\nu}=0$，其余满足 $(\square-m_g^2)\gamma_{\mu\nu}=0$。移到一粒子的静系，横向条件给 $\gamma_{0\nu}=0$，空间对称三阶方阵去掉迹恰有五个自由度，按一固定轴分解即 $m=0,\pm1,\pm2$。这五条约束来自场方程，只对 $m_g\ne0$ 的理论成立，不能在加质量项之前先把 $\gamma$ 置于 TT 规范。调谐只控制围绕平直背景的二次作用量，推广到相互作用时还需检查非线性约束是否保留。

也可从传播子中读出同一结果。对 $k^2=-m_g^2$ 定义 $\Pi_{\mu\nu}=\eta_{\mu\nu}+k_\mu k_\nu/m_g^2$，则 $k^\mu\Pi_{\mu\nu}=0$、$\Pi^\mu{}_\mu=3$。一个对称张量在这三维横向空间的无迹投影为
\[
 P^{(2)}_{\mu\nu,\rho\sigma}
 =\tfrac12(\Pi_{\mu\rho}\Pi_{\nu\sigma}+\Pi_{\mu\sigma}\Pi_{\nu\rho})
  -\tfrac13\Pi_{\mu\nu}\Pi_{\rho\sigma}.
\]
对一对守恒源，诸 $k_\mu$ 项与源缩并为零，迹项中的 $1/3$ 则留下。相比之下，无质量波实际只在二维屏幕上去迹，得到 $1/2$。这给出了第三章讨论静态偏折时所用两个交换振幅的来历。纵向空间的维数并不随 $m_g$ 连续改变，因此把两个投影算符中的 $m_g$ 同时趋于零并不合法。关于静力校准及非线性屏蔽，见第三章「曲率修正、有质量的引力与双度规」。

若只观测传播，直接测到的是色散。以一维窄带波包为例，写 $k(\omega)=\sqrt{\omega^2-m_g^2}$，对 $\omega_0>m_g$ 展开传播相位 $Dk(\omega)$。常数项改变相位，一次项使包络延迟 $Dk'=D/v_g$，其中 $v_g=k/\omega$。二次项 $Dk''=-Dm_g^2/k^3$ 使不同频率的包络到时不同。源在两个频率的发射次序本可不同，因此比较峰值时须减去源内时间差。若 $\omega<m_g$，固定频率的稳态响应含 $\exp[-D\sqrt{m_g^2-\omega^2}]$；在有限距离仍有近场，却不给出可传播到无穷远的平面波。对一个有质量标量也有这条色散关系，所以独凭群速度不能决定自旋或偏振数。

\subsection{纵向极化的耦合与极限次序}

线性传播子中的 $1/3$ 可以从各个偏振的贡献看得更清楚。以 $k^\mu=(\omega,0,0,k)$ 为例，定义单位类空纵向矢量 $s^\mu=(k,0,0,\omega)/m_g$ 及横向矢量 $e_1^\mu,e_2^\mu$。三者与 $k^\mu$ 正交且范数为一。令 $e_\pm=(e_1\pm\mathrm i e_2)/\sqrt2$，五个正交张量可分别取为
\[
 e_\pm e_\pm,\qquad
 \tfrac1{\sqrt2}(s e_\pm+e_\pm s),\qquad
 \tfrac1{\sqrt6}(e_1e_1+e_2e_2-2ss).
\]
这里并列表示张量积。它们的螺旋度依次为 $\pm2,\pm1,0$。固定能量令 $m_g/\omega\ll1$ 时，$s^\mu=k^\mu/m_g+\delta s^\mu$，而 $\delta s^\mu=O(m_g/\omega)$，因为 $\omega-k=m_g^2/(\omega+k)$。与守恒源缩并会杀掉 $k^\mu/m_g$。于是 $\pm1$ 极化的辐射振幅至少小 $m_g/\omega$，但零螺旋度张量中的 $e_1e_1+e_2e_2$ 并不消失。这就是从五维极化空间取小质量极限时，出现额外标量而未留下显著矢量耦合的原因。

具体检验这个标量所测的源。由 $k_\mu T^{\mu\nu}=0$ 有 $T^{0\nu}=kT^{3\nu}/\omega$。因此 $T=-T^{00}+T^{11}+T^{22}+T^{33}=T^{11}+T^{22}+(m_g^2/\omega^2)T^{33}$。再算 $s_\mu s_\nu T^{\mu\nu}=(m_g^2/\omega^2)T^{33}$，所以零螺旋度的源振幅为 $[T-3(m_g^2/\omega^2)T^{33}]/\sqrt6$。在小质量极限留下 $T/\sqrt6$，两源间多出的交换分子正是 $TT'/6$。这比单看静态势更能说明质量改动了什么：色散修正按 $m_g^2/\omega^2$ 变小，标量的迹耦合却可保持有限值。

小质量与弱场展开的极限还可能不交换。以一个在平直参考度规上补入一般非线性质量势的早期模型为例，设 $r_s=2GM/c^2$、$\lambda_g$ 为约化 Compton 波长。弱球对称解的某些高阶项同一阶解的比值是 $r_s\lambda_g^4/r^5$，故线性解只在 $r\gg (r_s\lambda_g^4)^{1/5}$ 可靠。其来源也能从去耦尺度估计。暂用 $\hbar=c=1$，以 $M_{\rm Pl}\sim G^{-1/2}$，纵向标量 $\pi$ 若有 $(\partial^2\pi)^3/\Lambda_5^5$ 相互作用，$\Lambda_5^5\sim M_{\rm Pl}m_g^4$。静态点源的线性纵向场约为 $\pi_{\rm lin}\sim M/(M_{\rm Pl}r)$，所以非线性场方程与线性项的量级比为 $\pi_{\rm lin}/(\Lambda_5^5 r^4)\sim r_s/(m_g^4r^5)$。这里是由一个非线性长度制约线性计算，即便 $r_s/r$ 已很小也可发生。

若理论的质量势经过特殊调谐，主导的纵向相互作用可以改为 $(\partial\pi)^2\square\pi/\Lambda_3^3$，$\Lambda_3^3\sim M_{\rm Pl} m_g^2$。此时只需用 $\partial^2\pi_{\rm lin}/\Lambda_3^3\sim r_s/(m_g^2r^3)$ 来计数，得到第三章所用的 $(r_s\lambda_g^2)^{1/3}$。因而这两个不同的幂次不矛盾，它们属于不同的非线性作用量。知道某区域需重求非线性解，尚不足以保证解是稳定的或内外支都可平滑衔接。可检验的步骤是在内区选定源和边界条件，再同外区衰减解匹配；线性的质量项自身定不出这些性质。

\subsection{平面螺旋度态的角动量和宇称}

给定动量后只许两种螺旋度，是否意味着引力场只能从角动量 $j=2$ 的源跃迁产生？绕动量方向和绕任一固定空间轴是两个不同的旋转。设 $|\hat{\bm z},\lambda\rangle$ 为动量沿 $z$ 的态，$\lambda=\pm2$，定义 $|\bm n,\lambda\rangle=R_z(\phi)R_y(\theta)|\hat{\bm z},\lambda\rangle$。把动量在整个方向球面上叠加，
\[
 |jm;\lambda\rangle
 =\sqrt{\frac{2j+1}{4\pi}}\int\mathrm d\Omega\,
 D^{j*}_{m\lambda}(\phi,\theta,0)\,|\bm n,\lambda\rangle ,
\]
固定了能量而未指定动量方向，$D^j$ 是旋转群的 Wigner 矩阵。对 $z$ 轴旋转可直接对 $\phi$ 作变量代换，得投影量子数 $m$；一般旋转则由 $D$ 矩阵的乘法律定出总角动量 $j$。因 $|\lambda|\le j$，第一项必须从 $j=2$ 开始，但 $j=3,4,\ldots$ 同样存在。把球上带方位变化的振幅当成只含 $\lambda=\pm2$ 两种整球角图，就会错误地删掉高多极辐射。

空间反演使 $\bm k\mapsto-\bm k$，而角动量为轴矢量，故螺旋度变号。取相容的横向基约定，可使 $\mathcal P|jm;\lambda\rangle=(-1)^j|jm;-\lambda\rangle$。于是真正的确定宇称态是 $[|jm;+2\rangle\pm|jm;-2\rangle]/\sqrt2$，宇称分别为 $(-1)^j$ 和 $(-1)^{j+1}$。后面球面张量的电型、磁型就是这两个组合在位置空间的代表。若初末态有确定总角动量 $j_i,j_f$ 与宇称 $P_i,P_f$，辐射阶数还须满足 $|j_i-j_f|\le j\le j_i+j_f$；质量型要求 $P_iP_f=(-1)^j$，流型则要求 $P_iP_f=(-1)^{j+1}$。例如两个无自旋、内部角动量皆为零的孤立定态之间不能单独放出一个任意多极的实引力子，同时发射多个量子则须另作耦合。这与第一章给出的原子质量四极跃迁属于同一条选择规则。

\section{孤立源的辐射多极矩}\label{sec:gw-a-multipoles}

给定一个任意运动的局域源，式 \eqref{eq:arbitrary-source-radiated-energy} 已直接给出远处的能谱。若要从一族波形认出哪些差别来自源的质量分布、哪些来自内部流动，就应把方向 $\bm n$ 显式分离。取源的质心为原点，辐射时间 $u=t-r/c$。本节恢复 $c,G$，使质量矩、流矩和它们的相对阶数易于辨认。

\subsection{对称去迹和源内推迟}

把一连串空间指标记作 $L=i_1\ldots i_l$，$x_L=x_{i_1}\cdots x_{i_l}$，$\hat x_L=x_{\langle L\rangle}$ 表示对称去迹（STF）。如 $\hat x_{ijk}=x_ix_jx_k-(x^2/5)(x_i\delta_{jk}+x_j\delta_{ki}+x_k\delta_{ij})$。全对称的 $l$ 阶张量有 $(l+1)(l+2)/2$ 个分量，迹是 $l-2$ 阶对称张量，相减得 $2l+1$。与 $x_L$ 缩并便成为 $l$ 次齐次多项式，取 Laplace 算符每次都缩并一对指标，故无迹恰使多项式为调和函数；限制在 $r=1$ 上就给普通球谐函数。这样同一角分布可用固定 Descartes 张量或 $Y_{lm}$ 的系数表示。

对 STF 张量 $A_L,B_K$，有
\[
 \int (A_Ln_L)(B_Kn_K)\,\mathrm d\Omega
 =\delta_{lk}\,\frac{4\pi l!}{(2l+1)!!}\, A_L B_L
 \qquad (k=|K|,\ l=|L|),
\]
其中右边的缩并只在 $k=l$ 时取。理由是奇数个 $n_i$ 积分为零；偶数个 $n_i$ 的积分只能是所有两两配对的 $\delta_{ij}$ 的对称和。逐对缩并并用 $n_i n_i=1$ 即定出 $4\pi/(2p+1)!!$（$2p$ 个指标）。同一 STF 张量内部的配对为零，唯余两张量一一相配的 $l!$ 项。稍后求总通量，就是依此剔除不同 $l$ 的交叉项。

先看标量模型 $\square\varphi=-4\pi\rho$。在源外，$F(u)/r$ 是无入射的球面解；再对其作任意空间导数仍是解，因 $\partial_i$ 与 $\square$ 对易。因此可设
\begin{equation}
 \varphi(t,\bm x)=\sum_{l\ge0}\frac{(-1)^l}{l!}
       \partial_L\!\left(\frac{F_L(u)}{r}\right),\qquad
 F_L(u)=\int \mathrm d^3y\,\hat y_L\!\int_{-1}^{1}\mathrm dz\,
       w_l(z)\rho(u+z|\bm y|/c,\bm y) ,
 \label{eq:gw-a-stf-ret}
\end{equation}
其中 $w_l(z)=(2l+1)!!(1-z^2)^l/(2^{l+1}l!)$。对静态源即得通常的 Newton 多极势。对一般源，时间参数不是源处的单一时刻；$z$ 积分保留了波由源的前后两侧发出时的光行时差。这一点亦可从 Fourier 空间检验。单色源中，辐射场带 $\exp(-\mathrm i k\bm n\cdot\bm y)$ 或其共轭；该平面波的 $l$ 次角分量是球 Bessel 函数 $j_l(k|\bm y|)$。令 $p=k|\bm y|$，由 Rodrigues 公式对 $\int P_l(z)\exp(\mathrm i p z)\,\mathrm dz$ 分部积分 $l$ 次得
\[
 \int_{-1}^{1}w_l(z)\mathrm e^{\mathrm i pz}\,\mathrm dz
       =(2l+1)!!\,\frac{j_l(p)}{p^l}.
\]
右边的 $p^l$ 由 $\partial_L[F_L(u)/r]$ 在远区的 $l$ 次推迟时间导数补上。因每个 $l$ 的出射径向方程只有一个出射解，在远区对上的这些系数便定出了源外全解。

$w_l$ 是偶函数，$\int w_l\,\mathrm dz=1$。用 Beta 积分还得 $\int z^2w_l\,\mathrm dz=(2l+3)^{-1}$，四阶矩为 $3/[(2l+3)(2l+5)]$。于是源内加权推迟对任意光滑 $f$ 的展开为
\[
 \int_{-1}^1 w_l(z)f(u+zy/c)\,\mathrm dz
 =f(u)+\frac{y^2f^{(2)}(u)}{2c^2(2l+3)}
    +\frac{y^4f^{(4)}(u)}{8c^4(2l+3)(2l+5)}+\cdots .
\]
没有一次时间导数是前后两个方向相消的结果。若把完整传播时间 $r/c$ 也如此在 $t$ 处展开，远区展开立即失去意义。只对大小为源半径的 $y$ 展开，正是上式能与远区 $1/r$ 极限并用的原因。

\subsection{质量矩与流矩}

张量场比标量场多了一道约束。对推迟解的每个分量可先如式 \eqref{eq:gw-a-stf-ret} 展开，得到每个多极阶的能量、动量、应力矩。由于 $\partial_\nu T^{\mu\nu}=0$，它们并不独立。例如考虑对称应力的一次矩 $S_{ij,k}=\int T^{ij}y_k\,\mathrm d^3y$。定义未去迹的质量矩 $D_{ijk}=\int T^{00}y_iy_jy_k\,\mathrm d^3y/c^2$，两次用守恒律分部积分给 $\ddot D_{ijk} =2(S_{ij,k}+S_{ik,j}+S_{jk,i})$。另定义 $J_{ij}=\int y_{\langle i}(\bm y\times \bm T^{0}/c)_{j\rangle}\,\mathrm d^3y$。对其作时间导数，引号中的动量密度换成应力散度，分部积分后 $\epsilon_{jab}T^{ab}=0$，余下 $\dot J_{ij}={\rm STF}_{ij}\,\epsilon_{jab}S_{bi,a}$。故 $S_{ij,k}$ 的全对称部分给质量八极矩，混合对称部分给流四极矩的导数。剩下的迹经 TT 投影为零或因守恒化入低阶矩。除了给出八极修正，这是最低阶便可看见的电型、磁型两族辐射的由来。

精确地说，在平直背景上解有守恒紧支撑源的线性方程时，令 $\sigma=(T^{00}+T^{aa})/c^2$、$\sigma_i=T^{0i}/c$、$\sigma_{ij}=T^{ij}$。用线性规范自由消去纯规范矩后，源矩可由如下积分生成（被积的各密度均在 $(u+z y/c,\bm y)$ 求值，$y=|\bm y|$）：
\begin{align}
 I_L={}&\int\mathrm d^3y\int_{-1}^{1}\!\mathrm dz\,
 \bigg[w_l\hat y_L\sigma
 -\frac{4(2l+1)w_{l+1}}{c^2(l+1)(2l+3)}\hat y_{aL}\dot\sigma_a
 +\frac{2(2l+1)w_{l+2}}{c^4(l+1)(l+2)(2l+5)}
       \hat y_{abL}\ddot\sigma_{ab}\bigg],\label{eq:gw-a-sourceI}\\
 J_L={}&{\rm STF}_{L}\,\epsilon_{ab i_l}\!
 \int\mathrm d^3y\int_{-1}^{1}\!\mathrm dz\,
 \bigg[w_l\hat y_{L-1\,a}\sigma_b
 -\frac{(2l+1)w_{l+1}}{c^2(l+2)(2l+3)}
      \hat y_{L-1\,ac}\dot\sigma_{bc}\bigg].
 \label{eq:gw-a-sourceJ}
\end{align}
这里 $L-1$ 表示去掉末指标，$\epsilon_{123}=1$。在应力项中 $c$ 作指标时只表示求和的空间指标。质量矩的单位是质量乘长度 $l$ 次方，流矩比它多一个速度。按慢运动计数，$\sigma_i\sim\sigma v$、$\sigma_{ij}\sim\sigma v^2$，时间导数给 $v/y$。例如式 \eqref{eq:gw-a-sourceI} 中流密度项为相对 $(v/c)^2$，恰与 $w_l\sigma$ 中的首个推迟修正同阶，不能只保留一个。以四极矩为例，展开到这一级即
\[
 I_{ij}=\int\mathrm d^3y\left[
 \hat y_{ij}\sigma+\frac{y^2\hat y_{ij}}{14c^2}\ddot\sigma
              -\frac{20\hat y_{aij}}{21c^2}\dot\sigma_a\right]
 +O(c^{-4}) .
\]
这里 $O(c^{-4})$ 指相对于同量纲最低项的 PN 计数。$\sigma$ 本身含应力迹以及 $T^{00}$ 的动能修正，对自引力系统还须计入引力场的有效源。因此这条公式应用于守恒的线性源时是完整的，若将一对外加速度已知的粒子直接作源，它们的物质应力却未必满足所用守恒律。

以两族辐射矩 $U_L,V_L$ 分别记远区的质量型、流型系数。在线性理论中它们是 $I_L^{(l)},J_L^{(l)}$，非线性时还会有传播修正。置 $P_{ij}=\delta_{ij}-n_i n_j$、$\Lambda_{ij,ab}=P_{i(a}P_{b)j}-P_{ij}P_{ab}/2$。远区波用这两族 STF 矩可统一写为
\begin{equation}
 h^{\rm TT}_{ij}=\frac{4G}{c^2r}\Lambda_{ij,ab}
 \sum_{l=2}^{\infty}\frac{1}{c^l l!}
 \left[n_{L-2}U_{abL-2}
  -\frac{2l}{c(l+1)}n_{dL-2}\epsilon_{de(a}V_{b)eL-2}\right].
 \label{eq:gw-a-multipolar-wave}
\end{equation}
$1/c^l$ 来自源区时间导数与空间导数的转换。第一项代入 $l=2$ 就是第一章的质量四极波形。下两项相对同属 $v/c$ 阶，具体为 $2G\Lambda_{ij,ab} [n_k I_{abk}^{(3)}/3-(4/3)n_d\epsilon_{de(a}J^{(2)}_{b)e}]/(c^5r)$。这两项均随 $\bm n\to-\bm n$ 变号，与四极矩的干涉是净线动量流的最低阶来源，见第三章的反冲通量。反过来，求全天总能量时没有额外的方向矢量，这个干涉项恰为零。

\subsection{由波面分量求总功率}

对式 \eqref{eq:gw-a-multipolar-wave} 求一次时间导数，用 $\Lambda_{ij,ab}\Lambda_{ij,cd}=\Lambda_{ab,cd}$，再以 STF 角积分公式收缩，得向无穷远的功率
\begin{equation}
 \mathcal F=\sum_{l\ge2}\frac{G}{c^{2l+1}}
 \left[\frac{(l+1)(l+2)}{l(l-1)\,l!\,(2l+1)!!}\dot U_L\dot U_L
 +\frac{4l(l+2)}{(l-1)(l+1)!\,(2l+1)!!\,c^2}\dot V_L\dot V_L\right].
 \label{eq:gw-a-multiflux}
\end{equation}
可对窄带波包按波周期平均，暂态信号则在需要的时段积分。这里给出一个系数的直接检验。置 $A_{ij}=I_{ij}^{(3)}$，有
\[
 A_{ij}\Lambda_{ij,ab}A_{ab}
       =A_{ij}A_{ij}-2n_iA_{ij}A_{jk}n_k
                        +\tfrac12(n_iA_{ij}n_j)^2.
\]
最后两项的球面平均分别是 $-2A_{ij}A_{ij}/3$ 与 $A_{ij}A_{ij}/15$，相加为 $2A_{ij}A_{ij}/5$。连同应变的 $2G/(c^4r)$、能流的 $c^3r^2/(32\pi G)$ 及立体角 $4\pi$ 恰给 $G/(5c^5)$。若把 $\Lambda$ 中的 $1/2$ 误作三维去迹的 $1/3$，这个检查立即失败。

依次取 $(l,\text{型})=(3,I),(2,J)$，得 $G I_{ijk}^{(4)}I_{ijk}^{(4)}/(189c^7)$、$16G J_{ij}^{(3)}J_{ij}^{(3)}/(45c^7)$。它们是振幅中 $v/c$ 修正在功率中的平方。同在相对 1PN 阶的还有四极矩自身 1PN 修正与最低阶四极矩的干涉。把上面两个正平方单独相加不可能给出完整的 1PN 功率，后者的修正系数完全可以为负。

为算一个可检验的例子，取无自旋两体，$M=m_1+m_2$、$\mu=m_1m_2/M$、$\Delta=(m_1-m_2)/M$，轨道间距 $a$，$\bm e=(\cos\phi,\sin\phi,0)$。在 Newton 质心系，两星位置分别是 $m_2a\bm e/M$ 和 $-m_1a\bm e/M$。取 $\dot\phi=\Omega$，$\bm l$ 为单位轨道法向，则最低阶矩为
\[
 I_{ij}=\mu a^2 e_{\langle i}e_{j\rangle},\quad
 I_{ijk}=-\mu\Delta a^3e_{\langle i}e_je_{k\rangle},\quad
 J_{ij}=-\mu\Delta a^3\Omega e_{\langle i} l_{j\rangle}.
\]
质量八极系数来自 $m_1(m_2/M)^3-m_2(m_1/M)^3=-\mu\Delta$。流矩中两个轨道角动量同向，但与位置相乘时后者异号，亦给 $-\mu\Delta$。等质量极限两项一起消失。为看时间导数的作用，可先做两时刻的缩并再微分。STF 加法公式给
\begin{align*}
 e_{\langle ijk\rangle}(t)e_{\langle ijk\rangle}(t')
 &=\tfrac25 P_3(\cos\delta)
 =\tfrac{3}{20}\cos\delta+\tfrac14\cos3\delta,\\
 e_{\langle i}(t)l_{j\rangle}e_{\langle i}(t')l_{j\rangle}
 &=\tfrac12\cos\delta,\qquad \delta=\Omega(t-t').
\end{align*}
分别对 $t,t'$ 求四次或三次导数，再取 $t'=t$，得到前式的因子 $\Omega^8(3/20+3^8/4)$ 与后式的 $\Omega^6/2$。质量八极的高次谐波虽在矩中只占有限比例，在能流中却受高阶导数的 $3^8$ 放大。由此例还可理解，仅依据某个质量矩在空间上的大小排各次谐波的功率，往往会错判主要频带。

\subsection{球面上电型与磁型的张量谐波}

对探测器而言，真正能区分的是每个方向上的两个波面分量。取单位二球度规 $q_{AB}$、体积形式 $\varepsilon_{AB}$、协变导数 $D_A$。由归一化标量球谐 $Y=Y_{lm}$ 造
\begin{align*}
 E^{lm}_{AB}&=N_l\left(D_A D_B-\tfrac12q_{AB}D^2\right)Y,\\
 B^{lm}_{AB}&=\tfrac{N_l}{2}
      (\varepsilon_A{}^C D_BD_C+\varepsilon_B{}^C D_AD_C)Y,\qquad
 N_l=\sqrt{\frac{2(l-2)!}{(l+2)!}} .
\end{align*}
二者都切于球面且无迹。$l=0$ 时梯度为零，$l=1$ 时二阶导数全为迹，因此只从 $l=2$ 起有这两族。电型的宇称为 $(-1)^l$，磁型的体积形式在球面反演下多变一次号，宇称为 $(-1)^{l+1}$。两族名称沿自电磁多极辐射，其构造只涉及同一引力波的张量场。

归一系数可由一次换序得到。写 $\lambda_l=l(l+1)$，二球的 Ricci 张量为 $q_{AB}$，故 $D^2D_AY=D_AD^2Y+D_AY=-(\lambda_l-1)D_AY$。连续两次分部积分给
\[
 \int |D_AD_BY|^2\,\mathrm d\Omega=\lambda_l(\lambda_l-1),\qquad
 \int |D_{\langle A}D_{B\rangle}Y|^2\,\mathrm d\Omega
   =\tfrac12\lambda_l(\lambda_l-2)=\frac{(l+2)!}{2(l-2)!}.
\]
同阶两族的内积因宇称为零，不同 $l,m$ 的正交性由自伴随的旋转生成元给出。于是保留辐射面切分量 $h_{AB}$，可逐模读出 $u_{lm}=r\int h_{AB}E_{lm}^{*AB}\,\mathrm d\Omega$ 及 $v_{lm}=r\int h_{AB}B_{lm}^{*AB}\,\mathrm d\Omega$，并由完备性重建 $rh_{AB}=\sum_{lm}(u_{lm}E^{lm}_{AB}+v_{lm}B^{lm}_{AB})$。若原场实，系数满足 $u_{l,-m}=(-1)^m u_{lm}^*$，$v$ 同理。在这个归一化下，能流就是 $\mathcal F=c^3\sum_{lm}(|\dot u_{lm}|^2+|\dot v_{lm}|^2)/(32\pi G)$。与式 \eqref{eq:gw-a-multiflux} 比，STF 归一化中那些双阶乘已计入球谐系数的定义。

作为例子，$Y_{20}=\sqrt{5/(16\pi)}(3\cos^2\theta-1)$。在单位正交切基上，$D_{\hat\theta}D_{\hat\theta}Y=Y''$、$D_{\hat\phi}D_{\hat\phi}Y=\cot\theta\,Y'$，两者之差为 $6\sqrt{5/(16\pi)}\sin^2\theta$。所以 $m=0$ 电型四极仅有加偏振，沿对称轴的振幅为零。将切基转动 $45^\circ$ 会把加偏振写成叉偏振，却不能使其成为磁型四极；电型、磁型由整个球面上的宇称确定。一台仪器在一条视线上测得的一个偏振分量无法代替全天的这种区分。

有时把 $h=h_+-\mathrm i h_\times$ 展为自旋权 $-2$ 球谐。取复切矢 $m^A=(e_{\hat\theta}^A+\mathrm ie_{\hat\phi}^A)/\sqrt2$，即 $h=h_{AB}\bar m^A\bar m^B$。若 $m^A\mapsto\exp(\mathrm i\psi)m^A$，这个缩并获因子 $\exp(-2\mathrm i\psi)$。这只记载偏振基的局域选择；角动量量子数仍取 $l\ge2,-l\le m\le l$。波形重建中常把不同 $m$ 称为不同谐波。只有在源具有匀速转动的对称性时，空间量子数 $m$ 才同时对应时域频率 $m\Omega$；在偏心运动或受扰进动中，一个固定 $(l,m)$ 模本身还能含有多个时间频率。

\subsection{从一套 STF 矩直接读出整球波形}

前面先从源积分得 $U_L,V_L$，又在天球上用积分定义了 $u_{lm},v_{lm}$。实际计算毋须先生成一幅高分辨率的整球波形再积分。设 $Y_{lm}=\mathcal Y_L^{lm}n_L$，$\mathcal Y_L^{lm}$ 为常数 STF 张量，按 $\int |Y_{lm}|^2\mathrm d\Omega=1$ 归一。由 STF 角积分公式，对一个给定的 $U_L$ 有
\begin{align*}
 U_L n_L&=\sum_{m=-l}^l U_{lm}^{s}Y_{lm},&
 U_{lm}^{s}
 &=\frac{4\pi l!}{(2l+1)!!}\mathcal Y_L^{lm*} U_L ,\\
 \sum_m|U_{lm}^{s}|^2&=\int |U_L n_L|^2\,\mathrm d\Omega
 =\frac{4\pi l!}{(2l+1)!!}\, U_L^*U_L .
\end{align*}
上标 $s$ 仅表示暂时作成一个标量角函数的系数；还须由它恢复切球张量。取 $e_A^i=D_A n^i$，在单位球上 $D_A D_B n^i=-q_{AB}n^i$。对上式左边两次作球面导数，有
\[
 D_A D_B(U_L n_L)=l(l-1)U_{ijL-2}e_A^i e_B^j n_{L-2}
                        -l q_{AB}U_L n_L .
\]
故去掉切面上的迹后只剩第一项的切面无迹部分。它正是式 \eqref{eq:gw-a-multipolar-wave} 给出的 $\Lambda$ 作用在质量矩上。于是标量、张量二种球谐只差一个明确的微分，质量型振幅为 $u_{lm}=\frac{4G\,U_{lm}^{s}}{c^{l+2}l!\,N_l\,l(l-1)}$。同样置 $V_L n_L=\sum_m V_{lm}^{s}Y_{lm}$，只须在其中一个球面切指标上作 $\varepsilon_A{}^C$ 旋转，再将两个指标对称化。以 $\varepsilon_{\hat\theta\hat\phi}=+1$ 和外向 $\bm n$ 的定向，$-n_d\epsilon_{de(a}V_{b)eL-2}$ 在屏幕上的作用正是这个旋转。相应 $v_{lm}$ 与上式相差 $V_{lm}^{s}$ 代 $U_{lm}^{s}$ 和一个 $2l/[c(l+1)]$ 的因子。反转对球面的定向会改变 $v_{lm}$ 的符号，但不改变其对整球通量的正平方。

现可不作高阶 $n_i$ 积分而核对式 \eqref{eq:gw-a-multiflux} 的任意阶系数。对 $u_{lm}$ 的时间导数平方求和，乘 $c^3/(32\pi G)$ 后，去掉 $G/c^{2l+1}$，所余质量型系数为
\[
 \frac{2}{l!(2l+1)!!\,N_l^2 l^2(l-1)^2}
   =\frac{(l+1)(l+2)}{l(l-1)\,l!\,(2l+1)!!}.
\]
流型的系数再乘 $4l^2/[c^2(l+1)^2]$，正是另一项。这一推导依赖球面无边界以及正交归一的张量谐波，对任意取向的有限观测网络却不能使用球面的正交性。只看到天空中的几个方向时，不同 $l,m$ 给出的时间序列可以高度相关。

以四极为例，写 $U_{xx}+U_{yy}=-U_{zz}$，将 $U_{ij}n_in_j$ 按方位 Fourier 次数排开即
\begin{align*}
 U_{ij}n_i n_j={}&\tfrac12 U_{zz}(3\cos^2\theta-1)\\
 &+2\sin\theta\cos\theta(U_{xz}\cos\phi+U_{yz}\sin\phi)\\
 &+\sin^2\theta[\tfrac12(U_{xx}-U_{yy})\cos2\phi+U_{xy}\sin2\phi].
\end{align*}
这三行分别只有 $m=0$、$|m|=1$、$|m|=2$。若取一对沿 $x$ 相向振动的等质量小球，$m=0,2,-2$ 都可能出现，它们可有同一个时间频率。若主轴绕 $z$ 转动，$m=\pm2$ 的复系数可获得相反的时间相位。若改绕倾斜轴转动，当改用沿新轴的源标架时只需用 $D^l$ 矩阵在同一个 $l$ 内混合这些系数。旋转坐标不可能由一个真实的八极矩产生四极矩。

有时场方程的自然基先固定的是轨道角动量，而边界条件给的是径向或切向分量。取恒定的五个三维 STF 极化张量 $t^{m_s}_{ij}$（$m_s=-2,\ldots,2$），可先写
\[
 {\cal T}^{Lj m}_{ij}(\bm n)
 =\sum_{m_L,m_s}\langle Lm_L,2m_s|jm\rangle
                   Y_{Lm_L}(\bm n)\,t^{m_s}_{ij}.
\]
这组谐波同时对角化 $J^2,J_z,L^2$，固定 $j$ 时 $L$ 从 $|j-2|$ 到 $j+2$。单个 $L$ 的张量一般有径向指标。要使辐射横向，须在同一 $j,m$ 下把几个允许的 $L$ 作特定组合。另一种等价作法是在每个方向先分成 $n_i n_j-\delta_{ij}/3$ 的标量、$n_{(i}v_{j)}$ 的切矢量、切面 STF 张量三种。它们分别有一、二、二个局域分量，合计五个。TT 投影杀掉前两种，余下两种恰是 $E^{jm}_{AB},B^{jm}_{AB}$。因此 $Y_{lm}$ 上单乘一个不变偏振张量一般不会得到同一 $l$ 的物理横波；源计算中采用 STF 投影与角动量耦合，正是为避免此误认。

\subsection{应力一次矩的不可约分解}

前文以 $S_{ij,k}$ 说明了八极和流四极的由来。现在可把分解写全，以便对有内部应力的源逐项检查系数。先从 $S_{ij,k}=S_{ji,k}$ 减去迹。经与 $n_k\Lambda_{ab,ij}$ 缩并，所有 $\delta_{ij}V_k,\ \delta_{ik}W_j+\delta_{jk}W_i$ 均给零。余部有七维的完全对称 STF 张量和五维的混合对称张量，这与 $18-3-3=12$ 相符。设 $H_{ijk}=S_{\langle ij,k\rangle}$、$C_{ij}={\rm STF}_{ij}\epsilon_{jab}S_{bi,a}$。则忽略迹项时 $S_{ij,k}=H_{ijk}+\tfrac13(\epsilon_{ki a}C_{ja}+\epsilon_{kj a}C_{ia})$。检验右边对三个指标对称化，含 $\epsilon$ 的项相消，故第一项等于 $H$。再用 $\epsilon_{jab}S_{bi,a}$ 投影，在 $S_{bi,a}$ 的第一项中 $\epsilon_{jab}\epsilon_{ab c}$ 给 $2\delta_{jc}$，另一项给 $C_{ij}$ 及一个为零的 $C$ 的迹，系数 $1/3$ 恰由此定出。用守恒律 $H_{ijk}=I^{(2)}_{ijk}/6$、$C_{ij}=\dot J_{ij}$，推迟展开的 $4G n_k\dot S_{ij,k}/(c^5r)$ 便同时重现两个 $v/c$ 振幅系数。这样无需先知道全阶多极公式，也能从三个源指标算出第一批高阶修正。

更一般地，若 $T_{iL}$ 只对末 $l$ 个指标 STF，它属于 $1\otimes l=(l+1)\oplus l\oplus(l-1)$，维数恒等式是 $3(2l+1)=(2l+3)+(2l+1)+(2l-1)$。三个分量可分别从 $T_{\langle iL\rangle}$、$\epsilon_{abi_l}T_{a,bL-1}$ 的 STF 部分及 $T_{a,aL-1}$ 读出。对守恒源，后一族常与时间分量的低阶矩关联。张量源还有第二个空间指标，构造的各族虽更多，但无穷远只传送横向的两族。这也解释了为何源计算中出现应力迹、流密度的导数等许多项，而辐射多极的最终类型并未随计算阶数增加。

\section{从质量分布求波形}\label{sec:gw-a-sources}

\subsection{刚体的两条谱线}

刚体的质量矩在本体坐标中为常数，随时间变化的是它的主轴在天空中的投影。设惯性张量 $\mathcal I_{ij}=\int\rho(y^2\delta_{ij}-y_iy_j)\,\mathrm d^3y$，质量四极矩正好是 $I_{ij}=-\mathcal I_{\langle ij\rangle}$。绕第三主轴匀速转动时，仅本体 1、2 方向的惯量差参与辐射。相对于固定坐标的 $\mathcal I_{xx}-\mathcal I_{yy}$、$2\mathcal I_{xy}$ 分别为 $(\mathcal I_1-\mathcal I_2)\cos2\Omega t$、$(\mathcal I_1-\mathcal I_2)\sin2\Omega t$。于是沿轴应变幅为 $4G|\mathcal I_1-\mathcal I_2|\Omega^2/(c^4r)$。即使物体转得很快，只要关于转轴严格轴对称，这两个量仍然为零。

轴对称体可否只靠转动发波？若其角动量 $\bm J$ 与对称轴 $\bm s$ 不平行，问题就不同了。令两个相等主惯量为 $A$，第三主惯量为 $C$，$\Delta I=C-A$。于是 $\mathcal I_{ij}=A\delta_{ij}+\Delta I s_is_j$。无外力矩时 $\bm J$ 固定，$\bm\omega=\bm J/A+(1/C-1/A)(\bm J\cdot\bm s)\bm s$，由 $\dot{\bm s}=\bm\omega\times\bm s$ 得
\[
 \dot{\bm s}=\frac{\bm J}{A}\times\bm s,\qquad
 \bm s=(\sin\alpha\cos\phi,\sin\alpha\sin\phi,\cos\alpha),\qquad
 \dot\phi=\Omega=J/A .
\]
这里 $z$ 轴沿 $\bm J$，$\alpha$ 为两轴夹角。$\bm s$ 在惯性系中的进动角频率是 $\Omega$；在本体坐标中角速度矢量绕第三轴的慢进动另有频率 $(C-A)\omega_3/A$，不可把后者直接当作波的基频。

取观测方向 $\bm n=(\sin\iota,0,\cos\iota)$，偏振基 $\bm p=(\cos\iota,0,-\sin\iota)$、$\bm q=(0,1,0)$。用四极波的 $h_+=G(\ddot I_{pp}-\ddot I_{qq})/(c^4r)$ 和 $h_\times=2G\ddot I_{pq}/(c^4r)$，只需计算 $(\bm p\cdot\bm s)^2-(\bm q\cdot\bm s)^2$ 与 $2(\bm p\cdot\bm s)(\bm q\cdot\bm s)$ 的二阶导数。记 $H=G\Delta I\Omega^2/(c^4r)$、$\phi=\Omega u$，便得
\begin{align*}
 h_+/H&=-\sin2\alpha\sin\iota\cos\iota\cos\phi
        +2\sin^2\alpha(1+\cos^2\iota)\cos2\phi,\\
 h_\times/H&=-\sin2\alpha\sin\iota\sin\phi
        +4\sin^2\alpha\cos\iota\sin2\phi .
\end{align*}
基频来自 $s_xs_z,s_ys_z$，二倍频来自 $s_x^2-s_y^2,2s_xs_y$。若沿 $\bm J$ 观察，基频恰消失。若在赤道面观察，二倍频只在加偏振，基频只在叉偏振。足够小的摆角使基频振幅为 $O(\alpha)$，二倍频为 $O(\alpha^2)$；因此缺少可见的二倍频不等于缺少自由进动。两个偏振中同一条谱线的幅比限制 $\iota$，再用两个不同频率的幅比限制 $\alpha$，尚须独立的距离才可把 $\Delta I$ 从总振幅中分开。

第三阶时间导数的 $xz,yz$ 分量给基频功率，$xx,yy,xy$ 分量给二倍频功率。分别平方并权计非对角元的两次出现，有
\begin{equation}
 \mathcal F_1=\frac{2G}{5c^5}(\Delta I)^2\Omega^6\sin^2\alpha\cos^2\alpha,\qquad
 \mathcal F_2=\frac{32G}{5c^5}(\Delta I)^2\Omega^6\sin^4\alpha .
 \label{eq:gw-a-precessflux}
\end{equation}
二者相等时 $\tan\alpha=1/4$（限 $0<\alpha<\pi/2$）。这些平均功率针对全空间；单台仪器中的两线功率比还受上面的 $\iota$ 和天线响应影响。

辐射反作用也改变摆角。源的转动能为 $E=J^2[\sin^2\alpha/A+\cos^2\alpha/C]/2$。两个模的 $(m,\omega)$ 分别是 $(1,\Omega)$ 和 $(2,2\Omega)$，故各自都给 $\dot J=\dot E/\Omega$。从 $\dot E=E_{,J}\dot J+E_{,\alpha}\dot\alpha$ 消去 $\dot E$ 得 $(1/A-1/C)J\cos\alpha (\cos\alpha\,\dot J-J\sin\alpha\,\dot\alpha)=0$。所以 $J\cos\alpha$ 保持不变。特别地 $\dot\alpha=(\dot J/J)\cot\alpha$，$\dot J=-(\mathcal F_1+\mathcal F_2)/\Omega$。在小摆角处 $J$ 的损失率随 $\alpha^2$ 下降，摆角本身近似指数衰减。它可以对齐到几乎不辐射而仍保留很大的轴向自转能。若有内部摩擦，$J$ 在短时间内几乎不变，此时应在固定 $J$ 下极小化 $E$：扁体 $C>A$ 趋向 $\alpha=0$，长体 $C<A$ 趋向 $\alpha=\pi/2$。内部耗散和外辐射可以驱动不同的长期演化。

\subsection{三个主惯量不等时的自由进动}

若 $A_1<A_2<A_3$，对称轴已不存在。刚体在无外力矩时仍有两个积分 $2E=\sum_i A_i\omega_i^2$、$J^2=\sum_i A_i^2\omega_i^2$，本体坐标中的 Euler 方程为 $A_1\dot\omega_1=(A_2-A_3)\omega_2\omega_3$ 及其循环置换。因在 $A_1$ 或 $A_3$ 轴附近的扰动振荡，在中间轴附近的扰动增长，这里取靠近 $A_3$ 轴的一支，$2EA_2<J^2<2EA_3$。在两个积分中消去 $\omega_1^2,\omega_3^2$，从而
\begin{align*}
 \omega_1^2&=\frac{2EA_3-J^2-A_2(A_3-A_2)\omega_2^2}
                   {A_1(A_3-A_1)},\\
 \omega_3^2&=\frac{J^2-2EA_1-A_2(A_2-A_1)\omega_2^2}
                   {A_3(A_3-A_1)} .
\end{align*}
$\omega_2$ 的方程平方后，右边即为两个关于 $\omega_2^2$ 的一次因子。令 $\omega_2=B y$，$B^2=(2EA_3-J^2)/[A_2(A_3-A_2)]$，化成
\[
 \dot y^2=\nu_p^2(1-y^2)(1-k_e^2 y^2),\qquad
 \nu_p^2=\frac{(A_3-A_2)(J^2-2EA_1)}{A_1A_2A_3},\quad
 k_e^2=\frac{(A_2-A_1)(2EA_3-J^2)}{(A_3-A_2)(J^2-2EA_1)} .
\]
范围 $0\le k_e^2<1$ 由选支保证。取坐标原点使 $y(0)=0,\dot y(0)>0$，直接求逆椭圆积分便是 $y=\mathrm{sn}(\nu_p t,k_e)$。其余分量为
\[
 \omega_1=\sqrt{\frac{2EA_3-J^2}{A_1(A_3-A_1)}}\,\mathrm{cn}(\nu_pt,k_e),
 \qquad
 \omega_3=\sqrt{\frac{J^2-2EA_1}{A_3(A_3-A_1)}}\,\mathrm{dn}(\nu_pt,k_e).
\]
可用 $\mathrm{cn}^2=1-\mathrm{sn}^2,\ \mathrm{dn}^2=1-k_e^2\mathrm{sn}^2$ 检查两个能量积分；各分量的初始符号另由 Euler 方程锁定。$k_e\to0$ 时三函数退成正弦、余弦和常数，回到近对称体；$k_e\to1$ 时完整椭圆积分 $K(k_e)$ 对数发散，表示刚体在接近不稳定中间轴时逗留很久。以很小的摆角作展开不能描述后一种情形。

这还只求得本体系中 $\bm\omega$ 的周期 $T=4K(k_e)/\nu_p$。要求固定空间基上的四极矩，尚须确定整个刚体绕 $\bm J$ 的转角。取第三主轴的单位矢量 $\bm s$，其与 $\hat{\bm J}$ 的夹角有 $\cos\alpha=A_3\omega_3/J$。若其方位为 $\phi$，用 $\dot{\bm s}=\bm\omega\times\bm s$ 和 $\hat{\bm J}\cdot(\bm s\times\dot{\bm s})=\dot\phi\sin^2\alpha$ 得
\[
 \dot\phi=\frac{J(A_1\omega_1^2+A_2\omega_2^2)}
                {A_1^2\omega_1^2+A_2^2\omega_2^2}.
\]
分子、分母均周期变化。将其写成平均角速度 $\bar\Omega$ 与零均值的周期函数之和，有 $\phi(t)=\bar\Omega t+\delta\phi(t)$。因此经过一整个本体周期，角速度在本体内回到初值，物体在固定系中却转过 $\bar\Omega T$。若把三个主轴的一周期都误当成惯性系中质量矩的周期，就会把两个本来不可公度的基频错并为一个。

记 $R(t)$ 为本体到惯性系的转动，$Q= -[R\,\mathrm{diag}(A_1,A_2,A_3)R^{\mathsf T}]_{\rm STF}$。经过一个 $T$ 有 $Q(t+T)=R_z(\bar\Omega T)Q(t)R_z^{\mathsf T}(\bar\Omega T)$。于是绕 $\bm J$ 的 $m$ 分量乘以 $\exp(-\mathrm i m\bar\Omega T)$（采用被动分量约定），除去连续因子 $\exp(-\mathrm im\bar\Omega t)$ 后是 $T$ 周期函数。Fourier 线上可出现 $m\bar\Omega+2\pi n/T$，其中 $|m|\le2$、$n$ 为整数。刚体的对称性会消去其中一些系数，靠近主轴时多数高边带极弱。这一构造只需先求 Euler 运动，再作一次转动与两次导数，无须为每条谱线重新求引力波方程。

比较两种小偏离的作用。若三轴性 $A_2-A_1$ 极小而摆角有限，主线是上小节的自由进动双线。若摆角趋零而 $A_2-A_1$ 有限，体轴绕自身转动产生 $2\omega_{\rm rot}$ 一条线。两者均小时常用 $\omega_{\rm rot}$ 记轴向自转、$\omega_p$ 记本体系慢进动，保留的主要正频角线为 $2\omega_{\rm rot}$、$\omega_{\rm rot}+\omega_p$ 与 $2(\omega_{\rm rot}+\omega_p)$。最后一项与摆角平方同阶，前两项另分别受三轴性、一次摆角支配。不能把任何三条相近谱线无条件解释为这三个频率；它们的相对间隔、两种偏振及随观测方向的消失条件应相容。

\subsection{从偏近点角逐项计算辐射谐波}

第三章给出了每个谐波的功率 $g_n(e)$。要重建有两个独立偏振的时域波形，仅有功率还不够，需要各矩分量在同一近心点时刻的相位。令 $\ell=\Omega(t-t_p)=u_b-e\sin u_b$ 为平近点角，其中 $u_b$ 是偏近点角。取轨道 $x$ 轴指向近心点，$x=a(\cos u_b-e),\ y=a\sqrt{1-e^2}\sin u_b$。作无量纲组合 $X=(x^2-y^2)/a^2$、$Y=2xy/a^2$、$Z=(x^2+y^2)/a^2$。将三角函数相乘，在 $u_b$ 中只需二次谐波，
\begin{align*}
 X&=(1-e^2/2)\cos2u_b-2e\cos u_b+3e^2/2,\\
 Y&=\sqrt{1-e^2}(\sin2u_b-2e\sin u_b),\\
 Z&=1+e^2/2-2e\cos u_b+(e^2/2)\cos2u_b .
\end{align*}
但探测器均匀记录的是 $t$，故应在 $\ell$ 上作 Fourier 展开。以近心点为零相的选择使 $X,Z$ 为偶、$Y$ 为奇。对整数 $k,n>0$ 记 $\cos ku_b$ 的第 $n$ 个余弦系数为 $C_{kn}$，$\sin ku_b$ 的第 $n$ 个正弦系数为 $S_{kn}$。例如
\[
 C_{kn}=\frac1\pi\int_0^{2\pi}\cos(ku_b)\cos(n\ell)\,\mathrm d\ell
       =\frac{k}{n\pi}\int_0^{2\pi}\sin(ku_b)
                \sin[n(u_b-e\sin u_b)]\,\mathrm du_b .
\]
第二步是对 $\mathrm d(\sin n\ell)$ 分部积分，故 Jacobian $1-e\cos u_b$ 不应再乘一次。同样对 $\mathrm d(\cos n\ell)$ 分部积分求 $S_{kn}$，积化和差就得到
\[
 C_{kn}=\frac{k}{n}[J_{n-k}(ne)-J_{n+k}(ne)],\qquad
 S_{kn}=\frac{k}{n}[J_{n-k}(ne)+J_{n+k}(ne)] .
\]
余弦的常数项须单算。$\langle\cos ku_b\rangle_\ell=-e\delta_{k1}/2$，正弦均值为零。常数项从矩取二、三阶导数后消去，求系数时却不能以 $\ell$ 平均等同于以 $u_b$ 平均。

记 $J_s=J_s(ne)$，用 $J_{s-1}+J_{s+1}=2sJ_s/(ne)$，在 $e>0$ 时可逐步简化 $X_n=(1-e^2/2)C_{2n}-2eC_{1n}$。例如
$e^2(J_{n-2}-J_{n+2})=2e(J_{n-1}-J_{n+1})-2e(J_{n-1}+J_{n+1})/n$，后一和为 $2J_n/e$。同样处理 $Y_n,Z_n$，有
\begin{align*}
 X_n&=\tfrac2n A_n,\quad
 A_n=J_{n-2}-2eJ_{n-1}+\tfrac2n J_n+2eJ_{n+1}-J_{n+2},\\
 Y_n&=\tfrac2n\sqrt{1-e^2}\, B_n,\quad B_n=J_{n-2}-2J_n+J_{n+2},\qquad
 Z_n=-\tfrac4{n^2}J_n .
\end{align*}
$X_n,Y_n,Z_n$ 是有符号系数。对给定的观测方向先把这些系数作 TT 投影，乘 $-(n\Omega)^2$ 便给出应变；若先逐项平方便不能重建偏振。特别地 $Z_n$ 虽是平面内的径向二次矩，除去的是三维迹，其贡献并不整个消失。

具体地，$I_{xx}/(\mu a^2)=Z/6+X/2$、$I_{yy}/(\mu a^2)=Z/6-X/2$、$I_{zz}/(\mu a^2)=-Z/3$、$I_{xy}/(\mu a^2)=Y/2$。于是第 $n$ 分量的三阶导缩并周期平均正比于 $(n\Omega)^6[Z_n^2/6+(X_n^2+Y_n^2)/2]/2$。除以同半长轴圆轨道的通量便得第三章所列的 $g_n$，其中 $4J_n^2/(3n^2)$ 恰来自 $Z_n^2/6$，前述三维去迹在这里可以检查。若取相位原点偏离 $t_p$，每个系数只多一个 $n$ 倍的相位转动，$g_n$ 不改变。

还有两种敏感的极限检验。先取 $e\ll1$，用 $J_s(z)=(z/2)^s/s!+O(z^{s+2})$（$s\ge0$）及 $J_{-s}=(-1)^s J_s$。保留功率到 $e^2$，得 $g_1=29e^2/192$、$g_2=1-5e^2$、$g_3=729e^2/64$，其余至少从 $e^4$ 起。相加为 $1+(157/24)e^2$，恰与总通量的偏心增强因子展开相同。三倍频比基频强很多，不仅因系数不同，也因三次时间导数的 $n^3$ 在功率中要平方。$e=0$ 应直接取极限，避免在前面的 Bessel 递推中除以零。

再取 $1-e\ll1$。近心点速度 $v_p^2=GM(1+e)/[a(1-e)]$，单位时间扫过的真近点角为 $v_p/[a(1-e)]=\Omega\sqrt{1+e}/(1-e)^{3/2}$。大部分辐射由很短的一段近心点运动给出，因此谱可延伸到 $n$ 约为 $(1-e)^{-3/2}$ 的谐波。这个尺度还出现在 Bessel 积分本身。$J_n(ne)$ 的相位在近心点邻域为 $n[(1-e)u_b+eu_b^3/6+\cdots]$，把 $u_b$ 缩放为 $(2/(ne))^{1/3}s$ 后，主部近似为
\[
 J_n(ne)\simeq (2/(ne))^{1/3}\,
 \mathrm{Ai}\!\left[(2/e)^{1/3}n^{2/3}(1-e)\right].
\]
当 $n^{2/3}(1-e)$ 增大，Airy 函数转为指数衰减。谱峰的精确位置尚受 $n^4$ 和各阶 Bessel 组合影响，以上估算定出的是适用的谐波范围。这也解释了为何对高偏心系统先固定少数 $n$ 再取 $e\to1$，无法得到正确的总能量。

\subsection{偏心旋近的寿命与每个频段的能量}

偏心率随半长轴的关系、Peters--Mathews 谐波功率和 $\dot P_b$ 的观测意义已见 \S\ref{sec:pn-matching} 后的双星讨论。这里把同一组衰减率积分，求一个高偏心系统在给定半长轴还可生存多久。记 $\beta=64G^3\mu M^2/(5c^5)$，$q=121/304$，第三章给出的轨迹常数为 $C_0$。由 $a=C_0 e^{12/19}(1+qe^2)^{870/2299}/(1-e^2)$ 代入 $\dot e$ 并取倒数，得到
\[
 t(e_0\to 0)=\frac{12 C_0^4}{19\beta}
   \int_0^{e_0}
     \frac{e^{29/19}(1+qe^2)^{1181/2299}}{(1-e^2)^{3/2}}\,\mathrm de .
\]
幂次可以逐项检验。$a^4/e$ 给 $e^{48/19-1}$，$a^4(1-e^2)^{5/2}$ 给 $(1-e^2)^{-3/2}$，余下 $(1+qe^2)^{3480/2299-1}$。下端 $e\simeq0$ 时积分近似为 $19e_0^{48/19}/48$，恰与 $C_0^4\simeq a_0^4 e_0^{-48/19}$ 抵消，得圆轨道寿命 $a_0^4/(4\beta)$。因此 $C_0$ 的形式奇点不会使圆形系统的合并时间发散。

在 $e_0\to1$ 处，$\int(1-e^2)^{-3/2}\mathrm de=e/\sqrt{1-e^2}$。保留奇异端的其余因子，且固定初始半长轴 $a_0$，有
\[
 \frac{t(e_0\to0)}{a_0^4/(4\beta)}
 \simeq\frac{48}{19(1+q)}(1-e_0^2)^{7/2}
       =\frac{768}{425}(1-e_0^2)^{7/2}.
\]
这是很陡的缩短。在固定近心距 $r_p=a_0(1-e_0)$ 而非固定 $a_0$ 时，反而有 $t\propto r_p^4(1-e_0)^{-1/2}$。因为每次在近心点损失的能量虽较大，相继两次经过近心点的等待时间也随半长轴增加。这一极限需在每次能量损失远小于轨道束缚能的区间内取；若一个近心点脉冲就能显著改变轨道，需逐圈演化，不能无限外推轨道平均寿命。

取圆轨道的一段绝热旋近，单位源系频带辐射的能量为 $\mathrm d E_{\rm GW}/\mathrm df_s=\mathcal F/\dot f_s=-\mathrm dE_{\rm orb}/\mathrm df_s$。以 $E_{\rm orb}=-\mu(GM\pi f_s)^{2/3}/2$ 化简便是
\[
 \frac{\mathrm dE_{\rm GW}}{\mathrm df_s}
    =\frac{(\pi G)^{2/3}}{3}\mathcal M^{5/3} f_s^{-1/3}.
\]
此谱已把源在每一频率停留的时间计入，所以指数不等于瞬时功率的指数。以偏心轨道为例，设轨道周频为 $F_b$，第 $n$ 次谐波 $f_s=nF_b$；当轨道缓慢演化时，任一指定 $f_s$ 得到的总能谱是 $\sum_n \mathcal P_n/(n\dot F_b)|_{F_b=f_s/n}$。各项所用 $e$ 应由 $C_0$ 求在该半长轴的数值。把同一时刻计算的一组谐波占比当作所有 $f_s$ 的修正因子，会遗漏这一差别。

再从能量算观测振幅。采用正频 Fourier 振幅 $\widetilde h(f)=\int h(t)\mathrm e^{-2\pi\mathrm i ft}\mathrm dt$，波穿过单位面积的谱能量为 $\pi c^3 f^2(|\widetilde h_+|^2+|\widetilde h_\times|^2)/(2G)$，由式 \eqref{eq:gw-effective-stress} 及 Parseval 恒等式可得。一个源系频带 $\mathrm df_s$ 中的能量到观者处少 $1+z$，而频带本身也少同一因子。于是对发射方向作平均后有
\[
 \frac{\pi c^3}{2G} f_o^2\!\left\langle|\widetilde h_+|^2+|\widetilde h_\times|^2\right\rangle
     =\frac{1}{4\pi D_G^2}
       \left.\frac{\mathrm dE_{\rm GW}}{\mathrm df_s}\right|_{f_s=(1+z)f_o}.
\]
这里 $D_G=D_L/(1+z)$，故左右量纲皆是单位面积、单位频率的能量。代入旋近的 $f_s^{-1/3}$ 即得 $|\widetilde h|\propto \mathcal M_z^{5/6}f_o^{-7/6}/D_L$。这是从能量输运对频域波形所作的独立检查，也便于把许多不能逐一分辨的双星谱相加。

\subsection{径直落下的点质量与扩展天体}

径向落体没有轨道角动量，却能使质量四极矩改变。考虑质量 $m\ll M$ 的粒子自远处静止径直落向质量 $M$ 的黑洞，先在 $r\gg GM/c^2$ 区域用 Newton 轨道，$\dot r=-\sqrt{2GM/r}$。记观者距离为 $D$，落体轴为单位矢量 $\bm s$，则 $I_{ij}=m r^2s_{\langle i}s_{j\rangle}$。相继求导，
\[
 (r^2)^{(2)}=2(\dot r^2+r\ddot r)=2GM/r,\qquad
 (r^2)^{(3)}=-2GM\dot r/r^2 .
\]
所以 $h_+=2G^2Mm\sin^2\iota/(c^4Dr)$、$h_\times=0$（偏振基的一轴取为 $\bm s$ 的投影），$s_{\langle i}s_{j\rangle}s_{\langle i}s_{j\rangle}=2/3$ 给
\[
 P=\frac{16G^4M^3m^2}{15c^5r^5},\qquad
 E(r>r_0)=\int_{r_0}^{\infty}\frac{P}{|\dot r|}\,\mathrm dr
      =\frac{16\sqrt2}{105}\frac{m^2}{M}c^2
         \left(\frac{GM}{c^2r_0}\right)^{7/2}.
\]
最后一个积分中的 $r^{-9/2}$ 使能量主要出自最靠近 $r_0$ 的几段。把截止半径取在若干引力半径，便可读出强场极限的量纲 $m^2c^2/M$，但不能由这个 Newton 计算定出强场的数值系数。靠近光子球时传播势垒和黑洞振铃都会改变波形。

在弱场段积分 $\dot r$ 得 $r=[9GM(t_*-t)^2/2]^{1/3}$，其中 $t_*$ 是形式上落到原点的时刻。故应变早期部分正比于 $(t_*-t)^{-2/3}$。在低频极限对远早于强场终止的一段作 Fourier 积分，$\int_0^\infty s^{-2/3}\mathrm e^{\mathrm i\omega s}\mathrm ds=\Gamma(1/3)\mathrm e^{\mathrm i\pi/6}\omega^{-1/3}$。于是谱能量 $\mathrm dE/\mathrm d\omega\propto\omega^2|\widetilde h|^2\propto\omega^{4/3}$。这条幂律取决于从无穷远静止落下的早期边界条件。若起点有限或初速度不为零，要以实际轨道替换，尤其不能以这一无限历时的幂律代表任意一次短落体。

若落下的是半径 $R_*$ 的恒星，强场中的点粒子解还需与星体形变比较。由 $GM R_*/r_t^3\sim Gm/R_*^2$ 得潮汐半径 $r_t\sim R_*(M/m)^{1/3}$。与黑洞引力半径的比为 $r_t/(GM/c^2)\sim [R_*c^2/(Gm)](m/M)^{2/3}$，故增大黑洞质量虽会扩大 $r_t$，却更快扩大视界，相同恒星最终可能在瓦解前就已穿过视界。若在视界外瓦解，星的不同部分受激时间不同，需比较到时宽度与波周期。

具体地，以许多小质量元各按相同的点粒子波形落入，但第 $A$ 个质量元的到时迟了 $\tau_A$。记单位质量波形为 $\widetilde h_1$，线性叠加得 $\widetilde h=m\widetilde h_1\sum_A w_A\exp(-2\pi\mathrm i f\tau_A)$，$w_A=m_A/m$。若 $\tau_A$ 独立取自密度 $p(\tau)$，平均能量的形状因子为
\[
 \left\langle\left|\sum_A w_A\mathrm e^{-2\pi\mathrm i f\tau_A}\right|^2\right\rangle
 =\sum_A w_A^2+\left(1-\sum_A w_A^2\right)|\widetilde p(f)|^2 .
\]
第一项是各元自项，第二项才随相位散布而减小。低频时 $\widetilde p(0)=1$，恢复总质量的相干能量；高频对等质量元只剩 $1/N$。若连续物质在宽 $\Delta\tau$ 的区间均匀落入，则相干因子为 $[\sin(\pi f\Delta\tau)/(\pi f\Delta\tau)]^2$，其零点反映特定的均匀分布，并非所有被瓦解星的共同性质。可以用整个星的到时分布同点源强场响应作卷积；若不同质量元各有不同角动量和轨道，单一卷积核也不够。

\subsection{潮汐瓦解的形状因子}

上例的均匀到时只是说明相干性。若要把临界频率同星的物理半径联系起来，可先取一均匀球形星，在质心跨过 $r_t$ 的时刻释放各质元，随后都在同一个中心势中下落。设星内沿径向的位移为 $\delta r$。从无穷远静止的点粒子轨迹可写为 $t=t_A-2r^{3/2}/(3\sqrt{2GM})$；在同一释放时刻分别令 $r=r_t+\delta r$，可得
\[
 t_A=t_c+\sqrt{\frac{r_t}{2GM}}\,\delta r
                 +O\!\left(\frac{(\delta r)^2}{\sqrt{GM r_t}}\right).
\]
所有质元处于半径为 $R_*$ 的均匀球中。置 $s=\delta r/R_*$，等 $s$ 截面的质量比例正比于 $1-s^2$，归一化后 $p(s)=3(1-s^2)/4$。于是令 $w=\omega R_*\sqrt{r_t/(2GM)}$，相干因子为
\[
 F(w)=\tfrac34\int_{-1}^1(1-s^2)\,\mathrm e^{\mathrm iws}\,\mathrm ds
     =3\frac{\sin w-w\cos w}{w^3} .
\]
只须对 $2\sin w/w$ 再作两次 $w$ 微分便得最后一式。小 $w$ 时 $F(w)=1-w^2/10+O(w^4)$，其中 $1/10$ 也可由 $\langle s^2\rangle/2$ 核对；大 $w$ 时 $F\sim-3\cos w/w^2$，故能谱乘上 $\sim w^{-4}$ 的振荡包络。与均匀到时模型的 $w^{-2}$ 能谱包络相比，此处较快的衰减来自球截面在两端变为零。已知点粒子能谱后，形状因子提供了一种把几何尺寸与强场传播近似分开的办法。

临界约化波长并非总为未受扰星的 $R_*$。由 $r_s=2GM/c^2$，条件 $w\sim1$ 等于 $\omega^{-1}\sim(R_*/c)\sqrt{r_t/r_s}$，可解释为等落入时刻切片上星的径向拉长。若黑洞振铃最强的 $\omega_{\rm peak}\sim c^3/(GM)$ 对应 $w_{\rm peak}\gg1$，瓦解会特别压低本来最亮的一段。若 $w_{\rm peak}\ll1$，有限大小主要压高频尾，总能量仍可近似点质量的估计。若星在临界面附近就有不可略的压强或自转，各质元释放时的速度也会有分布，这个只含一个 $w$ 的表达式须重建。

还可把相干性问题反过来。对于已知点波形，以 $F(\omega)$ 的相位保留到时分布信息，其模平方只给到时分布的自相关；一个分布和整体平移、时间反演后的分布都有相同模平方。因此即使完整测得归一化的能谱，也不能由上式唯一重建被拉长星的密度剖面。若同时测得复波形并有可靠的点源相位模型，则在 $F$ 的若干零点外才有反卷积的可能。

\subsection{加速粒子的软谱与永久位移}

式 \eqref{eq:soft-graviton-radiated-energy} 给出了入射、出射各腿的干涉。其 $\omega\to0$ 极限为常数，表示有限的每单位频率能量，却允许应变在过程前后有不同的常数极限。为何如此？令 $h(u)$ 从 $h_i$ 变到 $h_f$，其变化率在有限历时内可积，则 $\widetilde{\dot h}(0)=h_f-h_i$。对于正的非零频率，按上一节的 Fourier 约定有 $\widetilde{\dot h}=2\pi\mathrm i f\widetilde h$；故 $f^2|\widetilde h|^2\to |h_f-h_i|^2/(4\pi^2)$。在点粒子模型中，由推迟积分在各段匀速轨迹上的 Jacobian 可把这个差直接写为
\begin{equation}
 \Delta h_{ij}^{\rm TT}
   =\frac{4G}{c^4D}\Lambda_{ij,kl}(\bm n)
     \left[\sum_A\frac{E_A\beta_{Ak}\beta_{Al}}
                            {1-\bm n\cdot\bm\beta_A}\right]_{\rm in}^{\rm out},
 \qquad \bm\beta_A=\bm v_A/c .
 \label{eq:gw-a-memory}
\end{equation}
称为普通或线性\textbf{记忆}。变化后局域曲率可已归零，相邻质点的固有间距却改变了 $\Delta h_{ij}^{\rm TT}\xi^j/2$ 的相应分量。只有无界运动确实改变远前与远后速度分布时，才可直接以式 \eqref{eq:gw-a-memory} 估计。若系统的所有质量最终回到原来的静止分布，这一粒子项消失，仍可能有被辐射本身携走的能量所生的非线性记忆。

举一条匀速腿为例，其角因子是 $\sin^4\theta/(1-\beta\cos\theta)^2$，其中 $\theta$ 是视线与速度的夹角。分子是把 $\beta_i\beta_j$ 投在屏幕上再取无迹所留的四次正弦，分母来自推迟 Jacobian 的平方。这可代表轻粒子突然从静止加速的软频角图，给粒子反冲的重物虽也应列入守恒源，其波幅在重质量极限受到抑制。令 $y=\cos\theta$，非零极值满足 $\beta y^2-2y+\beta=0$，物理解 $y=(1-\sqrt{1-\beta^2})/\beta$。当 $\Gamma=(1-\beta^2)^{-1/2}\gg1$，峰角约为 $\sqrt{2/\Gamma}$。然而在任意固定的 $0<\theta<\pi$，角因子趋于 $(1+\cos\theta)^2$，没有把辐射都压到 $1/\Gamma$ 小锥内。前向处反因横向投影出现一个很窄的缺口。场幅在不同散射腿间还可相消，不能只给每条腿单独的能谱再相加。

\subsection{两体碰撞中各条腿的干涉}

碰撞过程中只有总应力能动张量守恒。在一种极限中可以只算一条被加速的轻粒子；若两粒子的质量可比，必须把所有入射段和出射段相干相加。用式 \eqref{eq:soft-graviton-radiated-energy} 还能在不求轨迹细节的条件下作完全的软频角积分。暂取 $G=c=1$，第 $A$ 条腿的能量、动量为 $E_A,\bm p_A$，约定未来指向 $p_A$，入、出射由 $\eta_A=\mp1$ 标记。需积的角量为
\[
 \mathcal I_{AB}=\int\!\frac{\mathrm d\Omega}
 {(E_A-\bm n\cdot\bm p_A)(E_B-\bm n\cdot\bm p_B)} .
\]
两个分母都正。用 $1/(ab)=\int_0^1 [sa+(1-s)b]^{-2}\,\mathrm ds$，置 $P^\mu=sp_A^\mu+(1-s)p_B^\mu$，对 $\bm P$ 为极轴积分得
\[
 \int\!\frac{\mathrm d\Omega}{(P^0-|\bm P|\cos\theta)^2}
   =\frac{4\pi}{(P^0)^2-\bm P^2},\quad
 \frac{\mathcal I_{AB}}{4\pi}
 =\int_0^1\!\frac{\mathrm ds}
 {s^2m_A^2+(1-s)^2m_B^2+2s(1-s)m_Am_Bw_{AB}} ,
\]
其中 $w_{AB}=-p_A\cdot p_B/(m_Am_B)\ge1$ 是相对 Lorentz 因子。再以 $t=sm_A/[(1-s)m_B]$ 换元，剩下 $\int_0^\infty\mathrm dt/(t^2+2w_{AB}t+1)$，前因子为 $1/(m_Am_B)$。置 $w_{AB}=\cosh\chi_{AB}$，分母等于 $(t+\mathrm e^\chi)(t+\mathrm e^{-\chi})$，部分分式给出 $\mathcal I_{AB}=4\pi\chi_{AB}/(m_Am_B\sinh\chi_{AB})$。在 $\chi=0$ 处按极限取 $\chi/\sinh\chi=1$。于是
\[
 \left.\frac{\mathrm dE}{\mathrm d\omega}\right|_{\omega\to0}
 =\frac{2}{\pi}\sum_{A,B}\eta_A\eta_Bm_A m_B
 \left(\cosh^2\chi_{AB}-\tfrac12\right)
 \frac{\chi_{AB}}{\sinh\chi_{AB}} .
\]
原式的一项虽含看似参考系相关的 $E_A,\bm p_A$，角积分后却只依赖于惯性质量和相对快速度。整体系若不改变，入射腿与对应出射腿处处抵消，总结果为零。若在求模平方前就把不同腿的振幅当成统计独立量，连这个零都无法恢复。

再看一次非相对论弹性碰撞。由于在质心系 $\sum m_A\bm v_A=0$，整个系统的速度二次矩为 $\sum m_A v_A^i v_A^j=\mu v^iv^j$，$v$ 是相对速度，弹性使其大小在前后相同。令 $B_{ij}=\mu(v_i^{\rm out}v_j^{\rm out}-v_i^{\rm in}v_j^{\rm in})$，则 $B_{ii}=0$。式 \eqref{eq:gw-a-memory}（仍取几何单位）变成 $\Delta h^{\rm TT}_{ij}=4\Lambda_{ij,ab} B_{ab}/D$，而 $B_{ij} B_{ij}=2\mu^2v^4\sin^2\Theta$，$\Theta$ 为质心散射角。由式 \eqref{eq:arbitrary-source-radiated-energy} 的正频约定和前面的 STF 角积分，得
\[
 \left.\frac{\mathrm dE}{\mathrm d\omega}\right|_{0}
 =\frac{4}{5\pi} B_{ij}B_{ij}
 =\frac{8}{5\pi}\mu^2v^4\sin^2\Theta .
\]
可直接从快速度表达式将各腿的 $O(v^0),O(v^2)$ 依守恒消掉，首次非零也在 $v^4$。前向不散射当然为零；正对碰后两条相对速度只反号，速度张量积也未改变，故 $\Theta=\pi$ 的这个软极限同样为零。这并不妨碍碰撞在有限频率发射脉冲，因为有限历时内的四极矩仍会改变。一个作用时间为 $\Delta t$ 的碰撞若只知其初末速度，则上述常数只控制 $\omega\Delta t\ll1$，把它从零积到任意高频得出的总能量没有依据。

\section{源矩到达无穷远之前}\label{sec:gw-a-pn}

\subsection{外部真空解的逐阶构造}

\S\ref{sec:pn-matching} 已把近区和远区的展开域区分开来，\S\ref{sec:exp-pn} 给出保守轨道所需的近区势。现固定一种匹配算法，以确定远处波形上的系数。取谐和坐标，设 $\mathfrak h^{\mu\nu}=\sqrt{-g}g^{\mu\nu}-\eta^{\mu\nu}$（其线性项为负的迹反转扰动），满足 $\partial_\nu\mathfrak h^{\mu\nu}=0$。在物质以外作后 Minkowski 展开 $\mathfrak h=\sum_{n\ge1}G^n h_n$。不先展开 $u=t-r/c$，逐阶方程具有三角结构
\[
 \square h_1^{\mu\nu}=0,\qquad
 \square h_n^{\mu\nu}=\Lambda_n^{\mu\nu}[h_1,\ldots,h_{n-1}],\qquad
 \partial_\nu h_n^{\mu\nu}=0 .
\]
$\Lambda_n$ 取自非线性 Einstein 张量的 $G^n$ 系数。二阶已有 $h_1\partial^2h_1$ 和 $\partial h_1\partial h_1$，所以即便物质完全局限于一个小球，二阶有效源也延伸到无穷远。先不按 $1/c$ 展开可在一次迭代中保留同一 $G$ 阶的全部传播时延。

在线性阶，出射解可由式 \eqref{eq:gw-a-stf-ret} 中的 $\partial_L[F_L(u)/r]$ 作基。取适当坐标，它只需质量、流两族矩 $M_L,S_L$；一般谐和坐标则另有一个满足 $\square\varphi_1^\mu=0$ 的规范矢量。将 $\varphi_1^\mu$ 的时空分量各作 STF 分解，通常记其中的四族函数为 $W_L,X_L,Y_L,Z_L$。$M$ 是单极，$M_i$ 在质心系可为零，$S_i$ 是守恒的总角动量。较高矩随 $u$ 变化，可用它们来标记源外解。若在线性阶直接删掉这四族规范矩，再用不同的 $h_1$ 构造 $\Lambda_2$，所得二阶项会有差别。要在高阶保持同一物理解，必须同时变换各阶场并重新定义 $M_L,S_L$。故称两族 $M_L,S_L$ 为\textbf{正则多极矩}，而匹配到近区积分所得的记作\textbf{源矩} $I_L,J_L$；线性阶二者相等。

若 $h_1$ 只有单极，$h_1\propto1/r$，已会有 $\Lambda_2\propto r^{-4}$。将只在源外成立的式子延拓到原点再直接作三维推迟积分，会遇到 $\int_0 r^{-2}\mathrm dr$。这个发散来自不适用的延拓，并非天体中心真的有这种场。Blanchet--Damour 构造先给源乘 $(r/r_0)^B$，对复数 $B$ 定义推迟积分，解析延拓回 $B=0$ 后取 Laurent 展开的常数项，称为\textbf{有限部}，记作 $\mathrm{FP}_{B=0}$。当原点和无穷远的收敛实部条件不相容，可先用一个中间半径把积分分成两段，各在收敛半平面定义，再延拓相加。固定 $r_0$ 使复幂的底无量纲。以最简单的对数径向积分为例，
\[
 \int_0^R(r/r_0)^B\frac{\mathrm dr}{r}
       =\frac{1}{B}\left(\frac{R}{r_0}\right)^B
       =\frac1B+\ln(R/r_0)+O(B).
\]
由此外部的有限部出现任意尺度 $r_0$。其在可测量中或与匹配内解的同尺度项抵消，或等效为辐射时刻的原点重选；若一条最终可观测预言仍随任意 $r_0$ 改变，就少算了相同阶的另一项。

现在令 $u_n^{\mu\nu}=\mathrm{FP}_{B=0}\,\square_{\rm ret}^{-1}[(r/r_0)^B\Lambda_n^{\mu\nu}]$。它解对了波动方程，却还不能直接作为 $h_n$，因为 $\partial_\nu$ 会作用到正则因子上。利用低阶方程可得 $\partial_\nu\Lambda_n^{\mu\nu}=0$（源外），从而
\[
 w_n^\mu=\partial_\nu u_n^{\mu\nu}
 =\mathrm{FP}_{B=0}\,\square_{\rm ret}^{-1}
       \left[B(r/r_0)^B(n_i/r)\Lambda_n^{\mu i}\right],
 \qquad \square w_n^\mu=0 \quad(r>0).
\]
显式因子 $B$ 并不足以使它为零，被积函数的延拓可有 $1/B$ 的极点。从 $w_n^\mu$ 的齐次多极展开造另一个出射齐次解 $v_n^{\mu\nu}$，使 $\partial_\nu v_n^{\mu\nu}=-w_n^\mu$，则 $h_n=u_n+v_n$ 同时满足两条方程。高阶规范的这一步在解普通标量波方程时并不存在。

还可直接看见某些对数的来历。冻结缓变时间系数，在三维有 $\Delta(r^a\hat n_L)=[a(a+1)-l(l+1)]r^{a-2}\hat n_L$。源若为 $r^{l-2}\hat n_L$，欲反演时会需 $a=l$，方括号为零。对 $a$ 微分再取 $a=l$ 得特解 $r^l\ln(r/r_0)\hat n_L/(2l+1)$。改变 $r_0$ 只多一个齐次的 $r^l\hat n_L$，恰需匹配定出。外部解向近区展开一般既有 $r^p$ 又有 $(\ln r)^q$，不能预先要求单纯的 $r$ 次幂级数。

\subsection{内区正则解怎样选定}

对扩展、慢运动的物质源，近区把谐和场方程看成 $\Delta\mathfrak h=c^{-2}\partial_t^2\mathfrak h+16\pi G\tau/c^4$，$\tau^{\mu\nu}$ 含物质和引力场。这相当于依次反演 Poisson 方程。若只写它的特解，算符形式为
\[
 \mathcal I^{-1}[s]=\mathrm{FP}_{B=0}
   \sum_{k\ge0}\frac{1}{c^{2k}}\partial_t^{2k}
        \Delta^{-k-1}[(r/r_0)^B s],\qquad
 \Delta^{-k-1}s=-\frac{1}{4\pi(2k)!}\int
      |\bm x-\bm y|^{2k-1}s(\bm y)\,\mathrm d^3y .
\]
正则因子在积分中取在 $\bm y$ 处。检验第二式只需反复用 $\Delta R^p=p(p+1)R^{p-2}$，最后用 $\Delta(1/R)=-4\pi\delta^{(3)}(\bm R)$。$k=0$ 就是 Newton 反演。$k$ 增加时核反随 $|\bm y|$ 增长，场应力的远处积分一般不再收敛，这一次有限部处理的是把近区展开延伸到无穷远带来的困难。

近区有一族在原点正则的齐次波，标量例子可写成
\[
 \hat\partial_L\!\left[
   \frac{A_L(t-r/c)-A_L(t+r/c)}{2r}\right].
\]
单独延拓出射项会有 $1/r$，一减一加后分子为奇函数，除以 $r$ 在原点只余偶次幂。以 $l=2$ 为例，
\[
 \frac{A(t-r/c)-A(t+r/c)}{2r}
  =-\frac{A^{(1)}}{c}-\frac{r^2A^{(3)}}{6c^3}
                    -\frac{r^4A^{(5)}}{120c^5}+\cdots .
\]
对第一项作 $\hat\partial_{ij}$ 为零，第二项因纯迹也为零，第三项用 $\hat\partial_{ij}r^4=8\hat x_{ij}$ 首次给非零。这正是质量四极反作用中第五时间导数和 $c^{-5}$ 的阶数由来。与之相比，上面的瞬时算符本身只有 $c$ 的偶次逆幂，单做逐阶 Poisson 反演永远选不出这种耗散的边界条件。这里近区的超前函数是一个正则齐次解的表示；匹配完成后，总场在过去零无穷处仍可严格无入射波。

设外展开记为 $\mathcal M(h)$，上横线表示 PN 展开。两解在源尺度 $a$ 与约化波长 $c/\Omega$ 之间的交叠区满足 $\overline{\mathcal M(h)}=\mathcal M(\bar h)$。先作外解近区展开，用它定出内区各 $A_L(t)$；反过来，外解各出射矩由内解作源外多极展开给出。在满足过去充分平稳、孤立的边界条件下，这两个方向共同决定自由函数。其结果之一是源矩仍可写作式 \eqref{eq:gw-a-sourceI}、\eqref{eq:gw-a-sourceJ} 的形式，但把 $T^{\mu\nu}$ 换为 $\bar\tau^{\mu\nu}$，并在整个空间积分上取有限部。$T^{\mu\nu}$ 的物质支撑有界，$\tau^{\mu\nu}$ 的引力部分却没有边界，不能把积分只算到星体表面。

对点质量的替代还需另一种正则化。上面针对的是把渐近展开用在错误区域所造成的发散；把单颗致密星缩成点以后，在粒子位置又有自势的紫外发散。可先在两星各自的缓冲区把内部解和外部 PN 度规相接，使质量包含固有结合能，再取点粒子极限。计算高 PN 阶的点源积分时也可把空间维数延拓为 $d=3+\varepsilon$，以解析延拓处理自场，最后再回到三维。以任意截止球粗暴删去粒子附近一块，在相同 PN 阶可留下有限项差异。最后须以同一可观测量（如圆轨道的 $E(\Omega)$）比较，才可分清坐标或世界线位置的重定义与真正物理差别。

另一种办法 DIRE 直接在过去零锥上分区积分。内区 $|\bm y|<\mathcal R$ 用慢运动展开算，外区用推迟时间和方向改写积分。分界半径 $\mathcal R$ 在波长量级选择，单独两区可各含 $\mathcal R$ 的幂或对数。一个简单检查是对精确积分作 $\mathrm d/\mathrm d\mathcal R$，内区的边界面元与外区下限面元等大反号。截断到某一阶后，在两侧都展开到该阶仍须有同样消去。\S\ref{sec:pn-matching} 所举的任意分界球即属此法；它与前面用解析延拓的匹配给出同一孤立源的波形。

\subsection{波尾的实部、相位与记忆}

近源的 $I_L,J_L$ 要先换成正则矩，再传播成式 \eqref{eq:gw-a-multipolar-wave} 中的 $U_L,V_L$。单极 $M$ 与某一辐射多极混合的效应特别直观。弱静态外场中径向零线近似满足 $\mathrm dt/\mathrm dr=(1+2GM/(c^2r))/c$，积分给 $t=r/c+(2GM/c^3)\ln(r/b)+\mathrm{const}$。因而谐和坐标的平直推迟时刻须改成 $U=t-r/c-(2GM/c^3)\ln(r/b)+\cdots$，才使同一波前在 $r\to\infty$ 仍有有限标记。$b$ 是辐射坐标的参考长度。这个坐标修正消掉的是场中随 $\ln r$ 增长的相位；散射造成的延迟分布仍真实地留在波中。

令 $M_L$ 在远过去的高阶导数消失。领头的质量型尾项为
\begin{align}
 U_L(U)&=M_L^{(l)}(U)
 +\frac{2GM}{c^3}\int_0^\infty M_L^{(l+2)}(U-\tau)
             [\ln(c\tau/(2b))+\kappa_l]\,\mathrm d\tau+\cdots ,
 \label{eq:gw-a-tails}\\
 \kappa_l&=\sum_{k=1}^{l-2}\frac1k+
       \frac{2l^2+5l+4}{l(l+1)(l+2)} .\notag
\end{align}
流型的同阶尾将 $M_L$ 换成 $S_L$，常数换成 $\pi_l=\sum_{k=1}^{l-1}1/k+(l-1)/[l(l+1)]$。例如 $\kappa_2=11/12$、$\pi_2=7/6$。若把 $b$ 乘以 $\exp\delta$，积分增加 $-2GM\delta M_L^{(l+1)}(U)/c^3$，等价于将直接项的时间参数平移。它是检验尾积分和坐标变换是否采用相容约定的简便办法。

要从尾积分读出实际的能流修正，取 $M_L^{(l)}(U)\propto\exp(-\mathrm i\omega U)$、$\omega>0$，在过去加一个绝热因子 $\exp(-\varepsilon\tau)$。利用
\[
 \int_0^\infty \tau^{s-1}\mathrm e^{(\mathrm i\omega-\varepsilon)\tau}\,\mathrm d\tau
     =\Gamma(s)(\varepsilon-\mathrm i\omega)^{-s},
\]
在 $s=1$ 求导，再令 $\varepsilon\to0^+$，得到
\[
 \int_0^\infty\mathrm e^{\mathrm i\omega\tau}
     [\ln(c\tau/(2b))+\kappa_l]\,\mathrm d\tau
 =-\frac{\pi}{2\omega}
 +\frac{\mathrm i}{\omega}[\kappa_l-\gamma_E-\ln(2b\omega/c)] .
\]
乘 $M_L^{(l+2)}/M_L^{(l)}=-\omega^2$，振幅的实修正即 $\pi GM\omega/c^3$，虚部只改相位。功率取模平方使实修正加倍；对主模 $\omega=2\Omega$ 得 $4\pi GM\Omega/c^3=4\pi x^{3/2}$，恰好是 \S\ref{sec:pn-wave-energy} 的半整数阶。这里的 $\pi$ 来自取出射 Green 函数的边界值，与某个局域保守力的系数无关。更高次的 $M\times M\times M_L$ 散射还会产生尾的尾，须与源矩及项间干涉合并计算。

另取一束向方向 $\bm n'$ 传播的辐射。其 4-动量 $p^\mu=(E/c)(1,\bm n')$，可在式 \eqref{eq:gw-a-memory} 的角核中取 $\bm\beta=\bm n'$。把每束带走的能量以 $\mathrm dE/(\mathrm du'\mathrm d\Omega')$ 表示，就得零质量辐射的永久位移
\begin{equation}
 h_{ij}^{\rm mem}(u)=\frac{4G}{c^4D}\Lambda_{ij,kl}(\bm n)
   \int_{-\infty}^u\!\mathrm du'\int\mathrm d\Omega'\,
   \frac{n'_k n'_l}{1-\bm n\cdot\bm n'}
                 \frac{\mathrm dE}{\mathrm du'\mathrm d\Omega'} .
 \label{eq:gw-a-nonlinearmem}
\end{equation}
能量若由波本身提供，这一项即为非线性引力波记忆。积分核的分母在 $\bm n'=\bm n$ 消失，但经 TT 投影的分子正比于两方向夹角的正弦平方，可积。各向同性辐射积分只给 $\delta_{kl}$ 和 $n_kn_l$，两者都被投影消去。记忆所量的是整个先前辐射的各向异性能量；尾则是过去源矩受到曲率散射的线性组合，二者虽都为历史积分，计算所需的资料不同。

用非进动圆轨道检验这个角核。令 $\bm l$ 为轨道法线，第一个四极波的角能量分布正比于 $1+6z'^2+z'^4$，其中 $z'=\bm l\cdot\bm n'$。对全天平均给 $1+6/3+1/5=16/5$，故一段时间释放总能量 $E$ 时 $\mathrm dE/\mathrm d\Omega'=5E(1+6z'^2+z'^4)/(64\pi)$。取观测方向为 $z$ 轴，$\bm l$ 落在 $xz$ 面，与视线夹角为 $\iota$。令 $w=n'_z$，则 $z'=w\cos\iota+\sqrt{1-w^2}\sin\iota\cos\phi'$。对加偏振缩并，角核化为 $(1+w)\cos2\phi'$，对叉偏振则是 $(1+w)\sin2\phi'$。后者积分为零，前者只须在多项式中留下 $\cos^2\phi',\cos^4\phi'$，二者乘 $\cos2\phi'$ 后的积分均为 $\pi/2$。执行 $w\in[-1,1]$ 的积分得 $\Delta h_+^{\rm mem}=\frac{GE}{24c^4D}\sin^2\iota\,(17+\cos^2\iota)$。例如边缘方向 $\iota=\pi/2$，剩下的角积分正比于 $\int[3(1-w^2)+(1-w^2)^2/2]\,\mathrm dw=68/15$，可检查归一化。$E$ 在绝热旋近中可同轨道束缚能损失相比，因此长时积累能使非振荡波形达到与最低阶振荡波相同的应变量级，尽管每一周期中增加的记忆很小。

\subsection{一个可积的外部波源}

尾积分为何含有对数，可用二阶场方程里实际出现的一类项直接检验。暂置 $c=1$，在平直背景取一个有效源 $S(t,\bm y)=\hat y_{L}^{\,n}F(t-y)/y^2$，其中 $\hat y_L^{\,n}$ 表单位方向 $\hat{\bm y}$ 的 $l$ 阶 STF 积，$y=|\bm y|$。这类 $1/y^2$ 项可来自单极场的 $1/y$ 乘出射波导数的 $1/y$。令 $F$ 在过去有限时刻以前为零，以免另加初态项。所求出射解为
\begin{equation}
 \square^{-1}_{\rm ret}S(t,\bm x)
 =-\hat n_L\int_1^\infty Q_l(z)\,F(t-rz)\,\mathrm dz,\qquad
 Q_l(z)=\tfrac12\int_{-1}^1\frac{P_l(w)}{z-w}\,\mathrm dw,\quad z>1 .
 \label{eq:gw-a-legendre}
\end{equation}
$Q_l$ 为第二类 Legendre 函数。先在标量 $l=0$ 时检验。由推迟 Green 函数的 $-1/(4\pi R)$，以 $\bm x$ 为极轴，$R=|\bm x-\bm y|$，有 $R\,\mathrm dR=-ry\,\mathrm d(\hat{\bm x}\cdot\hat{\bm y})$。将 $v=y+R$ 固定，三角条件 $|r-y|\le R\le r+y$ 给 $(v-r)/2\le y\le(v+r)/2$、$v\ge r$，故
\[
 \square^{-1}_{\rm ret}S
 =-\frac1{2r}\int_r^\infty F(t-v)\,\mathrm dv
       \int_{(v-r)/2}^{(v+r)/2}\frac{\mathrm dy}{y}
 =-\frac1r\int_r^\infty F(t-v)\,
              \frac12\ln\!\left(\frac{v+r}{v-r}\right)\mathrm dv .
\]
这就是式 \eqref{eq:gw-a-legendre} 的 $Q_0$。积分中的同一个到达时刻由一整层椭球形散射位置 $y+R=v$ 贡献。最短路径 $v=r$ 附近的几何核为对数，但端点奇性可积。

若 $l>0$，先把方位角积分做掉，$n'_L$ 的 STF 投影在 $\hat{\bm x}$ 的轴上只余 $\hat n_L P_l(w)$，$w=\hat{\bm y}\cdot\hat{\bm x}$。由 $R=v-y$ 代回余弦定律，得 $w=v/r-(v^2-r^2)/(2ry)$，从而 $\mathrm dy/y=\mathrm dw/(v/r-w)$。这一步用内层 $y$ 积分便直接得到 $2Q_l(v/r)$。在近光锥处令 $z=1+\tau/r$，$Q_l(1+\tau/r)=\tfrac12\ln(2r/\tau)-H_l+O[(\tau/r)\ln(r/\tau)]$，调和数 $H_l=\sum_{j=1}^l1/j$。由 $Q_0$ 的显式对数和递推 $(l+1)Q_{l+1}=(2l+1)zQ_l-lQ_{l-1}$ 可依次验证此常数（$Q_1=zQ_0-1$）。因此固定 $u=t-r$ 而 $r$ 很大时，
\[
 \square^{-1}_{\rm ret}S
  =-\frac{\hat n_L}{r}\int_0^\infty\! F(u-\tau)
            [\tfrac12\ln(2r/\tau)-H_l]\,\mathrm d\tau+\cdots .
\]
其中 $\ln r$ 可同辐射时刻的重定义合并，$\ln\tau$ 却必须在积分内对所有历史加权。完整引力计算还须满足谐和条件并作张量投影，因而常数成为式 \eqref{eq:gw-a-tails} 中按多极而异的 $\kappa_l$。这个可积例子给出了对数和历史依赖同时出现的几何原因。

同一变元也说明 DIRE 在求外部积分时为何能既保留传播又利用近区 PN 结果。令 $u'=t-R-y$ 为从源心发出、到外部一点后再到观者的等效推迟时刻。固定 $u'$ 时用上面的 $v=t-u'$ 及方向 $w$ 表示半径 $y=(v^2-r^2)/[2(v-rw)]$。原先遍及无穷远的半径积分变成过去历史和有界立体角积分；对给定分界球 $\mathcal R$，允许的 $w$ 还需满足 $y>\mathcal R$。原点邻域在这一部分已被排除。若换元之后又把 $F(u')$ 收缩为 $F(t)$ 的有限时间 Taylor 级数，就会把最要保留的长时间传播项丢掉。这是外部场即使在源很慢时也必须保持推迟解的一个显式理由。

\subsection{在分界球两边消去任意半径}

在上例中可保留一个实际的内边界 $a$，使 $S=0$ 于 $y<a$，而在 $y>a$ 为 $F(t-y)/y^2$。求距原点为 $r$ 处的标量解，令 $l=0$，且选 $r<\mathcal R$，这里 $\mathcal R>a$ 是任意分界球。在固定 $v=y+R$ 的椭球上，上下半径为 $y_\pm=(v\pm r)/2$，还需 $y>a$。当该椭球和分界球相交，亦即 $2\mathcal R-r<v<2\mathcal R+r$，内外两积分的这一层分别给
\[
 -\frac{F(t-v)}{2r}\,\ln\frac{\mathcal R}{\max(a,y_-)}\,\mathrm dv,\qquad
 -\frac{F(t-v)}{2r}\,\ln\frac{y_+}{\mathcal R}\,\mathrm dv .
\]
若 $y_+\le a$ 整层为零。若 $y_+\le\mathcal R$ 且 $y_+>a$，整层属于内区；若 $y_-\ge\mathcal R$ 则整层属于外区。两项相加，$\ln\mathcal R$ 在每个混合层分别消去；换到极角变量 $w$ 时，同一分割成为 $w_{\mathcal R}=[v-(v^2-r^2)/(2\mathcal R)]/r$。这给出了外部零锥角积分的真实下限，而非一律从 $-1$ 到 $1$。分界条件随 $v$ 移动，若错用恒定角域，末尾即使想手动丢掉 $\mathcal R$ 也修不好时间核。

在中心观者 $r=0$ 的极限 $R=y$，可用一个单色分量检验两边各自近似时，半径究竟如何相消。暂恢复 $c$，取 $F(t)=\exp(-\mathrm i\omega t)$，在远过去绝热开通。精确解在中心等于 $-\exp(-\mathrm i\omega t)\int_a^\infty \exp(2\mathrm iky)\mathrm dy/y$，$k=\omega/c$。把正核的内区与外区积分分别记为 $I_<,I_>$。令 $C(z)=\mathrm{Ci}(z)+\mathrm i\,\mathrm{Si}(z)$，则 $C'(z)=\exp(\mathrm iz)/z$，无穷远的绝热边界值为 $\mathrm i\pi/2$。于是
\[
 I_<=C(2k\mathcal R)-C(2ka),\qquad
 I_>=\mathrm i\pi/2-C(2k\mathcal R).
\]
若 $k\mathcal R\ll1$，内区可先按 $k$ 展开，而非紧致的外区必须先作推迟积分再展开。利用 $C(z)=\gamma_E+\ln z+\mathrm iz-z^2/4+O(z^3)$，得
\begin{align*}
 I_< &=\ln(\mathcal R/a)+2\mathrm ik(\mathcal R-a)
                      -k^2(\mathcal R^2-a^2)+O((k\mathcal R)^3),\\
 I_> &=\mathrm i\pi/2-\gamma_E-\ln(2k\mathcal R)
                      -2\mathrm ik\mathcal R+k^2\mathcal R^2+O((k\mathcal R)^3).
\end{align*}
任意边界的对数、一次幂、二次幂依次抵消，留 $I_<+I_>=\mathrm i\pi/2-\gamma_E-\ln(2ka)-2\mathrm ika+k^2a^2+\cdots$。$\ln(2ka)$ 依赖实际内边界和频率，绝热边界值的虚部则记载推迟方向。若在外部积分内把 $\exp(2\mathrm iky)$ 截成多项式，每项的上限都会发散，既无法求出虚部，也不能得到正确的频率对数。这个标量例子中 $a$ 是人为规定的真实源孔；完整引力问题在物质内部还要解场，所得 $a$ 的作用由真源矩与传播共同代替。

\subsection{反作用的力矩与降阶}

在 \S\ref{sec:pn-matching} 中已用反作用力所做的功核对过能量损失。对同一力还须核对角动量。让 $I_{ij}$ 表 Newton 质量 STF 矩，质点 $A$ 所受的反作用加速度取第三章的规范。总力矩为
\[
 N_i=-\frac{2G}{5c^5}\sum_A m_A\epsilon_{iab}x_A^a x_A^j I^{(5)}_{bj}
    =-\frac{2G}{5c^5}\epsilon_{iab}I_{aj} I^{(5)}_{bj}.
\]
质量二次矩的迹经反对称缩并为零。在时间积分中两次分部积分得到
\[
 \epsilon_{iab}I_{aj} I^{(5)}_{bj}=
 \epsilon_{iab}\ddot I_{aj}I^{(3)}_{bj}
 +\frac{\mathrm d}{\mathrm dt}\left[
 \epsilon_{iab}(I_{aj}I^{(4)}_{bj}-\dot I_{aj}I^{(3)}_{bj})\right].
\]
正因为反作用势随 $x^ix^j$ 改变，各个质点受的力并不同向，才能给出这一个净力矩。对于周期近似，平均后剩的第一项乘以前因子，恰好是远处角动量通量的负值。对于散射，若初态或末态有近区储存角动量的变化，第二项不能只凭它是全导数而略去。

运动方程中出现第五阶导数，若按一个精确的高阶常微分方程求解，就会要求多给三阶乃至更高阶的初始数据。实际展开的是低阶方程的微扰解。设 $\ddot{\bm r}=\bm F_0(\bm r)+c^{-5}\bm F_5(\bm r,\dot{\bm r},\ldots)$，右边 $\bm F_5$ 内一切加速度及其导数先用 $\bm F_0$ 消去，误差才在所求的阶以外，称为\textbf{降阶}。举圆轨道为例，以 $\bm n$ 表示两星连线单位矢量，$\bm\lambda=\dot{\bm n}/\Omega$ 为切向单位矢。令 $r$ 在一周期内为常数，四极矩 $I_{ij}=\mu r^2 n_{\langle i}n_{j\rangle}$。依次微分给
\begin{align*}
 I_{ij}^{(1)} &=2\mu r^2\Omega\, n_{(i}\lambda_{j)},&
 I_{ij}^{(2)} &=2\mu r^2\Omega^2(\lambda_i\lambda_j-n_i n_j),\\
 I_{ij}^{(5)} &=(2\Omega)^4 I_{ij}^{(1)},&
 n_j I_{ij}^{(5)}&=16\mu r^2\Omega^5\lambda_i .
\end{align*}
于是相对加速度的耗散部分为 $-32G\mu r^3\Omega^5\bm\lambda/(5c^5)=-\zeta\bm v$，其中 $\zeta=32G^3M^3\eta/(5c^5r^4)$，$\eta=\mu/M$。这一操作把快时间 $\Omega^{-1}$ 上的运动用于求慢时间 $\zeta^{-1}$ 上的漂移；在最后几个轨道若两时间相近，就不能依旧令 $r$ 为常数。

为何切向减速仍使轨道角速度增大？把矢量运动方程写成极坐标，切向式为 $r\dot\Omega+2\dot r\Omega=-\zeta r\Omega$，故 $\mathrm d(r^2\Omega)/\mathrm dt=-\zeta r^2\Omega$。而在准圆序列上 $\Omega\propto r^{-3/2}$，代入得 $\dot r/r=-2\zeta$、$\dot\Omega/\Omega=3\zeta$。能量和角动量一起减小，系统落到角速度更高的圆轨道上。由于 $\dot r=O(c^{-5})$，$\dot r^2$ 相对于径向保守式要到 $O(c^{-10})$ 才出现。至 3.5PN 可把径向关系看成缓慢参数化的一族保守圆轨道，同时在切向方程中保留耗散；不能在求切向力之前就把 $\dot\Omega,\dot r$ 也都置零。

\subsection{次阶反作用势与体内尺度}

领头的 Burke--Thorne 势只用远处的质量四极。内区若要求到相对这项的下一阶，不仅 $I_{ij}$ 应取到 1PN，新的质量八极、流四极也会出现。用近区正 Newton 势的记号，一种反作用规范中标量、矢量势的相关项为
\begin{align*}
 V_{\rm R}={}&-\frac{G}{5c^5}x_i x_j I_{ij}^{(5)}
    +\frac{G}{c^7}
       \left[\frac{x_i x_j x_k}{189} I_{ijk}^{(7)}
          -\frac{r^2x_i x_j}{70} I_{ij}^{(7)}\right]+\cdots,\\
 V^i_{\rm R}={}&\frac{G}{c^5}
       \left[\frac{\hat x_{ijk}}{21} I_{jk}^{(6)}
                    -\frac{4}{45}\epsilon_{ijk}x_j x_a J_{ka}^{(5)}\right]+\cdots .
\end{align*}
此处 $\bm x$ 是近区场点，$r=|\bm x|$，矩均取在 $t$。标量、矢量在度规中的领头因子分别含 $1/c^2$、$1/c^3$；矢量作用于慢运动还多一次速度或时间微分，故后者的首项也贡献于 3.5PN 反作用。两式依赖坐标选择，不能只把矢量势一项与观测通量比较。一致的 3.5PN 算法还须考虑 Newton 反作用势对 1PN 机械能、角动量的功以及 PN 轨道对矩的改变。

可在标量项中看见前述正则齐次波的痕迹。以 STF 的 $x_i x_j x_k$ 与 $I_{ijk}^{(7)}$ 缩并，近区称重为长三次方乘质量八极的七阶导。与 $I_{ij}^{(5)}x_i x_j$ 相比，若场点就在轨道附近，多出两个时间导数和两次轨道尺度，配上 $c^{-2}$，给 $(v/c)^2$。$r^2x_ix_j I_{ij}^{(7)}$ 来自同一角阶四极齐次波的再一项径向展开。两者虽用同一 $c^{-7}$ 因子，却有不同的角分布，不能用修正一个四极系数代替全部 3.5PN 反作用。

到再高半阶，质量单极使一部分出射波散射并重入近区。相应的标量势中有
\[
 -\frac{4G^2M}{5c^8} x_i x_j
       \int_0^\infty I_{ij}^{(7)}(t-\tau)
              \ln(\tau/\tau_b)\,\mathrm d\tau .
\]
时间尺度 $\tau_b$ 随谐和规范及与局域项的分配而变。与 $x_i x_j I_{ij}^{(5)}$ 比，两个额外导数经一次积分保留一幂频率，再乘 $GM/c^3$，故相对最低的反作用是 1.5PN，而相对 Newton 运动是 4PN。由此可区分常混淆的两个阶次：远处通量的 1.5PN 尾，反馈到轨道的耗散项位于 4PN。完整的 4PN 保守动力学还可用时间对称的非局域作用量表示，不能把这里只列出的推迟反作用势当作整个 4PN 近区度规。

双星中为何能一直到很高的 PN 阶才显式需要星的结构？\S\ref{sec:pn-finite} 与第三章的强自引力讨论说明质量参数本已含结合能。对无自旋、球对称平衡星，可以用匹配进一步定量判断何时需要新参数。在星外缓冲区 $R_*\ll s\ll r$，外加的潮汐势有 $U_{\rm ext}=U_0+a_i y_i+\mathcal E_{ij}y_iy_j/2+\cdots$，$s=|\bm y|$。常数与均匀重力可由时钟、共动中心选择收进度规；二阶项是不能去掉的潮汐。该角阶的真空 Laplace 方程还允许 $Q_{ij}n_i n_j/s^3$ 的响应。内外边界条件决定两解之比，写为 $Q_{ij}=-\lambda\mathcal E_{ij}$，符号按电型曲率定义。无量纲外场乘 $s^2$，响应乘 $s^{-3}$，故 $\lambda\propto R_*^5/G$；无量纲的比例常数才记录物态。

外伴星的潮汐 $\mathcal E\sim Gm'/r^3$，由此诱导的 $Q$ 回到轨道能量时还有一次外部潮汐，$\delta E_{\rm tide}\sim\lambda(Gm'/r^3)^2$。相对 $Gmm'/r$ 约为 $(m'/m)(R_*/r)^5$ 乘 Love 数。若紧致度 $Gm/(R_*c^2)$ 在缩小轨道速度的序列中保持有限，此比随 $x^5$，即称 5PN；对蓬松的普通恒星却不能只按 $x$ 推断其可忽略。固有自旋四极并不由外场诱导，强度可在较低阶出现。若模频与轨道驱动相近，静态比例 $Q=-\lambda\mathcal E$ 也应替换成含相位的频率响应。这些是内部结构能够重新影响源矩与运动的三个不同条件。

\subsection{用频率消去坐标半径}

以下取无自旋、非吸收的点质量准圆双星。令 $\gamma=GM/(c^2r)$，$x=(GM\Omega/c^3)^{2/3}$。以谐和坐标写广义 Kepler 律为 $\Omega^2=(GM/r^3)K(\gamma)$。为明确哪些对数只是坐标选择，可列出其到 3PN 的系数：
\begin{align*}
 K(\gamma)&=1+k_1\gamma+k_2\gamma^2+k_3\gamma^3+O(\gamma^4),\\
 k_1&=-3+\eta,\qquad k_2=6+\tfrac{41}{4}\eta+\eta^2,\\
 k_3&=-10+\left(22\ln(r/r'_0)-\tfrac{75707}{840}
                  +\tfrac{41\pi^2}{64}\right)\eta
                    +\tfrac{19}{2}\eta^2+\eta^3 .
\end{align*}
这里的 $r'_0$ 来自近区的谐和坐标，与辐射时刻的 $b$、源多极的正则化尺度 $r_0$ 应分开。取三次根得 $x=\gamma K^{1/3}$。设 $\gamma=x(1+c_1x+c_2x^2+c_3x^3)$，那么最前两级有 $c_1=-k_1/3=1-\eta/3$、$c_2=k_1^2/3-k_2/3=1-65\eta/12$。$k_3$ 中的 $22\eta\ln(r/r'_0)$ 在 $c_3$ 中以 $-22\eta\ln(r/r'_0)/3$ 出现；代换该对数中的 $r$ 时只需保留 $r=GM/(c^2x)$ 的领头值，多出的相对 $O(x)$ 属下一阶。

谐和半径下的束缚能有 $E=-\mu c^2\gamma[1+e'_1\gamma+e'_2\gamma^2+e'_3\gamma^3]/2$。其中 $e'_1=-7/4+\eta/4$，$e'_2=-7/8+49\eta/8+\eta^2/8$，$e'_3$ 的对数项为 $+22\eta\ln(r/r'_0)/3$。代入前式，$E(x)$ 中 $x^3$ 括号系数里两个对数恰相消。低两级的系数已在 \S\ref{sec:pn-wave-energy} 给出，新的 3PN 能量系数为
\[
 e_3=-\frac{675}{64}
    +\left(\frac{34445}{576}-\frac{205\pi^2}{96}\right)\eta
    -\frac{155}{96}\eta^2-\frac{35}{5184}\eta^3,
 \qquad E=-\tfrac12\mu c^2x(1+e_1x+e_2x^2+e_3x^3).
\]
这些系数依展开的频率所用的总质量而定。若误把 $M$ 换成减缩质量，连测试粒子极限的自变量也会失去意义。

此极限有一个不依赖 PN 展开的核对。Schwarzschild 圆测地线上的单位静质量能量为 $(1-2x)/\sqrt{1-3x}$，其中 $x=GM/(c^2r_{\rm areal})$。取 $\eta\to0$ 后应令 $E/(\mu c^2)=(1-2x)(1-3x)^{-1/2}-1$。二项式展开的各阶依次为 $-x/2+3x^2/8+27x^3/16+675x^4/128$，最后一项等于 $-e_3(0)x^4/2$。精确函数在 $x=1/3$ 处有光子圆轨道的奇点，稳定圆轨道序列却已在 $x=1/6$ 结束。仅因有限阶多项式在 $x>1/6$ 仍为实数，并不能证明那里还有绝热稳定旋近。

\subsection{在一个完整 PN 阶内合并各个辐射矩}

把 $E(x)$ 与总通量直接代进平衡方程非常方便，却看不出一个总通量系数需要多少种矩。以下以相对 1PN 为例，把每一项真正求和。暂在谐和坐标下写 $\bm r=r\bm n$、$\bm v=r\Omega\bm\lambda$。把匹配后的质量、动量及应力密度代入式 \eqref{eq:gw-a-sourceI} 的 1PN 展开，对无自旋圆轨道，源质量四极矩为
\[
 I_{ij}=\mu\left[(1+a\gamma)r_{\langle i}r_{j\rangle}
                 +b\frac{r^2}{c^2}v_{\langle i}v_{j\rangle}\right]+O(c^{-4}),
 \quad a=-\frac1{42}-\frac{13}{14}\eta,\quad
 b=\frac{11}{21}-\frac{11}{7}\eta .
\]
在一般偏心轨道还有与 $\dot r\,r_{\langle i}v_{j\rangle}$ 等组合有关的项。此处圆轨道的 $a$ 已包含引力场能量对源矩的修正，不能只把各粒子静质量乘上 Lorentz 因子来重得。

恒有 $n_i n_j+\lambda_i\lambda_j=\delta_{ij}-l_il_j$，$\bm l$ 固定。因此纯振荡的 $\lambda_{\langle i}\lambda_{j\rangle}$ 和 $n_{\langle i}n_{j\rangle}$ 只差负号。又因该速度项本已 $O(c^{-2})$，其中可用 $v^2/c^2=\gamma+O(\gamma^2)$，故 $I_{ij}^{(3)}$ 的振幅相对改为 $1+(a-b)\gamma$。功率中这一因子须平方。在频率 $\Omega$ 固定时，量纲因子 $r^4\Omega^6$ 也有 1PN 修正，因 $r=(GM/\Omega^2)^{1/3}(1+k_1x/3)+O(x^2)$（余项相对于首项）。故仅由质量四极矩给出的修正为
\[
 \frac{\mathcal F_{I_2}}{\mathcal F_N}
 =1+\left[2(a-b)+\tfrac43 k_1\right]x
 =1+\left(-\frac{107}{21}+\frac{55}{21}\eta\right)x .
\]
若边计算边令 $r=(GM/\Omega^2)^{1/3}$，会漏掉 $4k_1/3$，其大小并不比真正的源矩修正小。

质量八极、流四极的波幅比主模小 $v/c$，与质量四极的干涉对整个球面无贡献，但其各自平方正好到 1PN 阶。由前文圆轨道的两时刻缩并及式 \eqref{eq:gw-a-multiflux}，有
\begin{align*}
 \frac{\mathcal F_{I_3}}{\mathcal F_N}
 &=\frac{5}{32\cdot189}\left(\frac{3}{20}+\frac{3^8}{4}\right)
      \Delta^2 x=\frac{1367}{1008}\Delta^2 x,\\
 \frac{\mathcal F_{J_2}}{\mathcal F_N}
 &=\frac{5}{32}\frac{16}{45}\frac12\Delta^2 x
   =\frac1{36}\Delta^2 x .
\end{align*}
相加得 $(155/112)\Delta^2 x$，再用 $\Delta^2=1-4\eta$，总 1PN 系数为
\[
 -\frac{107}{21}+\frac{155}{112}
       +\left(\frac{55}{21}-\frac{155}{28}\right)\eta
       =-\frac{1247}{336}-\frac{35}{12}\eta .
\]
这实际算出了 \S\ref{sec:pn-wave-energy} 中的一个系数。等质量系统的高奇宇称辐射在这一级全消，修正完全来自四极波的改变；质量比趋零却仍辐射高多极，两个极限可分别检查求和的质量因子。虽然高多极通量恒正，固定轨道频率时总 PN 修正为负，因为质量四极振幅本身及 $r(\Omega)$ 均改变了。

这一核算还给出任意阶的取舍办法。$I_l$ 相对于四极振幅约为 $(v/c)^{l-2}$，$J_l$ 还多一幂 $v/c$。所以求全天通量到相对 $n$ PN 阶，最低阶的 $I_{n+2}$ 和 $J_{n+1}$ 都要，质量四极须算到 $n$ PN，质量八极与流四极须算到 $n-1$ PN，以此类推（这里 $n$ 为非负整数）。到半整数阶另有尾和反作用等非线性效应。若要的是方向分辨的波形或带 $\bm n$ 权重的动量通量，不同 $l$ 间干涉不再全消，选取的最大 $l$ 也应按振幅的相对阶数重新计数。

\subsection{从 3.5PN 功率求波相位}

将功率以 Newton 值 $\mathcal F_N=32c^5\eta^2x^5/(5G)$ 归一，写 $\mathcal F/\mathcal F_N=1+\sum_{k=2}^7 f_k x^{k/2}+O(x^4)$。$f_2,f_3,f_4$ 已列在 \S\ref{sec:pn-wave-energy}。其余的系数为
\begin{align*}
 f_5={}&-\pi\left(\frac{8191}{672}+\frac{583}{24}\eta\right),\\
 f_6={}&\frac{6643739519}{69854400}+\frac{16\pi^2}{3}
           -\frac{1712\gamma_E}{105}-\frac{856}{105}\ln(16x)\\
      &+\left(-\frac{134543}{7776}+\frac{41\pi^2}{48}\right)\eta
           -\frac{94403}{3024}\eta^2-\frac{775}{324}\eta^3,\\
 f_7={}&\pi\left(-\frac{16285}{504}+\frac{214745}{1728}\eta
                                +\frac{193385}{3024}\eta^2\right).
\end{align*}
与 $K$ 中的 $\ln(r/r'_0)$ 不同，$\ln(16x)$ 只有可测频率和质量，不能用重选谐和半径消去。在同一 3PN 阶，将源矩修正、尾与尾的散射以及尾的绝对值平方都计入后，才剩这个无量纲对数。功率总式还有一个极限检验：除去 $\eta^2$ 再令 $\eta\to0$，应给测试粒子在 Schwarzschild 背景上以出射边界条件求得的通量展开。只核对 $E(x)$ 并不足以验证耗散系数。

这些系数如何进入可用的相位？记 $v=\sqrt{x}$（无量纲），$f=c^3v^3/(\pi GM)$ 为主模的周频。用
\[
 Q(v)=\frac{1+2e_1v^2+3e_2v^4+4e_3v^6}
                {1+f_2v^2+f_3v^3+\cdots+f_7v^7}
      =\sum_{k=0}^7 q_k(v)v^k+O(v^8)
\]
写 $\mathrm dt/\mathrm dv=5GM Q(v)/(32\eta c^3 v^9)$、$\mathrm d\phi/\mathrm dv=5Q(v)/(32\eta v^6)$。分子来自 $\mathrm dE/\mathrm dx$ 的乘幂，分母来自功率。对分母作级数除法足以避免在长分式中漏掉混合项。若记分子系数为 $d_k$（$d_0=1,d_2=2e_1,d_4=3e_2,d_6=4e_3$，其余零），那么 $q_0=1$ 且 $q_k=d_k-\sum_{j=2}^k f_jq_{k-j}$。例如 $q_3=-f_3$，$q_5=-f_5-f_2q_3-f_3q_2$，其中两项低阶交叉不能略去。$q_6$ 才第一次出现 $\ln v$。

令正频振幅为 $|\widetilde h|\exp(-\mathrm i\Psi)$，则与第三章相同，$\Psi=2\pi f t(v)-2\phi(v)-\pi/4$。若 $q_k$ 是常数且 $k\ne 5,8$，把前面的两条微分式分别积分，到括号内第 $k$ 阶相位系数的净贡献为 $40q_k/[(k-8)(k-5)]$。因为两项都含同一 $v^{k-5}$，若只对 $\phi(v)$ 积分，得到的系数会错。$k=5$ 时 $\mathrm d\phi/\mathrm dv$ 含 $1/v$，因此出现 $\ln(v/v_*)$；改变 $v_*$ 只移动任意的合并初相。$k=6$ 时的对数要按 $\int v^p\ln v\,\mathrm dv=v^{p+1}[\ln v-1/(p+1)]/(p+1)$ 处理，常数部分也被改动。由此可在 \S\ref{sec:pn-wave-energy} 的 2PN 频域相位旁直接续写
\begin{align*}
 \Psi={}&2\pi f t_c-\Phi_c-\tfrac\pi4
  +\frac{3}{128\eta v^5}\left[1+\sum_{k=2}^7\psi_k v^k\right],\\
 \psi_5={}&\pi\left(\frac{38645}{756}-\frac{65}{9}\eta\right)
                       [1+3\ln(v/v_*)],\\
 \psi_6={}&\frac{11583231236531}{4694215680}-\frac{640\pi^2}{3}
 -\frac{6848}{21}[\gamma_E+\ln(4v)]\\
 &+\left(-\frac{15737765635}{3048192}+\frac{2255\pi^2}{12}\right)\eta
 +\frac{76055}{1728}\eta^2-\frac{127825}{1296}\eta^3,\\
 \psi_7={}&\pi\left(\frac{77096675}{254016}+\frac{378515}{1512}\eta
                  -\frac{74045}{756}\eta^2\right).
\end{align*}
$\psi_5$ 中的非对数常数同样可归入 $\Phi_c$。这个表示已在指定阶数重新展开。另一种积分方法直接以截断的 $E$ 和 $\mathcal F$ 作商进行数值积分，两种方法到所列的 PN 阶完全相合，商的未展开部分在并合前却可给出不同的高阶外推。因此比较波形时须说明是对哪个商展开以及终止频率如何选取。

作频域振幅的驻相积分时也有一项可检验的相消。设 $h(t)=a(t)\cos\Phi(t)$，$a$ 与 $\dot f$ 在相位驻点附近的 $\Delta t\sim\dot f^{-1/2}$ 内缓变。则 $|\widetilde h|\simeq a(t_f)/(2\sqrt{\dot f(t_f)})$。对准圆主模 $a\propto v^2$、$\dot f\propto v^{11}$，故得 $v^{-7/2}$，与前面能谱所得一致。对第 $m$ 谐波，驻点却取在 $2\pi f=m\Omega(t_m)$，曲率由 $m\dot\Omega$ 给出；多模之间的相位差不能靠对同一驻点处的极化系数求和得到。在频带最末端，若驻点距终止边界只剩一个 $\Delta t$，完整的 Gaussian 积分变成有边界的 Fresnel 积分，$-\pi/4$ 的常数相位和 $1/\sqrt{\dot f}$ 的幅度均不足以给出精确形状。

\subsection{幅度的半整数阶与累积周期}

对非进动准圆双星，将振幅写成 $h_A=2G\mu x[H_A^{(0)}+\sqrt x H_A^{(1/2)}+x H_A^{(1)}+\cdots]/(c^2D)$。观测基的一轴取为两平面的交线，选相位使 $H_+^{(0)}=-(1+\cos^2\iota)\cos2\psi$、$H_\times^{(0)}=-2\cos\iota\sin2\psi$。取与此相容的轨道、质量差约定，前述质量八极与流四极在波面上的合成为
\begin{align*}
 H_+^{(1/2)}&=-\frac{\Delta\sin\iota}{8}
     [(5+\cos^2\iota)\cos\psi-9(1+\cos^2\iota)\cos3\psi],\\
 H_\times^{(1/2)}&=-\frac{3\Delta\sin\iota\cos\iota}{4}
         [\sin\psi-3\sin3\psi].
\end{align*}
这里 $\psi$ 是吸收尾相位的轨道角，到首个尾阶可写成 $\phi-(2GM\Omega/c^3)\ln(\Omega/\Omega_0)$。式中奇数谐波随交换两星变号，同时旋转轨道相位 $\pi$ 后，物理波形当然不变。正对轨道极点时奇数模也消失。若同时测到奇偶谱线，质量差和视线倾角便可由不同的组合进入参数估计；只有主模时两种信息更易混在总振幅和偏振比中。

相位为何常算到比振幅更高的阶？在固定的相对带宽内，未计入的幅度 $v^k$ 仍是 $O(v^k)$，每个周期的相位误差却会相干地累积。用上面的 $\mathrm d\phi/\mathrm dv$ 计算从 $v_1$ 到 $v_2$ 的主模周期数，
\[
 N=\frac{\phi(v_2)-\phi(v_1)}{\pi}
   =\frac{v_1^{-5}-v_2^{-5}}{32\pi\eta}
   +\frac{5}{32\pi\eta}\int_{v_1}^{v_2}
                        v^{-6}[Q(v)-1]\,\mathrm dv .
\]
当 $v_1\ll v_2$，第 $k$ 阶误差约按 $v_1^{k-5}$ 计，$k=5$ 还会积成对数。把振幅限于主四极而相位保留至较高阶称为限制型 PN 波形。它适于主要靠周期比对的某些搜索，但会丢掉上述奇数模以及 $m=4$ 等更高的幅度信息。例如在相同观测低频 $f_{\min}$，$m=3$ 分量对应的轨道频率只是 $m=2$ 的 $2/3$。在 Newton 近似下到同一末态的剩余时间按 $\Omega^{-8/3}$ 缩放，故两个谐波入带的时刻相差一个可计算的提前量；最低可测轨道频率并不由主模单独决定。

\subsection{计时轨道中的三个离心率}\label{sec:gw-b-timing}
第三章 \S\ref{sec:pn-peri} 已从二体运动求得每周进动。\hyperlink{subsection.3.13.1}{「双星脉冲星的计时」}又给出由发射钟到地面钟的时延。若要在一整圈而不仅在近心点上追踪脉冲相位，怎样把后 Newton 轨道放进这些时延？以下只求保守动力学到 1PN，取束缚而非圆形的轨道。利用守恒能量和角动量消去速度，径向与角向方程在常用的谐和坐标下可写成
\[
 \dot r^2=A+\frac{2B}{r}+\frac{C}{r^2}+\frac{D}{r^3},
 \qquad \dot\varphi=\frac{H}{r^2}+\frac{I}{r^3}.
\]
这里 $A,B,C,H$ 各含 Newton 和 1PN 项，$D,I$ 从 1PN 起。$A<0,C<0$，系数由同一组守恒量决定。写 $r=\bar r+s$，则 $1/r^3$ 的系数变为 $D-2Cs$。选 $s=D/(2C)$ 便使之消失，并将 $1/\bar r^2$ 的系数改成 $\bar C=C-2Bs$。这一步只在截断至 1PN 后成立。

置 $\bar a=-B/A$、$e_t^2=1-A\bar C/B^2$。将新的径向方程乘 $\bar r^2$ 再配方，得 $\bar r^2\dot{\bar r}^2=(-A)[\bar a^2 e_t^2-(\bar r-\bar a)^2]$。令 $\bar r=\bar a(1-e_t\cos u)$，沿去程选正平方根，随后绕全轨道连续延拓，便有 $\d t/\d u=\bar a(1-e_t\cos u)/\sqrt{-A}$。从近心点 $T_0$ 起积分，得到
\[
 n(t-T_0)=u-e_t\sin u,\qquad
 r=a_r(1-e_r\cos u),\qquad
 n=\frac{\sqrt{-A}}{\bar a},\quad a_r=\bar a+s,\quad e_r=\frac{\bar a}{a_r}e_t .
\]
控制钟的 $e_t$ 和控制半径的 $e_r$ 只在 Newton 极限相同。用一个经验离心率同时拟合钟与形状，相当于漏掉了一项周期时延。

角向方程可写为 $\dot\varphi=H/[r-I/(2H)]^2$。再令 $a_q=a_r-I/(2H)$、$e_q=\bar a e_t/a_q$，便有 $\d\varphi/\d u=[H/(na_q^2)](1-e_t\cos u)/(1-e_q\cos u)^2$。因 $e_t-e_q$ 只有 1PN 大小，最后的比值可写成 $1/(1-e_\varphi\cos u)$，其中 $e_\varphi=2e_q-e_t$，误差为 2PN。定义连续展开的真近点角 $v_e(u)$，其半角满足 $\tan(v_e/2)=\sqrt{(1+e)/(1-e)}\tan(u/2)$。由 $\d v_e/\d u=\sqrt{1-e^2}/(1-e\cos u)$，有
\[
 \varphi-\varphi_0=(1+k)v_{e_\varphi}(u),\qquad
 1+k=\frac{H}{n a_q^2\sqrt{1-e_\varphi^2}} .
\]
因 $u$ 增 $2\pi$ 时 $v_e$ 也增 $2\pi$，径向一周后方位角多出 $2\pi k$。代入广义相对论系数应还原第三章的 $k=3GM/[a(1-e^2)c^2]$；在已是 1PN 的修正项里可取 Newton 的 $a,e$。$e_r,e_t,e_\varphi$ 是选定坐标与参数化后的三个系数，不能把三者之差都当作独立的引力检验。真正可辨的是它们在钟读数上形成的不同谐波。

还有一种时移即使沿用 Newton 椭圆轨道也会出现。随星运动的观者看见的视线随轨道速度横分量改变，波束要稍早或稍晚转向这个新方向才能被接收。设 $\hat{\bm s}$ 为自转轴、$\bm n$ 为脉冲星指向地球的视线，$\zeta$ 为两者夹角。自转相位改变 $\delta\phi$ 使波束沿 $\hat{\bm s}\times\bm n$ 移动。将一阶像差 $\delta\bm n$ 投至此方向，纵向相位变化为 $\delta\phi=(\hat{\bm s}\times\bm n)\cdot\delta\bm n/\sin^2\zeta$。故时移正比于 $\bm v_{\rm orb}/(c\Omega_{\rm spin})$ 的一固定投影，符号由自转和视线的约定决定。

以近心点为轨道平面的一根轴，比角动量记为 $j$。速度分量为 $(GM/j)(-\sin v,e+\cos v)$。转过近心点角 $\omega$，投到同一波束方向上，像差时延必有
\[
 \Delta_A=A_s[\sin(\omega+v)+e\sin\omega]
          +B_s[\cos(\omega+v)+e\cos\omega] .
\]
$A_s,B_s$ 有时间量纲，随星的自转方位改变。在圆轨道极限，这只在原几何时延上加了同频的正弦与余弦，可吸收到投影半长轴和参考相位。长期自旋进动或脉冲轮廓随轨道相位改变，才可能另给几何约束。横于自转扫掠的像差还可改变观者切过射电束的位置，主要改变轮廓；用固定轮廓拟合所有脉冲时，这也会混入到达时刻。

\subsection{把观测站的钟换到太阳系质心}
实际计时的两端各有一座钟。前面的轨道方程联系源处的发射事件，同一个脉冲到达不同地面站的时刻还需换至共同标准。取太阳系质心到观测站的位移为 $\bm r$，脉冲星位于 $D\bm n$。先略引力和介质，在相同发射事件下以到达时间之差直接算得
\[
 t_{\rm obs}-t_{\rm bar}
 =\frac{|D\bm n-\bm r|-D}{c}
 =-\frac{\bm n\cdot\bm r}{c}
   +\frac{r^2-(\bm n\cdot\bm r)^2}{2cD}
   +O\!\left(\frac{r^3}{cD^2}\right).
\]
第一项中的年周振幅至多约一日地距离的光行时，即约 $500\,$s。方向若误了 $\delta\bm n$，折算的质心时刻便有 $\bm r\cdot\delta\bm n/c$ 的周年残差。第二项是波前曲率的效应。令地球公转为半径 $R$ 的圆，源的黄纬为 $\beta$，则 $\bm n\cdot\bm r=R\cos\beta\cos(\Omega_\oplus^{\rm orb}t-\lambda)$。减去常量，余下半年的项为 $-R^2\cos^2\beta\cos[2(\Omega_\oplus^{\rm orb}t-\lambda)]/(4cD)$。在黄极方向，圆轨道模型中这项变为不可测的恒定距离增量。故计时视差同时依赖距离和天空位置，所得距离还可用来估算第三章讨论的横向运动运动学修正。

地面钟给的是固有时。取地球中心为试验钟，只看太阳弱场的随年变化部分，在质心坐标时 $t$ 下有 $\d\tau/\d t=1-[GM_\odot/r+v_\oplus^2/2]/c^2$。代入活力公式 $v_\oplus^2=GM_\odot(2/r-1/a)$，以偏近点角 $u$ 写 $\d t=(1-e\cos u)\d u/n$。时间平均 $\langle1/r\rangle=1/a$，故去掉恒定速率后的质心时差为
\[
 \Delta_{\rm E}^{\rm per}
 =\frac{2GM_\odot}{c^2}\int(1/r-1/a)\d t
 =\frac{2GM_\odot e}{a n c^2}\sin u .
\]
取 $GM_\odot/(ac^2)\simeq10^{-8}$、$e\simeq0.017$、$n=2\pi/{\rm yr}$，振幅约 $1.7\,\mathrm{ms}$。这项来自接收端绕日运动，与脉冲星在其双星内产生的 Einstein 时延有相同的动力学来源，但频率和轨道参数不同。太阳系的传播时延则由 \S\ref{sec:pn-light} 的零线积分按实际太阳及行星位置计算。将地面钟时刻换到质心并外推到无穷高射电频率之后，才与双星自身的计时模型比较。

色散外推还可从一束宽带脉冲本身检查。冷的无磁电子等离子体中，平面电磁波满足 $\omega^2=\omega_p^2+c^2k^2$，其中在 SI 制下 $\omega_p^2=n_e e^2/(\epsilon_0m_e)$。由 $v_g=\partial\omega/\partial k$ 展开单程传播时间，在射电频率 $\nu=\omega/(2\pi)$ 有
\[
 \Delta t_{\rm pl}=\frac{e^2}{8\pi^2\epsilon_0m_e c\,\nu^2}
       \int n_e\,\d l+
       \frac{3}{8c(2\pi\nu)^4}\int\omega_p^4\,\d l+\cdots .
\]
最低阶只量出电子柱密度 $\mathrm{DM}=\int n_e\d l$。下一项却依赖 $\int n_e^2\d l$，穿过稠密局部云团时二者并非按相同比例增加。通常射电计时中 $\omega_p^2/\omega^2$ 足够小，可以只对 $\nu^{-2}$ 作线性拟合。即使最低阶完全适用，一个宽度 $\delta\nu$ 的频率通道内也会有约 $2|\Delta t_{\rm pl}|\delta\nu/\nu$ 的展宽。通道中心的时移能对齐不同通道，却不能恢复在通道内平均时丢掉的脉冲相位；精确测时可在平均强度之前先对电场解除色散。

还应检查源钟是否真正数准了所有转数。若在一无记录区间 $\Delta T$ 内自转频率偏差为 $\delta\nu_s$、频率导数偏差为 $\delta\dot\nu_s$，累积圈数的不确定量级为 $|\Delta T\,\delta\nu_s+\Delta T^2\delta\dot\nu_s/2|$。一旦可达半圈，前后段各自都可有极小的拟合残差，连接两段的脉冲编号却有不同的整数解。这类离散歧义须逐个尝试或借助中间观测消除，在同一整数解附近缩小连续参数的协方差并不能排除它。它与后面用多颗星同时寻找共同红色计时残差，分属源钟与传播的两层推断。

\section{含噪声的应变测量}\label{sec:gw-b-spectrum}
\subsection{有限观测的功率谱}
设仪器已用一已知位移标定，输出为实序列 $d(t)=s(t)+n(t)$，其中 $s$ 已含仪器的方向响应。先截取长为 $T$ 的一段，定义 $\tilde n_T(f)=\int_{-T/2}^{T/2}n(t)\e^{-2\pi\i ft}\d t$。同取两个不同频率，乘积的期望会如何依赖 $T$？设噪声均值为零且二阶平稳，即 $\langle n(t)n(t')\rangle=C_n(t-t')$。在两个时间积分中先固定时差 $\tau$，则可供两点同时落入窗内的时长为 $T-|\tau|$。特别地，
\[
 \frac{1}{T}\langle|\tilde n_T(f)|^2\rangle
 =\int_{-T}^{T}\left(1-\frac{|\tau|}{T}\right) C_n(\tau)\e^{-2\pi\i f\tau}\d\tau .
\]
若相关时间远小于 $T$，右端趋于双边功率谱。以下对 $f>0$ 用其两倍，即单边谱 $S_n(f)$，于是 $C_n(0)=\int_0^\infty S_n(f)\d f$。仪器曲线上以 $\mathrm{Hz}^{-1/2}$ 为单位的量是 $\sqrt{S_n}$。白噪声在宽 $\Delta f$ 的带内有均方振幅 $S_n\Delta f$；给一条谱线报振幅时须区别峰高、积分功率与频率分辨率。

窗的 Fourier 变换为 $\delta_T(f)=T\operatorname{sinc}(\pi fT)$，其中 $\operatorname{sinc}x=\sin x/x$。在谱平滑的一小段内，$\langle\tilde n_T^*(f)\tilde n_T(f')\rangle\simeq\tfrac12 S_n(f)\delta_T(f'-f)$。矩形窗内相隔 $1/T$ 整数倍的 Fourier 基函数正交。有限观测时长的有色噪声谱系数，只在上述局部平谱近似下不相关。若按间隔 $\Delta t$ 取样，只须保留 $0<f<1/(2\Delta t)$ 的独立正频率。在取样前必须除去更高频率，否则 $f$ 与 $f+k/\Delta t$ 在样点处有完全相同的相位。减小 Fourier 变换的频率步长可以使一条峰的位置读得更细，分辨两条近线的能力则仍取决于 $T$。

怎样从一次实验估计 $S_n$？对 Gaussian 噪声，某一非端点频格上 $\tilde n_T$ 的实、虚部为独立的零均值正态量，方差各约 $TS_n/4$。故周期图 $I=2|\tilde n_T|^2/T$ 满足 $2I/S_n\sim\chi^2_2$。于是 $\langle I\rangle=S_n$，但标准差也等于 $S_n$。增长单段时长使频格变窄，同一格上的相对误差仍为一。将 $K$ 段独立且谱相同的记录取平均，$2K\bar I/S_n\sim\chi^2_{2K}$，相对误差降为 $1/\sqrt K$。这说明为什么要在时长、谱的局部光滑性及噪声漂移之间选择估谱区间。

矩形窗遇到不恰在频格中心的强窄线会怎样？单频复指数的窗内变换是 $\delta_T(f-f_0)$，故旁瓣按距离的倒数衰减。采样一条很弱的近邻谱线时，强线的旁瓣可能比它高。改乘平滑窗 $w(t)$ 后须用 $T_w=\int w^2\d t$ 代替周期图分母 $T$。反之，对一个已知频率的确定信号，峰值振幅正比于 $\int w\d t$。两种标定之比给等效噪声带宽 $B_w=\int w^2\d t/(\int w\d t)^2$。例如在窗内取 $w=(1+\cos(2\pi t/T))/2$，所得 $B_w=3/(2T)$。抑制旁瓣的代价是主峰加宽。邻窗重叠所给的谱估计相互相关，不能把段数直接当独立平均次数。

测到的若是光电流 $j$，写 $\tilde j=H(f)\tilde s+\tilde j_n$。把同一数据换算为应变给 $S_n=S_j/|H|^2$。在 $H$ 的零点附近相除会显著放大噪声。若直接在电流空间拟合预测 $H\tilde s$，权重取 $1/S_j$，所得内积相同。白化则是除以噪声振幅谱，使各频率的噪声方差相同。校准需要传递函数的相位；白化对噪声功率只用到模。用受信号本身支配的一小段粗估噪声也会压低显著性，实际应在可比较的无信号时段估谱并检验其稳定性。

\subsection{给定一个模板会有多少假警报？}\label{sec:gw-b-match}
第三章由似然给出匹配滤波。依上面的单边谱约定，两实波形的内积是 $(a|b)=4\operatorname{Re}\int_0^\infty \tilde a^*\tilde b\,\d f/S_n$。固定一个单位模板 $u$，则 $X=(d|u)$ 在纯噪声中有方差一，在有 $s=A u$ 时均值移到 $A$。因而阈值 $x_*$ 给单次虚警概率 $\tfrac12\operatorname{erfc}(x_*/\sqrt2)$，检出率为 $\tfrac12\operatorname{erfc}[(x_*-A)/\sqrt2]$。这两个数各自的条件不同。若要由一次超过阈值的记录求它来自真实源的概率，还要知道真实源与噪声事件出现的先验比。

相位未知时，取正交且单位归一的两个模板 $u,v$，令 $s=A(u\cos\phi+v\sin\phi)$，并记 $X=(d|u),\,Y=(d|v)$。似然比的对数为 $A(X\cos\phi+Y\sin\phi)-A^2/2$。先对 $\phi$ 极大，再对非负的 $A$ 极大，得 $2\ln\Lambda_{\max}=R^2=X^2+Y^2$。纯噪声时 $P(R>r)=\e^{-r^2/2}$，有信号时 $R^2$ 服从非中心 $\chi^2_2$，非中心参数为真实信噪比平方。对相位取均匀先验作边缘化则须积分，似然比变为 $\e^{-A^2/2}I_0(AR)$，其中 $I_0$ 为零阶修正 Bessel 函数。即使对同一组记录，极大值与积分值也只在一定条件下给出相似的排序。

若这一个模板平移许多到达时刻，则 $X(t)$ 自身是一个随机过程。不同延迟的协方差为 $(u_t|u_{t'})$。若有 $N$ 次独立尝试，任一统计量越阈的概率是 $1-(1-p)^N$，小 $Np$ 下约为 $Np$。密集的时移和相邻质量模板并不独立，实际总虚警率可由整套搜索在错时数据中的结果估计。例如两个独立台站，单台越阈事件率为 $\lambda_1,\lambda_2$，要求到达差在预定时差两侧各 $\Delta t$ 内，稀疏事件的偶然符合率为 $2\Delta t\,\lambda_1\lambda_2$。若还要质量及频段相符，应在估背景时施加同样的条件。把一台相对另一台移过大于物理时差的时间，再反复计算符合，可保留单台偶发扰动的真实形状。共有的环境干扰若恰在零时差相关，则错时背景不能估计这种污染。

分辨不同的波形参数还需考察模板间距。取归一化波形 $\hat h(\lambda)$，由 $(\hat h|\hat h)=1$ 得 $(\hat h|\partial_a\hat h)=0$，再微分一次消去二阶导数在 $\hat h$ 上的投影，便有
\[
 1-(\hat h(\lambda)|\hat h(\lambda+\Delta\lambda))
 =\frac{1}{2}(\partial_a\hat h|\partial_b\hat h)\,
              \Delta\lambda^a\Delta\lambda^b+O(\Delta\lambda^3).
\]
对大多误差来自相位的模板，这个度量也可写成相位导数的协方差。设归一权重 $W(f)\propto|\tilde h(f)|^2/S_n(f)$，则对常相位极大后，匹配率是 $|\langle\e^{\i\Delta\Psi}\rangle_W|$。展开指数并取模，损失为 $[\langle(\Delta\Psi)^2\rangle_W-\langle\Delta\Psi\rangle_W^2]/2$。再对到达时刻极大，等于从 $\Delta\Psi(f)$ 中扣掉最佳线性函数 $a+bf$。这一投影解释了为何给每个频点都增添同一相位不需要增添一个质量模板。

取一个容易检验的极限。对时长 $T$ 的矩形窗单色信号，频率误差 $\Delta f$ 给相位极大后的匹配率 $|\operatorname{sinc}(\pi T\Delta f)|$，小误差时损失 $(\pi T\Delta f)^2/6$。若所有源都在 Euclid 空间均匀分布，振幅损失 $m$ 使探测距离缩为 $1-m$ 倍，可见体积缩为 $(1-m)^3$ 倍。质量网格的密度要由可接受的损失与计算代价共同确定。若波形族本身少了一项相位，任意加密网格也未必能补偿这项模型误差。

\subsection{短脉冲、到达时刻和带宽}\label{sec:gw-b-burst}
对持续时间甚短而形状未知的波，先定义两偏振合计的 $h_{\rm rss}^2=\int(h_+^2+h_\times^2)\d t$，量纲为时间。观测方向上的能流由式 \eqref{eq:gw-plane-wave-stress} 给出。若频谱集中在 $f_0$ 附近，Parseval 等式给
\[
 \int {\cal F}_{\rm gw}\d t
 =\frac{c^3}{16\pi G}\int(\dot h_+^2+\dot h_\times^2)\d t
 \simeq \frac{\pi c^3}{4G}f_0^2h_{\rm rss}^2 .
\]
乘 $4\pi D^2$ 得附近源的各向同性等效能量；若源辐射有方向性，总能量须先对真实方向图积分。若两偏振相位或方向未知，需要用投影后的 $s$ 才能算信噪比。平坦噪声带内，$\rho^2\simeq2\int s^2\d t/S_n(f_0)$。因此相同的峰值应变在持续时间变短时较难检测，固定能量、距离与辐射图样时，高频脉冲又比低频脉冲具有更小的 $h_{\rm rss}$。

未知波形怎样搜索？把白化记录在一个时频区域内展开为 $k$ 个独立的实正交系数 $d_i$，对区域内任意信号 $s_i$，有 $2\ln\Lambda=2\sum_i d_i s_i-\sum_i s_i^2$。逐项极大给 $s_i=d_i$，所得统计量 $E=\sum_i d_i^2$ 在零假设下服从 $\chi_k^2$，有信号时均值由 $k$ 加至 $k+\rho^2$，纯噪声标准差为 $\sqrt{2k}$。故把太多无信号的频格圈入会抬高门槛。宽 $B$、长 $T$ 的实带通信号约有 $2BT$ 个实自由度。不同尺度的窗口、小波或若干相邻显著频格的聚类，都是为了用较少的自由度包括较多的信号能量。筛选窗口本身也带来额外的尝试次数。

取复指数乘 Gaussian 包络 $g(t)\propto\exp[-(t-t_0)^2/(4\sigma_t^2)+2\pi\i f_0 t]$。则 $|g|^2$ 的时间方差为 $\sigma_t^2$，$|\tilde g|^2$ 的频宽为 $\sigma_f=1/(4\pi\sigma_t)$。固定 $Q=f_0/\sigma_f$ 等于在高频使用短窗、低频使用长窗。窄谱振铃在长窗里分辨较好，脉冲边沿在短窗里定时较好。不能把相邻且重叠的 Gaussian 窗都当独立的 $\chi^2$ 试验。

有大信噪比的已知波形时，求到达时刻的精度并非简单地取 $1/f_0$。频域平移乘以 $\e^{-2\pi\i ft_0}$。把常相位作为伴随参数消去，得到关于 $t_0$ 的信息 $4\pi^2\rho^2(\langle f^2\rangle_W-\langle f\rangle_W^2)$，从而
\[
 \sigma_{t_0}\simeq \frac{1}{2\pi\rho\sqrt{\langle f^2\rangle_W-\langle f\rangle_W^2}} .
\]
单色波即使很强，时移也与相位变换不可分。多个频率共同确定包络，才可给相对到达时刻。两地基线长度 $b$，与传播方向夹角为 $\vartheta$，则在已确定时差环附近有 $\sigma_\vartheta\simeq c\,\sigma_{\Delta t}/(b|\sin\vartheta|)$。接近基线端射方向时这一线性误差式失效，应使用完整的 $\cos\vartheta$ 似然。

\subsection{误差区间与模型的权重}
一次事件给出最佳拟合波形后，怎样给质量或距离报误差？第三章由波形导数构造了 Fisher 矩阵。这里先保留一个常被省掉的步骤。设真实参数为 $\theta_*$，数据 $d=h(\theta_*)+n$。极值方程 $(d-h|\partial_a h)=0$ 在线性近似下给 $\Gamma_{ab}\Delta\theta^b=(n|\partial_a h)$。对噪声平均，右端均值为零、协方差为 $\Gamma_{ab}$，于是 $\langle\Delta\theta^a\Delta\theta^b\rangle=(\Gamma^{-1})^{ab}$。这一步平均的是同一真实源在不同噪声中的估计值。若以参数的先验密度 $\pi(\theta)$ 乘似然，再将所得乘积归一化，就得到给定这一份记录的后验。远离边界且局域先验近乎常数时，两种计算给相似的误差椭球。

设只关心一个相位参数 $\alpha$，同时要估计另一个参数 $\beta$。两者信息矩阵的对角元为 $A,B$，非对角元为 $C$。若固定 $\beta$，$\alpha$ 的方差近似于 $1/A$；若 $\beta$ 也由数据拟合，则
\[
 \sigma_\alpha^2=\frac{B}{AB-C^2}
 =\frac{1}{A(1-r^2)},\qquad r=\frac{C}{\sqrt{AB}} .
\]
这是波形导数夹角的余弦进入误差的方式。当 $r\to1$ 时，两种物理效应可有很大的单独导数，同一条记录却主要测到二者之和。若有独立资料给 $\beta$ 方差 $\sigma_{\beta,0}^2$，可将 $B$ 替换为 $B+\sigma_{\beta,0}^{-2}$，其它两项不变。到达时刻的已知电磁对应体、连续波的射电自转相位，便可通过这种途径改善引力波参数精度。若先验有截断或多峰，则应直接积分，简单相加信息并不充分。

取一个带物理边界的可算例子。设一单位模板的输出为 $X$，观测误差标准差为一，振幅 $A\ge0$。若对非负振幅采用平坦密度，并在远大于似然宽度处截断，则 $p(A|X)\propto\exp[-(A-X)^2/2]$ 在 $A=0$ 截去一侧。用标准正态累积分布 $\Phi_{\rm N}$，有
\[
 P(0<A<a\mid X)=
 \frac{\Phi_{\rm N}(a-X)-\Phi_{\rm N}(-X)}{1-\Phi_{\rm N}(-X)} .
\]
若 $X=0$，后验的 $90\%$ 单侧上限满足 $\Phi_{\rm N}(a)=0.95$，即 $a\simeq1.645$。若换成对 $A^2$ 平坦的先验，振幅密度多一因子 $A$，同一记录给 $1-\e^{-a^2/2}=0.9$，上限为 $\sqrt{2\ln10}$。这两种区间都以具体先验为条件。频率学派的置信区间则须说明重复实验的构造规则，使给定真实 $A$ 时覆盖它的概率达到所报比例。

先验有时由源群给出。近邻均匀球内的距离密度为 $3D^2/D_{\max}^3$。若方向和本征强度已知而 $A=K/D$，则振幅先验正比于 $A^{-4}$，只从 $K/D_{\max}$ 起算。将平坦振幅先验直呼为「对距离无先验」会改变弱信号的推断。对方向未知的双星还要同时积分倾角；把振幅似然压缩为单一对称的距离误差，常会损失面向观者系统的长后验尾。

两个假设的后验赔率可写成先验赔率乘以 Bayes 因子。以一单位模板为例，假设 $H_0$ 固定 $A=0$，而 $H_1$ 令实振幅取均值零、方差 $\sigma_A^2$ 的正态先验，用符号的两向性代表已知波形及其反号。对 $A$ 在实轴积分，$X$ 在 $H_1$ 下的边缘分布方差为 $1+\sigma_A^2$。于是
\[
 \frac{p(X|H_1)}{p(X|H_0)}
 =\frac{1}{\sqrt{1+\sigma_A^2}}\,
 \exp\!\left[\frac{X^2}{2}\frac{\sigma_A^2}{1+\sigma_A^2}\right].
\]
当先验宽度远大于噪声尺度，极大似然比给出的指数增益几乎都在，但要乘上一个宽度的倒数。若备择假设允许许多数据根本不支持的模板，模型证据要将这些位置也平均进去。与单次虚警概率直接取倒数会缺失这一积分，而且还未给出仪器触发来自天体的先验率。

\subsection{波形误差怎样偏移参数}
有时观测噪声已很小，错误来自用于反演的波形族。将真实应变写成 $h_{\rm true}=h(\theta_*)+\delta h$，用同一极值方程线性化，得到
\[
 \Delta\theta^a=(\Gamma^{-1})^{ab}(\partial_b h|n+\delta h).
\]
其中含 $n$ 的项随数据改变，含 $\delta h$ 的项在重复观测同一个源时不消失。若波形振幅整体放大 $g$ 倍，而相对模型误差固定，$\Gamma\propto g^2$、$(\partial_bh|\delta h)\propto g^2$。于是统计误差按 $1/g$ 下降，系统偏移约恒定。更近的源有更高信噪比，也对模板精度提出更严格的要求。

定义到模板切空间的投影 $P v=\partial_a h\,(\Gamma^{-1})^{ab}(\partial_b h|v)$。拟合能吸收 $P\delta h$，无噪声最佳残差为 $(1-P)\delta h$。因此小的拟合残差本身并不能证明质量的偏移很小。一个恰好沿啁啾质量导数的相位误差，在一阶可被全部拟合掉。若关心的是一族模板对叠加在噪声中的真实源的检出效率，剩余正交分量较为重要；若关心的是比较啁啾质量和电磁测得的质量，还须估计沿切空间的偏移。

在未知瞬时干扰很多的频带，单一的高匹配滤波输出也未必可信。取一个单位模板，分为 $p$ 个互不重叠的子频带 $u_i$，使 $(u_i|u_i)=1/p$。若波形真是 $A\sum_i u_i$，则各实输出 $X_i=(d|u_i)$ 的均值均为 $A/p$，方差 $1/p$，带间独立。写 $X=\sum_i X_i$，则
\[
 \chi^2=p\sum_i(X_i-X/p)^2
\]
在理想 Gaussian 模型中服从 $\chi^2_{p-1}$，且与 $X$ 独立。这是将一个 $p$ 维独立正态矢量沿 $(1,\dots,1)$ 及正交方向分解的结果。对常相位未知的两个正交分量作相同计算，自由度加倍为 $2p-2$。若强瞬态只把一小段频带抬高，虽然 $X$ 很大，差分统计量也会很大。若模板本身不准，真实信号同样能令此统计量变大，因此门槛需要用注入实验标定。

计算上限时还要区分一次事件的振幅上限和一类源的事件率上限。例如观测了若干不同时段，各时段的可见体积与探测效率给有效体积时间 $\mathcal V=\int \epsilon(\lambda,\bm x,t)p(\lambda)\d\lambda\,\d V\,\d t$。设事件是均匀独立的 Poisson 过程，且阈值以上无事件，$P(N=0|\mathcal R)=\exp(-\mathcal R\mathcal V)$。一个单侧 $90\%$ 零计数界为 $\mathcal R_{90}=\ln10/\mathcal V$。拟合噪声中波形振幅时的 $1.645$ 不能替代这里的 $\ln10$。若背景触发数不可忽略，应同时给信号与背景写计数或样本似然。

\subsection{从能流到频谱的两个算例}
先取单一方向、单一偏振的 Gaussian 正弦包 $h(t)=h_0\e^{-t^2/(2\tau^2)}\cos\omega_0t$。正负载频的两个峰并非总可独立处理。记 $b=\omega_0^2\tau^2$，直接用 $\cos^2 u=(1+\cos2u)/2$ 积分得 $h_{\rm rss}^2=h_0^2\sqrt\pi\tau(1+\e^{-b})/2$。再把 $\dot h$ 的包络导数也保留。利用 $\int t^2\e^{-t^2/\tau^2}\cos(2\omega_0t)\d t=-\tfrac14\partial_{\omega_0}^2(\sqrt\pi\tau\e^{-b})$，各交叉项合并后有
\[
 \frac{\int\dot h^2\d t}{\int h^2\d t}
 =\frac{1}{2\tau^2}+\frac{\omega_0^2}{1+\e^{-b}} .
\]
在 $\omega_0\tau\gg1$ 时，领项为 $\omega_0^2$，还留包络的 $1/(2\tau^2)$ 修正。对 $\omega_0=0$，本式仍给非零能流，只依赖升降边沿。用中心频率代替全部时间导数在这个极限会错误地预言零能量。

再取并合后简化的单模 $h(t)=H\,\Theta(t)\e^{-t/\tau}\sin\omega_0t$。它在 $t=0$ 连续，导数不产生 $\delta$ 脉冲。正频与负频两极点的差给
\[
 \tilde h(f)=\frac{H\omega_0}{(\tau^{-1}+2\pi\i f)^2+\omega_0^2},
 \qquad \int h^2\d t=\frac{H^2\tau}{4}
       \frac{(\omega_0\tau)^2}{1+(\omega_0\tau)^2}.
\]
设噪声在信号带内近乎平坦且 $Q=\omega_0\tau/2\gg1$，则最优的 $\rho^2\simeq H^2\tau/(2S_n)$。衰减慢使功率压入窄线并增加可积累的周期数。若在任意相位处用阶跃突然将非零振幅的波形接入，模型的导数却多一 $\delta$ 项，不能直接对它求平方来估总发射能量。求信噪比可对有限观测作适当窗选；归属于物理源的辐射能量还要求合理的启动或与前段波形平滑匹配。

\subsection{周期源的相位钟}
考虑一颗固定在天球方向 $\bm n$ 的孤立星，这里 $\bm n$ 从太阳系质心指向源。去掉与方向无关的常量，同一波面在探测器时刻 $t$ 所对应的质心时刻为 $t_b=t+\bm n\cdot\bm r(t)/c$，其中 $\bm r$ 是测站相对于质心的位置。精密计时还须作 Einstein 和 Shapiro 修正，其来源见第三章\hyperlink{subsection.3.13.1}{「双星脉冲星的计时」}。令 $\tau=t_b-t_{\rm ref}$，缓慢减速的星体给 $\Phi(t)=\phi_0+2\pi(f_0\tau+\dot f_0\tau^2/2+\ddot f_0\tau^3/6+\cdots)$。最低阶的瞬时频率是 $f(t_b)[1+\bm n\cdot\bm v(t)/c]$。若在每段记录中忽略调频，频率漂移须小于一格，即 $|\dot f_{\rm obs}|T^2\lesssim1$。地球运动给 $\dot f_{\rm obs}$ 一项 $f_0\bm n\cdot\dot{\bm v}/c$。几小时内看来静止的谱线，在半年记录中可扫过许多格。对已知方向可按 $t_b$ 等间隔重新采样，消去这一扫频，再在质心系内修正自转导数。

设探测器处在半径 $R$、角频率 $\Omega$ 的圆轨道上，源方向在轨道平面的投影长 $\sin\vartheta$。由传播时差得相位调制 $\beta\cos(\Omega t-\varphi)$，其中 $\beta=2\pi f_0R\sin\vartheta/c$。将 $\exp(\i\beta\cos u)=\sum_k\i^k J_k(\beta)\exp(\i ku)$ 代入，原来的单线分裂为间隔 $\Omega/(2\pi)$ 的边带，显著边带大抵止于 $|k|\sim\beta$。由 $\sum_kJ_k^2=1$，总功率守恒。对极轴上的源 $\sin\vartheta=0$，这项调制消失。日周 Doppler 相位和年周 Doppler 相位可以分别计算，其频谱是两组边带的卷积。

探测器还会随地球自转而改变指向。其长波响应张量在惯性系中为 $D(t)=R(t)D(0)R^{\mathsf T}(t)$。每个矩阵元最多含两个自转角的正弦或余弦相乘。故扣掉 Doppler 相位后，振幅中只含 $0,\pm\Omega_\oplus,\pm2\Omega_\oplus$ 五个日周分量。令 $a(t),b(t)$ 为偏振角为零时的两列方向响应，将未知偏振角、倾角、初相位和总振幅合并为四个常数，信号可写成 $s=\sum_{\mu=1}^4 A^\mu h_\mu$，四个 $h_\mu$ 即 $a,b$ 分别乘 $\cos(\Phi-\phi_0),\sin(\Phi-\phi_0)$。取 $x_\mu=(d|h_\mu)$、$M_{\mu\nu}=(h_\mu|h_\nu)$，极大化线性参数给
\[
 M_{\mu\nu}\widehat A^\nu=x_\mu,\qquad
 2\mathcal F=x_\mu(M^{-1})^{\mu\nu}x_\nu .
\]
若 $M$ 满秩，无信号时 $2\mathcal F\sim\chi_4^2$。观测太短而 $a,b$ 比值近乎固定时，实际可辨的振幅组合少于四个。此时用可测子空间中的逆，或把其它台站的数据一并加入内积。这个分离让搜索只须在天区、频率和频率导数上布网。参见 \href{https://arxiv.org/abs/gr-qc/9804014}{Jaranowski、Królak 与 Schutz}。

把方向偏开 $\Delta\bm n$，长基线上遗留的相位是 $2\pi f_0\,\Delta\bm n\cdot[\bm r(t)-\bm r(t_{\rm ref})]/c$。同一记录若还拟合 $f_0,\dot f_0$，则位置的常数、速度及部分加速度投影也能被这些参数吸收。故短弧天区精度不宜用整条地球公转直径估算。长时观测扫过完整轨道后，要求残留相位小于一弧度，给出量级 $\Delta\theta\lesssim c/(2\pi f_0R)$，在几何退化方向须单独计算。另一方面，频率导数的相位误差正比于 $T^2\Delta\dot f$，故网格步长要按 $T^{-2}$ 缩小。

当相干搜索总代价过高，可先将 $T=N T_s$ 分成 $N$ 段，分别求短时谱，沿给定源的频率轨迹相加功率。段内相位已被利用，段间相位则不需对准。设每段的白化功率在纯噪声下均值、标准差均为一，弱信号带来的均值增量正比于 $h_0^2T_s/S_n$。相加后的显著性正比于 $h_0^2T_s\sqrt N/S_n$，故固定虚警率和检出率所需 $h_0\propto(S_n/T_s)^{1/2}N^{-1/4}$。相干相加则为 $(S_n/T_s)^{1/2}N^{-1/2}$。把短段加长虽然改善灵敏度，却迫使段内频率及位置网格变密。

另一取法是只记每段的功率是否超过阈值 $u$。若功率在噪声下服从均值一的指数分布，单格被选概率为 $p=\e^{-u}$。对一个试探源沿轨迹数峰，纯噪声下计数的均值为 $Np$，方差 $Np(1-p)$。反过来，每个已选时频格给所有可能经过该格的源参数投一票，便是 Hough 搜索。幅度异常高的一格也只投一票，对单次强干扰较稳健。若不同段灵敏度不同，可给票加权，理想权重正比于该段在弱信号下的选中概率增量除以其噪声方差。分级搜索先用计数缩小参数空间，再对候选区作长时相干跟踪。

\subsection{双星的可见距离}
第三章 \S\ref{sec:pn-wave-energy} 得到准圆旋近的驻相相位。对探测率还需要它的振幅。设源轨道轴和视线夹角为 $\iota$，记 $u=\cos\iota$。主谐波的两偏振分别带 $(1+u^2)/2$ 与 $u$，相差四分之一周期。对一台直角干涉仪，定义 $w^2=F_+^2(1+u^2)^2/4+F_\times^2u^2$。在最优天区、面向观者时 $w=1$。取观者所量的啁啾质量 $\mathcal M_z$，用光度距离 $D_L$，限于领阶振幅及频率缓变区间，有
\[
 \rho^2=\frac{5}{6\pi^{4/3}}\frac{c^2w^2}{D_L^2}
       \left(\frac{G\mathcal M_z}{c^3}\right)^{5/3}
       \int_{f_{\min}}^{f_{\max}}\frac{f^{-7/3}}{S_n(f)}\,\d f .
\]
这里的 $f_{\max}$ 同时受仪器带宽与旋近近似的终止频率限制。时间平移只给模板乘一个线性相位，所有到达时刻的内积可由同一次快速 Fourier 变换得到。振幅及常相位可如前解析极大化。用到的只有主谐波；进动、潮汐和高谐波均可能使这个分离不再精确。

令 $\rho=\rho_*$ 且 $w=1$ 所得 $D_{\rm hor}$ 称最优视距。随即有三个不同的方向平均。第一，偏振基旋转不改 $F_+^2+F_\times^2$。用第三章的方向函数，对方位角先平均，剩下 $(1+6\cos^2\theta+\cos^4\theta)/8$，再对 $\cos\theta$ 从 $-1$ 至 $1$ 平均得 $2/5$。第二，平均轨道倾角给 $\langle(1+u^2)^2/4\rangle=7/15$、$\langle u^2\rangle=1/3$。再平均偏振角给 $\langle w^2\rangle=(7/15+1/3)/5=4/25$。这给的是平方信噪比平均。第三，均匀空间中方向 $w$ 的可见体积正比于 $w^3$。故等体积视距为 $D_{\rm vol}=D_{\rm hor}\langle w^3\rangle^{1/3}$。对各向同性天区、均匀的 $u$ 和偏振角，直接作上述三重角积分约得 $D_{\rm vol}=D_{\rm hor}/2.26$。

若把一次探测换成总体事件率，还须按源时刻计时。令 $\mathcal R(z,\lambda)$ 为单位共动体积、单位源固有时间和单位内参 $\lambda$ 的事件率，$p_{\rm det}(z,\lambda)$ 为方向和噪声平均后的检出率，则观测率为
\[
 \dot N_{\rm obs}=\int \mathcal R(z,\lambda)\,\frac{1}{1+z}
 \frac{\d V_c}{\d z}\,p_{\rm det}(z,\lambda)\,\d z\,\d\lambda .
\]
这里积分波形用红移后的质量，而 $\mathcal R$ 是源系质量的函数。在低红移、带宽固定且全为旋近信号时，$D_{\rm hor}\propto\mathcal M^{5/6}$，故可见体积正比于 $\mathcal M^{5/2}$。一组检出事件中大质量系统会偏多，即使其本征数密度并不高。若总质量已大至使末段移出灵敏带，上述幂律即失效。

\subsection{用台站网分开距离和轨道倾角}
同一双星在多台仪器中的信号应共用源方向、到达时刻、距离和两偏振的相位关系。在第三章的零流检验中，消去两个张量分量便可检验一致性。这里反过来，求处于张量平面内的那两份信号以及它们的误差。设天区已知，按几何时差移齐各台记录并依自谱白化。在一个频格上取 $F$ 为白化后的两列响应，按约定包含频率传递函数。独立噪声协方差已化为单位阵，最小二乘方程为 $(F^\dagger F)\hat{\bm h}=F^\dagger\bm d$。其误差协方差为 $(F^\dagger F)^{-1}$；若噪声尚有台间相关，应先对白化使用完整的谱矩阵。可辨精度取决于 $F$ 的两个奇异值。当最小奇异值很小时，已知方向和大总信噪比仍不足以精确测出较弱的那个偏振组合。

在长波实响应且各台共用一个相位模板时，可以把频率积分并入正权重 $w_I$。写
\[
 A=\sum_I w_I F_{I+}^2,\qquad B=\sum_I w_I F_{I\times}^2,\qquad
 C=\sum_I w_I F_{I+}F_{I\times}.
\]
则
\[
 AB-C^2=\sum_{I<J}w_Iw_J(F_{I+}F_{J\times}-F_{J+}F_{I\times})^2 .
\]
证法是展开 $AB$，同站项与 $C^2$ 相消，再将 $(I,J)$ 与 $(J,I)$ 合并。同一方向上有两台完全平行的仪器时，额外的台站虽可提高同一应变组合的信噪比，却不能使这个行列式非零。若各站取向使许多二阶子式显著，则由 $F$ 的两个左、右主方向可分开两种偏振。转动偏振基，角为 $2\psi$ 的正交变换会使 $F^\dagger F$ 对角化，其本征值为 $[A+B\pm\sqrt{(A-B)^2+4C^2}]/2$。这两个数对选取偏振角的零点不变。

假设已重构两个偏振，并在适当的偏振主轴与相位零点下求得主谐波的椭圆半轴。非进动圆双星给半轴 $H(1+u^2)/2$ 与 $H|u|$，其中 $H\propto D_L^{-1}$，$u=\cos\iota$。两半轴之比 $\kappa=2|u|/(1+u^2)$ 与总距离无关。要由比值反推倾角，解二次方程得
\[
 |u|=\frac{\kappa}{1+\sqrt{1-\kappa^2}},\qquad
 H=\frac{2H_{\rm major}}{1+u^2}.
\]
保留相位随时间转动的方向，才能确定 $u$ 的符号，只有两半轴长度时互为镜像的倾角仍等价。对边缘朝向的轨道 $u\simeq0$，弱轴按 $|u|$ 线性长出；但若它落在测不准的偏振通道，实际仍很难给出 $\kappa$。面对观者时写 $|u|=1-\epsilon$，则 $\kappa=1-\epsilon^2/2+O(\epsilon^3)$。两偏振的比直到 $\epsilon^2$ 阶才偏离圆，振幅却已在 $\epsilon$ 阶随方向变小，故微小的距离改变即可在一阶掩盖倾角变化。

同一个退化也可从左右圆偏振看出。把相位差 $\pi/2$ 的两偏振组合，强弱两种螺旋度的幅度正比于 $H(1+u)^2/4$ 与 $H(1-u)^2/4$。当 $u\to1$ 时，反向螺旋度与主项之比为 $(1-u)^2/(1+u)^2$。用 $\iota\ll1$ 近似，这个比为 $\iota^4/16$。即使两个偏振方向的仪器噪声完全相等，也需要很高信噪比才能由波形本身量出这样小的椭率偏差。若并合激发的高次谐波可测，因其自旋加权球谐在天球上的角形不同，可能提供独立于主谐波比值的倾角信息。与电磁喷流联合分析依靠的则是另一假定，即喷流和轨道轴的相对取向已有模型。

再来看时差决定的方向是否唯一。以 1 号站为原点，两个基线矢量为 $\bm b_2,\bm b_3$，测得 $\bm b_i\cdot\bm N=c\Delta t_i$。令 $G_{ij}=\bm b_i\cdot\bm b_j$，解平面内的投影为 $\bm N_\parallel=\sum_{ij}\bm b_i(G^{-1})_{ij}c\Delta t_j$。若两基线不共线且 $|\bm N_\parallel|\le1$，余下两个解为
\[
 \bm N=\bm N_\parallel\pm
   \sqrt{1-|\bm N_\parallel|^2}\,
   \frac{\bm b_2\times\bm b_3}{|\bm b_2\times\bm b_3|}.
\]
三个站只给两个独立的时差，故两解关于台站平面互为镜像。相对振幅、相位或另一条不共面的基线常可选择其中之一。若求出的 $|\bm N_\parallel|>1$，不应强行把平方根改为零；应以有限的时差误差在单位球面上重新求似然。对来向靠近台站平面，时差的一阶变化在某方向上很小，天区误差通常呈细长弧形，局域 Gaussian 椭圆只是近似。

作一个量纲例。两条主要基线的光行时若为 $10\,\mathrm{ms}$，一个有效均方带宽 $100\,\mathrm{Hz}$、信噪比为十的台站，理想时刻标准差约为 $0.16\,\mathrm{ms}$。两站精度相同时，时差误差为其 $\sqrt2$ 倍。在与基线成直角附近，这相当于约 $0.023$ 弧度的角厚度。若将它换算为天空面积，还须乘另一方向的宽度并考虑两基线的投影夹角。把时差的精度直接当作每个方向上相同的角精度，会特别高估只有两台仪器的定位能力。

\subsection{背景的应变谱与能量谱}\label{sec:gw-b-background}
一颗星的初相位未知，仍可将它当作待拟合参数。许多独立源各有自己的相位时，可靠的预测却常是不同方向和频率之间的二点函数。设背景平稳、各向同性、无偏振，写实场为
\[
 h_{ij}(t,\bm x)=\sum_A\int_{-\infty}^\infty\d f\int\d^2\bm N\,
  h_A(f,\bm N)e^A_{ij}(\bm N)\,
  \exp[2\pi\i f(t-\bm N\cdot\bm x/c)] .
\]
这里 $\bm N$ 沿波的传播方向，$e^A_{ij}e^{Bij}=2\delta^{AB}$，$h_A(-f,\bm N)=h_A^*(f,\bm N)$。定义 $S_g(f)$ 为正频率上、对全天积分且将两个偏振相加的单边谱。于是
\[
 \langle h_A^*(f,\bm N)h_B(f',\bm N')\rangle
 =\frac{S_g(|f|)}{16\pi}\,
 \delta_{AB}\delta(f-f')\delta^2(\bm N,\bm N') .
\]
因子 $16\pi$ 中包含全天 $4\pi$、两偏振和双边频率各一个因子二。令波场在同一点缩并，得 $\langle h_{ij}h^{ij}\rangle=2\int_0^\infty S_g(f)\d f$。原书以每偏振的单边谱 $S_h$ 表示同一背景，故 $S_g=2S_h$。以上对角形式还用了不同传播方向彼此不相关。只有在 Gaussian 假设下，二点函数才决定全部统计性质。稀疏爆发叠加出的背景可能有相同的谱而有完全不同的高阶矩。

将同一展开代入式 \eqref{eq:gw-effective-stress} 的能量密度 $\epsilon_{\rm gw}=c^2\langle\dot h_{ij}\dot h^{ij}\rangle/(32\pi G)$。每个时间导数带 $2\pi f$，再把单位频率间隔换成单位对数间隔，故
\begin{equation}\label{eq:gw-b-omega}
 \frac{\d\epsilon_{\rm gw}}{\d\ln f}=\frac{\pi c^2}{4G}f^3 S_g(f),
 \qquad \Omega_{\rm gw}(f)=\frac{2\pi^2}{3H_0^2}f^2h_c^2(f),
 \qquad h_c^2(f)=fS_g(f).
\end{equation}
其中 $\Omega_{\rm gw}$ 以今日临界能量密度 $3H_0^2c^2/(8\pi G)$ 归一化，式中频率是在今日测得的，波长须处在可作短波平均的范围内。平的能量谱对应 $S_g\propto f^{-3}$；平的应变功率谱所含的能量则向高频迅速增大。对确定性的啁啾也有人用「特征应变」称 $2f|\tilde s(f)|$，因为它的平方除以 $fS_n$ 给每对数区间的 $\rho^2$。这两个同名量一个由实测时长内的波形定义，一个由单位时间的统计平均定义，比较时应先写清积分。

若把背景当作引力子气体，每个偏振每单位体积的模数为 $\d^3 k/(2\pi)^3$，单量子能量为 $\hbar c k$。两个等占据的偏振在每个相空间格平均有 $n(k)$ 个粒子，扣除真空项后 $\d\epsilon_{\rm gw}/\d\ln k=\hbar c k^4 n(k)/\pi^2$。同式 \eqref{eq:gw-b-omega} 相等且取 $k=2\pi f/c$ 就可将能量密度换为占据数。固定 $\Omega_{\rm gw}$ 时 $n\propto f^{-4}$，低频段即使只占总能量的一小部分，也可有极大的模占据数。这给将传播的随机波作为经典场处理的一个尺度判据。

\subsection{两台仪器的互谱}
对长波直角干涉仪，$s_I=D_I^{ij}h_{ij}$。由上述二点函数，单台背景的单边谱为 $(S_g/2)\langle F_+^2+F_\times^2\rangle=S_g/5$。若只有一台仪器，所见自谱还包括本底 $P_I(f)$。没有对噪声的独立了解就难以区分两个任意的正谱。考虑两个位置为 $\bm x_1,\bm x_2$ 的仪器。其随机信号互谱含有同一波面的传播时差。在长波臂近似下定义
\begin{align}
 \gamma_{12}(f)&=\frac{5}{8\pi}\int \d^2\bm N\,
 \e^{2\pi\i f\bm N\cdot(\bm x_1-\bm x_2)/c}
 \sum_A F_{1A}F_{2A},\nonumber\\
 S_{12}^{\rm gw}(f)&=\frac{1}{5}\gamma_{12}(f)S_g(f)
 =\frac{3H_0^2}{10\pi^2}\,\frac{\gamma_{12}(f)\Omega_{\rm gw}(f)}{f^3}.
 \label{eq:gw-b-cross}
\end{align}
这里 $S_{12}$ 对应 $\tilde d_1^*\tilde d_2$。若响应本身为复数，要代入 $F_{1A}^*F_{2A}$。因 $\bm N$ 与 $-\bm N$ 的贡献共轭，以上长波响应对各向同性背景给出实的 $\gamma$。同址、同向时 $\gamma=1$；绕共同法向转动角 $\alpha$ 时 $\gamma=\cos2\alpha$。故一同址仪器转过 $90^\circ$ 会把应变反号，互谱同样变号。若两台离得很远，取 $f|\bm x_1-\bm x_2|/c\gg1$，在天空积分中相位迅速变化，绝大部分频率上的相关便很小。

这一积分可怎样计算？记 $P_{ij}=\delta_{ij}-N_iN_j$。偏振求和给 $P_{ik}P_{jl}+P_{il}P_{jk}-P_{ij}P_{kl}$。与两对臂张量缩并后，角积分化为常数、$N_iN_j$、$N_iN_jN_kN_l$ 三类。令 $\bm y=2\pi f(\bm x_1-\bm x_2)/c$，则 $\int\e^{\i\bm N\cdot\bm y}\d^2\bm N=4\pi j_0(y)$，每个 $N_i$ 可由 $\partial/(\i\partial y_i)$ 生成。故精确的长波臂结果只含球 Bessel 函数 $j_0,j_2,j_4$。例如两仪器平行放置，基线沿共同法向，方位角平均后便有
\[
 \gamma(y)=\frac{5}{32}\int_{-1}^{1}
 (1+6\mu^2+\mu^4)\e^{\i y\mu}\d\mu
 =\frac{5}{16}\big[j_0(y)-6j_0''(y)+j_0^{(4)}(y)\big].
\]
取 $y=0$ 得一，展开至二次得 $\gamma=1-11y^2/42+O(y^4)$。即使两台方向相同，也不能在整个带宽内都取 $\gamma=1$。基线与仪器臂平行时，上面多项式的系数也会改变。

在每个正频格 $f_k=k/T$ 上取 $C_k=(2/T)\operatorname{Re}[\tilde d_1^*(f_k)\tilde d_2(f_k)]$。设两台噪声统计独立，且远大于信号，测量均值为 $S_{12}^{\rm gw}$，方差为 $P_1P_2/2$。后二分之一可从两个独立复 Gaussian 变量的实虚部直接相乘得到。令待检谱为 $\Omega_{\rm gw}=A\,\omega(f)$，$\omega$ 形状已知。记 $K_k=S_{12}^{\rm gw}(f_k)/A$，则对这个单参数作加权最小二乘，有
\begin{align*}
 \widehat A&=\frac{\sum_k K_k C_k/(P_1P_2)_k}{\sum_k K_k^2/(P_1P_2)_k},&
 \sigma_A^{-2}&\simeq 2T\int_0^\infty\frac{K^2(f)}{P_1(f)P_2(f)}\,\d f .
\end{align*}
这就是互相关的最优滤波，见 \href{https://arxiv.org/abs/gr-qc/9710117}{Allen 与 Romano}。方向不利的频格通过 $\gamma$ 自动减权，反相关的频格符号反过来后仍增加灵敏度。互谱中的随机量实际为振幅的平方，因而给定检出率时可辨的谱幅 $A\propto T^{-1/2}$，可辨的特征应变却按 $T^{-1/4}$ 缩小。单色信号相位可追踪时的 $T^{-1/2}$ 振幅规律不能用于此处。

若噪声随时段而变，将各段估计用各自的 $\sigma_A^{-2}$ 加权。若信号本身不可忽略，联合 Gaussian 模型对单个频格给 $\operatorname{var} C_k=[S_{11}S_{22}+\operatorname{Re}(S_{12}^2)]/2$，这里的谱都含信号。这是有限样本对背景本身的涨落，不能靠降低仪器本底无限压低。若台站间存在非零噪声互谱 $P_{12}$，则估计的偏差正比于 $\sum K_k\operatorname{Re}P_{12}/(P_1P_2)_k$。加长观测只会更精确地测到这份偏差。可用环境监测通道测出共用磁场等扰动，再按各站的传递函数估其互谱。

对于方向性背景，把 $S_g(f)/(4\pi)$ 推广为 $S_g(f,\bm N)$。取 $S_g(f,\bm N)=H(f)\sum_\alpha a_\alpha Y_\alpha(\bm N)$，在不同时间段仪器随地球转动，得到不同的响应 $K_{k\alpha}(t)$。实际测得的是一组线性方程 $\langle Z_k(t)\rangle=\sum_\alpha K_{k\alpha}(t)a_\alpha$，其中 $Z_k=(2/T)\tilde d_1^*\tilde d_2$ 保留复互谱的两部分。带噪声求解时，法矩阵为 $K^\dagger W K$，$W$ 由各频格方差决定。某角模几乎落在该矩阵零空间中时，单纯求逆会给出巨大的噪声图样。分辨角尺度由基线的波长数、响应方向图和独立台站数共同决定。

\subsection{如何合并时段、频带和基线}
用两台仪器的相关限制能量密度时，常只给一个参考频率的数。这个数实际怎样取自宽带记录？以式 \eqref{eq:gw-b-cross} 写 $\Omega_{\rm gw}(f)=\Omega_\alpha(f/f_*)^\alpha$，令 $K_\alpha(f)=3H_0^2\gamma(f)(f/f_*)^\alpha/(10\pi^2f^3)$。互谱的最优权重正比于 $K_\alpha/(P_1P_2)$。若 $\alpha$ 改变，权重连同对 $\Omega_\alpha$ 的上限都改变。选择 $f_*$ 只是对同一条幂律重新标记振幅；将 $\alpha$ 固定为零则是关于源的物理假设。若完全不愿预设谱形，须在独立频带内各自估谱，并承担较少频格所带来的方差。

例如假定一带内 $P_1=P_2=P_0$、$\gamma=1$，测的是平的 $\Omega_*=\Omega_{\rm gw}$。令 $\mathcal B=3H_0^2/(10\pi^2)$，则
\[
 \sigma_{\Omega_*}^{-2}
 =\frac{2T\mathcal B^2}{P_0^2}
 \int_{f_1}^{f_2}f^{-6}\d f
 =\frac{2T\mathcal B^2}{5P_0^2}(f_1^{-5}-f_2^{-5}) .
\]
尽管噪声谱平坦，若 $f_2\gg f_1$，信息主要来自靠近 $f_1$ 的频率；原因是相同能量密度在低频有较大的应变谱。若假定平的是 $S_g$，各频格才近乎等权。若 $\gamma$ 在此处过零，对单个频格求 $\Omega=C_k f^3/(\mathcal B\gamma)$ 的方差发散。先将所有频格乘 $\gamma$ 再作加权和则保持有限，完全没有信息的格自然退出。

信号在两台的互谱也可能有两种分量。例如星系中的源群与宇宙学背景在所测频带的形状分别为 $K_1,K_2$，写 $S_{12}=aK_1+bK_2$。取 $\langle K_i,K_j\rangle=2T\int K_iK_j/(P_1P_2)\,\d f=W_{ij}$。欲估 $a$ 而允许任意 $b$，只需把 $K_1$ 中与 $K_2$ 平行的部分去掉，取 $K_{1\perp}=K_1-K_2 W_{12}/W_{22}$。则 $\langle K_{1\perp},K_2\rangle=0$，关于 $a$ 的有效信息为 $W_{11}-W_{12}^2/W_{22}$。若两源谱在当前带内几乎同形，再长的观测也不如增加能分开斜率的第二个频带有效。这样的投影也可用于已测得频谱形状的环境相关噪声；它只能去掉模板张成的那部分噪声。

对随时间改变的探测器，本底在每段 $\Delta T_j$ 内可近似稳定，各段最优估计为 $\widehat A_j$，方差 $\sigma_j^2$。求唯一的无偏最小方差线性估计，令 $\widehat A=\sum w_j\widehat A_j$ 并在 $\sum w_j=1$ 下极小化 $\sum w_j^2\sigma_j^2$，得 $w_j=\sigma_j^{-2}/\sum_i\sigma_i^{-2}$。若一段噪声的振幅谱增大为原来的两倍，两台都受影响，则 $P_1P_2$ 增为十六倍，这段资料在相同时间内的权重只及原来的十六分之一。先把噪声功率作算术平均再代入最优滤波，与逐段处理通常不同。若用了相互重叠的窗，权重应由完整的段间协方差求得。

在同一频率，四个台站给六对互谱。这六个数是否都可按独立的实测量合并？令单边谱矩阵为 $P_{ij}$，含噪声及信号，取频格内的 Fourier 振幅 $z_i$ 已归一化，使 $\langle z_i^*z_j\rangle=P_{ij}$。对正频率的一组圆对称复 Gaussian 变量，Wick 分解给
\[
 \operatorname{cov}\big(\operatorname{Re} z_i^*z_j,\operatorname{Re} z_k^*z_l\big)
 =\frac{1}{2}\operatorname{Re}(P_{ik}P_{lj}+P_{il}P_{kj}).
\]
例如 $(i,j)=(k,l)$ 时得 $[P_{ii}P_{jj}+\operatorname{Re}(P_{ij}^2)]/2$，重现前面的单基线方差。式中若谱为复须按共轭次序处理；上式按 $\langle z_i^*z_j\rangle$ 定义，第一项的两个因子对应 $z_i^*$ 接 $z_k$ 与 $z_l^*$ 接 $z_j$。纯独立噪声、弱信号时，不同的基线即便共用一个台站，只要另两个台站不相关，上式的互协方差也为零。若背景本身已很强，多对基线却在测同一随机实现，忽略互协方差会低估误差。

作一对方向相同、共址但噪声独立的理想仪器。令信号为 $h$，两台噪声独立同分布，自谱均为 $P$。将输出改为和、差 $d_\pm=(d_1\pm d_2)/\sqrt2$，则和通道的谱为 $2S_{11}^{\rm gw}+P$，差通道为 $P$。二者谱之差除以二就估计出 $S_{11}^{\rm gw}$。用原变量展开，$|d_+|^2-|d_-|^2=2\operatorname{Re}(d_1^*d_2)$，恰为互谱。若存在共同的地面运动，它也只出现在和通道；若两台取向不同，把差通道全当作纯噪声则会丢掉真实引力波。零信号的通道需由真实几何响应决定。

最后看背景的 Gaussian 假定何时有物理根据。设一频带内独立爆发的观者率为 $\mathcal R$，每个在带内持续 $\tau$，任意时刻重叠事件数的均值为 $D=\mathcal R\tau$。若过程为 Poisson 且各波包同分布，则零个事件的概率为 $\e^{-D}$。$D\ll1$ 时大部分记录没有信号，少数短爆发却给长时均方；$D\gg1$ 且没有单一源支配方差时，中心极限定理才使振幅趋于 Gaussian。例如把同一时刻的应变写成 $s=\sum_{i=1}^N a_i$，$N$ 服从均值 $D$ 的 Poisson 分布，$a_i$ 独立、零均值。由 $\ln\langle\e^{\i t s}\rangle=D[\langle\e^{\i t a}\rangle-1]$，二阶累积量为 $D\langle a^2\rangle$，四阶累积量为 $D\langle a^4\rangle$。超额峰度按 $D^{-1}$ 衰减，但若振幅分布有很重的尾部，所需的重叠数会很大。互谱仍测二点函数，依赖 Wick 分解的方差则应补这类四点关联，或由真实数据估计。

\subsection{棒、球与干涉仪能否测到同一背景？}
式 \eqref{eq:gw-b-cross} 中的 $1/5$ 专用于长波直角干涉仪。两种不同仪器若放在一起，该怎样定标互谱？先把任意一个线性读数写成 $s_D=D^{ij}h_{ij}$，回复或储光的频率增益已另行除去。张量偏振无迹，故 $D$ 的各向同性部分完全不参加。置 $\bar D=D-\tfrac13\delta\,\operatorname{tr}D$，另一仪器同样记为 $\bar E$。对同址仪器，须求的角平均为 $\mathcal A_{DE}=\langle\sum_A F_D^AF_E^A\rangle_{\bm N}$。用偏振完备式逐项展开，有
\[
 \sum_A F_D^AF_E^A
 =2\bar D:\bar E-4(\bar D\bm N)\cdot(\bar E\bm N)
       +(\bm N\cdot\bar D\bm N)(\bm N\cdot\bar E\bm N).
\]
这里 $D:E=D_{ij}E^{ij}$。对均匀球面有 $\langle N_iN_j\rangle=\delta_{ij}/3$，而四阶矩为 $(\delta_{ij}\delta_{kl}+\delta_{ik}\delta_{jl}+\delta_{il}\delta_{jk})/15$。后一式与两个无迹张量收缩，只剩两个相等的项。故三项平均后的总系数为 $2-4/3+2/15=4/5$，即 $\mathcal A_{DE}=4\bar D:\bar E/5$。按本章两偏振总谱约定，测得的单边互谱为 $\mathcal A_{DE}S_g/2$。这个算法不必为每一种仪器重新旋转两套偏振基。

例如两等长臂夹角为 $\zeta$ 时，$D=(\bm u\bm u-\bm v\bm v)/2$ 自身无迹，$D:D=[1-(\bm u\cdot\bm v)^2]/2=\sin^2\zeta/2$。其自谱为 $S_g\sin^2\zeta/5$。若将三角形空间仪器的一个简单 Michelson 通道先在长波近似下单独看，$60^\circ$ 夹角使功率比直角臂少 $1/4$；该三角形中不同合成通道仍有共有链路噪声，不能据此当作若干噪声独立的直角干涉仪。若两臂几乎同向，响应按夹角的一次方减小，信号功率按二次方减小，这与此时两个投影几乎相消一致。

对一根谐振棒，以轴向应变 $h_\parallel$ 归一，$E=\bm w\bm w$，则 $\bar E:\bar E=2/3$。若两根棒的轴夹角为 $\beta$，交叉收缩为 $\bar E_1:\bar E_2=\cos^2\beta-1/3$。用两自谱几何平均归一化的相关为 $(3\cos^2\beta-1)/2$。因而两根相互垂直的棒仍有负相关；在 $\cos^2\beta=1/3$ 时，虽然每根棒都有信号方差，等时的全天平均相关却为零。对于轴在直角干涉仪平面内的棒，若与第一臂成 $\beta$ 角，有 $D:\bar E=\cos2\beta/2$，归一化相关为 $\sqrt3\cos2\beta/2$。就算棒恰沿一条臂，同一源总体在两个仪器中的响应也并非完全相同。把这个相关再乘棒的复机械响应，才是光学相位与机械位移的原始互谱。

这些取向关系还提供一种区分天体背景与固定环境干扰的办法。设可转向的棒先沿第一臂、再沿第二臂各测一段，与同一干涉仪作互谱。若两段灵敏度相同，背景信号的均值分别为 $+K S_g,-K S_g$，而一项与棒转向无关的附加互谱 $P_c$ 在两段中相同。则互谱之差除 $2K$ 给 $S_g$，半和给 $P_c$。若任意一段包含不同的环境条件，或转动同时改变了磁场对读出电路的耦合，这个消去便不成立。对不同地点的仪器，天区平均还须带传播相位，不能继续用刚才的单一 $\cos2\beta$ 因子。

球形天线的五个四极模可取相互正交的无迹 $E_m$。由上述收缩，在相同径向阶且各模机械增益已除去时，各向同性背景的模协方差正比于 $\delta_{mn}$。单个方向的一列波却能在几个模之间产生确定的相位关系。由此一个球即使测齐五模，也不能简单地将两个正交模的互谱当作各向同性背景的主信号；同址相关所需的是重叠的潮汐投影，以及对仪器噪声相关的独立约束。

\section{弹性体怎样吸收引力波}\label{sec:gw-b-bar}
\subsection{由连续介质求有效长度}
第三章\hyperlink{subsection.3.12.11}{「把潮汐力转换成读数」}给出了一个受迫模的方程。现在从棒的质量分布计算方程中留下的两个常数。取均匀细棒，质量 $M$、全长 $L$、纵波声速 $v_s$。以质心为原点，令 $u(x,t)$ 表示沿轴的位移，$h_\parallel=h_{ij}n^in^j$ 为棒轴投影。长波潮汐力密度为 $\rho x\ddot h_\parallel/2$，自由端的纵向应力为零，故 $\ddot u-v_s^2u''=x\ddot h_\parallel/2$，且 $u'(\pm L/2)=0$。由于左右两端受力反号，质心不动，偶宇称伸缩模式也不被激发。可把受激部分写成 $u=\sum_{j\ge0}q_j(t)\sin k_jx$，其中 $k_j=(2j+1)\pi/L$、$\omega_j=v_sk_j$。共振时 $\omega_jL/c=(2j+1)\pi v_s/c$，对低阶弹性模，长波条件自然满足。

取 $U_j=\sin k_jx$，则模质量 $m_j=\int U_j^2\d m=M/2$。广义力由虚功 $\int (\ddot h_\parallel/2)x\,\delta u\,\d m$ 给出。除以模质量可得
\[
 L_j=\frac{1}{m_j}\int xU_j\,\d m
 =\frac{2}{L}\left[-\frac{x\cos k_jx}{k_j}
             +\frac{\sin k_jx}{k_j^2}\right]_{-L/2}^{L/2}
 =\frac{4(-1)^j L}{(2j+1)^2\pi^2}.
\]
所以常说的有效长度依赖于如何定义读出位移。取 $q_0$ 为一端的位移时 $L_0=4L/\pi^2$，取端到端的伸长为坐标就要同时换算质量和有效长度。振子与波的耦合强度涉及的组合 $m_jL_j^2$ 与模函数的任意归一化无关。对有限截面、非均匀密度或复合棒，只须把两个积分换成真实位移图样的质量积分。

\subsection{脉冲沉积的能量与吸收截面}
单取一模，记 $m=m_j,\ell=L_j,\gamma=\omega_0/Q$。机械响应函数为 $\chi(\omega)=[m(\omega_0^2-\omega^2+\i\gamma\omega)]^{-1}$，沿用前节的 Fourier 符号。于是 $\tilde q=-(m\ell/2)\omega^2\chi\,\tilde h_\parallel$。瞬态脉冲长度若远短于衰减时间，可在激发期间令 $\gamma=0$。脉冲过后，平衡位置两侧的自由振幅为
\[
 q_{\rm pk}=\frac{|\ell|}{2\omega_0}
 \left|\int \ddot h_\parallel(t')\e^{-\i\omega_0t'}\d t'\right|
 =\frac{|\ell|\omega_0}{2}|\tilde h_\parallel(f_0)|,\qquad
 E=\frac{m\ell^2\omega_0^4}{8}|\tilde h_\parallel(f_0)|^2 .
\]
分部积分已用到脉冲前后的导数消失，这里先取本身也回到零的波包。所得能量只取决于一个复谱点，棒内随后持续许久的正弦振荡是本征模的形状。两个总能量相同的入射波包，若一个恰在该频率附近有零点，沉积的能量可以十分不同。若加上阻尼，衰减时间对振幅为 $2/\gamma$，对能量为 $1/\gamma$。

对正入射、与棒取向最佳的单偏振谐波 $h_\parallel=h_0\cos\omega t$，稳态耗散功率为 $m\gamma\langle\dot q^2\rangle$，入射平均能流为 $c^3\omega^2 h_0^2/(32\pi G)$。二者之比即吸收截面，代入上式有
\[
 \sigma(\omega)=\frac{4\pi Gm\ell^2}{c^3}
 \frac{\gamma\omega^4}{(\omega^2-\omega_0^2)^2+\gamma^2\omega^2}.
\]
近峰时以 $4(\omega-\omega_0)^2+\gamma^2$ 代替分母除以 $\omega_0^2$ 的结果，便见峰高正比于 $Q$，半高全宽为 $\gamma$。对覆盖整个窄峰而不延伸到其它模的频带积分，$\int_{\rm peak}\sigma(\omega)\d\omega\simeq2\pi^2Gm\ell^2\omega_0^2/c^3$。峰下的面积与 $Q$ 无关。这个积分公式只用于窄峰，用单模黏滞阻尼模型一直积分到无穷大没有意义。代入基模的 $m,\ell,\omega_0$，面积正比于 $GMv_s^2/c^3$。改变棒长可移动敏感频段，同种材料、同质量下却很难靠增大 $L$ 提高积分截面。

与棒成夹角 $\theta$ 的来波，棒轴在横平面上的投影长 $\sin\theta$，响应含两个投影，故振幅正比于 $\sin^2\theta$。再对未知线偏振角平均，最佳截面乘 $\sin^4\theta/2$。对天区平均得 $4/15$，说明引用峰值截面时须说明方向。这里把能量的两偏振之和固定，与每偏振分别给单位功率时的方向平均有一因子二之差。

\subsection{用小振子读出大振子}
设棒的基模等效质量为 $M_b$，广义位移为 $x$；将质量 $m_t$ 的传感器通过弹簧装在读出端，绝对位移为 $y$。写 $k_b=M_b\omega_0^2$、$k_t=m_t\omega_t^2$，暂略损耗及传感器自身受到的直接潮汐力，则 $M_b\ddot x+k_b x+k_t(x-y)=F$、$m_t\ddot y+k_t(y-x)=0$。线性代数消去 $y$，自由振荡的频率由
\[
 (\omega_0^2-\omega^2)(\omega_t^2-\omega^2)
       -\mu\omega_t^2\omega^2=0,\qquad \mu=m_t/M_b
\]
决定。调至 $\omega_t=\omega_0$ 时，$\omega_\pm^2/\omega_0^2=1+\mu/2\pm\sqrt{\mu+\mu^2/4}$。小质量比给出频差 $\omega_+-\omega_-\simeq\omega_0\sqrt\mu$。若一开始只有棒在动，则两正常模以不同频率积累相位，能量可在约 $\pi/(\omega_+-\omega_-)$ 后转至小振子。相同的机械能对应 $|y|/|x_{\rm initial}|\sim\sqrt{M_b/m_t}$。放大的是读出的位移，其代价是交流能量所需的时间。耦合太弱或频率失谐超过分裂宽度时，小振子来不及从棒接受足够的能量。

直接可读的是相对位移 $z=y-x$。由第二个方程可得 $\tilde z=\omega^2\tilde x/(\omega_t^2-\omega^2)$；同第一个联立，整个响应的极点仍只有上面两个。单看这一比值似乎在裸频率 $\omega_t$ 发散，实际该频率的 $\tilde x$ 同时趋零。对两个通道作校准时应处理完整的耦合系统。弹性端的运动还可转为电学量。例如以平行板间隙 $d-z$ 作电容，$C(z)\simeq C_0(1+z/d)$；恒偏压 $V$ 时，信号电流为 $I_{\rm sig}=VC_0\dot z/d$。同一偏压下电场力 $F=C'(z)V^2/2$ 也随电压扰动改变。增加 $V$ 同时增强信号及电路反作用，这限制了仅靠提高机电增益所能达到的精度。

\subsection{热涨落、反作用与两个正交分量}
对黏滞阻尼的单模，在经典温度 $T_m$ 下涨落耗散关系给单边热力谱 $S_F^{\rm th}=4k_{\rm B}T_m m\gamma$。位移谱为 $S_q=|\chi|^2 S_F^{\rm th}$。以 $\d f=\d\omega/(2\pi)$ 积分，近峰处用 Lorentz 曲线代替，便得 $\langle q^2\rangle=k_{\rm B}T_m/(m\omega_0^2)$。降低损耗会把热峰变窄并提高峰值，总均方位移由均分定理固定。要判断一段短脉冲能否引起可辨的增量，还应看观测前后振子的差。在 $\gamma\Delta t\ll1$ 时，慢变正交分量在 $\Delta t$ 内的热扩散只占平衡方差的 $O(\gamma\Delta t)$，连续读出并不需要每次从未知的完整热能重新开始。

若位移读出另有独立的测量噪声 $S_q^{\rm imp}$ 和反作用力谱 $S_F^{\rm ba}$，将输出全折算到输入的棒轴应变，可得
\[
 S_{h,\parallel}^{\rm eq}=\frac{4}{m^2\ell^2\omega^4}
 \left[S_F^{\rm th}+S_F^{\rm ba}
             +\frac{S_q^{\rm imp}}{|\chi|^2}\right].
\]
读出噪声在共振处折算得最小。由同一机械力通道驱动的信号和热噪声则被同一 $\chi$ 放大，二者的比值并无共振增益。实际可用的频宽可大于本征线宽，边界约由折算后的读出噪声重新高于力噪声的位置决定。给一次爆发探测报「有效温度」$T_{\rm eff}$，通常把某一模板、滤波和门槛下的最小能量写成 $k_{\rm B}T_{\rm eff}$，不等于棒本身热力学温度。

用量子力学看这一限制。自由模可写 $q(t)=X\cos\omega_0 t+Y\sin\omega_0t$，其中 $X=q(0)$、$Y=p(0)/(m\omega_0)$，对易子为 $[X,Y]=\i\hbar/(m\omega_0)$。若对任意未知相位都要求同一精度，两个正交分量的噪声便不能同时任意小。只在 $t_k=2\pi k/\omega_0$ 测量自由振子，所得全为同一个 $X$，两次自由演化位置的对易子在这些时刻为零。测量踢给 $Y$ 的扰动在理想的下一次采样不进入 $X$。这就是一种避开反作用的量子非破坏测量。若信号与所选分量正交，或振子在两次读出间耗散、失谐，该方案也会损失信号或重新混入噪声。

\subsection{球的五个四极模}
对均匀、各向同性且自由的弹性球，潮汐体力 $f_i=\rho\ddot h_{ij}x^j/2$ 是标量势 $\rho\ddot h_{ij}x^ix^j/4$ 的梯度。横向张量波无迹，故 $h_{ij}x^ix^j$ 在球面上是纯 $\ell=2$ 的函数。模形可分为径向与切向梯度组成的球状模，以及切向旋转型的环状模。写任意模为 $\bm U_\alpha$，其驱动正比于 $\ddot h_{ij}\int\rho U_\alpha^ix^j\d^3x$。球谐正交性只留下球状四极模。对给定径向阶它有五重简并，每个读数给一对称无迹张量的一个分量；纯横向张量波仍只有两偏振。一台球形仪器由此可以同时估来向和偏振。

例如取五个正交的实无迹基 $E_m^{ij}$，从已校准的模振幅读出 $H_m=E_m^{ij}h_{ij}$，重构 $H_{ij}=\sum_m H_mE_{m,ij}$。波在同一时间沿 $\bm N$ 传播，必有 $H_{ij}N^j=0$，另两个特征值为 $\pm\sqrt{\operatorname{tr}(H^2)/2}$。零特征值的特征矢给出传播轴，但 $\bm N$ 和 $-\bm N$ 的差别须由别的时差或外部源位置来定。若只测到一个特定瞬间且波形恰过零，则三根轴都不能确定。将一段数据的 $H^{\mathsf T}H$ 相加后再取最小特征方向，可用多个时刻的约束；这样的检验以该时段内只有一个主导平面波为前提。

球面第 $a$ 处的径向传感器，方向 $\bm n_a$，对模 $m$ 的几何响应正比于 $B_{am}=n_a^i n_a^j E_{m,ij}$。若矩阵 $B$ 满列秩，五个模可由加权逆 $(B^{\mathsf T}W B)^{-1}B^{\mathsf T}W$ 重建。六个适当对称放置的传感器可使五个通道在理想情形等灵敏，另外一个通道不含四极信号。例如选择截角二十面体上六个适当五边形中心所定的轴，即 TIGA 布置。真实球面上换能器本身有质量，打破简并并产生频率分裂，几何矩阵须和耦合振子的频率响应一同标定。

\subsection{把换能器和放大器匹配起来}
两正常模的分裂本身并不等于已得到相应的灵敏带宽。可把传感器的热运动与棒的热运动分别折算为输入力，看看带宽还受什么限制。仍取棒、传感器的绝对坐标为 $x,y$，读数为 $z=y-x$。在角频率 $\omega$ 写两弹簧的复刚度为 $K_b=k_b+\i\omega b_b$、$K_t=k_t+\i\omega b_t$，$b_b,b_t$ 为黏滞系数。令 $F_b$ 包含潮汐力和基模的热力，$F_t$ 为小振子弹簧内部的热力，并约定后者对小振子用正号，对棒用负号。Newton 方程变为
\[
 (K_b-M_b\omega^2)x-K_t z=F_b-F_t,\qquad
 -m_t\omega^2x+(K_t-m_t\omega^2)z=F_t .
\]
消去 $x$，记行列式为 $\mathcal D=(K_b-M_b\omega^2)(K_t-m_t\omega^2)-K_t m_t\omega^2$，则
\[
 z=\frac{m_t\omega^2 F_b+[K_b-(M_b+m_t)\omega^2]F_t}{\mathcal D}.
\]
这个结果比直接把两个热谱乘上相同放大率多了一步。内弹簧的随机力总和为零，但恰好作用在差模；两质量共同的位移则不被 $z$ 测到。若把同一个热力加到两个方程的右边而取同号，就会引入虚假的质心热激发。

温度同为 $T_m$ 时，两个独立黏滞元件各有 $4k_{\rm B}T_m b_i$ 的单边力谱。以输入的 $F_b$ 为基准，传感器附加的等效力谱为 $4k_{\rm B}T_m b_t|K_b-(M_b+m_t)\omega^2|^2/(m_t^2\omega^4)$。若近裸频率且调节恰当，分子小，小振子的热扰动折算到输入端受到抑制；稍远离该频率，小质量 $m_t$ 又会放大这项。读出系统若有独立的位移噪声 $S_z$，还须加 $|\mathcal D|^2 S_z/(m_t^2\omega^4)$。因此峰间的噪声可能高于两峰，完整可用频带须由全曲线与波形相乘后积分得到。取 $m_t\to0$ 虽可使理想位移变大，热噪声和读出噪声并未同样缩小。

电路可用类似的两个噪声源表示。以某端口的速度 $v$ 为输入、外力 $F$ 为其共轭量，机械阻抗为 $Z=F/v$。输入端等效的速度测量噪声为 $S_v$，仪器反推到端口的随机力为 $S_F$；先假定二者不相关。若端口为实阻抗 $R$，则折算到外力的附加谱为 $S_F+R^2S_v$。按热噪声 $4k_{\rm B}TR$ 定义放大器噪声温度，得
\[
 T_n(R)=\frac{1}{4k_{\rm B}}\left(\frac{S_F}{R}+RS_v\right).
\]
当 $R=\sqrt{S_F/S_v}$ 时两项相等，最小值为 $\sqrt{S_FS_v}/(2k_{\rm B})$。若源阻抗还有虚部，信号及电路调谐可先补偿电抗，再谈实部的匹配。这个最优值固定的是信号源呈给放大器的阻抗，与把棒和传感器的裸频率调到相等是两种条件。

可以怎样改变呈给放大器的阻抗？一个无损变换若将 $v$ 放大 $\eta$ 倍，则必须将共轭力缩小为 $F/\eta$ 以保持平均功率，故阻抗变成 $Z/\eta^2$。谐振小质量相对于大质量的位移放大起类似作用。被动电容读出还可接调谐回路和低噪声电流放大器。若给电容加上交流泵浦，间隙变化调制泵的载频，在 $\omega_{\rm pump}\pm\omega_0$ 产生边带。外泵提供能量，变换不再受纯被动传感器的功率增益约束，泵浦相位涨落也会混进边带。选择泵、机械模和电学模时要同时检查反作用、稳定性及两个边带的噪声。

怎样把谱曲线归结为一次爆发的阈值？令短脉冲写成 $h_\parallel(t)=A p(t)$，取 $A$ 为待测振幅。棒的最优振幅方差满足 $\sigma_A^{-2}=4\int_0^\infty|\tilde p|^2/S_{h,\parallel}^{\rm eq}\,\d f$。若波形谱在窄机械带内平坦，可把它提到积分外，将有效带宽定义成
\[
 B_{\rm eff}=S_{\min}\int_{\rm band}\frac{\d f}{S_{h,\parallel}^{\rm eq}(f)},
 \qquad
 \sigma_A^2\simeq \frac{S_{\min}}{4B_{\rm eff}|\tilde p(f_0)|^2}.
\]
对于噪声水平相同、灵敏峰更宽的仪器，可辨到更弱的短脉冲；对于频率已知的稳态线，则要使用该线上实际的噪声水平及相干观测时长。把 $A=\sigma_A$ 代入前面的单模沉积能量，在此近似下得 $E_{\min}\simeq m\ell^2\omega_0^4S_{\min}/(32 B_{\rm eff})$。若阈值取 $\rho_*$，最小能量再乘 $\rho_*^2$。报出一个有效温度同时隐含了这个阈值、方向、波形及有效带宽。

例如在 $\omega_0/(2\pi)=1\,\mathrm{kHz}$、$Q=10^6$ 的理想单模中，振幅自由衰减时间约 $318\,\mathrm{s}$，本征功率全宽仅 $10^{-3}\,\mathrm{Hz}$。若换能器质量比 $m_t/M_b=10^{-4}$，理想分裂却约 $10\,\mathrm{Hz}$，能量第一次交换需约 $0.05\,\mathrm{s}$。三个数字刻画不同的时间和带宽。长振铃使到达后的取样从容；到达时刻能定多准，还取决于整个被读出频段的相位与信噪比。

\section{以光行时读出应变}\label{sec:gw-b-optic}
\subsection{任意来向的双程传递函数}
第三章已算出垂直入射时一次往返中的抵消。现设平面波沿 $\bm N$ 传播，分光镜在原点，一端镜在 $L\bm u$。取 $T_L=L/c$、$\mu=\bm N\cdot\bm u$。以光子回到分光镜的时刻 $t$ 标记一趟往返。在 TT 坐标中，自由镜的空间坐标恒定，把零线元沿未扰动的两段光路积分至一阶就够了。去程中，距分光镜 $s$ 处的事件发生于 $t-2T_L+s/c$，波的相位时刻是 $t-2T_L+(1-\mu)s/c$。回程中，从端镜走过 $s$ 处的坐标为 $(L-s)\bm u$，相位时刻为 $t-(1+\mu)T_L+(1+\mu)s/c$。故
\begin{align*}
 \delta T_{\bm u}(t)=\frac{u^iu^j}{2c}\int_0^L\d s\,\big[
 &h_{ij}(t-2T_L+(1-\mu)s/c)\\
 +{}&h_{ij}(t-(1+\mu)T_L+(1+\mu)s/c)\big].
\end{align*}
\begin{figure}[htbp]
 \centering
 \begin{tikzpicture}[x=1.6cm,y=0.64cm,>=Stealth]
  \draw[->] (-.05,0)--(3.5,0) node[right] {$s$};
  \draw[->] (0,-.3)--(0,6.9) node[above] {$ct'$};
  \draw[densely dashed] (3,0)--(3,6.4);
  \draw[thick,->] (0,.35)--(1.6,1.95);
  \draw[thick] (1.6,1.95)--(3,3.35);
  \draw[thick,->] (3,3.35)--(1.45,4.90);
  \draw[thick] (1.45,4.90)--(0,6.35);
  \fill (0,.35) circle (1.2pt) node[left] {$c(t-2T_L)$};
  \fill (3,3.35) circle (1.2pt) node[right] {反射};
  \fill (0,6.35) circle (1.2pt) node[left] {$ct$};
  \node[below] at (0,0) {$0$};
  \node[below] at (3,0) {$L$};
 \end{tikzpicture}
 \caption{以接收时刻 $t$ 标记一趟往返。图中纵坐标为事件的坐标时乘 $c$。波在各事件处的相位还须减去 $\omega\bm N\cdot\bm x/c$，所以光走过等长的两段时所跨的波相位一般不同。}
 \label{fig:gw-b-path}
\end{figure}

图~\ref{fig:gw-b-path} 给出去、回两段的取时。这里两段的卷积宽度不同。记 $x=\omega T_L$，令 $\delta\tilde T_{\bm u}=T_Lu^iu^j{\cal T}(x,\mu)\tilde h_{ij}$。将谐波在每段上积分，得
\begin{align}
 {\cal T}(x,\mu)=\tfrac12\big\{&
 \e^{-\i x(3+\mu)/2}\operatorname{sinc}[x(1-\mu)/2]\nonumber\\
 &+\e^{-\i x(1+\mu)/2}\operatorname{sinc}[x(1+\mu)/2]\big\}.
 \label{eq:gw-b-arm}
\end{align}
特别地 ${\cal T}(x,0)=\e^{-\i x}\operatorname{sinc}x$。随频率增加，第一个零点位于 $f=c/(2L)$，即一个往返时长等于波周期。另一个检验是 $\mu\to\pm1$，表面上换元的分母趋零，但 sinc 的极限存在，同时 $u^iu^jh_{ij}$ 由横性为零。以两个近似 $h_{ij}(t)$ 相减再除以很小的 $1-\mu$，数值上容易造成假发散。

两臂在相同接收事件处比较相位，激光角频率为 $\omega_{\rm L}$，读数为 $\omega_{\rm L}(\delta T_{\bm u}-\delta T_{\bm v})$（差分端口符号可互换）。换算成长波下校准的应变，即除以 $2\omega_{\rm L}T_L$。故完整的响应张量为 $D^{ij}(f)=[u^iu^j{\cal T}(x,\bm N\cdot\bm u)-v^iv^j{\cal T}(x,\bm N\cdot\bm v)]/2$。当 $x\ll1$ 才退化为第三章的实常数张量。对长基线高频测量，波形的相位和方向响应都随频率改变。在传播臂不等时，即使源很窄带，公共激光的频率噪声也会以两条延迟的差出现在测量端。

\subsection{局域自由落体系的校验}
光行时是分光镜上一只钟读出的量，应当与所选坐标无关。以垂直入射、$x$ 臂取 $h_{xx}=h$ 为例，在分光镜的局域自由落体系中，若只保留首个曲率项，$g_{00}=-1-R_{0x0x}X^2$，臂线上 $g_{0x}$ 和 $g_{xx}-1$ 的同阶项因 Riemann 的反称性为零。$R_{0x0x}=-\ddot h/(2c^2)$。端镜坐标移动约 $Lh(t)/2$；反射发生在 $t-T_L$，因此移动端点带来的双程时间约为 $T_L h(t-T_L)$。零线元还给光的坐标时增量 $\d t\simeq[1-R_{0x0x}X^2/2]\d X/c$（回程取路径正长度）。两段合计，在这一阶可把曲率取在中点时刻，额外贡献为 $T_L^3\ddot h(t-T_L)/6$。

同 TT 系的精确结果比较，垂直入射时任一波形均有 $\delta T=\frac12\int_{t-2T_L}^{t}h(t')\d t'$。绕中点 Taylor 展开，奇次项互相抵消，得 $\delta T=T_L h(t-T_L)+T_L^3\ddot h(t-T_L)/6+\cdots$。两种作法给出的时间一致。只追踪端镜而令光速严格等于其平直坐标值会漏掉后一项。反之，在 TT 坐标中若已沿光路积分，又给端镜重复添加同一物理臂长变化，会把信号算两次。

\subsection{光腔的自由光谱范围、线宽和反射相位}
两镜的振幅反射率取实数 $r_1,r_2$，相位约定使共振时往返因子 $a=r_1r_2$ 为正。透射振幅满足 $t_i^2=1-r_i^2$，暂不计额外损耗。令激光失谐引起的往返相位为 $\delta=2\omega_{\rm L}T_L-2\pi m$，第 $n$ 次往返带因子 $(a\e^{\i\delta})^n$。于是腔内行波振幅相对于入射振幅为 $t_1/(1-a\e^{\i\delta})$，腔的功率透射率为
\[
 {\cal P}_{\rm tr}=\frac{t_1^2t_2^2}{(1-a)^2+4a\sin^2(\delta/2)} .
\]
相邻共振频率间隔为 $\nu_{\rm FSR}=1/(2T_L)$。若 $1-a\ll1$，半高点满足 $|\delta|\simeq(1-a)/\sqrt a$，精细度即自由光谱范围除以功率全宽，为 $\mathcal F_{\rm cav}\simeq\pi\sqrt a/(1-a)$。停掉入射光后，每一往返把腔内场乘以 $a$，故振幅衰减时间为 $\tau_E=-2T_L/\ln a$，能量的衰减时间为 $\tau_E/2$。共振功率曲线在频率上的半高半宽约为 $1/(2\pi\tau_E)$。使用「储光时间」时先区分场和能量，否则会在极点频率中差一个二。

将从入射镜立即反射的场也计入，选端口相位使总反射振幅为 $r_{\rm cav}=(r_1-r_2\e^{\i\delta})/(1-a\e^{\i\delta})$。若 $r_1=r_2$，在共振时中央载频的反射完全抵消，理想光腔全透射。若端镜全反射，$r_2=1$，则任意失谐均有 $|r_{\rm cav}|=1$。取对数求导，在 $\delta=0$ 得 $\partial_\delta\arg r_{\rm cav}=(1+r_1)/(1-r_1)$。此时共振附近臂长变化几乎全转成反射相位。把激光频率乘以静态有效光程得到很大的相移，还须检验引力波在储光期间是否已反号。

对每次往返注入的微小相位 $\delta\phi(t)$，共振的稳态载频给各次相位边带以几何权重。去掉直流增益后，反射边带随调制频率的因子为
\[
 H_{\rm cav}(\omega)=(1-a)\sum_{j=0}^\infty a^j\e^{-2\i j\omega T_L}
 =\frac{1-a}{1-a\e^{-2\i\omega T_L}}
 \simeq\frac{1}{1+\i\omega\,2aT_L/(1-a)} .
\]
最后一步要求 $\omega T_L\ll1$；高精细度时其中的时间趋近于 $\tau_E$。信号先经式 \eqref{eq:gw-b-arm} 的单次光路平均，再经储光平均。这两个平均各有自己的时间尺度。把精细度提高一倍，低频增益近乎加倍，通带宽度约减半。接近自由光谱范围时应使用几何级数的精确形式，不能把单极点延拓到任意高频。

\subsection{有限光束与腔的几何稳定性}
先看衍射。将电场写成 $u(x,y,z)\exp[\i(kz-\omega_{\rm L}t)]$，若纵向包络变化缓慢，则 Helmholtz 方程化为 $2\i k\partial_z u+\nabla_\perp^2u=0$。取轴对称 Gaussian 解，其腰斑处强度取 $I(r,0)\propto\e^{-2r^2/w_0^2}$。离开腰斑后斑半径、波前曲率半径及额外的 Gouy 相位分别为
\[
 w^2(z)=w_0^2[1+(z/z_R)^2],\quad
 R(z)=z[1+(z_R/z)^2],\quad
 \psi(z)=\arctan(z/z_R),\qquad z_R=\pi w_0^2/\lambda_{\rm L}.
\]
例如以半径 $b$ 的圆镜截束，落在镜外的功率分数为 $\int_b^\infty I\,r\d r/\int_0^\infty I\,r\d r=\exp[-2b^2/w^2(z)]$。从而加大光斑虽可平均更多镜面热涨落，却增加衍射损耗。只由射线光学看镜面中心一条好轨迹，无法估算光腔的精细度。

令 $q$ 为复光束参量，取 $1/q=1/R-\i\lambda_{\rm L}/(\pi w^2)$。在一阶光学系统的光线矩阵 $\bigl(\begin{smallmatrix}A&B\\C&D\end{smallmatrix}\bigr)$ 下，它变为 $(Aq+B)/(Cq+D)$。自由传播长 $L$ 给矩阵 $\bigl(\begin{smallmatrix}1&L\\0&1\end{smallmatrix}\bigr)$，曲率半径 $R_i$ 的镜反射给 $\bigl(\begin{smallmatrix}1&0\\-2/R_i&1\end{smallmatrix}\bigr)$。将两传播、两反射矩阵相乘可得往返矩阵行列式为一、半迹为 $2g_1g_2-1$，其中 $g_i=1-L/R_i$。腔模须在一往返后恢复，故 $Cq^2+(D-A)q-B=0$。其判别式为 $(A+D)^2-4$；存在正虚部的不动点且小扰动有界，通常选 $0<g_1g_2<1$ 的非退化稳定区。边界需另查简并与对失准的敏感性。

取两镜等曲率半径 $R$，腰斑居中。要求 $R(L/2)=R$ 给 $z_R^2=(L/2)(R-L/2)$，故 $w_0^2=(\lambda_{\rm L}/\pi)\sqrt{(L/2)(R-L/2)}$。将此代回 $w(L/2)$ 即可按腔长和镜径估截光损失。入射光腰的位置和半径若与固有模不合，只会有部分功率耦合到基本模。例如两个共址、共轴且腰斑分别为 $w_1,w_2$ 的归一 Gaussian 场，重叠积分平方为 $4w_1^2w_2^2/(w_1^2+w_2^2)^2$，这也是只调总入射功率不能弥补的几何损失。

\subsection{如何在暗纹上取出线性信号}
两个臂返回的场若恰好等幅，则反对称端口场正比于 $E_0[\e^{\i\phi_x}-\e^{\i\phi_y}]/2$，输出功率为 $P_0\sin^2(\Delta\phi/2)$。暗纹处载频功率最小，可是功率对小相移的导数也为零。怎样既留在暗纹附近，又测得线性相位？一种办法是取小的固定失谐 $\phi_b$，使直流上有本地振荡场。写 $\Delta\phi=\phi_b+\delta\phi$，信号功率增量约为 $P_0\phi_b\delta\phi/2$，而暗端载波功率约为 $P_0\phi_b^2/4$。相干光的光子数涨落正比于功率平方根，故在载波主导的理想散粒噪声极限，两者的比值对小偏置近乎不变。偏置再变小至电子暗噪声或杂散光不可略时，这个结论不再成立。

另一办法给载频加射频相位调制。调制指数 $\beta\ll1$ 时，入射场 $\e^{\i\omega_{\rm L}t+\i\beta\sin\Omega_m t}$ 有载频及振幅分别为 $\beta/2,-\beta/2$ 的两边带。令 $r_0,r_+,r_-$ 为三频率在光腔的复反射率。光电流中正比于 $\exp(\i\Omega_m t)$ 的项为 $P_0\beta(r_+r_0^*-r_0r_-^*)/2$。乘同频本振并低通，在两个解调相位分别得到这复量的实部、虚部。若两边带远离腔共振、反射近似相同且为实，虚部随载频反射相位过零而变号，正好给出指示应向哪边调腔长的误差信号。只量反射功率，失谐的两侧反而可能一模一样。此即 Pound--Drever--Hall 锁定的基本计算。

将误差信号反馈到致动镜。用频域传递函数写 $e=C(\delta L-Au)$、$u=K e$，其中 $\delta L$ 是无反馈时的等效变化，$C$ 是光学及读出增益，$A$ 为致动器的米每控制单位。则 $e=C\delta L/(1+CAK)$，控制量为 $u=KC\delta L/(1+CAK)$。高环路增益使误差趋近零，但被测信息并未消失，重构应使用 $\delta L=e/C+Au$。环路的相位在增益超过一的频段若接近反号，反馈会转成激发。多腔仪器通常要分别约束臂长之和与差、中央 Michelson 长度差以及回收腔长度，各传感端口组成响应矩阵；只把某一误差通道除以静态增益，不能替代这一频域标定。

理想等臂时大部分载频返回激光端，在这端加一面部分透射镜可使载波再入仪器，构成功率回收腔。令臂与分光镜合起来对公共模的振幅反射率为 $r_{\rm com}$，回收镜为 $r_p,t_p$。共振时回收功率相对输入的增益为 $G_p=t_p^2/|1-r_p r_{\rm com}|^2$。若 $r_{\rm com}$ 实且其模小于一，令 $t_p^2=1-r_p^2$ 求极值得 $r_p=r_{\rm com}$，即输入耦合与等效损耗匹配。反对称端口的信号边带可再经一面镜子回收，形成差模耦合腔。改变这面镜子的透射率或失谐量，可交换探测带宽与窄带响应。公共端功率回收主要增大载波，差模端的信号回收主要改变边带的传播。两者须用各自的场相位判断共振条件。

\subsection{从光子计数到位移的量子极限}
设有效入射光子通量为 $\dot N=P/(\hbar\omega_{\rm L})$，输出理想、可见度为一。在一段 $\tau$ 内计数，暗口均值 $\bar n=\dot N\tau\sin^2(\Delta\phi/2)$。Poisson 涨落给 $\operatorname{var} n=\bar n$。小偏置时线性误差传播给 $\sigma_\phi\simeq1/\sqrt{\dot N\tau}$。若将这方差视为白相位噪声在长为 $\tau$ 的平均值上的方差，单边相位谱便为 $S_\phi=2/\dot N$。因矩形时间平均对白噪声的单边等效带宽为 $1/(2\tau)$，两式相符。折算到应变还要除以相位响应平方。臂腔低频提高响应，超过腔极点后，等效应变的振幅谱便随频率升高。在此计算中 $\dot N$ 是按指定读出模型定义的有效光通量，不能同时用臂内循环功率和已有腔增益重复计数。

同一束光还给镜子动量。以单个正入射完美反射面为例，每个光子冲量为 $2\hbar k_{\rm L}$，独立到达的冲量流单边力谱为 $S_F=2(2\hbar k_{\rm L})^2\dot N$。这随功率增大；相反，测量噪声随功率减小。实际差模中两臂的光子涨落与分光镜暗端进入的真空场有关，必须用同一个广义位移定义两项。可用理想无相关线性位置测量说明最优功率的尺度。记单边测量噪声和反作用噪声为 $S_x^{\rm imp},S_F^{\rm ba}$，取理想极限 $S_x^{\rm imp}S_F^{\rm ba}=\hbar^2$。对于机械响应 $\chi_m$，总位移噪声至少为
\[
 S_x^{\rm imp}+|\chi_m|^2S_F^{\rm ba}
 \ \ge\ 2\sqrt{S_x^{\rm imp}|\chi_m|^2 S_F^{\rm ba}}
 =2\hbar|\chi_m| .
\]
在悬挂共振以上的自由质量带，$|\chi_m|\simeq1/(m_{\rm eff}\omega^2)$。若四面镜子质量同为 $M$，差模 $x=(x_{\rm EX}-x_{\rm IX})-(x_{\rm EY}-x_{\rm IY})$，广义外力 $F$ 在四镜上分别施加 $\pm F$，得 $|\chi_m|=4/(M\omega^2)$。长波下 $x=Lh$，于是单边自由质量标准量子极限为 $S_h^{\rm SQL}=8\hbar/(M\omega^2L^2)$。若只允许端镜运动，或用半个差模作读数，须重新计算 $m_{\rm eff}$ 与 $x/h$。同一表观因子四可分别来自不一致的坐标或谱定义。

对某一固定功率，等效噪声形如 $a(f)/P+b(f)P$。求极值得 $P_{\rm opt}(f)=\sqrt{a/b}$，而随频率变化的最小值才是上述包络。一台固定功率的装置通常只在窄频段触及它。若读出的光场正交分量与辐射压所来自的分量有相关，总谱还含 $2\operatorname{Re}[\chi_m S_{xF}]$（互谱取与此式相应的次序）。适当选同相读出角可令这项在一段频带部分抵消反作用；压缩入射真空也可在两个光学正交分量间重分配方差。有光学损耗时未测的真空重新进入，不能把理想相关的抵消直接当作真实仪器的性能。

\subsection{哪类噪声随镜子运动？}
设悬摆支点有位移 $y$、镜的绝对位移为 $x$。取回复和黏滞力均依赖于相对坐标，则 $m\ddot x=-m\omega_p^2(x-y)-m\gamma_p(\dot x-\dot y)+F$。从支点位移到镜位移的传递函数为
\[
 \frac{\tilde x}{\tilde y}=
 \frac{\omega_p^2+\i\gamma_p\omega}{\omega_p^2-\omega^2+\i\gamma_p\omega}.
\]
在 $\omega\gg\omega_p$ 且损耗项仍小于弹性项的频带，幅度约按 $(\omega_p/\omega)^2$ 下降。串接且适当隔开共振的多级悬摆可进一步衰减地面噪声。非常低频时整套镜子跟随地面，须用低频控制维持锁定。地面密度涨落所产生的时变 Newton 引力则直接驱动测试质量。若一块质量相对于镜子运动 $\delta r$，引力梯度的量级为 $GM/r^3$，在自由质量带的位移量级为 $GM|\delta r|/(r^3\omega^2)$。它不经过悬挂连接，在当地地震场中增加悬摆级数不能按相同传递率消除它。

热噪声可由受扰物块对外力的耗散求出。对读出坐标 $x$，取前面 Fourier 约定下正频率的耗散部为 $-\operatorname{Im}\chi$，经典涨落耗散关系为 $S_x(f)=4k_{\rm B}T[-\operatorname{Im}\chi(2\pi f)]/(2\pi f)$。若回复刚度为 $k[1+\i\phi(\omega)]$，损耗角 $\phi\ll1$，则
\[
 S_x(f)=\frac{4k_{\rm B}T}{\omega}
 \frac{k\phi(\omega)}{(k-m\omega^2)^2+k^2\phi^2(\omega)} .
\]
即使在本征共振之外也有热谱尾。黏滞阻尼对应 $k\phi=m\gamma\omega$，若 $\phi$ 近乎为常数，远离峰时的频率依赖便不同。悬丝内耗、反射镀膜及基底的弹性内耗、热膨胀与热扩散之间的耦合都可贡献镜面坐标的等效涨落。读出坐标实际上是光斑强度加权的表面位移。按同一空间权重对镜面施加周期试探压力、计算一周期内的耗散，就能应用涨落耗散关系，不能只把镜子当单一质点。

两臂长度若差 $\Delta L$，激光频率噪声也可能伪装成镜子运动。静态差相位为 $4\pi\nu_{\rm L}\Delta L/c$，其变化等价于差臂位移 $\delta x_{\rm eq}=\Delta L\,\delta\nu_{\rm L}/\nu_{\rm L}$。这说明光源稳频与臂长匹配须同时考虑。腔极点不等时，即使几何长度相等也会有频率相关的公共模泄漏。束斑偏离镜心时，角抖动还可转成几何路程扰动；失准损耗则降低与主模的重叠。给出一条应变噪声曲线时，要用每类噪声进入的实际端口分别折算，独立项加功率，相关项另加互谱。

由这些频率依赖可辨认地面干涉仪灵敏带的三段。低频端要控制地震、悬挂与辐射压，中频常由热噪声等限制，高频信号边带受储光时间抑制而散粒噪声占优。提高循环功率可以改善高频灵敏度，同时会增强低频的辐射压和镜体吸热。选择更重的镜子减小自由质量的力响应，扩大光斑降低局部表面涨落，信号回收调节对目标频带的响应。这些取舍说明地面各代仪器为何在臂长相近时仍能有不同的探测距离。扩大基线到空间还改变了式 \eqref{eq:gw-b-arm} 中的 $x$，不能只按 $1/L$ 缩放地面灵敏度。

\subsection{空间基线上的延迟相减}
空间仪器中的两颗航天器不能像实验室反射镜那样把足量入射光直接送回。可用远端激光同收到的光相锁，再发出相位相关的新束。先算一段从 $i$ 到 $j$ 的单程链路。设方向为 $\bm u_{ij}$、时长为 $T_{ij}$，接收事件在 $(t,\bm x_j)$，令 $u_r=t-\bm N\cdot\bm x_j/c$。把式 \eqref{eq:gw-b-arm} 的单程积分对接收时刻微分，可得光行时的相对变化率
\[
 \frac{\d\delta T_{ij}}{\d t}
 =\frac{u_{ij}^m u_{ij}^n}{2(1-\bm N\cdot\bm u_{ij})}
 [h_{mn}(u_r)-h_{mn}(u_r-(1-\bm N\cdot\bm u_{ij})T_{ij})] .
\]
到达频率的相对变化与此反号。式中传播差与几何分母必须一起取极限。低频时，差分给大约一个光行时乘 $\dot h$，所以 Doppler 读数随 $fT_{ij}$ 减小；把它再积分成相位，便还原有限的长度响应。频率和相位是同一信息的两种读数，比较两种装置时须在同一种读数下折算噪声。

考虑三站的简化模型，各站激光噪声只由一个相位过程 $p_i(t)$ 表示。$i$ 发 $j$ 收的相位读数为 $y_{ij}=D_{ij}p_i-p_j+s_{ij}$，其中 $D_{ij}a(t)=a(t-T_{ij})$，$s_{ij}$ 包括波和其它噪声。先在中心站 1 合成两次往返，令 $U=y_{21}+D_{21}y_{12}$、$V=y_{31}+D_{31}y_{13}$。来自远端站 2、3 的激光噪声分别精确相消，剩下
\[
 U=(A-1)p_1+u_s,\qquad V=(B-1)p_1+v_s,\qquad
 A=D_{21}D_{12},\quad B=D_{31}D_{13}.
\]
若两臂等长，则 $U-V$ 不含 $p_1$。若两臂不等且长度恒定，构造 $X=(1-B)U-(1-A)V$。代入上式，中心激光的残项为 $(AB-BA)p_1=0$，因为此时纯时间平移可交换。引力波沿不同方向送入各 $s_{ij}$，一般不会同时消失。此为第三章所述时间延迟干涉的一个显式构造。相位噪声被消去所凭的是推迟后的发射事件重合；以同两条臂的瞬时长度作一次普通差分并不能消去它。

长度随时间改变时，两算符为何不再可交换？将 $Aa(t)=a[t-T_A(t)]$ 等代入，$AB$ 的有效延迟为 $T_A(t)+T_B(t-T_A(t))$，近似 $T_A+T_B-T_A\dot T_B$。交换顺序则末项变成 $-T_B\dot T_A$。对一阶缓变臂，前述消除留下的激光噪声正比于 $(T_A\dot T_B-T_B\dot T_A)\dot p_1$。可将路径加长。例如先由 $U,V$ 获得 $U+AV$ 与 $V+BU$，其激光部分为 $(AB-1)p_1$ 与 $(BA-1)p_1$。再依前式比较这两条合成路径，余下 $[AB,BA]p_1$。$AB$ 和 $BA$ 都各含一次 $A$ 和一次 $B$，故它们到零阶的总光行时及其一阶导数相同，此交换子的缓变一阶项消失。精确实现须分别处理各单程的光行时及其变化。参见 \href{https://arxiv.org/abs/gr-qc/0310017}{Tinto、Estabrook 与 Armstrong}。

即使臂长真是常数，所用延迟也得校准。若误差为 $\delta T$，则 $p(t-T-\delta T)-p(t-T)\simeq-\delta T\dot p(t-T)$。对频率 $f$ 的相位噪声，残余功率约按 $(2\pi f\delta T)^2$ 乘原噪声。测距和授时的误差因而会直接影响激光噪声消除。在已去掉激光噪声的低频读数中，测试质量的非引力残余加速度仍经 $1/(2\pi f)^2$ 转成位移，高频则有光链路散粒及读出噪声。若测试质量的单边加速度谱为 $S_a$，位移测量谱为 $S_x$，两者的交界量级由 $S_a/(2\pi f)^4=S_x$ 定出。实际的方向平均、传递函数和台间相关须再由各链路组合给出。

\section{恒星怎样激发引力波}\label{sec:gw-cstars}

在球对称塌缩中，即使一颗恒星在一个自由落体时间内释放了相当于太阳静能十分之一的结合能，外部真空仍可一直取 Schwarzschild 度规。观者若只收引力波，看见的是什么？是质量分布偏离球对称的那一小部分如何随时间改变。由此出发，爆炸的总能量、辐射的时间尺度和波的偏振便成了三个不同的问题。

\subsection{铁芯失稳及回弹}

天文学家依光谱区分超新星。有氢线的是 II 型，无氢线的是 I 型。后一类中 Ia 型有显著的硅吸收线，通常源于白矮星热核燃烧；Ib、Ic 型则可由失去氢包层、甚至氦包层的大质量恒星核塌缩而来。光谱主要看外层，决定引力波的却是内核。两颗外观分属不同类型的星，内核可以经历相近的回弹。反过来，中心形成黑洞而包层未成功抛出时，光学爆发很弱，中心仍可能发波。

为何铁芯会突然失去支撑？以质量 $M$、半径 $R$ 做一次同调压缩。若压强 $p=K\rho^\gamma$，内能随 $R^{-3(\gamma-1)}$ 变，自引力能随 $-R^{-1}$ 变。记 $q=3(\gamma-1)$，势能为 $U=A R^{-q}-B R^{-1}$。在平衡处 $q A R^{-q}=B/R$，故 $R^2 U''=(q-1)B/R$。$\gamma=4/3$ 正好在稳定性分界上。极端相对论简并电子给 $p\propto(Y_e\rho)^{4/3}$，其中 $Y_e$ 是每个重子的电子数。用 $p/R\sim GM\rho/R^2$ 消去 $R$，得可支撑质量 $M\propto Y_e^2$。电子俘获一面减小 $Y_e$，一面以逃逸中微子带走能量；高能光子分解重核又消耗热能。故实际芯核可在质量几乎未变时越过稳定界。

内芯声波来得及往返，速度近似正比于半径；外芯落速超过声速，尚不知中心已发生了什么。密度接近核物质密度时物态方程变硬，内芯制动，压缩波在音速面附近聚成激波。若只看反弹动能，容易高估爆炸。把 $0.1M_\odot$ 的铁分解到核子，每个核子约需 $8$\,MeV，共约 $1.5\times10^{51}$\,erg，已与一次普通爆炸的抛射动能相当。激波外行时还受持续吸积压迫，可能停滞于原中子星以外。中微子从更热的深处扩散出来，在激波后的一层物质中被吸收。对这一层，设比结合能为 $|e_{\rm b}|$、净加热率为 $\dot q_\nu$、厚度为 $\Delta r$，则要比较的是 $t_{\rm heat}=|e_{\rm b}|/\dot q_\nu$ 与穿流时间 $t_{\rm adv}\sim\Delta r/|v_r|$。若前者较短，有机会把持续向内的流改成向外的流。非径向对流可以延长某些质元的停留时间，故一维计算会遗漏影响爆炸成败的运动。

取原中子星 $M=1.4M_\odot$、$R=12$\,km，$GM^2/R$ 约 $4\times10^{53}$\,erg；实际结合能还要乘结构因子。中微子在致密物质内的平均自由程若为 $\lambda_\nu\ll R$，扩散时间量级为 $R^2/(c\lambda_\nu)$，比自由传播时间大 $R/\lambda_\nu$ 倍。这就容许毫秒动力学与持续数秒的冷却共存。留在中心的质量还可因物质回落增加。是否最终成黑洞，需把遗留质量与相应温度、转速、物态方程下的最大支撑质量比较，不能只用主序星质量判别。前身星的金属丰度通过星风改变它在塌缩前保留下的质量及角动量；它影响的是初值。

内外芯的分界怎样随状态改变？塌缩初期若电子近于极端相对论，内芯质量可用上述 $M\propto Y_e^2$ 的质量尺度估计，$Y_e$ 从 $0.46$ 减至 $0.32$ 可使该尺度缩小到原值的约一半。密度高到中微子的扩散时间超过落体时间后，总轻子数近似随流携带，局部产生的中微子可参与压强与弱平衡。此后只用自由逸出的电子俘获率估落体加速度会过大。若用 $p=K\rho^\gamma$ 作简化的演化物态，核密度以下选略小于 $4/3$ 的 $\gamma_1$ 触发不稳，以上选较大的 $\gamma_2$ 模拟核力。为使冷压强在 $\rho_n$ 连续，须有 $K_2=K_1\rho_n^{\gamma_1-\gamma_2}$；冷比内能由 $\d e_{\rm cold}=p_{\rm cold}\d\rho/\rho^2$ 积出，第二支还须加积分常数使内能连续。只改幂指数不匹配积分常数，会在穿过核密度时人为注入能量。激波带来的热压则可另取 $p_{\rm th}=(\gamma_{\rm th}-1)\rho(e-e_{\rm cold})$。这样的模型能隔开刚度与转动的影响，细致的中微子输出仍须有温度、成分和输运。

在激波本身可用局部静止激波系比较上下游。记上游自由落入的速率为 $v_1>0$、下游为 $v_2$，忽略跨越很薄的激波时的引力势差。质量、法向动量和比总焓的跳跃条件是
\begin{align*}
 j&=\rho_1v_1=\rho_2v_2,\qquad
 p_2-p_1=j(v_1-v_2),\\
 \tfrac12v_1^2+\frac{\gamma}{\gamma-1}\frac{p_1}{\rho_1}
 &=\tfrac12v_2^2+\frac{\gamma}{\gamma-1}\frac{p_2}{\rho_2}
   +e_{\rm diss}.
\end{align*}
$e_{\rm diss}$ 为单位质量耗于分解重核的能量。强激波取 $p_1\ll\rho_1v_1^2$、$\beta_s=\rho_2/\rho_1$，代入得到
\[
 \left(1-\frac{2e_{\rm diss}}{v_1^2}\right)\beta_s^2
 -\frac{2\gamma}{\gamma-1}\beta_s+\frac{\gamma+1}{\gamma-1}=0 .
\]
没有分解时，一根 $\beta_s=1$ 是无激波状态，另一根 $(\gamma+1)/(\gamma-1)$ 是有压缩的激波。对近相对论气体 $\gamma=4/3$，后者为七。损耗正且小于入射比动能时，物理压缩比还会增加，下游速度进一步降低。同一吸积率 $\dot M=4\pi r_s^2\rho_1 v_1$ 给冲压 $\rho_1v_1^2=\dot M v_1/(4\pi r_s^2)$。若取 $v_1\simeq(2GM/r_s)^{1/2}$，此项随 $r_s^{-5/2}$ 增长，已开始后退的激波会面临更强冲压。要判断是否复苏，得把热输入、核分解和对流一同纳入下游能量，而不能只比较中心总中微子能量与外层结合能。

有个观测上的区别也可用时间尺度说明。若扩散面近似热球，给定味的中微子光度量级为辐射常数乘 $R_\nu^2 T_\nu^4$，通量的总量只限制这个组合。弱相互作用的吸收截面在相关能区大致随中微子能量平方变，因此相同光度的硬谱一般有较高加热效率。取另一极端，若回落物质以 $\dot M$ 沉积到半径 $R$ 的致密星表面，每秒可释放 $GM\dot M/R$ 的引力能。光度曲线晚期的一段增强可能包含新的吸积；只有联用中微子平均能量、持续时间和核心振荡频率，才可区分冷却面收缩与外界补给。

\subsubsection{光学样本怎样标定这类源的数目}

一次塌缩在光学上能有多亮，与内芯有多不对称是两种测量。若抛射层中合成 $N_0$ 个 $^{56}\mathrm{Ni}$，经 $^{56}\mathrm{Co}$ 到 $^{56}\mathrm{Fe}$ 的两步衰变可使光变在爆炸后持续很久。取两个衰变常数 $\lambda_1,\lambda_2$，初始无子核，解 $\dot N_1=-\lambda_1N_1$、$\dot N_2=\lambda_1N_1-\lambda_2N_2$ 得
\[
 N_1=N_0e^{-\lambda_1t},\qquad
 N_2=N_0\frac{\lambda_1}{\lambda_1-\lambda_2}
             (e^{-\lambda_2t}-e^{-\lambda_1t}).
\]
若每次衰变可沉积的电磁能为 $Q_1,Q_2$，放射加热率为 $Q_1\lambda_1N_1+Q_2\lambda_2N_2$ 再乘各自沉积比例。在短寿命的母核基本耗尽而抛射层对 $\gamma$ 射线仍厚时，晚期光度可追随子核的指数尾，估合成镍的质量。若包层重且富氢，氢复合使光球退入较深的质量层，光变又可出现近平台。尽管这几种光变能区分外层结构，它们对毫秒内芯的方位分布只给间接约束。出现一次亮超新星并不蕴含有某个固定能量进入四极波。

若从星系的恒星形成率密度 $\dot\rho_\star(t)$ 估核塌缩次数，需先说明它按形成的总质量计，不能直接除以一颗典型中子星的质量。设初始质量函数按星数为 $\xi(m)\,\d m$，归一任意，前身寿命为 $\tau_\star(m)$。最简估计为
\[
 R_{\rm cc}(t)=
  \frac{\int_{m_{\rm min}^{\rm cc}}^{m_{\rm max}^{\rm cc}}
   \dot\rho_\star[t-\tau_\star(m)]\,\xi(m)\,\d m}
       {\int_{m_{\rm min}}^{m_{\rm max}}m\,\xi(m)\,\d m}.
\]
分子只数能塌缩的高质量星，分母却含所有出生质量。若形成率在几十兆年内几乎不变，寿命延迟可略。举例取 $\xi\propto m^{-2.35}$，全范围为 $0.1$ 至 $100M_\odot$，塌缩区间暂取 $8$ 至 $100M_\odot$。以太阳质量为单位，比例为 $[(8^{-1.35}-100^{-1.35})/1.35]/[(0.1^{-0.35}-100^{-0.35})/0.35]\simeq7.4\times10^{-3}$，即形成约 $135M_\odot$ 的恒星伴随一次潜在的核塌缩。真正可见的爆炸率还要乘保留包层、成功推出激波、巡天能找到等概率。这些区间不能用来判定个别星最后是中子星还是黑洞。白矮星热核爆炸依双星演化，延迟时间可长得多，故由老年恒星族看到的 Ia 型率不能直接代入这一核塌缩估计。

对给定引力波源类的探测率，还须将 $R_{\rm cc}$ 与内部转动、距离和观测取向相折叠。以纯轴对称回弹为例，若同样的星在赤道方向可见到距离 $d_0$，则从任意 $\iota$ 可及 $d_0\sin^2\iota$。暂不计仪器天线图，在低红移均匀空间，取向平均的体积因子为 $\langle\sin^6\iota\rangle=(1/2)\int_{-1}^1(1-\mu^2)^3\,\d\mu=16/35$。这次平均按振幅的三次方，因为是问能看到多少颗；由赤道能流推总能量用的是二次方。对同一物理源这两个角平均不能交换。只有加上台天线、实际噪声和恒星空间分布，才能把它化成探测次数。

\subsection{由流体的一张快照求波}

模拟输出的量通常是 $\rho,\bm v,p$ 及引力势 $\Phi$。取 $\bm a=-\nabla\Phi-\nabla p/\rho$，并以尖括号表示一对指标的对称无迹部分。令 $I_{ij}=\int\rho x_{\langle i}x_{j\rangle}\,\d^3x$。由连续性方程，对随流质元的任意函数 $F$ 有 $\d(\int\rho F\,\d^3x)/\d t=\int\rho\,\d F/\d t\,\d^3x$。于是第一次求导得到 $2\int\rho x_{\langle i}v_{j\rangle}\,\d^3x$，再求一次为
\begin{align}
 \ddot I_{ij}
 &=2\int\!\left[\rho v_{\langle i}v_{j\rangle}
 -\rho x_{\langle i}\partial_{j\rangle}\Phi
 -x_{\langle i}\partial_{j\rangle}p\right]\d^3x \notag\\
 &=2\int\!\rho\left[v_{\langle i}v_{j\rangle}
                  -x_{\langle i}\partial_{j\rangle}\Phi\right]\d^3x .
 \label{eq:gw-cstress}
\end{align}
末步将压强项分部积分，$\partial_jx_i=\delta_{ij}$ 的无迹部分为零，外表面压强取零。若只截取一小块吸积气体，边界上的动量和压强通量均须保留。式 \eqref{eq:gw-cstress} 使每张快照都能给出一次波幅读数。二次时间差分会放大网格噪声，在暴烈流动中尤其不利。此法适于慢运动、弱场的波源估计；一旦时空自身参与强烈振荡，应从演化出的曲率提波。

以 $z$ 为对称轴，定义 $D=I_{zz}-I_{xx}=\int\rho r^2P_2(\cos\theta)\,\d^3x$。取视线与轴夹角 $\iota$，偏振基矢沿天球子午线和纬线。收缩四极波得到 $h_+=G\ddot D\sin^2\iota/(c^4 d)$、$h_\times=0$，$d$ 为距离。具体地，$e_+^{ij}I_{ij}=D\sin^2\iota$；故正沿对称轴看不到这一回弹信号。代 $I_{zz}=2D/3$、$I_{xx}=I_{yy}=-D/3$，有 $I^{(3)}_{ij}I^{(3)}_{ij}=2[D^{(3)}]^2/3$。相应总功率为 $2G[D^{(3)}]^2/(15c^5)$。处理单一方向波形时须先恢复这个角分布，不能把赤道方向的通量直接乘 $4\pi d^2$。

柱坐标 $(s,\phi,z)$ 中有 $D=\int\rho(z^2-s^2/2)\,\d^3x$。用式 \eqref{eq:gw-cstress} 先做 $\phi$ 积分得
\[
 \ddot D=\int\rho\big[2v_z^2-v_s^2-v_\phi^2
                   -2z\,\partial_z\Phi+s\,\partial_s\Phi\big]\,\d^3x .
\]
这还给一个算例。对均匀轴对称椭球，赤道和极向半轴各为 $a,b$，$D=M(b^2-a^2)/5$。若同调变化，则 $\ddot D=2M(\dot b^2+b\ddot b-\dot a^2-a\ddot a)/5$。两方向都在加速塌缩时，决定正负号的是哪一方向先制动。若 $a=b$ 始终相等，无论径向加速度多大，所有项恰好抵消。轴对称信号出现在回弹，并不要求恒星在两条赤道轴间也有不对称。

用动力学时间 $t_{\rm dyn}=(R^3/GM)^{1/2}$ 和形变 $\epsilon\sim|D|/(MR^2)$ 估算。若形变在一个 $t_{\rm dyn}$ 内改变，$h\sim \epsilon(GM/c^2d)(GM/Rc^2)$，功率约为 $(c^5/G)\epsilon^2(GM/Rc^2)^5$。两颗质量与形变相近的星，半径缩小一半会使这个功率估计增大 $32$ 倍。只有参与运动的质量才该算入 $D$。若取振荡分量的尺度太大，例如把整个外包层都算作与内芯同相回弹，质量因子和时间尺度就会相互矛盾。

\subsection{有哪些不同的非球对称运动？}

设转动动能为 $T$、引力能为 $W<0$。固定角动量 $J$ 做近似同调收缩时，$T=J^2/(2I)\propto R^{-2}$，故 $\beta=T/|W|\propto R^{-1}$。先随内芯收缩而自旋加快，随后核压强或离心力先后制动，产生的轴对称波形只持续几个动力学时间。若回弹主要靠离心支撑，发生制动的密度低，典型频率也较低。核物质压强、电子俘获和相对论引力会改变是否发生多次回弹，故从一条计算波形读出的频率并非物态方程的唯一指标。

另一种可能是内芯由近轴对称变作一根转棒。对均匀不可压 Maclaurin 椭球，沿平衡序列在 $\beta\simeq0.14$ 处有长期不稳定性，在 $\beta\simeq0.27$ 处有动力学不稳定性。前者须靠耗散逐渐离开原平衡，后者即使流体无粘也能在周转时间内生长。这两个数只用于该模型。差速转动流体中，模式图样角速度若与某层流体共转，该处能交换角动量，较低的 $\beta$ 也可出现非轴对称不稳定性。若用两团各为 $m$ 的物质、到轴的距离为 $b$ 作转棒估计，则 $I_{xx}-I_{yy}=2mb^2\cos 2\Omega t$、$2I_{xy}=2mb^2\sin2\Omega t$。观者可见持续许多周的窄频信号，峰频为图样转频的两倍。振荡寿命和参与质量比瞬时最大 $\beta$ 更直接地决定可积累多少信噪比。

回弹后的对流、吸积柱撞击原中子星和停滞激波的整体摆动，也改变四极矩。前两者有一串相位不相干的激发，不能把相邻波包按确定相位累加。激波的驻立吸积激波不稳定性可由一次回路理解：激波形变产生熵和涡量扰动，被向内平流到强制动区，那里激起向外传播的声波，再作用于激波。若一次回路的复放大率为 ${\cal Q}$、时间为 $\tau$，则 $\exp[(\sigma+\mathrm i\omega)\tau]={\cal Q}$，增长率 $\sigma=\ln|{\cal Q}|/\tau$、振荡频率由回路相位确定。改变停滞半径会改变 $\tau$；这与核心弹性和密度设置的回弹频率可分辨。原中子星逐渐收缩时表面重力增强，受撞击激发的浮力振荡也会沿时间向较高频率移动。

若快速差速转动的芯核带有极向磁场，感应方程给 $\partial_t B_\phi\simeq s(\bm B_{\rm p}\cdot\nabla)\Omega$。缠绕增加环向磁压 $B_\phi^2/(8\pi)$，并以 Maxwell 应力 $-B_s B_\phi/(4\pi)$ 输运角动量。超过外流动压时可沿轴推出喷流。轴对称喷流即使结构逼近定速膨胀也可产生迟迟不回到原值的波幅：对喷出质量 $m$、轴向速度 $v$，$D\propto m v^2 t^2$，故 $\ddot D$ 趋于常数，直至物质遭受下一次显著制动。它与高频物质振铃可以叠加。还有一种模型设塌缩物质碎裂成两个能公转的致密块。此时可直接引用 \S\ref{sec:pn-wave-energy} 的双星关系，应用前需确认碎块在若干轨道内都不被潮汐撕碎。几种机制分别需要怎样的转动、磁场和冷却，决定了哪一估计有物理根据。

\subsection{由中微子光度求记忆}

设 $\bm n$ 为中微子的逃逸方向，$\bm N$ 指向观者，$u=t-d/c$。中微子极相对论时，在式 \eqref{eq:gw-a-nonlinearmem} 的角积分中以单位立体角的中微子光度 $\d L_\nu/\d\Omega_{\bm n}$ 代入能量流即可。
在 $\bm N$ 作极轴的球坐标，令 $\mu=\cos\theta$，则 $e_+^{ij}n_i n_j/(1-\mu)=(1+\mu)\cos 2\phi$，另一偏振为 $(1+\mu)\sin2\phi$。前向的表观奇点已消失。若辐射各向同性，方位积分给零。定义 $\alpha_+(u)=\int(1+\mu)\cos2\phi\,\d L_\nu/\big(L_\nu\d\Omega\big)\,\d\Omega$，便有 $\dot h_+=2G\alpha_+L_\nu/(c^4d)$。这说明总共发出多少中微子只能给波幅的一个尺度，还须知道随时间变化的各向异性。

设 $E_\nu$ 在 $\tau$ 内以恒光度放出，且各向异性 $\alpha_+$ 不变。波幅随时间线性上升，净增量为 $\Delta h_+=2G\alpha_+E_\nu/(c^4d)$。记 $x=\pi f\tau$。略去直流分量，Fourier 振幅模为 $|\Delta h_+/(2\pi f)|$ 乘 $|\sin x/x|$。在 $f\tau\ll1$ 时呈 $1/f$，但可测与否还依赖仪器的低频截止。若后半段各向异性的符号反转，初末值可能相同，中途仍有可观测的瞬变。可见“记忆”所保留的是有方向权重的时间积分。

怎样从轴对称的中微子输运结果取出这个方向权重？设以恒星轴 $\bm k$ 表示一个最低阶的各向异性模型，
$\d L_\nu/\d\Omega=L_\nu[1+b P_2(\bm k\cdot\bm n)]/(4\pi)$。令 $\bm k$ 在观者天球的投影沿 $x$，取 $\bm k\cdot\bm n=\cos\iota\,\mu+\sin\iota\sqrt{1-\mu^2}\cos\phi$。代回积分，含常数、$\cos\phi$ 的项都和 $\cos2\phi$ 正交，故只余 $\tfrac32\sin^2\iota(1-\mu^2)\cos^2\phi$。而 $\int_0^{2\pi}\cos^2\phi\cos2\phi\,\d\phi=\pi/2$，$\int_{-1}^1(1+\mu)(1-\mu^2)\,\d\mu=4/3$，得 $\alpha_+=b\sin^2\iota/4$、$\alpha_\times=0$。在该模型中轴向与赤道光度密度之比为 $(1+b)/(1-b/2)$，但这一比值不在一般角分布上唯一确定记忆。譬如给光度加一个保持轴向与赤道值相同的更高多极，积分可能仍变。

还可算上述模型的波能。用双偏振能流 $c^3(\dot h_+^2+\dot h_\times^2)/(16\pi G)$，积分 $d^2\d\Omega$ 并用角均值 $\langle\sin^4\iota\rangle=8/15$，得 $P_{\rm GW}(t)=\frac{G}{30c^5}\, b(t)^2L_\nu(t)^2$。因而同样的永久位移若来自不同的光度历程，所放能量可以不同。定 $b$、总能量和持续时间 $\tau$ 时，Cauchy 不等式给 $\int L_\nu^2\,\d t\ge E_\nu^2/\tau$，等号恰为恒光度发射。若 $b$ 半途变号，末记忆是有符号量，它可小到近零，能量积分却始终为正。这个例子也解释从引力波和中微子同时观到的时间结构可以恢复什么信息：引力波在此测各向异性的加权光度，近球对称的大部分释放仍须由粒子探测本身测量。

\subsection{转动中子星上的两种频率尺度}

脉冲星自转在 $P$--$\dot P$ 平面的位置可用来估算长期能量收支。设 $I$ 近于常数，角速度 $\Omega=2\pi/P$ 满足 $\dot\Omega=-K\Omega^n$。若 $K$ 不变，则制动指数 $n=\Omega\ddot\Omega/\dot\Omega^2$，从出生周期 $P_0$ 转到 $P$ 所需时间为
\[
 t=\frac{P}{(n-1)\dot P}
       \left[1-\left(\frac{P_0}{P}\right)^{n-1}\right]\quad(n\ne1).
\]
磁偶极辐射给 $n=3$，只有在 $P_0\ll P$ 时常用的 $P/(2\dot P)$ 才近似年龄。真空斜偶极的辐射功率随磁矩垂直分量平方及 $\Omega^4$ 变，配平 $I\Omega|\dot\Omega|$ 得表面偶极场 $B\propto\sqrt{P\dot P}$，系数依磁轴倾角、半径及磁层模型而变。磁星的爆发还可消耗内部磁场能 $E_B\sim B^2R^3/6$（Gaussian 制），此时电磁光度不能直接用来标定自转损耗。

宽轨双星中的老星若吸积了伴星的角动量，可以被重新加速至毫秒周期。设吸积率为 $\dot M$、盘内缘为 $r_{\rm in}$，角动量输入量级 $\dot M\sqrt{GMr_{\rm in}}$。当星的自转接近盘内缘的 Kepler 角速度，物质与磁场交换角动量的符号可能改变。故出生很快、后来减速的孤立星，与先减速再被加速的毫秒星即使周期相同，其年龄与内部热状态也可大不相同。计时测得的 $\dot P$ 还受运动学影响，横向速度 $v_\perp$ 使视向距离的二阶导数含 $v_\perp^2/d$，带来 $\dot P/P=v_\perp^2/(cd)$ 的正贡献。用极小的 $\dot P$ 推断本征转矩时必须先处理它。

对绕主轴稳定自转的星，定义赤道椭率 $\epsilon=(I_1-I_2)/I$，$I$ 取自转惯量。四极功率为 $32 G I^2\epsilon^2\Omega^6/(5c^5)$，故单此机制给 $n=5$。将取向因子分开后，连续波振幅尺度为 $h_0=4GI|\epsilon|\Omega^2/(c^4d)$。若所有本征自转能损失皆由此波带走，消去 $\epsilon$ 得
\begin{equation}
 h_0^{\rm sd}=\frac1d\left(\frac{5GI}{2c^3}
                      \frac{|\dot\Omega|}{\Omega}\right)^{1/2}.
 \label{eq:gw-csd}
\end{equation}
测得上限 $h_0<h_0^{\rm sd}$ 的某一比例 $b$ 时，这种四极波至多消耗 $b^2$ 的自转损耗功率。若同时有常系数的磁偶极和质量四极转矩，写 $-\dot\Omega=K_3\Omega^3+K_5\Omega^5$，求导得 $n=3+2K_5\Omega^5/|\dot\Omega|$。若 $K_3$ 随磁轴倾角而变，后一种分账立即失效，而由能量守恒所得的上限仍可在没有外部补给时使用。弹性壳的非轴对称山、内部环形磁场引起的形变、吸积的不均匀沉积都可造成 $\epsilon$；同一个 $h_0$ 并不能挑出其中之一。



\subsection{星体的自由振荡}

自转给出一只频率很稳定的钟，振荡模给出另一只。先略去自转，用第 \ref{chpt:BH} 章的球对称静态平衡星作背景。在 Newton 极限，令质元的 Lagrange 位移为 $\bm\xi$、Euler 扰动为 $\delta$、随流扰动为 $\Delta=\delta+\bm\xi\cdot\nabla$。由小体积的 Jacobian 得 $\Delta\rho=-\rho\nabla\cdot\bm\xi$，从而 $\delta\rho=-\nabla\cdot(\rho\bm\xi)$。若移动快于热交换与成分调整，取 $\Delta p=\Gamma_1 p\Delta\rho/\rho$。时间依赖取 $\exp(-\mathrm i\omega t)$，运动方程和场方程给
\[
 -\omega^2\rho\bm\xi=-\nabla\delta p-\delta\rho\nabla\Phi
                      -\rho\nabla\delta\Phi,\qquad
 \nabla^2\delta\Phi=4\pi G\delta\rho .
\]
中心正则、表面随流压强 $\Delta p=0$、外部势衰减，三者合在一起选出离散频率。把度规或势的扰动冻结的做法称 Cowling 近似；它可分清流体的恢复力，不能求纯时空的模。

把球谐阶数记为 $\ell$，局域水平波数平方 $k_h^2=\ell(\ell+1)/r^2$。令 $g=-p'/\rho>0$、声速 $c_s^2=\Gamma_1p/\rho$。一小团物质上移 $\xi_r$ 后与周围压强平衡，团块密度和环境密度相差 $\rho[(\ln p)'/\Gamma_1-(\ln\rho)']\xi_r$，故浮力加速度为 $-N^2\xi_r$，其中 $N^2=g[(\ln p)'/\Gamma_1-(\ln\rho)']$。若 $N^2<0$，位移越大上浮越快，便发生对流。若 $N^2>0$，在短径向波长、忽略尺度高度修正时，声学及浮力两支色散关系可写成 $k_r^2\simeq\frac{(\omega^2-c_s^2k_h^2)(\omega^2-N^2)} {c_s^2\omega^2}$。高于两个特征频率时可有声压维持的 $p$ 模。低于两者时可有浮力维持的 $g$ 模；中间区域 $k_r^2<0$，两边的振荡只通过衰减解相接。$g$ 模频率因此对温度及成分梯度敏感。频率在声模和重力模之间而径向无节点的基模称 $f$ 模。以星为整体运动，$\omega_f^2$ 应与 $GM/R^3$ 同量级。

均匀不可压星给一个能算到底的例子。取无旋位移 $\bm\xi=A\nabla(r^\ell Y_{\ell m})$，因 $\nabla^2(r^\ell Y_{\ell m})=0$ 自动满足体内不可压。表面位移 $\xi_R Y_{\ell m}$ 等效于薄面密度 $\rho\xi_R Y_{\ell m}$，其中 $\xi_R=A\ell R^{\ell-1}$。扰动势在内外分别正比于 $r^\ell$ 和 $r^{-\ell-1}$；势连续，径向导数外减内为 $4\pi G\rho\xi_R Y_{\ell m}$，故 $\delta\Phi(R)=-4\pi G\rho R\xi_R Y_{\ell m}/(2\ell+1)$。在表面用 $\delta p/\rho=g\xi_R Y_{\ell m}$，积分运动方程 $\omega^2 A r^\ell Y_{\ell m}=\delta p/\rho+\delta\Phi$，得到
\[
 \omega_f^2=\frac{\ell g}{R}-\frac{4\pi G\rho\ell}{2\ell+1}
            =\frac{2\ell(\ell-1)}{2\ell+1}\frac{GM}{R^3}.
\]
$\ell=1$ 是整星平移，故频率为零；$\ell=2$ 才是可辐射的第一支。若把 $\delta\Phi$ 丢掉，连平移也会有频率，这就暴露了 Cowling 近似在最低多极上的不足。相对论星还可激发主要由曲率承担的 $w$ 模，须联立流体、度规并在无穷远施加出射条件，频率与衰减率同由一套边值问题决定。

自由振荡的波能有两个不同的衰减时间。若振幅 $q=q_0e^{-t/\tau}\cos\omega t$ 且 $\omega\tau\gg1$，能量与振幅平方成正比，故能量的 $e$ 倍衰减时间为 $\tau/2$。对上述均匀不可压的轴对称四极基模，可用一个无量纲坐标 $q$ 写 $a=R(1-q/2)$、$b=R(1+q)$，一阶体积恰守恒。位移速度为 $(\dot x,\dot y,\dot z)=\dot q(-x/2,-y/2,z)$，用 $\int\rho x_i^2\,\d^3x=MR^2/5$ 得
\[
 T_{\rm mode}=\frac{3}{20}MR^2\dot q^2,\qquad
 D=\frac35MR^2q+O(q^2).
\]
过零时位移消失，全模能 $E$ 就是此时的动能。波功率则随时间变化，须平均一周。将 $q=q_0\cos\omega t$ 代入得到
\[
 E=\frac{3}{20}MR^2\omega^2q_0^2,\qquad
 \langle P\rangle=\frac{3G}{125c^5}M^2R^4\omega^6q_0^2,\qquad
 \tau=\frac{2E}{\langle P\rangle}=\frac{25c^5}{2GMR^2\omega^4}.
\]
这里的 $\omega_f^2=4GM/(5R^3)$，故 $\tau\propto c^5R^4/(G^3M^3)$，与 $f_f\propto\sqrt{M/R^3}$ 的密度依赖不同。实际物态不是均匀不可压，强场也会改系数；两种尺度的不同正使得同时测到频率、寿命比只测一个频率多了一项结构信息。

将 \S\ref{sec:gw-b-burst} 的窄带等效能量用于轴对称 $m=0$ 星震，波幅随 $\sin^2\iota$ 改变。对能流作全天平均，$\langle\sin^4\iota\rangle=8/15$，故以赤道观者波幅算出的各向同性等效能量还应乘 $8/15$ 才是总辐射能。若释放相同能量却由较高频的模承担，$h_{\rm rss}$ 随 $1/f_0$ 降低；若同时仪器噪声随频率升高，检测还将进一步困难。这是即使千赫频段的星震固有频率很明确，能否实际测到模寿命仍取决于激发振幅的理由。

星体具有弹性壳和磁场时，除流体模还可出现切向剪切与 Alfvén 振荡。取薄壳剪切模量 $\mu_{\rm s}$，位移在球面上为无散的 $X_A$ 型，切向恢复力的量级给 $\omega_\ell^2\simeq(\mu_{\rm s}/\rho)(\ell-1)(\ell+2)/R^2$，$\ell=1$ 刚性转动不使壳发生剪切。若耦合到磁化流体核心，波还能沿磁力线以 $v_A\simeq B/\sqrt{4\pi\rho}$ 传播，不同场线长度给不同的往返时间。在磁星爆发后的电磁余辉中看见准周期调制，只能说明辐射区或耦合区域在振荡；要把同一峰认作强引力波，还需证明该模确有足够的质量或流多极矩。纯切向、近不可压模式可能强烈影响磁层，却只有很弱的质量四极。

\subsection{转动怎样使辐射成为激发机制}

令相位为 $m\phi-\omega t$，取 $\Omega>0$、$m>0$。变换到共转坐标 $\phi'=\phi-\Omega t$ 后频率为 $\omega'=\omega-m\Omega$。若 $0<\omega/m<\Omega$，惯性观者看图样顺转，流体却看见它逆转。向外的单色波若以速率 $L>0$ 带走能量，角动量流率为 $mL/\omega$；这是同一个模式在时间平移与绕轴转动下相位变化的比。对固定重子数的均匀转动平衡序列有 $\d E_0/\d J=\Omega$。星偏离新平衡的能量变化率遂为
\[
 \frac{\d(E-E_0)}{\d t}=-L+\Omega\,\frac{m}{\omega}L
                        =-\frac{\omega'}{\omega}L .
\]
在上述频率范围内它为正。波带走系统的正能量，却取走了足够多的角动量，使新平衡的能量下降更快，余下可供振荡的能量增加。这便是 Chandrasekhar--Friedman--Schutz (CFS) 不稳定性。它要求连续比较两个参考系中的图样转向，仅知道波为右旋或左旋不能做出判断。

在缓转、近乎正压的流体中，有一类位移近似沿等密度面，恢复力为 Coriolis 力。将线性旋转流体方程投影到轴型矢量球谐上，径向涡量方程给 $\omega'=-2m\Omega/[\ell(\ell+1)]$。对 $\ell=m=2$ 的 $r$ 模，惯性系中 $\omega=4\Omega/3$，既在 CFS 区间内，又不同于转动椭球质量四极的 $2\Omega$。最低阶 Euler 密度扰动很小，辐射来自流四极。令模速率扰动量级为 $\alpha\Omega R$，则模能量随 $\alpha^2\Omega^2$ 变，流四极随 $\alpha\Omega$ 变，功率随其三阶时间导数平方、即 $\alpha^2\Omega^8$ 变。因此增长率随 $\Omega^6$ 变，快速自转能使很弱的辐射反作用积累起来。

真实中子星还有剪切粘滞、体粘滞、壳与流体边界耗散。以振幅演化 $\dot\alpha=-\alpha/\tau$ 定义符号，则 $1/\tau=1/\tau_{\rm GR}+1/\tau_{\rm shear}+1/\tau_{\rm bulk}+\cdots$，第一项在失稳区为负。中微子透明、由修正 Urca 过程控制体粘滞的一个模型中，剪切阻尼率随温度约按 $T^{-2}$ 变，体阻尼率在常用的高振荡频率极限随 $T^6\Omega^2$ 变。低温与高温各有阻尼，中间温度可能打开不稳定窗口；实际边界由成分与超流性决定。线性计算只给启动条件，不能给饱和振幅。向子模的非线性能量转移、差速转动与磁场的耦合都可能限制可见波幅。双中子星并合后的热星转得快，物态又偏离冷星，常以最强的流体四极峰估计平均密度；若迅速塌成黑洞，谱便转入下节的曲率模。

还可从局部耗散积分算不稳定窗口的低温边。令速度扰动的复振幅为 $\delta v_i$，无迹对称应变率为 $\delta\sigma_{ij}=\tfrac12(\partial_i\delta v_j+\partial_j\delta v_i)-\delta_{ij}\partial_k\delta v^k/3$。单位体积的粘滞耗能在周期平均后正比于 $2\eta_s|\delta\sigma_{ij}|^2+\zeta|\nabla\cdot\delta\bm v|^2$，$\eta_s,\zeta$ 为剪切、体粘滞系数。能量约为 $\int\rho|\delta\bm v|^2/2\,\d^3x$，故模的体积积分决定阻尼率，不能取中心粘滞系数代入一个对全星一样的衰减时间。对慢转星的 $\ell=m$ 轴型模，速度随 $r^\ell$ 变且最低阶无散，体粘滞须从更高阶的压缩得到，剪切粘滞则已有领头贡献。具体地，在此近似下
\[
 \frac{1}{\tau_{\rm shear}}=(\ell-1)(2\ell+1)
   \frac{\int_0^R \eta_s r^{2\ell}\,\d r}
        {\int_0^R \rho r^{2\ell+2}\,\d r}.
\]
若 $\rho,\eta_s$ 均匀，分式为 $(2\ell+3)\eta_s/[(2\ell+1)\rho R^2]$。对于 $\ell=2$ 得 $7\eta_s/(\rho R^2)$，与粘性扩散时间 $R^2/(\eta_s/\rho)$ 只差模结构系数。

对同一支 $r$ 模，把速度角分布代入流多极功率并除以两倍模能，惯性频率 $m\Omega[1-2/\ell(\ell+1)]$ 给出
\[
 \frac{1}{\tau_{\rm GR}}=-\frac{32\pi G\Omega^{2\ell+2}}{c^{2\ell+3}}
 \frac{(\ell-1)^{2\ell}}{[(2\ell+1)!!]^2}
 \left(\frac{\ell+2}{\ell+1}\right)^{2\ell+2}
 \int_0^R \rho r^{2\ell+2}\,\d r .
\]
对于均匀星、$\ell=2$，积分为 $3MR^4/(28\pi)$；增长率随 $M R^4\Omega^6$ 变。质量与半径固定时，周期加倍，增长时间便变为原来的 $64$ 倍。把剪切边界简化为 $A\Omega^6=B T^{-2}$，临界角频率随 $T^{-1/3}$；在体粘滞边界 $A\Omega^6=D T^6\Omega^2$，则随 $T^{3/2}$。若 $\Omega$ 固定，星随中微子冷却从高温一侧进入窗口，可能先激发模，后因角动量损失或降温离开。初始很冷、转得慢的星即使测得极小的电磁制动，也不能不经此判断便认作正在发生 CFS 辐射。

吸积星还有另一个归一振幅的办法。若近于长期转矩平衡，设盘供给 $N_{\rm acc}\simeq\dot M\sqrt{GMr_{\rm in}}$，四极波的角动量流为 $L_{\rm GW}/\Omega$。由平衡得 $L_{\rm GW}\simeq\Omega\dot M\sqrt{GMr_{\rm in}}$。设恒定转动的质量四极给前述 $h_0$，消去椭率可写 $L_{\rm GW}=2c^3d^2\Omega^2 h_0^2/(5G)$，从而 $h_0^2\simeq\frac{5G}{2c^3d^2\Omega} \dot M\sqrt{GMr_{\rm in}}$。若吸积光度 $4\pi d^2 F_X\simeq GM\dot M/R$，距离可从表达式消掉。这里假定盘将上述角动量全传给星，X 射线近各向同性且来自吸积，故估计与 \eqref{eq:gw-csd} 对孤立星的能量上限有不同前提。有吸积补给时即使观测 $\dot\Omega\simeq0$，引力波转矩仍可有限。

\section{黑洞对扰动的响应}\label{sec:gw-cbh}

以下在黑洞邻近采用 $G=c=1$。孤立黑洞的平稳外场见 \S\ref{sec:KN}。设在这一外场投入一颗小质量物体，初生畸变先直接传到观者，随后是带衰减的振荡，再后只剩一个弱尾巴。这三段信号分别由光锥、曲率势垒和远处弯曲时空中的散射规定，宜从一个初值问题同时求出。

\subsection{从标量场到两种宇称}

先取无质量标量 $\varphi=\sum_{\ell m}r^{-1}\psi_{\ell m}(t,r)Y_{\ell m}$，在 Schwarzschild 背景上记 $F=1-2M/r$。波算符作用于单一分量时，径向部分是 $r^{-2}\partial_r[r^2F\partial_r(\psi/r)]$。展开括号，原来会作用在 $1/r$ 上的导数留下一项 $-F'\psi/r^2$。引入乌龟坐标 $\d r_*/\d r=F^{-1}$，得
\[
 (-\partial_t^2+\partial_*^2-V_\ell^{(0)})\psi_{\ell m}=0,\qquad
 V_\ell^{(0)}=F\left[\frac{\ell(\ell+1)}{r^2}+\frac{2M}{r^3}\right],
 \quad \partial_*=\partial/\partial r_* .
\]
$r_*=r+2M\ln(r/2M-1)$ 将视界送到 $-\infty$。势在两端消失，在中间形成势垒。平直极限仍有角动量势，曲率造成的额外无穷远散射不能简单当作一个有限宽度的方势垒。

对度规扰动，利用球面对称性可先在角向分解各个分量。把坐标指标分成轨道空间 $a,b=t,r$ 与单位球面 $A,B$，球面度规记 $\Omega_{AB}$。令 $Y=Y_{\ell m}$、$Y_A=D_A Y$、$X_A=-\varepsilon_A{}^B D_BY$。偶型球面张量为 $Y_{AB}=[D_AD_B+\ell(\ell+1)\Omega_{AB}/2]Y$，奇型为 $X_{AB}=D_{(A}X_{B)}$。这两项均无迹。球面反演使一组变号因子为 $(-1)^\ell$、另一组为 $(-1)^{\ell+1}$。背景没有空间手征性，所以即使没有任何轴对称假设，两组方程也独立。

对 $\ell\ge2$ 可将偶型扰动分成 $h_{ab}=H_{ab}Y$、$h_{aB}=j_aY_B$ 和 $h_{AB}=r^2(K\Omega_{AB}Y+G_{\rm e}Y_{AB})$；奇型只含 $h_{aB}=h_aX_B$、$h_{AB}=h_2 X_{AB}$。共七个偶变量、三个奇变量。一个无穷小坐标位移在偶域有三个自由函数，在奇域有一个。Regge--Wheeler 规范取 $j_a=G_{\rm e}=h_2=0$，其余变量仍有场方程中的约束。最后每一宇称、每组 $\ell,m$ 只需一个主变量。$\ell=0$ 的真空变化可改变总质量，奇 $\ell=1$ 可给黑洞增添角动量，偶 $\ell=1$ 涉及质心选择。这些低阶没有两个独立角向剪切，不能将 $\ell\ge2$ 的规范条件照搬给它们。

奇型消元可直接看出势垒从何而来。在上述规范写 $h_t=h_0$、$h_r=h_1$，无源角向方程给 $\partial_t h_0=F\partial_r(Fh_1)$，另一方程为
\[
 \partial_t^2h_1-\partial_r\partial_t h_0
 +\frac2r\partial_t h_0+\frac{F}{r^2}(\ell-1)(\ell+2)h_1=0 .
\]
令 $Q=Fh_1/r$。第一式右侧为 $F(Q+r\partial_r Q)$，代入第二式，乘 $F/r$ 并合并各阶导数，得
\begin{equation}
 (-\partial_t^2+\partial_*^2-V_{\rm o})Q=0,\qquad
 V_{\rm o}=F\left[\frac{\ell(\ell+1)}{r^2}-\frac{6M}{r^3}\right].
 \label{eq:gw-crw}
\end{equation}
其 $-6M/r^3$ 与标量的 $+2M/r^3$ 不同。可用实际的 $h_1$ 或规范不变量计算 $Q$。举例，奇位移生成元取 $\xi_B=\Lambda X_B$ 且 $h\mapsto h-{\cal L}_\xi g$，则 $\delta h_1=-\partial_r\Lambda+2\Lambda/r$、$\delta h_2=-2\Lambda$；组合 $h_1-\tfrac12\partial_r h_2+h_2/r$ 在任意规范下都等于上述规范中的 $h_1$。

偶型中令 $h_{rr}=H_2Y/F$、$h_{tt}=F H_0Y$，$\lambda=(\ell-1)(\ell+2)/2$。在无源区 $H_0=H_2$。一个常用的 Zerilli--Moncrief 归一化为
\[
 Z=\frac{r}{\lambda+1}\left[
 K+\frac{F}{\lambda+3M/r}(H_2-r\partial_rK)\right].
\]
代入约束消去 $H_2$ 与 $h_{tr}$ 后，$Z$ 满足与 \eqref{eq:gw-crw} 形同的方程，但势换为
\[
 V_{\rm e}=\frac{2F}{r^3(\lambda r+3M)^2}
 \left[\lambda^2(\lambda+1)r^3+3\lambda^2Mr^2
                         +9\lambda M^2r+9M^3\right].
\]
两者远处的首项均为 $\ell(\ell+1)/r^2$。近视界处 $F$ 乘有限函数，势随 $r_*$ 指数趋零，$r$ 坐标下很短的一段却是无限长的传播区。

为何不同的势给同一振铃谱？设 $W=W_0+3MF/[r(\lambda r+3M)]$，$W_0=\lambda(\lambda+1)/(3M)$。直接求导可验 $V_{\rm o}=W^2-\partial_*W-W_0^2$，$V_{\rm e}=W^2+\partial_*W-W_0^2$。故两个一维算符可分别分解成 $(-\partial_*+W)(\partial_*+W)-W_0^2$ 与次序颠倒的乘积。若 $Q$ 为频率 $\omega$ 的奇解，$(\partial_*+W)Q$ 即同频的偶解。在两端 $W\to W_0$，微分只给进出射波各乘常数，所以一般的进出射边界条件也保留。这一论证在某个乘数为零的代数特殊频率需要单独处理。两组的谱相同并不要求某个物体以相同强度激发它们；源对两种宇称的投影可以全然不同。

\subsection{怎样从主变量还原度规？}

上述计数中减去规范自由度后仍多于两个主变量。其余信息到哪里去了？考虑偶域在轨道空间的一般变换 $\xi_a=\xi_a(t,r)Y$、$\xi_A=r^2\xi(t,r)Y_A$。以 $D_a$ 表示二维轨道度规的协变导数，$r_a=D_a r$，逐项求 $-2\nabla_{(\mu}\xi_{\nu)}$ 得
\begin{align*}
 \delta H_{ab}&=-D_a\xi_b-D_b\xi_a,&
 \delta j_a&=-\xi_a-r^2\partial_a\xi,\\
 \delta G_{\rm e}&=-2\xi,&
 \delta K&=\ell(\ell+1)\xi-\frac2r r^a\xi_a .
\end{align*}
$\delta K$ 第一项是 $D_AD_BY$ 含迹的部分，第二项来自球面半径随二维坐标变化。令 $p_a=j_a-\tfrac12r^2\partial_aG_{\rm e}$，则 $\delta p_a=-\xi_a$。所以 $\tilde H_{ab}=H_{ab}-2D_{(a}p_{b)}$、$\tilde K=K+\ell(\ell+1)G_{\rm e}/2-(2/r)r^a p_a$ 在一般坐标扰动下不变。在 RW 规范两者恰退到 $H_{ab},K$。这给数值提波一种做法：先将扰动投影到球谐，取规范不变量，再代入主变量；不要求演化本身就在 RW 规范。

投影时沿用 \S\ref{sec:gw-a-multipoles} 的张量球谐归一。这里为便于解径向方程，用未归一的 $Y_{AB},X_{AB}$，其模平方在整球积分均为 $(\ell+2)!/[2(\ell-2)!]$。在正交球标架测量时还须除 $r^2$，因坐标分量 $h_{\theta\theta}$ 本身不是应变。

对真空偶模，有效的第一步消元是把度规迹系数写成 $Z$ 与其一阶导数。仍取本节的归一，有
\[
 K=g(r)Z+\partial_* Z,\qquad
 g(r)=\frac{\lambda(\lambda+1)r^2+3\lambda Mr+6M^2}
                  {r^2(\lambda r+3M)} .
\]
由 $Z$ 的定义反解即给
\[
 H_2=r\partial_r K+\frac{\lambda+3M/r}{F}
           \left[\frac{\lambda+1}{r}Z-K\right].
\]
因此 $H_2$ 涉及 $Z$ 的径向二阶导数。再用主方程 $\partial_*^2 Z=\partial_t^2 Z+V_{\rm e}Z$，即可换成时间二阶导数与较低导数。无源角向无迹方程给 $H_0=H_2$，剩余动量约束再定 $H_1$。这说明度规系数全随同一个传播变量而动。若源有张量角应力，$H_0-H_2$ 可直接受驱，便不能把无源重建式一路用到物质内部。

在 $r\gg M$ 的出射真空区，设 $Z\simeq Z(u)$、$u=t-r_*$，上式给 $K\simeq-\dot Z+O(r^{-1})$，$H_2$ 中甚至有 $r\ddot Z$。这里似乎不见 $1/r$ 衰减。原因是 RW 规范为分离径向方程而设，其坐标度规在辐射区未必是渐近辐射规范。还原到横向正交基时，角向的纯张量项才留下 $1/r$ 的应变。取奇变量为与偶 $Z$ 同量纲的 Cunningham--Price--Moncrief 变量 $\Psi_{\rm o}$，约定其符号使远处 $H_{\ell m}$ 为
\[
 H_{\ell m}=\frac12\sqrt{\frac{(\ell+2)!}{(\ell-2)!}}\,
             (Z_{\ell m}+\mathrm i\Psi_{{\rm o},\ell m}).
\]
对实度规的全部 $m$ 求和时偶奇交叉项相消，式 \eqref{eq:gw-cmodes} 给每个 $\ell$ 的总能流为 $(\ell+2)!/[64\pi(\ell-2)!]$ 乘 $\sum_m(|\dot Z_{\ell m}|^2+|\dot\Psi_{{\rm o},\ell m}|^2)$。在真空，$\partial_t\Psi_{\rm o}$ 与 $Q$ 成固定常数比例，符号依奇球谐取向；$Q=Fh_1/r$ 本身与 $Z$ 不同量纲。若把 $Q$ 直接加在 $Z$ 后面并对两者同求一次时间导数，奇模能谱会多出一个频率平方。

还有一个有用的检查是主方程的正能量。设实的 $X$ 代表任一宇称的主变量，$V$ 为相应实势，乘齐次波动方程以 $\partial_t X$ 后积分可得
\[
 \frac{\d}{\d t}\frac12\int_{x_1}^{x_2}
  [(\partial_tX)^2+(\partial_xX)^2+V X^2]\,\d x
    =\left[\partial_t X\,\partial_x X\right]_{x_1}^{x_2}.
\]
外区 $r>2M,\ell\ge2$ 时两个势均非负。左端纯向内、右端纯向外的边界使能量只减少。初值有限能量便不容许在固定范围内生长的无源不稳定模。这是 Schwarzschild 线性稳定的一项直接证据，其假设与“准正规模本征函数可以平方归一”无关。对旋转黑洞不能直接沿用此一维正定论证，因背景时间平移在能层内不再处处类时。

\subsection{散射、极点及晚时尾}

令单频分量为 $\exp(-\mathrm i\omega t)u(r_*)$。在视界，沿先进坐标 $v=t+r_*$ 的波 $\exp(-\mathrm i\omega v)$ 向内；在无穷远，沿 $u_{\rm ret}=t-r_*$ 的波向外。对实频和实势，从右入射且入射振幅为一，右端波为 $\exp(-\mathrm i\omega r_*)+{\cal R}\exp(\mathrm i\omega r_*)$，左端为 ${\cal T}\exp(-\mathrm i\omega r_*)$。用常数流 $j=(u^*u'-u u'^*)/(2\mathrm i)$ 比较两端得到 $|{\cal R}|^2+|{\cal T}|^2=1$，撇号指 $\partial_*$。若要求无入射源同时两端只许物理流出，该等式就禁止非零的实频解，振铃频率必然延拓到复平面。

具体地，从左端只取向内的齐次解 $u_{\rm in}$，从右端只取向外的 $u_{\rm up}$。对 $x=r_*$，频域 Green 函数为
\[
 G_\omega(x,x')=\frac{u_{\rm in}(x_<)u_{\rm up}(x_>)}{W(\omega)},\quad
 W=u_{\rm in}u_{\rm up}'-u_{\rm up}u_{\rm in}',\quad
 x_<=\min(x,x'),\quad x_>=\max(x,x').
\]
它在 $x=x'$ 连续，一阶导数跳跃为一。$u_{\rm in}$ 在远处若含 $A_{\rm in}\exp(-\mathrm i\omega x)+A_{\rm out}\exp(\mathrm i\omega x)$，$u_{\rm up}$ 的出射振幅取一，则 $W=2\mathrm i\omega A_{\rm in}$。$W(\omega_n)=0$ 时，同一非零解恰好满足两端辐射条件，称\textbf{准正规模}(quasi-normal mode)。$\omega_n=\omega_R-\mathrm i/\tau$ 的负虚部使扰动以 $\tau$ 衰减；此解在实 $x$ 的两端反而呈指数增长，故不能像有限弦的本征函数那样用平方可积条件归一化。源和初值决定留数，频率才是黑洞性质。

反演可先取 $\operatorname{Im}\omega>0$ 作 Laplace 积分，保证推迟性，再对晚时向下移积分路径。高频部分给直接到达的脉冲；孤立极点贡献衰减指数；无穷远长程势令 $\omega=0$ 附近有非解析部分，其分支积分给逆幂尾。若割线两侧的跃度低频随 $\sigma^p$ 变，$\omega=-\mathrm i\sigma$，则 $\int_0^\infty e^{-\sigma t}\sigma^p\,\d\sigma=\Gamma(p+1)t^{-p-1}$，所以晚到任何指数项之后，只要系数不为零，微弱的尾都会占优。在 Schwarzschild 外部的固定半径，紧支撑、一般初值的一支 $\ell$ 模有 $t^{-(2\ell+3)}$ 的 Price 尾；在零无穷远取的极限或有初始静态多极时幂数可变。环绕势峰的局域展开只给极点近似，不能由它推出此尾。

算一个势垒近似。若 $V$ 峰值在 $x_0$，写 $V\simeq V_0+V_0''(x-x_0)^2/2$，$V_0''<0$。倒抛物线的出射条件给 $\omega^2\simeq V_0-\mathrm i(n+\tfrac12)\sqrt{-2V_0''}$，$n=0,1,\ldots$。在 $\ell\gg n$ 时取 $V\simeq L^2F/r^2$、$L=\ell+1/2$。使 $\partial_r(F/r^2)=0$ 得 $r_0=3M$；峰处 $\partial_*^2 V=F^2\partial_r^2V=-2L^2/(729M^4)$，$V_0=L^2/(27M^2)$。取平方根即得 $M\omega_{\ell n}\simeq\frac{L-\mathrm i(n+1/2)}{3\sqrt3}$。实部反映不稳定光子圆轨道的角速度，虚部反映附近光线的散开速率。$\ell=2$ 不在大 $\ell$ 极限内，精确基模约为 $M\omega=0.374-0.0890\mathrm i$。恢复单位后，对 $10M_\odot$ 黑洞约为 $1.21$\,kHz、振幅衰减时间 $0.55$\,ms。质量同比放大，频率反比缩小、寿命同比延长，所以这个频率不能直接用于不同质量的黑洞。

\subsection{怎样从初值或频谱辨认准正规模？}

设在 $t=0$ 给空间函数 $X_0(x),\Pi_0(x)$，$x=r_*$。对 $t\ge0$ 定义 $\hat X(s,x)=\int_0^\infty e^{-st}X(t,x)\,\d t$，取 $\operatorname{Re}s>0$。两次分部积分把方程变为 $[-\partial_x^2+V+s^2]\hat X=sX_0+\Pi_0$。若 $G_s$ 为这个算符的逆，其两端物理解在右端按 $e^{-s x}$、左端按 $e^{s x}$ 变。于是
\[
 X(t,x)=\int_{a-\mathrm i\infty}^{a+\mathrm i\infty}
 \frac{\d s}{2\pi\mathrm i}e^{st}
 \int G_s(x,x')[sX_0(x')+\Pi_0(x')]\d x',\qquad a>0 .
\]
孤立简单极点的幅值含初值与该极点径向函数的重叠，也含 $1/W'(s_n)$。若 $X_0,\Pi_0$ 的组合与某个极点的投影恰好为零，该模的频率仍存在，观者却不见它。在实测波形上将频率可辨直接等同于频谱中存在本征值，会高估可测的模式数。

一只阻尼振子的推迟响应给一个可解的对照。取 $\ddot q+2\gamma\dot q+\omega_0^2 q=F(t)$，$\gamma<\omega_0$。其推迟 Green 函数为 $\Theta(t)e^{-\gamma t}\sin(\omega_d t)/\omega_d$，$\omega_d=(\omega_0^2-\gamma^2)^{1/2}$。极点在 $\omega=\pm\omega_d-\mathrm i\gamma$；稳态受迫振幅的峰却在 $\omega^2=\omega_0^2-2\gamma^2$（正根存在时）。因此频谱实轴峰、无阻尼固有频率、极点实部本来就可以不相等。黑洞的源还在运动，使 $S_{\ell\omega}$ 非平，观到的能谱峰与 $1/W$ 的峰又差一层。读图得一个峰值不足以决定寿命，必须从完整时域或复响应同时拟合。

在黑洞势峰附近可用一个连完整谱都可写下的模型检查符号。取 $V(x)=V_0/\cosh^2(\alpha x)$，令 $z=[1+\tanh(\alpha x)]/2$。两端出射因子可剥为 $X=[z(1-z)]^{-\mathrm i\omega/(2\alpha)}F(z)$。代入径向方程，$F$ 满足超几何方程，参数可取
\[
 a=\tfrac12-\frac{\mathrm i\omega}{\alpha}
               +\sqrt{\tfrac14-\frac{V_0}{\alpha^2}},\quad
 b=\tfrac12-\frac{\mathrm i\omega}{\alpha}
               -\sqrt{\tfrac14-\frac{V_0}{\alpha^2}},\quad
 c_0=1-\frac{\mathrm i\omega}{\alpha}.
\]
取 $F={}_2F_1(a,b;c_0;z)$ 可使左端正则因子恰留所需行波。将此解延拓到 $z=1$，不需要的右端入射因子系数正比于 $1/[\Gamma(a)\Gamma(b)]$，因而 $a=-n$ 或 $b=-n$ 给 $\omega_n=\pm\sqrt{V_0-\alpha^2/4}-\mathrm i\alpha(n+\tfrac12)$。令 $\alpha^2=-V_0''/(2V_0)$ 就匹配了真实势的峰高、曲率。在大 $V_0/\alpha^2$ 时，上式与倒抛物线结果一致；但此模型在远处指数衰减，不含 Schwarzschild 的长程势，不能重现 Price 尾。调节一个势使它的基模和寿命都准确，仍不保证它给出同一个完整推迟响应。

真实势中可用级数法补足低阶 WKB 的精度。以 $z=1-2M/r$ 把外区压缩到 $0\le z<1$，先从径向函数抽出能在视界和无穷远产生所需行波的因子，再写为 $\sum_{j=0}^\infty a_j z^j$。视界的 Frobenius 指数可直接看出：$r_*\simeq2M\ln z+\text{常数}$，进视界支为 $z^{-2\mathrm i M\omega}$。无穷远 $r_*\simeq r+2M\ln(r/2M)$，出射支不仅含 $\exp(\mathrm i\omega r)$，还含 $r^{2\mathrm iM\omega}$。忘去对数相位会使远边界的“纯出射”随截断半径而变。代回后整理到三项递推 $\alpha_j a_{j+1}+\beta_j a_j+\gamma_j a_{j-1}=0$。取 $R_j=a_j/a_{j-1}$，得到 $R_j=-\gamma_j/(\beta_j+\alpha_jR_{j+1})$。在 $z=1$ 收敛须取最小解，从大 $j$ 反向递推，再与左端 $\beta_0+\alpha_0R_1=0$ 配合，所得就是连分式的本征频率条件。数值上可同时增大尾端截断并检查根是否稳定，比在实径向坐标上直接辨认微弱的次优势指数牢靠。

两个常用的渐近极限还应分清。先定泛音数 $n$，让角动量 $\ell$ 很大，频谱由光子轨道附近的势峰决定，已得到实部随 $\ell$ 增大。先固定 $\ell$ 而令 $n\gg1$，则 Schwarzschild 引力扰动的高阻尼分支有 $8\pi M\omega_n\simeq\ln3-2\pi\mathrm i(n+\tfrac12)$，取正实部。此时殊异点附近的解析延拓参与定谱，不能用上述倒抛物线公式把小 $n$ 外推到任意 $n$。高泛音衰减时间短，在晚时观测中通常很难分辨；低阶模才是稳定遗留外场的主要可见记录。

也可避免频域求根，直接在特征网格上解初值演化。用 $u=t-x,v=t+x$，齐次方程为 $4\partial_u\partial_vX+VX=0$。在边长 $\delta$ 的一个零边界小矩形上积分，若四角依时序记为南 $S$、东西 $E,W$、北 $N$，以两相邻点平均近似中心值，更新为
\[
 X_N=X_W+X_E-X_S-\frac{\delta^2}{8}V(x_{\rm mid})(X_W+X_E)
       +O(\delta^4)
\]
（平滑场的对称中心差分）。给两条初始零线上的紧支撑波包，就可连续得到直接传播、振铃和尾，而不必先假定一个模展开。固定 $x$ 的数值解若有尾 $X\sim t^{-p}$，用局地幂指数 $-t\,\partial_t\ln|X|$ 可检验它是否趋于 $p$。选择振铃和幂尾相近的时段去拟合单一指数，得到的寿命会随拟合窗移动，这正是多组时间窗和边界半径应当一并检查的原因。

\subsection{落入的物体激起什么波？}

令质量 $\mu\ll M$ 的质点沿极轴径向自由落下。其能动张量可从世界线积分写成 $T^{\alpha\beta}=\mu\,u^\alpha u^\beta\delta^{(3)}(\bm x-\bm x_p)/(\sqrt{-g}\,u^t)$。轴对称只留 $m=0$，径向运动无环向流，只激发偶模。若从无穷远静止落下，$u^t=1/F$、$u^r=-\sqrt{2M/r}$，故在频域将 $\delta[r-r_p(t)]$ 沿 $t$ 积掉时多出 $|\d r_p/\d t|^{-1}$ 和相位 $\exp[\mathrm i\omega t_p(r)]$。这一 Jacobian 与 $1/u^t$ 各来自不同的变量替换。消去 Einstein 方程中的非传播分量后，Zerilli 源还可含 $\delta'$。在实际积分中应以分部积分将导数移给径向核，不能把带导数的源当成一个常数冲量。

以 $S_{\ell\omega}$ 表示上述投影和消元所得源，出射振幅为
\[
 Z_{\ell\omega}^{\rm out}=\frac{1}{W(\omega)}
       \int_{-\infty}^{\infty}u_{\rm in}(x,\omega)S_{\ell\omega}(x)\,\d x .
\]
这既是求解法，也是对数值积分的检查。先从近视界按纯向内级数出发积分 $u_{\rm in}$，在远处拟合两个行波得 $W$；再作源投影，逆变换回波形。频谱的峰受 $1/W$ 强化，峰高还随投影而变。从无穷远静止落入 Schwarzschild 黑洞时，总辐射能数值约为 $0.0104\,\mu^2/M$，四极成分最大。质量比平方可预先看出：源为 $O(\mu)$，功率为波幅平方，时间尺度为 $M$；无量纲系数则必须求解散射问题。它不能用于两颗等质量黑洞，也不能直接乘以很多个碎块的数目。在同一波长内若各块同相落下，应先加振幅再平方，若落下时间展宽至数个周期则会相消。



还可把源积分中随半径迅速振荡的相位先算出。令 $y=\sqrt{r/(2M)}$，由 $\d t/\d r=-y^3/(y^2-1)$ 得
\[
 t_p(r)=t_0-4M\left[\frac{y^3}{3}+y+
                 \frac12\ln\!\left(\frac{y-1}{y+1}\right)\right],
 \qquad \frac{\d t_p}{\d r_*}=-y .
\]
在近视界 $y\to1$，$t_p\simeq-r_*+\text{常数}$，故质点先进时 $v_p=t_p+r_*$ 趋于有限值，远方的推迟时间 $u_p=t_p-r_*$ 却趋于无穷大。远处观者看到外部驱动趋弱，不能据此认为质点须等无限固有时才跨越视界。沿世界线 $\d r/\d\tau=-1/y$ 给从 $r_0$ 到 $2M$ 的固有时 $\tfrac{4M}{3}[(r_0/2M)^{3/2}-1]$，明确有限。晚时外场所见的准正规波因而不是质点在视界内不断向外传递消息，而是留在视界外的先前扰动被势垒散射。

轴上点源的角分布可用 $\delta^{(2)}(\bm n,\bm e_z)=\sum_{\ell m}Y_{\ell m}(\bm n)Y_{\ell m}^*(\bm e_z)$。因 $Y_{\ell m}(\theta=0)=0$ 对 $m\ne0$ 成立，四极 $m=0$ 的投影权为 $\sqrt{5/(4\pi)}$。若是扩展天体，把单粒子的强场波形乘上 \S\ref{sec:gw-a-sources} 的落入时间分布形状因子。即使势的极点谱完全相同，源在某个模频率上的激发也可相消。

\subsection{旋转黑洞为何要改用曲率？}

在 Kerr 背景，球面对称性消失，偶、奇度规分量的上述分离也随之失效。但其 Weyl 张量有两个重复主零方向。沿它们取实零矢 $l^\mu,n^\mu$，另取复共轭角向矢 $m^\mu,\bar m^\mu$，以 $l\cdot n=-1$、$m\cdot\bar m=1$ 归一。Weyl 张量在这一基底只需五个复分量 $\Psi_0,\ldots,\Psi_4$；对适配的平稳背景，只有 $\Psi_2$ 非零。在一阶扰动中，极端分量便能独立描述传入及传出的曲率。取无穷远的零标架与本文曲率号数，使出射分量满足
\begin{equation}
 \Psi_4=\ddot h_+-\mathrm i\ddot h_\times+O(r^{-2}).
 \label{eq:gw-cpsi}
\end{equation}
这个归一在无穷远给定。缩放 $l\to b l,n\to b^{-1}n$ 会改变 $\Psi_4$ 的射影因子；转动 $m\to e^{\mathrm i\zeta}m$ 时其相位为 $e^{-2\mathrm i\zeta}$。因此出射振幅应展开成自旋权 $s=-2$ 的球谐，而单独一个曲率坐标分量不能认作探测器读数。

给出实际用于计算频谱的方程。以 \eqref{eq:kn-metric} 的无电荷情形为背景，$a=J/M$、$\Delta=r^2-2Mr+a^2$、$r_\pm=M\pm\sqrt{M^2-a^2}$。对 Kinnersley 标架取出射变量 $(r-\mathrm ia\cos\theta)^4\Psi_4$，分离成 $R_s(r)S_s(\theta)e^{\mathrm i m\phi-\mathrm i\omega t}$。记 $K=(r^2+a^2)\omega-am$、$\lambda_s=A_s+a^2\omega^2-2am\omega$，则角方程与真空 Teukolsky 径向方程为
\begin{align*}
 0={}&\frac{1}{\sin\theta}\frac{\d}{\d\theta}
          \left(\sin\theta\frac{\d S_s}{\d\theta}\right)
 +\left[a^2\omega^2\cos^2\theta-2a\omega s\cos\theta
       -\frac{(m+s\cos\theta)^2}{\sin^2\theta}+s+A_s\right]S_s,\\
 0={}&\Delta^{-s}\frac{\d}{\d r}\left(\Delta^{s+1}\frac{\d R_s}{\d r}\right)
 +\left[\frac{K^2-2\mathrm is(r-M)K}{\Delta}
                     +4\mathrm is\omega r-\lambda_s\right]R_s .
\end{align*}
在极点要求 $S_s$ 正则选出 $A_s$；$a\omega=0$ 给 $A_s=\ell(\ell+1)-s(s+1)$。有限 $a\omega$ 时角函数为自旋加权扁球谐，故对固定的整数 $m$ 仍可解一维径向方程，但分离常数本身依赖待解的复频率。先随意取球谐本征值、只改径向方程，在快转时就不能求得正确频谱。

在视界附近用 $\d r_*/\d r=(r^2+a^2)/\Delta$。先进方位角同时含 $\int a\,\d r/\Delta$，故视界处随光线传播的频率为 $k_H=\omega-m\Omega_H$，其中 $\Omega_H=a/(r_+^2+a^2)$ 亦可由 \eqref{eq:kn-generator} 读出。物理解在左端为 $R_s\sim\Delta^{-s}\exp(-\mathrm i k_H r_*)$，在右端为 $r^{-2s-1}\exp(\mathrm i\omega r_*)$。$s=-2$ 时后一幂为 $r^3$，除以先前乘入的 $\sim r^4$ 恰给 $1/r$ 的辐射曲率。实频散射中，视界的 Killing 能量通量的号由 $\omega k_H$ 定，故 $0<\omega<m\Omega_H$ 的波可以反射得比入射还强。此时进视界的角动量与能量都可为负，所转移的 $E-\Omega_HJ$ 却非负，与 \eqref{eq:bh-first-law} 相容。不能将 Schwarzschild 情形的 $|{\cal R}|^2+|{\cal T}|^2=1$ 原样沿用。

Kerr 准正规模的缩放仍写作 $\omega_{\ell mn}=M^{-1}\hat\omega_{\ell mn}(a/M)$。旋转劈开原先对 $m$ 简并的频率；系数 $a/M$ 越大，顺转与逆转分支可相差越多。对于某一被辨认的模式，$Q=\operatorname{Re}\omega/[-2\operatorname{Im}\omega]$ 先给自旋的函数，再由频率给质量。要用多个模检验同一外部几何，必须先分清不同 $(\ell,m)$ 和同一 $(\ell,m)$ 的泛音 $n$；它们的激发振幅由并合过程给定，未必足够大。初始时刻若取太早，源区尚未接近平稳，多个阻尼指数也可能只是在拟合非线性的过渡。

\subsection{由标架确定旋转黑洞的辐射分量}

具体取 Boyer--Lindquist 坐标次序 $(t,r,\theta,\phi)$、$\Sigma=r^2+a^2\cos^2\theta$。一组适配主零方向的逆变矢可写成
\begin{align*}
 l^\mu&=\left((r^2+a^2)/\Delta,\,1,\,0,\,a/\Delta\right),\\
 n^\mu&=\left(r^2+a^2,\,-\Delta,\,0,\,a\right)/(2\Sigma),\\
 m^\mu&=\frac{(\mathrm i a\sin\theta,\,0,\,1,\,\mathrm i/\sin\theta)}
                      {\sqrt2(r+\mathrm i a\cos\theta)}.
\end{align*}
轴上的角坐标奇性与这种明确的标架选择有关。用它投影曲率，可分别构成 $C_{lmlm}$、$C_{lnlm}$、$C_{lm\bar mn}$、$C_{ln\bar mn}$、$C_{n\bar mn\bar m}$ 类型的分量，取一致曲率号数的五个 Newman--Penrose 标量。在主标架中 $\Psi_2$ 描述平稳背景的库仑曲率；$a=0$ 时其量级为 $M/r^3$，远离源即迅速衰减。出射辐射只按 $1/r$ 衰减，尽管在近区比背景小，在渐近区反而是主导的偏离。标架作绕 $l$ 的无穷小零旋转时，$\Psi_4$ 的一阶变化中会混入原有 $\Psi_3$；在背景 $\Psi_3=0$，来自 $\Psi_2$ 的混入须二阶小转角。因此围绕主标架的小旋转不会在一阶把强静态库仑场误当作向外辐射。有限半径、较大的标架失配仍可产生混合。

为何在 \eqref{eq:gw-cpsi} 之前规定无穷远归一？在渐近 Minkowski 空间以一组正交观者基 $e_{(0)},e_{(r)},e_{(\theta)},e_{(\phi)}$，辐射抽取常选 $l=(e_{(0)}+e_{(r)})/\sqrt2$、$n=(e_{(0)}-e_{(r)})/\sqrt2$、$m=(e_{(\theta)}+\mathrm i e_{(\phi)})/\sqrt2$。纯出射量依 $t-r$，故 $(\nabla_{e_{(0)}}-\nabla_{e_{(r)}})h=2\dot h$，两次作用给 $4\ddot h$，与线性曲率的 $1/2$ 和两个零矢的归一配合，得到所写的二阶时间导数。上列 Kinnersley 标架却在远处给 $l\sim e_{(0)}+e_{(r)}$，$n\sim(e_{(0)}-e_{(r)})/2$，相对这一数值抽取标架有一次常数 boost。$\Psi_4$ 因而还差一个常数；将 Teukolsky 的径向解化为实际能流前，必须和用于提波的标架配准。否则所有模的频率虽正确，能量可错一个固定因子。

角向分离也可用微扰检验。令 $q_a=a\omega$ 很小，球谐按自旋权正交归一，将 $S_s={}_sY_{\ell m}+q_a\sum_{\ell'\ne\ell}c_{\ell'}\,{}_sY_{\ell'm}+\cdots$ 代进角方程。线性算符的扰动为 $-2q_a s\cos\theta$，令 $A_s=A_s^{(0)}+q_a A_s^{(1)}$。乘未扰动球谐的共轭并积分，得 $A_s^{(1)}=2s\langle\cos\theta\rangle$。用球面对称性整理乘 $\cos\theta$ 的递推关系，
\[
 \cos\theta\,{}_sY_{\ell m}=c_{\ell+1}\,{}_sY_{\ell+1,m}
 -\frac{sm}{\ell(\ell+1)}\,{}_sY_{\ell m}
 +c_\ell\,{}_sY_{\ell-1,m},\quad
 c_\ell=\sqrt{\frac{(\ell^2-m^2)(\ell^2-s^2)}{\ell^2(4\ell^2-1)}} .
\]
于是 $A_s^{(1)}=-2s^2m/[\ell(\ell+1)]$。非对角元只连到同 $m$ 的相邻 $\ell$。因此数值演化若在球谐基读取某一 $(\ell,m)$，即使其 Kerr 扁球谐激发只含另一相邻 $\ell$，也会因旋转得到可算的一阶渗入。此项随 $a\omega$ 消失，不能误认成一个独立的泛音；对真正的复频模式角向问题也须随 $\omega$ 解析延拓。

还有一项可不求频谱便核对的量，即视界的能量平衡。设实频波从远处入射，$\dot E_H,\dot J_H$ 为穿视界的通量，二者以增加黑洞为正。同一个 $e^{\mathrm i m\phi-\mathrm i\omega t}$ 分量满足 $\dot J_H=(m/\omega)\dot E_H$，适用时 \eqref{eq:bh-first-law} 给
\[
 \frac{\kappa_H}{8\pi}\dot A_H
 =\dot E_H\left(1-\frac{m\Omega_H}{\omega}\right)\ge0 .
\]
右侧即使在超辐射区也为正。若在视界同转阈值 $\omega=m\Omega_H$ 两侧各取一个频率，可发生能量通量反号，而面积通量不必反号。用实标量波作最简例，若无穷远纯单位入射振幅、视界透射振幅 ${\cal T}$ 均按 Klein--Gordon 径向流归一，Wronskian 直接给 $1-|{\cal R}|^2=(\omega-m\Omega_H)|{\cal T}|^2/\omega$。在 $m=0$ 或非转极限还原通常的守恒式。自旋二的曲率径向函数本身有 $\Delta$ 与幂次，式中不能把未归一的 Teukolsky 系数直接当作标量透射系数。

\section{由两体轨道接到并合}\label{sec:gw-cbinary}

\subsection{有效单体的能量字典}

\S\ref{sec:pn-wave-energy} 给出了以圆轨道频率为变量的保守能量与能流。接近最后稳定轨道时，慢速级数的少数几项往往无法确定接下来是继续绕转还是快速跌入。已有一个确知任意轨道速度的极限：当一方质量趋于零时，它在另一方的背景上作测地运动。可否保留这个极限，再逐阶打开有限质量比的修正？

令 $M=m_1+m_2$、$\mu=m_1m_2/M$、$\nu=\mu/M$。有效单体法 (effective one body, EOB) 将真实相对运动的两个作用变量 $J_\phi=\oint p_\phi\,\d\phi/(2\pi)$、$J_r=\oint p_r\,\d r/(2\pi)$ 与质量为 $\mu$ 的有效粒子配对。它们在正则变换下不变，故可在相同的 $(J_r,J_\phi)$ 处比较能量，无需认同两套径向坐标。例如两个自由粒子在质心系的总能量为 $E$，相对 Lorentz 因子由不变量 $(p_1+p_2)^2=-E^2$ 给出 $\gamma_{\rm rel}=(E^2-m_1^2-m_2^2)/(2m_1m_2)$。令 $H_{\rm eff}/\mu$ 对应这一比值，解出
\begin{equation}
 H_{\rm real}=M\sqrt{1+2\nu(H_{\rm eff}/\mu-1)} .
 \label{eq:gw-ceob}
\end{equation}
延拓到相互作用问题后，还须和已知的 PN 作用量匹配以确定有效势。若记 $E_b=H_{\rm real}-M$，则反解正好是 $H_{\rm eff}-\mu=E_b+E_b^2/(2M)$。可见有限质量比的修正有一部分进了能量的定义，不能只改 Schwarzschild 度规中的质量。

对无自旋情形写 $\d s_{\rm eff}^2=-A\,\d t^2+B\,\d r^2+r^2\d\Omega^2$，以 $u=M/r$、$j=J_\phi/(M\mu)$、$\hat p_r=p_r/\mu$ 记无量纲变量。到有效运动仍可表为测地运动的阶数，有 $(H_{\rm eff}/\mu)^2=A(1+j^2u^2+\hat p_r^2/B)$。匹配到 2PN 得 $A=1-2u+2\nu u^3+O(u^4)$、$A B=1-6\nu u^2+O(u^3)$。3PN 的 $A$ 还含 $(94/3-41\pi^2/32)\nu u^4$，并须在一般偏心运动中引入四次动量项。圆轨道上径向动量为零，故先从圆轨道比较 $A$，却不能仅凭它确定所有非圆运动。有效度规用于编码双星动力学，未直接规定从源到望远镜的光路。

在圆轨道上固定 $j$ 极小化有效能量，$A'(1+j^2u^2)+2A j^2u=0$，此处撇号表示 $\d/\d u$。由此可消去角动量，得
\[
 j^2=-\frac{A'}{u(2A+uA')},\qquad
 \left(\frac{H_{\rm eff}}{\mu}\right)^2=\frac{2A^2}{2A+uA'}.
\]
取 $A=1-2u$，回到 $j^2=1/[u(1-3u)]$。沿圆轨道序列最后的稳定条件为 $\d j^2/\d u=0$，故 $u=1/6$，而 $2A+uA'=0$ 的 $u=1/3$ 是光子轨道。对有限质量比，两条件分离的结构可在保留 $A$ 整体形式时继续使用。频率须用真实 Hamilton 量，链式求导带来 $\partial H_{\rm real}/\partial H_{\rm eff}=M/H_{\rm real}$，所以 $(M\Omega)^2=-\frac{u^3 A'(u)}{2}\frac{M^2}{H_{\rm real}^2}$。这给了一个可检验的不变量：以 $\Omega$ 消掉 $u$，所得 $E_b(\Omega)$ 在慢速展开时须回到 \S\ref{sec:pn-wave-energy} 的结果。若直接用测地线的有效时钟求频率，有限质量比的这一因子会漏掉。

高阶近似有时用有理函数代替多项式。以 $F(x)=1+a x+b x^2+\cdots$ 为例，若 $a\ne0$，$[1/1]$ Padé 近似为 $[1+(a-b/a)x]/[1-(b/a)x]$。乘回分母即可验到 $x^2$ 的系数完全相同。有理式在 $x=a/b$ 可有极点，原函数却未必在那里奇异；因此只能用测试质量极限、有无伪极点、数值轨道的相位等独立条件选择重求和。非微扰写法保留了最后稳定轨道这样的结构，并未凭空给出新的 PN 系数。

\subsection{用轨道不变量核对有效势}

作一个实际的作用量积分。设 $E_{\rm N}<0$ 为 Newton 束缚能、$k=M\mu$、$J$ 为角动量。两个径向转折点为 $r_p,r_a$，则
\[
 J_r=\frac{1}{\pi}\int_{r_p}^{r_a}
    \sqrt{2\mu(E_{\rm N}+k/r)-J^2/r^2}\,\d r
    =\frac{k\mu}{\sqrt{-2\mu E_{\rm N}}}-J .
\]
可取 $r=a(1-e\cos\chi)$，先把根号写为 $\sqrt{-2\mu E_{\rm N}}\sqrt{(r-r_p)(r_a-r)}/r$。分子、$\d r$ 各给一个 $\sin\chi$，积分中出现 $\int_0^\pi\sin^2\chi\,\d\chi/(1-e\cos\chi)=\pi[1-\sqrt{1-e^2}]/e^2$，即得上式。因此 $E_{\rm N}=-\mu k^2/[2(J+J_r)^2]$。只依赖两作用量之和使 $\Omega_\phi=\partial E/\partial J$ 与 $\Omega_r=\partial E/\partial J_r$ 相等，故轨道闭合。相对论修正使能量分别依赖二者，$2\pi(\Omega_\phi/\Omega_r-1)$ 直接给每径向周期的进动。这个比值适于匹配两套坐标完全不同的动力学。

有效球对称度规也可直接给近圆轨道的频率比。令能量和角动量每单位有效质量为 ${\cal E},{\cal J}$，径向固有时满足 $B\dot r^2={\cal E}^2/A-1-{\cal J}^2/r^2$。圆轨时右侧和一阶导数都为零，消掉常数得
\[
 {\cal E}^2=\frac{2A^2}{2A-r A_r},\qquad
 {\cal J}^2=\frac{r^3 A_r}{2A-r A_r},
 \quad A_r=\frac{\d A}{\d r}.
\]
在固定 ${\cal E},{\cal J}$ 下将右侧展开到位移平方，再以 $\d t_{\rm eff}/\d\tau={\cal E}/A$ 化到有效坐标时，得
\[
 \Omega_{\phi,\rm eff}^2=\frac{A_r}{2r},\qquad
 \Omega_{r,\rm eff}^2=\frac{1}{2B}\left(A_{rr}
              -\frac{2A_r^2}{A}+\frac{3A_r}{r}\right).
\]
由 \eqref{eq:gw-ceob}，同一圆轨上两个真实频率均乘 $M/H_{\rm real}$，比值不变。取 Schwarzschild 势，$\Omega_r^2/\Omega_\phi^2=1-6M/r$，零点正给稳定性的丧失；若保留有限 $\nu$，频率比中显含 $B$。所以离心率虽小但能辨认出径向周期的轨道，提供圆轨结合能不具备的一项信息。引入更多非测地的动量项时，也可用离心率的不同阶数分别确定它们。

再亲自作一次频率的消元。以 2PN 的 $A=1-2u+2\nu u^3$ 代前式，$H_{\rm eff}/\mu=1-u/2+3u^2/8+(27/16-\nu/2)u^3+\cdots$。经能量映射，$E_b/\mu=-u/2+(3-\nu)u^2/8+(27-5\nu-\nu^2)u^3/16+\cdots$，尚含坐标半径。令 $x=(M\Omega)^{2/3}$，可测频率的关系为
\[
 x=u+\frac{\nu}{3}u^2+
       \left(\frac{2\nu^2}{9}-\frac{5\nu}{4}\right)u^3+\cdots,\qquad
 u=x-\frac{\nu}{3}x^2+\frac{5\nu}{4}x^3+\cdots .
\]
先在 $u^3$ 阶发生的 $\nu^2$ 相消才得到右式。把它代回 $E_b$，$x^2$ 的系数为 $3/8+\nu/24$，$x^3$ 的系数为 $27/16-19\nu/16+\nu^2/48$，逐项回到 \S\ref{sec:pn-wave-energy} 的能量。这说明即使真实与有效的波形相位最终以相同 $x$ 标记，也不能在中途便把 $u$ 与 $x$ 认同。

还可将上述截断的势当成独立模型看其最后稳定轨道的轻微移动。展开角动量得 $j^2=1/[u(1-3u)]+\nu\,u(4u-3)/(1-3u)^2+O(\nu^2)$。若 $u_{\rm LSO}=1/6+\nu\,\delta u$，以 $\partial_u j^2=0$ 求解，给 $\delta u=19/648$。这里仅采用 $2\nu u^3$ 的 2PN 形变；所有更高阶势在 $u\simeq1/6$ 上仍会给 $O(\nu)$ 的平移，故不能把这个有理数视作自力理论对精确末轨道频率的预言。而且径向坐标可以随模型定义改变，观者若要比较应先以上述能量映射转成频率。

同一算例说明重求和能保持哪些性质。对三次势取分子一次、分母二次并配到 $u^3$，得 $A_{[1/2]}(u)=\frac{1-(2-\nu/2)u}{1+\nu u/2+\nu u^2}$。当 $0\le\nu\le1/4,u>0$ 时分母无零点，却有一个正 $u$ 使分子为零；$\nu\to0$ 时正好回到 Schwarzschild。若改用同一系数的单纯 Taylor 多项式，在足够大的 $u$，$2\nu u^3$ 会再次压过 $-2u$，给出未受慢速信息约束的行为。检验一种有效波形时，应分开比较已知的展开系数、强场定性的边界，以及与数值相位和振幅的一致性。三者所利用的信息各不相同。

不必事先选择一种有效度规也可作有物理根据的重求和。取随粒子交换对称的能量变量 $\epsilon_E=(H_{\rm real}^2-m_1^2-m_2^2)/(2m_1m_2)$，再定义 $e_E=\epsilon_E^2-1$。测试质量极限已知有限速度的圆轨道结果为 $\epsilon_E=(1-2x)/\sqrt{1-3x}$，平方减一即 $e_E(x,0)=-x\,\frac{1-4x}{1-3x}$。这一步把光子轨道处的平方根奇性换成简单极点。对有限 $\nu$，以上面的 $E_b(x)$ 代回定义，得到 $e_E=-x[1-a x-b x^2+O(x^3)]$，其中 $a=1+\nu/3$，$b=3-35\nu/12$。只需以最简单的有理式重写方括号，便有 $e_E^{P}=-x\,\frac{1-(a+b/a)x}{1-(b/a)x}$。当 $\nu=0$ 时 $a=1,b=3$，已知到 $x^3$ 的系数竟能还原精确有理函数；原因是先用测试质量的光子轨道奇性选对了变量与阶数。有限 $\nu$ 时有理式预言极点 $x_p=a/b$，可由极值 $\d e_E^P/\d x=0$ 估末轨道。令 $d=b/a$、$c=a+d$，求导的分子为 $1-2cx+cdx^2$，较小正根为 $[1-\sqrt{1-d/c}]/d$。测试质量取 $c=4,d=3$ 得 $1/6$；有限质量比的极值只是这一近似预言，可用数值能量--频率曲线检验。故重求和的可信部分来自选变量前已知的结构，数学上任意一个 Padé 极点都不能独立证明强场确有奇性。

\subsection{辐射反作用及波形的缝合}

把切向耗散力矩 ${\cal F}_\phi$ 加到 $\dot p_\phi$，径向力记为 ${\cal F}_r$。Hamilton 方程给 $\dot H_{\rm real}=\dot r\,{\cal F}_r+\Omega{\cal F}_\phi$。准圆阶段 $\dot r$ 小，远区能流 $L_{\rm GW}$ 与角动量流之比逼近 $\Omega$，于是 ${\cal F}_\phi\simeq-L_{\rm GW}/\Omega$。这条近似还可从旋转模的 $\omega/m$ 得出；明显偏心或快速跌入时径向做功与近区能量储存均可参与。求出 $r(t),\phi(t)$ 后，对每一辐射模构造 $h_{\ell m}$。常把它分成最低阶波形、有效源、曲率尾、相位修正与剩余振幅的乘积，将剩余振幅写为 $(\rho_{\ell m})^\ell$。例如若真实的剩余因子为 $1+\ell a x+O(x^2)$，取 $\ell$ 次根后才是 $1+a x+O(x^2)$，按模阶数增加时较易控制。跌入阶段的非圆修正可由峰附近振幅、斜率及频率与数值解匹配。

引入 $H_{\ell m}=r h_{\ell m}$，并取 $h_+-\mathrm i h_\times$ 的自旋球谐展开。球谐正交及有效应力式 \eqref{eq:gw-effective-stress} 给
\begin{equation}
 L_{\rm GW}=\frac{1}{16\pi}\sum_{\ell m}|\dot H_{\ell m}|^2,\qquad
 {\cal G}_z=\frac{1}{16\pi}\sum_{\ell m}m\,
                   \operatorname{Im}(H_{\ell m}\dot H_{\ell m}^*) .
 \label{eq:gw-cmodes}
\end{equation}
若 $H_{\ell m}\propto e^{-\mathrm i m\phi}$ 且慢变振幅，立刻有 $L_{\rm GW}=\Omega{\cal G}_z$。轴向角动量流来自相位，不能只由总能量和一个跨越整个频带的常数频率反推。式中已对全部正负 $m$ 求和。对有轨道面反射对称性的双星，可从正 $m$ 重建负 $m$；若只用正 $m$，能流中须相应补另一半。

当外部时空已接近一个 Kerr 黑洞，以晚时振铃接替非线性两体波形。设匹配时刻 $t_0$、主模复振幅及其导数为 $H_0,\dot H_0$，取两个准正规模 $H(t)=C_1 e^{s_1(t-t_0)}+C_2 e^{s_2(t-t_0)}$，$s_j=-\mathrm i\omega_j$。由值和斜率连续得 $C_1=(\dot H_0-s_2H_0)/(s_1-s_2)$、$C_2=(s_1H_0-\dot H_0)/(s_1-s_2)$。只用一个模时二者之比被完全固定，难以在峰附近同时吻合。用更多模可在一小段内作最小二乘；若两个 $s_j$ 太接近，系数的分母小，微小数值噪声也会被放大。合成信号须经过数值解校验，以免把匹配处的人为误差当成额外的物理振荡。

对有自旋的两体，中心背景可改作形变的 Kerr，另以有效自旋与轨道、两自旋之间的耦合校准慢速极限。几何制下 $\bm\chi_i=\bm S_i/m_i^2$，投影组合 $\chi_{\rm eff}=(m_1\bm\chi_1+m_2\bm\chi_2)\cdot\hat{\bm L}/M$ 反映最低几阶相位中较易辨认的对准成分。自旋与轨道角动量同向时，最后稳定的顺转轨道可更深入，临并合前多转几圈；反向时相反。若有面内成分，$\bm L$ 进动，可在随轨道面转动的基中计算主辐射，再用旋转矩阵将各 $m$ 混合到天球基底。时域振幅和偏振的慢调制由此伴随进动而来；只把 $\chi_{\rm eff}$ 放入无进动相位不能制造这项可测信息。

\subsubsection{把尾部和自旋各放在哪一项}

对准圆轨道的一个 $\ell,m$ 辐射模，常写成
\[
 h_{\ell m}=h_{\ell m}^{({\rm N})}\,\hat S_{\rm eff}\,
                T_{\ell m}\,e^{\mathrm i\delta_{\ell m}}\,
                \rho_{\ell m}^{\,\ell}\,N_{\ell m}.
\]
首因子取最低阶幅相，$\hat S_{\rm eff}$ 表示相应宇称的有效源（偶型可取约化有效能量，奇型可用适当归一的角动量）。$T_{\ell m}$ 对曲率多次散射的领头对数作重求和，$\delta_{\ell m}$ 收剩余相位，$\rho_{\ell m}$ 收剩余振幅，$N_{\ell m}$ 给偏离圆轨道的修正。对正 $m$ 令 $k=m\Omega$、$\hat k=H_{\rm real}k$，一种尾部归一为
\[
 T_{\ell m}=\frac{\Gamma(\ell+1-2\mathrm i\hat k)}{\Gamma(\ell+1)}
       e^{\pi\hat k}\exp[2\mathrm i\hat k\ln(2kr_0)] .
\]
人为尺度 $r_0$ 的依赖与相位、时间原点的约定配合。用 $\Gamma(1+\mathrm i y)\Gamma(1-\mathrm i y)=\pi y/\sinh(\pi y)$ 逐次递推得
\[
 |T_{\ell m}|^2=\frac{4\pi\hat k}{1-e^{-4\pi\hat k}}
                  \prod_{j=1}^\ell\left(1+\frac{4\hat k^2}{j^2}\right).
\]
小 $\hat k$ 时右端为 $1+2\pi\hat k+O(\hat k^2)$。四极 $\hat k\simeq2x^{3/2}$，故能流的相对修正为 $4\pi x^{3/2}$，正与 \S\ref{sec:pn-wave-energy} 的尾项相同。将完整尾因子乘进波形后，剩余振幅不能再保留这一相同的尾修正，否则即使每个单项看似都有物理来源，合计也会重算一次。

非圆修正可从哪种信息定出？在径向乌龟动量 $p_{r_*}$ 及轨道角动量 $p_\phi$ 给出的试探波形上，令 $v_r=p_{r_*}/(\mu r\Omega)$，并写 $N_{22}=(1+c_A v_r^2)\exp(\mathrm i c_\phi v_r)$ 作最低示例。足够多的 PN 圆轨系数都不能限定 $c_A,c_\phi$，因为圆轨 $v_r=0$。在匹配点 $t_m$ 已知数值振幅 $A_m$、原试探振幅 $A_0$ 时，如 $v_r(t_m)\ne0$ 可求 $c_A=(A_m/A_0-1)/v_r^2(t_m)$；匹配瞬时频率则由 $\dot v_r(t_m)$ 而非 $v_r$ 决定相位系数。若把 $t_m$ 取在早期几乎完全圆轨的部分，分母小，校准便很不稳定。实际各模还可同时比较峰时、曲率或斜率，用一族独立模拟核对插值。

自旋匹配的具体字典可从最低阶自旋--轨道 Hamilton 量读出。在几何制定义 $\bm S=\bm S_1+\bm S_2$、$\bm S_*=(m_2/m_1)\bm S_1+(m_1/m_2)\bm S_2$，则领头耦合为 $H_{\rm SO}^{\rm LO}=\frac{\bm L}{r^3}\cdot (2\bm S+\tfrac32\bm S_*)$。第一组合可进入有效背景的 Kerr 参数 $\bm a=\bm S/M$，第二组合承担有效粒子自旋的作用。测试质量极限若大洞有自旋、轻粒子无自旋，$\bm S_*$ 中来自大洞的贡献按 $m_2/m_1$ 抑制，保留了 Kerr 拖曳；反之轻天体的 $\bm S_2$ 乘了 $m_1/m_2$，使其与轨道的测地进动保持正确的质量阶。有限 $\nu$ 的更高 PN 自旋字典还可含 $\hat{\bm p}^{\,2}$、$(\hat{\bm r}\cdot\hat{\bm p})^2$ 和 $M/r$，这些量受正则坐标选择影响，仅按两质量加权的一个常自旋一般不够。有效 Kerr 例如可取 $\Delta_t=r^2 A(M/r)+a^2$、$\Delta_r=\Delta_t/(AB)$ 来区分角向和径向势，两者在 $\nu\to0$ 时才都回到 Kerr 的 $\Delta$。

一个有意令 $\bm S$ 为零的算例能显示这种分工。设 $m_1=2m_2$ 且 $\bm S_1=-\bm S_2=\bm S_a$，则 $\bm S=0$ 而 $\bm S_*=-3\bm S_a/2$，仍有 $H_{\rm SO}^{\rm LO}=-9\bm L\cdot\bm S_a/(4r^3)$。若只依总自旋决定所有动力学，就会把这对星误当作无自旋。若两自旋等量反向且质量也相等，则两组合均消失于此阶，但从下文的反冲计算可见，面内反向自旋仍可通过辐射的动量通量影响遗留黑洞。保守轨道耦合为零与波的动量分布对称并非同一条件。

\subsection{从有限网格到远方波形}

在 \S\ref{sec:Hamilton} 的 $3+1$ 初值问题中，先在空间片上解约束再演化。以相同物理单位比较初始总能、末态黑洞、无穷远辐射，才可从两颗视界在网格上相遇的现象得到一个物理波形。给初值的一种做法是取 $\gamma_{ij}=\psi^4\tilde\gamma_{ij}$、$K=0$，并用 $\tilde A^{ij}=\psi^{10}K^{ij}$ 表示外曲率无迹部。在共形平直真空片，动量约束为 $\partial_j\tilde A^{ij}=0$，Hamilton 约束变成 $\nabla^2\psi=-\tfrac18\psi^{-7}\tilde A_{ij}\tilde A^{ij}$。时间反演对称时右端为零，$\psi=1+\sum_A m_A/(2|\bm x-\bm x_A|)$ 在去掉的穿刺点外就是解。这里的 $m_A$ 是初值参数，有限间距的每个视界质量还含伴星引力影响。非零动量和自旋使右端回到非线性 Poisson 方程，可将已知奇部剥出，只数值求剩下的正则函数。它给的是某个合法的空间片，未必已含远过去的旋近辐射；演化初期的无关辐射脉冲应与天体信号分开。

直接差分一套约束良好的方程未必能稳定地演化它。BSSN 形式另演化共形度规、外曲率的迹与无迹部，以及收缩的共形联络函数；用约束改写若干导数，有利于控制不希望的约束模。广义谐和形式则令坐标满足带给定源函数的波方程，将 Einstein 方程主部化为度规分量的波算符，并可给偏离谐和条件的自由度加阻尼。移动穿刺法选适当 lapse 与 shift 使坐标绕开奇点，切除法在视界内挖去一小区，选择使全部特征向洞内流出的内边界。两种做法均须查物理和规范的特征传播，不能仅凭网格内最黑的一点决定边界条件。

在大半径的多个球面采样 \eqref{eq:gw-cpsi}，先按角向分解，得到 $r\Psi_{4,\ell m}$。把各球面的时间换成同一推迟时间，固定该时间逐项按 $1/r$ 外推可估计零无穷远的波；若取半径过小或坐标时不合适，外推也含近区项。另一办法是在有限世界管上给特征初边值资料，沿出射零面演化到零无穷远。再把曲率积分两次还原应变，直接时间积分易把一个极小的直流误差变成二次漂移；频域除以 $-\omega^2$ 也会强烈放大极低频初始误差。因而需利用无辐射的初始条件、合适的低频处理及独立的距离外推核对，特别是关心记忆时，任意滤去低频就可能抹掉信号。

质量守恒给出独立检查。以第 \ref{chpt:BH} 章的视界力学由末态面积和角动量求 $M_{\rm irr}=\sqrt{A_H/(16\pi)}$、$M_f^2=M_{\rm irr}^2+J_f^2/(4M_{\rm irr}^2)$。若初始 ADM 三动量为零，辐射的能、动量为 $E_{\rm rad},\bm P_{\rm rad}$，末态满足 $(M_{\rm ADM}-E_{\rm rad})^2-|\bm P_{\rm rad}|^2=M_f^2$。故较大反冲时应与运动黑洞的总能比较；直接将两个质量相减略去了其动能。ADM 能量取在空间无穷远，始终计入尚在向外走的波。Bondi 能量取在出射零无穷远的各截面，随波通过而减小，截面间差即辐射通量。这两个量取极限的方向不同。

例如初始单黑洞的共形因子取 $\psi=1+M/(2r)$，用第一章给的 ADM 球面积分，$\gamma_{ij}=\psi^4\delta_{ij}$ 使被积项 $\partial_j\gamma_{ij}-\partial_i\gamma_{jj}=-2\partial_i\psi^4$。在大半径球上 $r^2\partial_r\psi^4\to-2M$，故能量恰得 $M$。若从作用量构造同一能量，固定边界诱导度规的变分会有法向导数项；加入外曲率边界项并扣除同一诱导几何的平直参考值才能定义有限的 Hamilton 量。约束在壳上为零，时间平移的值由外边界贡献，正与此例相符。平稳区若有类时 Killing 矢，还可由它构造 Komar 质量。

数值误差也可用可观测量检验。对格距 $\Delta,\Delta/2,\Delta/4$ 的波形，在同一频率处比较相位，若 $\phi_\Delta=\phi_0+C\Delta^p+\cdots$，两次相邻相位差之比趋近 $2^p$。在不同分辨率同时拟合任意的额外相位漂移，会掩盖误差；通常只消去时刻与整体相位的约定，再比较所关心频带。真空两体问题只含总质量一个有量纲尺度，一次给定质量比、自旋的无量纲演化可伸缩到许多源质量。中子星有独立的密度、温度、微观反应尺度，这个便利随即消失。

\subsection{固定边界几何时的能量}

上面的质量核对为何要到边界去算？考虑有一段有限外壁的引力作用量。令 $\epsilon=n^\mu n_\mu=\pm1$ 表示边界法矢的号数、$q_{\mu\nu}=g_{\mu\nu}-\epsilon n_\mu n_\nu$ 为诱导度规、$\Theta=q^{\mu\nu}\nabla_\mu n_\nu$。在 \S\ref{sec:Einstein} 的 Hilbert 作用量上加 Gibbons--Hawking--York 项，取外向法矢，得
\[
 S_g=\frac{1}{16\pi}\int_{\cal V} R\sqrt{-g}\,\d^4x+
        \frac{1}{8\pi}\int_{\partial{\cal V}}\epsilon\Theta
                                   \sqrt{|q|}\,\d^3x .
\]
这里先取光滑的非零边界，或固定各段相遇角的正交外壁。由 $\delta R_{\mu\nu}$ 的体积分变成的散度含法向导数 $\nabla_n\delta q$，因此只固定诱导度规不会使原 Hilbert 项的边变分消失。可是适配法向高斯坐标中 $\Theta_{ab}=\tfrac12{\cal L}_nq_{ab}$，变分 $\sqrt{|q|}\Theta$ 恰有含相同导数的项。二者相加后，剩下的边变分只正比于 $(\Theta_{ab}-\Theta q_{ab})\delta q^{ab}$，故固定 $q_{ab}$ 足够。可以另外加任何仅依赖固定边界资料的参考泛函来选能量零点；在渐近平直问题中取平直背景的那一份使空时能量为零。

在一个静态片令 shift 为零，选边上正交静止观者的固有时为演化时间。球面上定义 $k=D_i s^i$，以相同内禀二面几何嵌入平直三空间给参考 $k_0$。由边界 Hamilton 项得到此类观者的准局域能量 $E_{\rm BY}=(8\pi)^{-1}\oint(k_0-k)\,\d A$。以 Schwarzschild 面积半径 $R$ 的球作算例。空间径向度规为 $\gamma_{RR}=1/F$，故 $s^R=\sqrt F$，$k=2\sqrt F/R$，而参考球为 $2/R$。于是 $E_{\rm BY}(R)=R\big(1-\sqrt{1-2M/R}\big) =M+\frac{M^2}{2R}+O(R^{-2})$。这里每个有限球面的边界观者不同，不能据此断言静态外部有实际能量向外流过。取 $R\to\infty$ 时才与渐近时间平移的 Hamilton 量一样得到 $M$。若未作参考扣除，单个 $k$ 积分会按 $R$ 发散，连 Minkowski 空间也有同样的问题。对一般非凸二面如何选参考较细致，这个球对称例子没有相关歧义。

对没有波向外流的平稳时空，令 $\xi^\mu$ 为无穷远归一的时间 Killing 矢。用 $\nabla_\mu\nabla^\mu\xi^\nu=-R^\nu{}_\lambda\xi^\lambda$ 及 Stokes 定理，可将两个二面的 Komar 质量差化为中间区域的 Ricci 曲率积分。真空中两者相等，故 Schwarzschild 取半径 $R$ 可直接用 $M_K=R^2 F'(R)/2=M$。与刚才的 $E_{\rm BY}(R)$ 比较，即使对完全相同的二面，不同定义也可给不同的有限半径值；前者明确用了延伸到整个外区的 Killing 矢，后者用了边界静止观者和参考嵌入。数值并合一般无全局类时 Killing 矢，故不能把这个简单真空等式用于一串演化中的球面。可操作的能量核对是初片 ADM、遗留视界及零无穷远辐射三方的比较。

能量守恒也能分辨不同半径之间的波包。例如取一段持续时间 $\tau_b$ 的纯出射波，发射后的第一个空间片在外球 $r_2$ 尚未见波，内球 $r_1<r_2$ 已见到前半段。如果在这两个球相同 $t$ 的截面比较总能，会把两球之间正在传播的能量误作误差。应改以相同 $u=t-r_*$ 标记各截面，并把从初始到该截面的通量积分计入；用 Bondi 通量时还应统一零无穷远的时间切面与角向标架。不同提波半径在相同 $u$ 仍有 $O(1/r)$ 的差，可以作外推，随 $u$ 有限传播的差却不能当普通的拟合系数消去。

\subsection{在空间片上找视界和辐射}

怎样判断网格中已形成共同黑洞？取空间片的未来单位法矢 $n^\mu$，闭二面的外向单位法矢 $s^i$，沿用第一章 $K_{ij}=\tfrac12{\cal L}_n\gamma_{ij}$ 的正号。由该章的外向零膨胀公式 $\theta_{(+)}=D_i s^i+K-K_{ij}s^i s^j$，可在网格上确定闭面。在给定一片上搜索 $\theta_{(+)}=0$ 的最外闭面就是找表观视界。可把试探面写成 $r=h(\theta,\phi)$，由 $s_i\propto D_i[r-h(\theta,\phi)]$ 归一化，迭代这个二维方程。所得视界位置随切片变化，坐标半径本身不代表不变量。第一次发现共同表观视界也不是事件视界的定义；后者还需知道此后哪些光线始终到不了无穷远。

用上一节的时间对称单穿刺作检验。$K_{ij}=0$、$\psi=1+M/(2r)$，某一坐标球的外法矢只有 $s^r=\psi^{-2}$，且 $\sqrt\gamma=\psi^6 r^2\sin\theta$。于是 $D_i s^i=[\partial_r(\psi^4 r^2)]/(\psi^6 r^2)$，视界条件等价于 $1+2r\psi'/\psi=0$，解为 $r=M/2$。它的面积半径为 $\psi^2 r=2M$，面积 $16\pi M^2$ 正还原前面的不可约质量，而坐标半径差了四倍。若双黑洞已靠近，同一片可能仍找到两个内面和一个包住二者的外面，给遗留者算面积要明确取哪一个。

若二面近轴对称，可选择在该面上近似 Killing 的切向矢 $\varphi^i$，其闭合轨道参数周期归一为 $2\pi$。以本书的外曲率号数，面角动量取 $J_H=-(8\pi)^{-1}\oint (K_{ij}-K\gamma_{ij})\varphi^i s^j\,\d A$，方向由 $\varphi^i$ 约定。$\gamma_{ij}\varphi^i s^j=0$，迹项自动消失。空间无穷远的渐近转动生成元也给类似积分，多黑洞相距很远时可同时估每个洞的自旋和总角动量；总量还含轨道角动量及传播中的场。若两洞在强烈相互作用区没有明确的各自轴对称，局部自旋的定义近似性必须与数值误差分开。临近静态 Kerr 末态，视界给的 $J_H,M_f$ 可和无穷远振铃推得的量独立核对。

提取 $\Psi_4$ 不必在网格上对时空度规作两次时间差分。定义 Weyl 张量的电、磁部 $E_{ij}=C_{i n j n}$、$B_{ij}={}^\star C_{i n j n}$，其中 $n$ 表示沿片法矢的投影。真空时，Gauss 关系给 $E_{ij}={}^{(3)}R_{ij}+K K_{ij}-K_{ik}K^k{}_j$。其迹恰为 Hamilton 约束。磁部是 $K_{ij}$ 的空间旋度，具体正负依空间体积形式及左对偶的方向；对固定方向可由 Codazzi 关系求得。取切于抽取球面的两条正交角矢 $e_\theta,e_\phi$，先把 $E,B$ 投影到它们组成的横向无迹两分量，再按向外零矢的定义组合为 $\Psi_4$。用已知沿径向传播的 $+$、$\times$ 平面波可同时校准号数和标架归一。空间曲率和外曲率已在演化中求得，这样便避开了跨许多步再差分一个小二阶信号的误差。

网格上益处明显的另一处理是拆开体积与形状。取 $\gamma_{ij}=e^{4\phi}\tilde\gamma_{ij}$、$\det\tilde\gamma=1$，则 $\phi=\ln\det\gamma/12$。再令 $\tilde A_{ij}=e^{-4\phi}(K_{ij}-\gamma_{ij}K/3)$，得 $\tilde\gamma^{ij}\tilde A_{ij}=0$。若 $\partial_t\gamma_{ij}=2\alpha K_{ij}+{\cal L}_{\bm\beta}\gamma_{ij}$，求迹给
\[
 (\partial_t-\beta^k\partial_k)\phi
       =\frac{\alpha K+\partial_i\beta^i}{6}.
\]
求 $\partial_t(e^{-4\phi}\gamma_{ij})$ 时其中 $\alpha K$ 项恰消，剩下 $2\alpha\tilde A_{ij}$ 与去掉体积膨胀的 shift 导数。这就是共形拆分使平均压缩与剪切分开演化的代数。另演化 $\tilde\Gamma^i=-\partial_j\tilde\gamma^{ij}$ 时，其与度规导数的差本身是一个须收敛到零的约束；对演化变量施加 $\det\tilde\gamma=1$ 不能代替检查原有 Hamilton、动量约束。不同数值代码若用相反 $K$ 号数，lapse 的驱动项也须一并换号。

初值和边界产生的误差如何影响相位？设抽取球离人工外边界的光行时间为 $T_{\rm b}$，若边界不完全吸收，一次反射最早在相应的因果回路时间后污染波形。扩大外边界和提升空间分辨率检验的是两类误差，不能用相同网格上的守恒量替代。若以源的特征长度归一，将外推后的主模作 $H=A(t)e^{-\mathrm i\phi(t)}$，瞬时频率 $\omega=-\operatorname{Im}(\dot H/H)$ 在振幅接近零时失去稳定性，应避开模节点再比较。按 \S\ref{sec:gw-b-match} 的失配计算，比较两模型时应先消去频域相位差中能被常相位和时移吸收的部分。有些线性的差异只是所选零点，某一短时高频误差在一种总质量上无关紧要，伸缩到另一质量后可进入最敏感频带。

\subsection{反冲和遗留黑洞的去向}

由 \eqref{eq:gw-effective-stress}，各方向能流乘传播方向便给线动量流。在几何制有
\[
 \frac{\d P^i_{\rm rad}}{\d u}
 =\frac1{16\pi}\int n^i\left|
       \sum_{\ell m}\dot H_{\ell m}\,{}_{-2}Y_{\ell m}\right|^2\d\Omega .
\]
能量积分使不同模正交；线动量多了 $n^i$，角积分可使相邻多极相干。等质量、无自旋时，对称性使各方向与其相反方向恰好配对，故反冲为零。若交换两无自旋天体只是旋转半周，净反冲的有符号分量在 $m_1\leftrightarrow m_2$ 下变号。记 $\delta=(m_1-m_2)/M$，则 $\delta^2=1-4\nu$。又在小质量比极限波幅为 $O(\nu)$，强辐射阶段持续时间为 $O(M)$，反冲速度的最低拟合因子可以取 $\nu^2\sqrt{1-4\nu}$。其最大值条件 $2/\nu-2/(1-4\nu)=0$ 给 $\nu=1/5$。实际峰位依高阶动力学而移，等质量零点则是精确对称性的结果。

对等质量而自旋相反、位于轨道面内的双黑洞，轨道面对称性可破缺，主要反冲甚至沿轨道法线。其方向和大小依面内自旋在并合瞬间相对分离矢量的方位。瞬时通量方向持续旋转，不同周期间动量大部相消，最大累积通常发生在晚期；若振铃仍沿相反方向放出一部分动量，最后速度可低于积分途中达到的峰值。用瞬时最大功率除 $c$ 再乘整个旋近寿命，既漏掉方向也漏掉相消。对星团或小星系，保留遗留者的条件近似为 $v_{\rm kick}<v_{\rm esc}$，$v_{\rm esc}$ 取当地势阱；即使未逃逸，往返星系核的过程也会延迟下一次并合。


可把上述对称性化为一个具体的 PN 算例。令无自旋双星绕 $xy$ 面运行，$x=GM/(ac^2)$，$\hat{\bm r},\hat{\bm\lambda}$ 为径向和切向单位矢。把质量四极与质量八极、流四极的交叉项代入线动量通量，圆轨道的领头项沿 $\hat{\bm\lambda}$，其模为 ${\cal F}_P=\frac{464}{105}\frac{c^4}{G} \nu^2|\delta|\,x^{11/2}$。为何总功率主要是四极平方，反冲却要看高一阶？线动量通量比能流多一个反演为负的 $n^i$，质量四极的平方在相反方向相等。要得到非零积分，须乘一个宇称不同或角动量不同而能满足三角选择规则的角分量。质量八极的源矩在无自旋双星质心系包含 $\mu(m_2-m_1)/M$，与多极所预期的质量比因子一致。若准圆参数在一转内几乎不变，则 $\d\hat{\bm r}/\d t=\Omega\hat{\bm\lambda}$，通量的时间原函数近似为 ${\cal F}_P\hat{\bm r}/\Omega$。除 $M$ 并用 Kepler 关系，质心速度的振荡幅值为 $(464/105)c\nu^2|\delta|x^4$。它随轨道缩小迅速增加，但完整一转的矢量平均仍为零。到并合时端点截断、相位和振幅快速改变，才可能留下有限净值；把 $x$ 取到光子轨道只能作量级估计，不能据此固定最后反冲系数。

有自旋时还可由现成的矢量推断反冲方向。取差自旋组合 $\bm\Delta_s=M(\bm S_1/m_1-\bm S_2/m_2)$，圆轨道时最低阶自旋--轨道动量通量中几何部分正比于
\[
 \hat{\bm r}\mathbin{\times}\bm\Delta_s+
 (\hat{\bm r}\mathbin{\times}\hat{\bm\lambda})
                     (\hat{\bm\lambda}\cdot\bm\Delta_s).
\]
以 $\hat{\bm L}=\hat{\bm r}\times\hat{\bm\lambda}$ 展开 $\bm\Delta_s=\Delta_r\hat{\bm r}+\Delta_\lambda\hat{\bm\lambda}+\Delta_L\hat{\bm L}$，上式化为 $2\Delta_\lambda\hat{\bm L}-\Delta_L\hat{\bm\lambda}$。自旋差平行于轨道轴时，通量位于轨道面；自旋差位于面内时，通量沿法线，且只由当时的切向投影决定。设面内自旋差近似固定，$\Delta_\lambda$ 随轨道相位作正弦变化，故最后的法向反冲对并合相位有近似正弦的依赖。仅给两颗黑洞的自旋大小和相对夹角，仍缺少决定法向反冲的一项角度。

反冲不仅影响天体留在何处，也会留下可查的光谱。对一颗在并合前有薄吸积盘的遗留黑洞，取盘元原速度 $v_K=(GM_f/r)^{1/2}$，在猝然受踢的黑洞共动系内，盘元速度改为 $\bm v_K-\bm v_k$。其新比轨道能约为 $e'=-\frac{GM_f}{2r}+\frac{v_k^2}{2}-\bm v_K\cdot\bm v_k$。忽略并合质量损失，且 $v_k$ 垂直盘时，内于 $r_b=GM_f/v_k^2$ 的盘元仍束缚，外部可脱离；若踢在盘内，不同方位有不同界。取 $M_f=10^8M_\odot$、$v_k=10^3$\,km/s，$r_b\simeq0.43$\,pc。足够紧束缚的宽线区和吸积盘可随黑洞移动，尺度较大的窄线区仍停在星系核。于是宽线与窄线可能有速度偏移，若再有空间偏移可加强识别；双活跃核与盘风也可能产生类似速度分量，不能仅凭一条偏移线就量出反冲。黑洞若只被踢到星系核外围，还可在动摩擦下回落，这与逃离星系的结局在后续并合率上不同。

\subsection{中子星的潮汐响应及并合阈值}

\S\ref{sec:pn-wave-energy} 已把潮汐在相位中的领头贡献写成 $\widetilde\Lambda$。给定一个冷物态方程 $p(\rho)$，怎样实际算出其中每一颗星的 $\Lambda$？在几何制，先用第 \ref{chpt:BH} 章的静态平衡求 $m(r),p(r)$，表面 $p(R)=0$，$M=m(R)$。取 $\ell=2$ 静态偶扰动，将星内流体与度规联立，正则的度规振幅记为 $H(r)$。外部真空的两支解在远离星面但仍远小于伴星距离的匹配区内各随 $r^2$、$r^{-3}$ 变；前者描述外加潮汐，后者描述诱导四极矩。只有两支的比例有意义，因此星内演化对数导数 $y=r H'/H$ 即可。令 $F=1-2m/r$、$c_s^2=\d p/\d\rho$，有
\begin{align*}
 r y'+y^2+y\,\frac{1+4\pi r^2(p-\rho)}{F}&=-r^2{\cal U},\\
 {\cal U}&=\frac{4\pi[5\rho+9p+(\rho+p)/c_s^2]-6/r^2}{F}
   -4\left[\frac{m+4\pi r^3p}{r^2F}\right]^2 .
\end{align*}
中心 $H\propto r^2$ 规定 $y(0)=2$。若物态有一阶相变，密度会跳变，方程里的 $\d\rho/\d p$ 有分布项，不能用同一个 $y$ 无条件跨过界面。在压强连续的界面 $r=r_d$，记密度下降量为 $\Delta\rho=\rho_{\rm in}-\rho_{\rm out}>0$，将奇异项沿平衡压强积分得 $y_{\rm out}-y_{\rm in}=-4\pi r_d^3\Delta\rho/(m_d+4\pi r_d^3 p_d)$。表面密度非零时，也须先取其外侧值再与真空解配对。

定义 $C=M/R$、$y_R=y(R^+)$、$b=2+2C(y_R-1)-y_R$。在表面连续匹配两个真空基解，以 $Q_{ij}=-\lambda_{\rm t}{\cal E}_{ij}$ 规定潮汐响应，得
\begin{align*}
 k_2={}&\frac{8 C^5(1-2C)^2 b}{5{\cal D}},\\
 {\cal D}={}&2C[6-3y_R+3C(5y_R-8)]\\
 &+4C^3[13-11y_R+C(3y_R-2)+2C^2(1+y_R)]\\
 &+3(1-2C)^2 b\ln(1-2C).
\end{align*}
于是 $\lambda_{\rm t}=2 k_2 R^5/3$、$\Lambda=2 k_2/(3C^5)$。$k_2$ 是内部密度分布对表面潮汐的响应，$C^{-5}$ 则单由长度标度而来。以 $C\to0$ 作检验，外部静态解趋于 $H=A r^2+B r^{-3}$。设 $b_{\rm N}=B/(A R^5)$，有 $y_R=(2-3b_{\rm N})/(1+b_{\rm N})$，故 $k_{2,\rm N}=b_{\rm N}/2=(2-y_R)/[2(3+y_R)]$。均匀不可压星的内侧解 $y=2$，表面密度跳变给 $y_R=2-3=-1$，因而 $k_{2,\rm N}=3/4$。若忽略跳变，会对最易被外场拉长的均匀流体反而得到零响应。

上述求法把伴星的潮汐视为在流体调整时间上近乎静止。若某个星模的频率 $\omega_\alpha$ 接近潮汐驱动的 $m\Omega$，对模式坐标应解 $\ddot q_\alpha+2\gamma_\alpha\dot q_\alpha+\omega_\alpha^2q_\alpha=F_\alpha e^{-\mathrm i m\Omega t}$。远离共振时 $q_\alpha\simeq F_\alpha/\omega_\alpha^2$，各模求和还原静态 Love 数；接近共振时 $(\omega_\alpha^2-m^2\Omega^2-2\mathrm i m\Omega\gamma_\alpha)^{-1}$ 同时引进增强和相位滞后。若驱动频率在阻尼时间内快速扫过，共振来不及建立稳态，上式还须作非定频积分。由此可以看出，冷星静态的 $\Lambda$ 并未概括所有可在高频出现的物质效应。

相接前一刻的频率量级可由 $\Omega_c^2\simeq G M_{\rm tot}/(R_1+R_2)^3$ 估算。两颗 $1.4 M_\odot$、$12$\,km 的星给 $f_{\rm GW}\sim1.6$\,kHz，潮汐会在接触以前已经改变相位。相接后剪切层、激波、磁场以及中微子输运均可能改变角动量和温度分布，不能接一个冷球星的线性振荡来代替。若总质量超过给定物态、热状态和角动量能支撑的上限，短促的并合信号后便转成黑洞振铃；稍轻一些可先形成由差速转动支撑的超大质量星，再随角动量搬运而塌缩。另一类遗留物高于冷非转星的最大质量却可由均匀转动支撑，须待长期制动后才失稳。更轻的还可长期稳定。比较这些上限时，应在同一质量定义下计入抛出物质和辐射；二元的引力质量总和与遗留星的重子质量并不相等。并合后谱峰频率主要受遗留星密度影响，谱峰的寿命则又依赖何时失稳。

\section{波形告诉我们怎样的天体？}\label{sec:gw-cobs}

\subsection{从第一次双黑洞记录辨认源}

对一段随频率升高而增强的信号，先问它在源附近要求怎样的长度和时间，再问哪些天体能满足。以 GW150914 为例，可在图上取远离峰值的一个时刻，由 $f,\dot f$ 代入 \S\ref{sec:pn-wave-energy} 的最低阶啁啾关系，估得红移后的啁啾质量为几十个太阳质量。这个估计只利用相位，未用距离或仪器天线方位。由 $\mathcal M=\nu^{3/5}M$ 且 $\nu\le1/4$，可知总质量至少为 $4^{3/5}\mathcal M$；故在相同啁啾质量下，不对称的系统需要更大的总质量。再从较晚时的旋转频率估分离 $a\sim[GM/(\pi f)^2]^{1/3}$。取 $M=65M_\odot$、$f=150$\,Hz，$a\sim340$\,km，而该质量的 $2GM/c^2$ 约 $190$\,km。轨道已深处强场，Newton 数值只能说明两个物体极其紧致，最后的质量与自旋应由完整强场波形估计。

怎样知道这样一个形似双黑洞的极强触发并非两台仪器噪声的偶合？若只用平稳 Gaussian 噪声的尾概率，极少的非平稳尖峰便可使答案相差许多数量级。可对真实台间时间序列施加大于传播时差的多组相对平移，用同样的触发与一致性统计重做搜索，以估计无天体相关时的背景。在等效背景时间 $T_{\rm bg}$ 内若有 $k$ 个背景触发超过候选阈值，估计假警率约 $k/T_{\rm bg}$。若没有一个超过，能给的是与 $T_{\rm bg}^{-1}$ 同量级的上限，其置信系数取决于所用统计量。假警率为 $R_{\rm FA}$ 时，在零延迟观测时间 $T$ 内至少出现一个的概率为 $1-\exp(-R_{\rm FA}T)$。搜索不同的质量区间和波形类别会增加试验次数，所引阈值必须与实际搜索流程相配。假警概率本身没有用到天体发生率的先验，不能直接充当“源为黑洞双星”的后验概率。

第一批事件已能显示同类天体也给不同的测量。总质量较小的 GW151226 在地面带宽内积累更多旋近周期，故质量组合与自旋引起的相位漂移有较长时间显现。GW170814 时第三台探测器也参与，除了两个独立时差还可比较三台的响应大小。设已由时差把某几个天区留下，在每个天区三台读数应位于张量偏振响应矩阵两列张成的平面；与该平面正交的组合理应只含噪声。一台仪器若在某方向应有很大响应却几乎无信号，也可排除该方向。故第三台较弱也能显著缩小天区；排除程度须用实际噪声和校准误差计算，不能将弱响应当作绝对的零。

将一次双黑洞测量用于基本引力的办法已见第三章。这里的来源区分尤为要紧。用旋近拟合质量、自旋，将其代入数值关系预言末态，然后和并合后波形给的末态比较，是贯穿强弱场的一致性检验。若用晚时某一模的 $\omega_R,\tau$ 单独定出 $(M_f,a_f)$，再用另一模去核对，检验依赖的是 \S\ref{sec:gw-cbh} 的边值问题。以第一个模拟合的两个数再次“验证”同一个模没有独立自由度。偏振检验取决于台网几何，传播时间检验又需区分源内延迟；后者在 GW170817 上的速度界，已在第三章讨论。

还可用一个面积检验连接首次并合记录与第二章。设初始和遗留黑洞质量各为 $m_1,m_2,M_f$，无量纲自旋为 $\chi_1,\chi_2,\chi_f$。Kerr 面积 $A=8\pi m^2(1+\sqrt{1-\chi^2})$，故在满足经典面积定理的假设时应有
\[
 M_f^2\left(1+\sqrt{1-\chi_f^2}\right)
 \ge \sum_{i=1}^2 m_i^2\left(1+\sqrt{1-\chi_i^2}\right).
\]
同一源的每个质量都可用红移质量代入，因为两边共同乘 $(1+z)^2$。若初态与末态皆无自旋且两初始质量等同，允许的最高辐射比例为 $1-1/\sqrt2\simeq0.29$；真实并合的动力学不会自动饱和这个几何上限。非零末态自旋使给定 $M_f$ 的面积较小，应回到原式。若将含有面积增长先验的同一个完整广义相对论模板同时拿来“测”三项，不构成独立的面积检验；实际可用早期旋近推两初始参数、较晚振铃推末态，并在频带分割和先验选择后传播相关的不确定度。

以第一批事件的量级作另一检查。若总红移质量约 $70M_\odot$，末态旋转使 $M_f\omega_R$ 为约半，几何时间 $G M_{f,z}/c^3$ 为数个 $10^{-4}$\,s，由此给几百 Hz 和数毫秒的振铃。把 $A e^{-t/\tau}\cos(\omega_Rt+\phi_0)$ 当作一个复指数及其共轭的和，$\omega_R\tau$ 不太小时，单个正频峰的功率半高全宽为 $1/(\pi\tau)$，无量纲宽度为 $1/Q$。低信噪比下两个相近极点若相隔不足谱宽，不易独立估计寿命。看到一段衰减波说明有晚期响应，测准第二个独立极点才多出对 Kerr 家族的检验。对同一模若用 \eqref{eq:gw-cmodes} 积分单一复模 $H=C e^{-\mathrm i\omega t}$，则该模在起点之后带走能量 $|C|^2|\omega|^2\tau/(32\pi)$。它还依赖所取振铃起点和模重建，不能直接把整个系统的质量亏损都归入这个单模。

用波形作种群比较时，GW150914、GW151226、GW170104、GW170608 所处的质量范围和有效自旋后验并不相同，单台带宽也使可见旋近长度不同。以一组质量为例，若低总质量源在 $f_{\min}$ 处进入频带、从 $f_{\max}$ 处离开，积累周期数为 $\int f\,\d f/\dot f$；给定啁啾规律后，低频端尤重要。对总质量较大的源，仪器往往看到了并合与振铃，旋近周期却明显更少。对于 $\chi_{\rm eff}$ 的后验，排除较大的同向自旋不等于独立测准了每颗黑洞的自旋模或证明了它倾斜；一正一负的投影可以抵消。比较不同事件时要对相同的参数定义、先验与选择效应作比较，而非数一数误差条是否跨零。

\subsection{双中子星的三种光}

GW170817 的旋近波给出紧致双星的啁啾与距离，短伽马暴给相对论外流的到达方向和时刻，随后光学至近红外瞬变给抛射物的质量、速度和组分信息。三种信号来自不同区域。若把伽马光子离开喷流的时刻与两星表面相接强行视为同一事件，就会把发动喷流、穿出包层和辐射的时间全部混入传播速度。第三章给出的时差讨论之外，还可直接从光变理解喷流与抛射物各自起什么作用。

并合把低电子丰度物质抛进低压区，大量中子迅速被种子核俘获。只要平均中子俘获时间 $[n_n\langle\sigma v\rangle]^{-1}$ 短于 $\beta$ 衰变时间，核素即可先沿同位素链远离稳定线，待自由中子耗尽后再由衰变走向稳定核。许多核素的不同衰变时间叠加，给膨胀物质持续加热。这里的能量储库来自核结合能；快衰减的中子丰度、电子丰度、逃逸中微子和 $\gamma$ 射线均决定最后能热化多少。潮汐尾倾向于保留低 $Y_e$，受激波或强中微子照射的物质可有较高 $Y_e$，即使同一次并合也可发生不同的核合成。

令抛射质量为 $M_{\rm ej}$，特征膨胀速度 $v$，时刻 $t$ 的半径为 $R\simeq vt$，灰不透明度为 $\kappa$。均匀球的径向光深 $\tau\simeq3\kappa M_{\rm ej}/(4\pi R^2)$，随机行走使 $t_{\rm diff}\sim\tau R/c$。早期扩散比膨胀慢，到了 $t_{\rm diff}\sim t$ 光开始有效逃逸，故
\[
 t_{\rm pk}\sim\left(\frac{3\kappa M_{\rm ej}}{4\pi vc}\right)^{1/2},
 \qquad L_{\rm pk}\sim M_{\rm ej} f_{\rm th}\dot q(t_{\rm pk}),\qquad
 T_{\rm eff}^4\simeq\frac{L}{4\pi\sigma_{\rm SB}v^2t^2}.
\]
这些是均匀灰球的尺度估计，层状速度与密度剖面可改数因子。取 $M_{\rm ej}=10^{-2}M_\odot$、$v=0.2c$、$\kappa=1\,\mathrm{cm^2\,g^{-1}}$，得 $t_{\rm pk}$ 约两天。若质量速度不变而 $\kappa$ 增大十倍，达峰时间延后三倍多，物质在更大半径才放光，温度较低。镧系元素有密集的束缚态谱线，富含这些元素的抛射物通常更不透明。早期蓝色、随后出现较红的成分，因而可由组分和速度不同的抛射物解释。只凭单滤光片的峰时只得 $\kappa M_{\rm ej}/v$，还需谱线、颜色或速度测量解除简并。

伽马暴的可见度还受相对论性束射影响。设流速 $\beta c$、Lorentz 因子 $\Gamma$，视线与局部喷流速度夹角 $\vartheta$，Doppler 因子 ${\cal D}=[\Gamma(1-\beta\cos\vartheta)]^{-1}$。当 $\Gamma\gg1,\vartheta\ll1$ 时，${\cal D}\simeq2\Gamma/(1+\Gamma^2\vartheta^2)$。喷流外侧的观者即使错过很窄的瞬发主束，仍可在外激波扫入环境、$\Gamma$ 下降到约 $1/\vartheta$ 后看到上升的余辉。若喷流与极区抛射物相互作用，受激波加热的物质还可形成较宽的茧，其破出或对环境的外激波亦能辐射。在这一阶段，单次弱 $\gamma$ 信号不够区分离轴喷流与较宽外流，还须对余辉随时间的变化和天球像运动作比较。温热抛射物的近各向同性热辐射与非热同步余辉提供不同的几何信息。

若宿主有电磁红移，波形的红移质量可还原成源质量，波幅给的光度距离可与近邻速度作比较。低红移时 $cz_{\rm obs}\simeq H_0 d_L+v_{\rm pec}$。若只用一例，距离和轨道倾角在波幅中强烈相关，星系本动速度也不能忽略。喷流方向若与轨道轴有联系，电磁拟合可给额外倾角信息；若把两者严格对齐当作先验，就须将这个物理假设计入测量。选作对应体的概率同样偏好更亮的源，合并许多事件时还须考虑选样。

\subsection{从颜色、减速和视运动解除简并}

把快速演化的热瞬变与普通超新星分开，可从一个热光球模型着手。若物质大致同调膨胀，各层速度从光谱的宽吸收结构给出；用共动光度与连续谱色温可估黑体半径 $R_{\rm BB}=[L/(4\pi\sigma_{\rm SB}T_{\rm col}^4)]^{1/2}$。若初始半径很小且 $\tau>1$，不同时刻推回的 $R_{\rm BB}/t$ 应与发光层的速度相容。由于不透明度随波长变化和光球在速度坐标中后退，同一固定质量壳的 $R=vt$ 并不意味着观到的 $R_{\rm BB}$ 必须始终线性上升。快速变红既可反映同一层冷却，又可因蓝光层变薄后显出原来遮住的富镧物质，因此需要分光而不只是两次不同滤镜的测光。

统计许多不稳定核的衰变，为何能给近幂律的热源？作一简单模型，设每个对数平均寿命区间的核数近乎相同，$\d N=C\,\d\tau/\tau$，单核某次释放的可热化能量近似 $Q$。时刻 $t$ 的总功率为 $CQ\int_{\tau_{\rm min}}^{\tau_{\rm max}}\exp(-t/\tau)\,\d\tau/\tau^2$。令 $x=t/\tau$，在 $\tau_{\rm min}\ll t\ll\tau_{\rm max}$ 得 $\dot E\simeq CQ/t$。若可释放能量、核种分布与热化效率还随寿命缓慢改变，指数也会略离 $-1$。因此一条近幂律光变可以由许多寿命的叠加产生，无须假定一种恰有这个时间规律的核。相反若只由单一寿命 $\tau_0$ 支配，晚期应显指数。真实的光变还经过膨胀和辐射输运滤过，不能由指数的一次拟合唯一指认某个核素。

扩散近似还能指出峰值关系在哪些条件下失效。取抛射物热能 $E_{\rm th}$、膨胀功近似 $E_{\rm th}/t$、加热率 $\dot Q(t)$，以及 $L\simeq E_{\rm th}/t_{\rm diff}$，$t_{\rm diff}=t_d^2/t$。热量方程变为
\[
 \dot E_{\rm th}+\left(\frac1t+\frac{t}{t_d^2}\right)E_{\rm th}
       =\dot Q(t),\qquad
 L(t)=\frac{e^{-t^2/(2t_d^2)}}{t_d^2}
        \int_0^t t' e^{t'^2/(2t_d^2)}\dot Q(t')\,\d t' .
\]
已取起初的热能可略。微分第二式得到 $\dot L=(t/t_d^2)(\dot Q-L)$，故此近似在峰处恰给 $L_{\rm pk}=\dot Q(t_{\rm pk})$。若 $\kappa$ 快速变化或者另有相对论喷流加热，$t_{\rm diff}=t_d^2/t$ 不成立，不能仍以这条峰值关系精确反求抛射质量。若热源是早已储存的激波能，积分初值项也不可省。于是同样的光学达峰时间，须和光谱及尾部热源一起解释。

对非热余辉可用减速的时间判别外流几何。取均匀外介质数密度 $n_{\rm ext}$、绝热的相对论性激波在未显出边缘以前的各向同性等效能量 $E_{\rm iso}$。扫入的静质量约为 $4\pi n_{\rm ext}m_p R^3/3$，其激波能量量级再乘 $\Gamma^2 c^2$。对沿速度方向的观者，光行时差给 $t_{\rm obs}/(1+z)\sim R/(2c\Gamma^2)$。消去 $R$ 得
\[
 \Gamma\propto\left(\frac{E_{\rm iso}}{n_{\rm ext}m_pc^5}\right)^{1/8}
                 \left(\frac{t_{\rm obs}}{1+z}\right)^{-3/8}.
\]
若明亮喷流核的边缘距视线 $\Delta\theta$，当 $\Gamma\Delta\theta\sim1$ 时观者开始见到那部分以前被束射压低的辐射，尺度 $t_{\rm pk}/(1+z)\propto(E_{\rm iso}/n_{\rm ext}m_pc^5)^{1/3}(\Delta\theta)^{8/3}$。这一极强的角度依赖说明晚到的余辉不要求中心引擎也延迟数月。若外介质随半径下降、喷流有角结构或横向扩展，指数与峰形会变；一个均匀锥模型只是读出符号和量级的方便办法。

还可直接用天球像的形心运动。设亮区沿与视线夹角 $\theta_v$ 方向前进，在源系经过 $\Delta t$，视向距离缩短 $\beta c\,\Delta t\cos\theta_v$。观者所见时间间隔除以 $(1+z)$ 为 $(1-\beta\cos\theta_v)\Delta t$，横向位移为 $\beta c\,\Delta t\sin\theta_v$。故表观横速满足 $\beta_{\rm app}=\frac{\beta\sin\theta_v}{1-\beta\cos\theta_v}$。固定 $\beta$，在 $\cos\theta_v=\beta$ 达 $\beta_{\rm app,max}=\beta\Gamma$。超过光速的天球形心运动可由普通的亚光速发射体给出，且要求辐射来自局部明亮区域；完全球对称同时明亮的壳没有同样的净形心漂移。这类几何信息与同时容许多种距离与倾角的引力波后验互补，能帮助区分窄喷流核与较宽、较慢的茧，而二者在一次并合中也可同时存在。

\subsection{从个别事件推到种群}

相同观测频带内，即使两个质量区间的真实并合率相同，目录中的数目也可大异。在旋近主导、低红移且末频高于主要敏感带的近似下，可见距离随 $\mathcal M^{5/6}$ 变，可见体积随 $\mathcal M^{5/2}$ 变。过重系统的并合频率会落到仪器低频截止以下，幂律便不再适用。给定源参数 $\vartheta$ 和种群超参数 $\lambda$，令归一化源分布为 $p(\vartheta|\lambda)$、探测概率为 $p_{\rm det}(\vartheta)$，则 $\alpha(\lambda)=\int p_{\rm det}p(\vartheta|\lambda)\,\d\vartheta$ 才是可见比例。若期望合计发生 $\mu_0$ 次，资料 $d_i$ 的种群似然除与模型无关的常数外为
\[
 {\cal L}(\mu_0,\lambda)\propto e^{-\mu_0\alpha(\lambda)}\mu_0^{N}
         \prod_{i=1}^{N}\int p(d_i|\vartheta)\,p(\vartheta|\lambda)\,\d\vartheta .
\]
在已知探测数 $N$ 的条件下，质量分布的似然含 $\alpha^{-N}$。单源推断得到的后验已经用了单源先验，放入上式重加权时应先除去它。例如若两质量段的可见比例分别为 $1/4$ 与 $1$，目录中各有四例，同等观测时间的校正率之比约为四比一。少数高质量源的存在可证明这种天体确能形成，却不能只看原始数目便断言此类是多数。恒星金属丰度、共同包层演化和致密星团动力学都可改变质量与自旋分布，辨认形成机制需同时处理这些选择效应。



\section{星系核的双黑洞}\label{sec:gw-csmbh}

\subsection{从星系并合到毫秒差距}

银河中心的恒星轨道怎样同时给出距离与中心质量，第三章已有计算。把附近星系中心的动力学质量和核球性质比较，可发现中心黑洞质量与恒星速度弥散 $\sigma$ 有关。这个关系带有内禀散布，不能把每个星系的光度直接换成唯一的黑洞质量。黑洞引力占优势的尺度约为 $r_{\rm inf}=GM/\sigma^2$。若两个含黑洞的星系并合，两核最初相距远大于 $r_{\rm inf}$，还未构成可以直接套用辐射反作用的 Kepler 双星。

较重天体穿过恒星背景时，许多小角重力散射形成拖在身后的密度尾迹。对速度 $V$、远大于单星质量的天体 $m$，在局域均匀、各向同性的 Maxwell 速度分布中，Chandrasekhar 动摩擦加速度为
\[
 \frac{\d\bm V}{\d t}=-4\pi G^2m\rho_\star\ln\Lambda_{\rm C}
 \frac{\bm V}{V^3}\left[\operatorname{erf}(X)
                  -\frac{2X}{\sqrt\pi}e^{-X^2}\right],
 \quad X=\frac{V}{\sqrt2\sigma}.
\]
Coulomb 对数近似是有效最大、最小碰撞参数之比的对数。括号内为比该天体慢的恒星比例；小 $X$ 时低次相消后正比于 $X^3$，故拖曳加速度在低速近于对速度线性。若核仍带一团被它束缚的恒星，早期实际产生尾迹的质量比裸黑洞大。向内移动时这些恒星又逐步被潮汐剥走，所用 $m$ 需随阶段改变。

算在奇异等温球中的近圆衰减。令圆速 $V_c=\sqrt2\sigma$ 为常数，则 $\rho_\star=V_c^2/(4\pi G r^2)$、$X=1$。取括号值 $B(1)\simeq0.428$，切向拖曳为 $a_{\rm df}=-B(1)Gm\ln\Lambda_{\rm C}/r^2$。比角动量 $j=rV_c$，故 $V_c\dot r=r a_{\rm df}$；从 $r_i$ 积到很小半径得 $t_{\rm df}\simeq1.17 V_c r_i^2/(Gm\ln\Lambda_{\rm C})$。初始半径增加一倍，耗时增加四倍；伴星质量减半，则时间再增加一倍。靠近已被双黑洞清空的核时，均匀背景及圆轨近似同时变差，不能把此式一路延到并合。

在 $a\lesssim r_{\rm inf}$ 形成束缚双星后，达到 $a_h\sim G\mu/(4\sigma^2)$ 左右可称\textbf{硬双星}，其中 $\mu=m_1m_2/M$。近旁恒星只要进入半径 $\sim a$ 的相互作用区，就可能以较高的速度离开，把双星束缚能的绝对值增大。在引力聚焦主导时，截面约为 $\Sigma\sim \pi a^2[1+2GM/(a\sigma^2)]\propto a$；每次移给恒星的能量量级 $G\mu m_\star/a$，二者乘积与 $a$ 近乎无关。将角分布与实际三体遭遇的系数归入 $H$，常写
\[
 s_\star=\frac{\d(1/a)}{\d t}\simeq H\,\frac{G\rho_\star}{\sigma},
 \qquad \dot a_\star=-s_\star a^2 .
\]
因此在持续充满遭遇轨道的恒星库中，$1/a$ 近似线性增加。若这类轨道很快被用尽，恒星密度仍高也未必能维持常数 $s_\star$。

所谓损失锥指角动量小到近心点进入相互作用区的轨道。在质量 $M$ 的近 Kepler 势内，近抛物线轨道有 $J_{\rm lc}^2\simeq2GM a$。若能量固定的恒星近各向同性，$J^2$ 分布近均匀，从而首批可供使用的比例约为 $J_{\rm lc}^2/J_c^2$，$J_c$ 为同能量圆轨道角动量。两体弛豫、非球对称星系势的长期扭矩、后续星系并合及气体盘都可重新填入损失锥。如果填充太慢，从秒差距级到引力波接手所需间距的那一段会比宇宙年龄还长，这就是末秒差距问题。其条件是星系核供给机制不足，并非所有超大质量双星都必经一个固定的停滞半径。

圆轨道辐射使 $\dot a_{\rm GW}=-K/a^3$，$K=64G^3\mu M^2/(5c^5)$。在恒定 $s_\star$ 的模型，两机制相等的间距、该处源系波频分别为
\[
 a_{\rm b}=(K/s_\star)^{1/5},\qquad
 f_{\rm b}=\frac{\sqrt{GM}}{\pi}(s_\star/K)^{3/10}.
\]
比较时间也可看得直观。若只靠辐射，距并合还需 $a^4/(4K)$。取两颗 $10^8M_\odot$ 的黑洞和 $a=1$\,pc，所得约 $3\times10^{14}$ 年；到 $10^{-2}$\,pc 已缩小为约 $3\times10^6$ 年。一个很小的间距差异即可在是否有效并合上产生决定性影响。对于大偏心率，辐射每轨集中在近心点，$K$ 还应乘随 $e$ 迅速增大的损失因子，环境同时能激发或压低 $e$。

\subsection{质量尺度和观测频带}

用 Schwarzschild 最后稳定圆轨道作示例，观者所见的主模频率为 $f_{\rm ISCO}=c^3/[6^{3/2}\pi G(1+z)M]$。对 $M=10^6M_\odot$，源系得 $4.4$\,mHz；$10^9M_\odot$ 则为 $4.4\,\mu$Hz。Kerr、自旋有限质量比和真正的并合会移动最高频，$1/[(1+z)M]$ 的缩放依然决定什么质量的末期进入空间干涉仪的带宽。更重的一群系统在并合很久以前、轨道为数月到数年时，就可在脉冲星计时的频率范围内一同辐射。把一台仪器称为对超大质量源敏感，还须说明它实际记录的是哪一段演化。

若观测时间 $T$，单频 Fourier 箱宽度约 $1/T$。一个圆旋近的频率漂移 $\dot f T$ 若小于 $1/T$，在单次观测内近似单色；若 $\dot f T^2\gtrsim1$，需要啁啾模板。取低阶 $\dot f\propto\mathcal M_z^{5/3}f^{11/3}$，同一质量内越靠近并合越难作单色源。设源系发生率为 $R$，平均在某对数频带内可见数与 $R f/\dot f$ 成正比。低频可同时积累许多缓慢系统，高频即使单个波更强，停留的源却更少。这是群体背景转为个别可辨源的一个原因。

极端质量比旋近则不同。一颗质量 $\mu$ 的致密体绕质量 $M$ 的大黑洞，$\eta_e=\mu/M\ll1$。在 $r\sim$ 数个 $M$ 处，单转时间为 $O(M)$，轨道束缚能量级 $O(\mu)$、辐射功率量级 $O(\mu^2/M^2)$，故辐射反作用时间为 $O(M/\eta_e)$，可积累 $O(\eta_e^{-1})$ 个弧度。大质量星双星短暂的非线性并合与这种在强场绕转许多周的运动，对模型精度的要求正相反：后者每转很小的误差经过 $10^4$ 至 $10^5$ 个周期就可移走许多相位。

\subsection{致密星怎样进入长期旋近？}

在星系核中，小体最初往往由许多远距离恒星遭遇逐步改变轨道角动量。记两体弛豫时间 $t_{\rm rlx}$ 为角动量变化到同能量圆轨道尺度 $J_c$ 所需的时间，则让初始 $J\ll J_c$ 改变其自身量级只需 $t_J\sim(J/J_c)^2t_{\rm rlx}$。Kepler 椭圆有 $J^2/J_c^2=1-e^2$。偏心极大时，虽然引力波在小近心距处很强，轨道也更易被下一次遭遇推入直坠区。若整个辐射收缩时间 $t_{\rm GW}\ll t_J$ 且初始近心点在俘获边界之外，随后可与星团脱耦，形成长旋近；若先越过俘获边界，只余一次爆发。比较对象必须是角动量扩散时间，不能直接用较长的总弛豫时间。

另一条通道是大黑洞潮汐拆散一对近来的小双星。质量合计 $m_b$、内半长轴 $a_b$，其内自引力加速度 $Gm_b/a_b^2$ 与大洞的潮汐差 $GMa_b/r^3$ 相等时给 $r_t\sim a_b(M/m_b)^{1/3}$。若质心最初近抛物线，近心点速度 $V\sim(GM/r_t)^{1/2}$，内部速度 $v_b\sim(Gm_b/a_b)^{1/2}$，拆散后每颗星相对于大洞的比轨道能改变 $\sim\pm Vv_b$。一颗获得正能逃逸，另一颗可束缚到半长轴 $a_{\rm cap}\sim GM/(2Vv_b)\propto a_b(M/m_b)^{2/3}$ 的轨道。只要被留下者为致密天体并未被肢解，日后也可能进入信号频带。盘内形成、气体迁移等机制可给更小的初始倾角和偏心率。不同轨道族若能观到，便有望反推核内致密星的供应途径。

\S\ref{sec:KN} 给了 Kerr 测地线的一阶积分。以能量、轴向角动量和 Carter 常数 $(E,L_z,Q_C)$ 标记一般有界轨道，相对远方时间有三种频率 $(\Omega_r,\Omega_\theta,\Omega_\phi)$。径向、纬向作周期运动，轴向每转还可有进动，源在频域包含组合 $\omega_{mkn}=m\Omega_\phi+k\Omega_\theta+n\Omega_r$。只保留 $m=2$ 的一个圆轨道频率，会把轨道偏心、倾斜及多种进动误作相位噪声。若轨道接近径向势的不稳定转折，粒子可先从远处迅速接近近心点，再在小半径绕转数圈才返回，波形表现为间歇的高频波列。此时用单一孤立椭圆的瞬时长轴描述谱已不够。

\subsection{自力的奇异部与长时相位}

小体产生自己的场，若把推迟场 $h_{\mu\nu}^{\rm ret}$ 在自身位置直接代入测地线方程，梯度发散。这与远处用场算波幅并不矛盾，后者的观测点不在源上。把局部与小体的奇异库仑型结构相同的场 $h^S$ 减去，定义 $h^R=h^{\rm ret}-h^S$。对无自旋小体的一阶质量比近似，物理运动可表为在 $g+h^R$ 上的测地线；换回背景固有时，用 $P^{\alpha\beta}=g^{\alpha\beta}+u^\alpha u^\beta$ 投影，得
\[
 a^\alpha=-\tfrac12 P^{\alpha\delta}
 (2\nabla_\beta h^R_{\gamma\delta}-\nabla_\delta h^R_{\beta\gamma})
 u^\beta u^\gamma+O(\eta_e^2/M).
\]
投影保证 $u_\alpha a^\alpha=0$。$h^R$ 虽在世界线上光滑，仍包含沿过去世界线传播并被背景曲率散射回来的尾，通常不能只由该时刻的速度加速度构造。奇部与有限部的区分还要配合所用规范，小体的坐标半径移动并非单独的可观测自力。

一种实际正则化从球谐分解入手。每一 $\ell$ 模在从内侧或外侧逼近质点时有限，形式和 $\sum_\ell F_\ell^{\rm ret}$ 才发散。若局部奇场的同一模大 $\ell$ 有 $F_\ell^S=A_\pm L+B+C/L+O(L^{-2})$，$L=\ell+1/2$，从每模减去前三项使余项可求和；再减去剩余奇部的有限总和即可恢复 $h^R$ 给的力。$A_\pm$ 的号依从哪侧逼近，不能混合取用。数值误差既来自单模积分，又来自截断的高 $\ell$ 尾，检查余项是否如 $L^{-2}$ 衰减可区分两者。

对有时间反演对称结构的轨道，可用推迟及先进解的差、和分离耗散与保守贡献。辐射场 $(h^{\rm ret}-h^{\rm adv})/2$ 中局部相同的奇部抵消，给长期耗散的便捷算法；保守部仍须正则化。对有界、非共振轨道的绝热近似，可用到无穷远及视界的能、角动量通量求 $\langle\dot E\rangle,\langle\dot L_z\rangle$。$Q_C$ 源自二阶 Killing 张量，不能仅用一股普通的角动量流来替它。一般倾斜偏心轨道需算相应模式和或局部自力。Kerr 视界通量某些频率可为负，故不能永远把“吸收”当成给轨道的额外正阻尼。

为何一阶的保守修正也须担心？以慢时 $T=\eta_e t/M$ 表示一支相位，若 $M\Omega(T)=\omega_0(T)+\eta_e\omega_1(T)+\cdots$，则
\[
 \Phi(T)=\frac1{\eta_e}\int^T\omega_0(s)\,\d s+
                         \int^T\omega_1(s)\,\d s+O(\eta_e).
\]
短时的 $O(\eta_e)$ 频率偏差积累成 $O(1)$ 相位，已足以降低与精密模板的相关。更进一步，积分 $\omega_0$ 本身在一个耗散时间上也会受次阶长期演化修正，因此欲使总相位误差小于一个弧度，通常既需一阶保守与振荡自力，也需二阶耗散信息。在低阶有理比 $k\Omega_\theta+n\Omega_r=0$ 附近，平常逐轨平均为零的某些项不再相消，穿过瞬时共振可使轨道常数发生依初相而定的改变。这说明平均法并非在一切强场轨道上都一致有效。

\subsection{从一颗源的能谱到星系核背景}

设 $n(z,\vartheta)\,\d z\,\d\vartheta$ 为单位今日共动体积内、在红移和参数区间中发生的并合数，已把天体发生率对源系时间积分。每个源发射的能量红移为本来的 $1/(1+z)$，$\d\ln f=\d\ln f_r$，于是今日单位对数频率能量密度为
\begin{equation}
 \frac{\d\rho_{\rm GW}}{\d\ln f}
 =\int\frac{n(z,\vartheta)}{1+z}
     \left.\frac{\d E_{\rm GW}}{\d\ln f_r}\right|_{f_r=(1+z)f}
        \d z\,\d\vartheta .
 \label{eq:gw-cpopulation}
\end{equation}
此处 $\rho_{\rm GW}$ 按能量密度计。若给的是每源系时间、每共动体积的速率密度 ${\cal R}$，就用 $n={\cal R}/[(1+z)H(z)]$。若从全天实测的每观者时间事件数出发，已经另含共动体积元与钟慢，二者不能再重复乘入。定义特征应变 $h_c$ 使能量密度为 $(\pi c^2/4G)f^2 h_c^2$。于是只要知道每一源在每个对数频率区间真正交给引力波的能量，便可求群体谱，而无需先给每一源指定随机相位。

若圆旋近完全由引力波驱动，$E_{\rm orb}=-\mu(GM\pi f_r)^{2/3}/2$，故 $\d E_{\rm GW}/\d\ln f_r=-\d E_{\rm orb}/\d\ln f_r=(\pi G)^{2/3}\mathcal M^{5/3}f_r^{2/3}/3$。代入 \eqref{eq:gw-cpopulation} 得
\[
 h_c^2(f)=\frac{4G^{5/3}}{3\pi^{1/3} c^2 f^{4/3}}
   \int n(z,\vartheta)\frac{\mathcal M^{5/3}}{(1+z)^{1/3}}
      \,\d z\,\d\vartheta .
\]
因此 $h_c\propto f^{-2/3}$。幂律依赖圆轨道与辐射驱动，且假定积分中没有随 $f$ 变化的并合截断或质量区间。在星系动力学主导的阶段，轨道束缚能的差大部交给了恒星，以上消元不再成立。此时应直接写 $\d E_{\rm GW}/\d\ln f_r=L_{\rm GW}\,f_r/\dot f_r$。

在上面的硬化模型，$\dot f_r/f_r=-(3/2)\dot a/a$，故每个频率的能谱等于纯辐射能谱乘
\[
 \frac{|\dot a_{\rm GW}|}{|\dot a_\star|+|\dot a_{\rm GW}|}
 =\frac{1}{1+(a/a_{\rm b})^5}
 =\frac{1}{1+(f_{\rm b}/f_r)^{10/3}} .
\]
低频时 $\d E_{\rm GW}/\d\ln f_r\propto f_r^4$，故此简化模型中 $h_c\propto f$，而纯辐射段仍为 $f^{-2/3}$。低频翻转说明源在该处停留得短、放入波的能量少，不能只从某个给定间距的瞬时四极振幅推断。若气体迁移、恒星密度耗尽或偏心率改变，翻转位置及幂律都会变。在一族不同质量和星系环境的系统中，各自的 $f_{\rm b}/(1+z)$ 也会摊宽转折。

偏心轨道又把同一时刻的功率分配到不同频率。若 $f_r=n f_{\rm orb}$，第 $n$ 谐波功率占总辐射的比例为 $g(n,e)/F(e)$，其中 $F(e)=\sum_{n\ge1}g(n,e)$。在辐射驱动的阶段，给定 $f_r$ 时每个 $n$ 须在不同的轨道时刻代入 $e=e(f_r/n)$。与同质量圆轨道的能谱相较，得到因子
\[
 \sum_{n\ge1}\left(\frac2n\right)^{2/3}
      \left.\frac{g(n,e)}{F(e)}\right|_{f_{\rm orb}=f_r/n}.
\]
$(2/n)^{2/3}$ 来自轨道结合能中的 $f_{\rm orb}^{2/3}$，因 $\d\ln f_r=\d\ln f_{\rm orb}$ 不再多一个 $n$。当 $e=0$ 只有 $n=2$，精确还原圆轨道；高偏心率把低频能量搬到许多高次谐波。若给每个频率都代同一个“平均偏心率”，就破坏了同一源的轨道演化。

有限观测的每个频率箱实际含多少颗源？若频宽 $\Delta f$，源系轨道在箱内停留 $\Delta f_r/\dot f_r$，一个壳层内的期望数为 $\int{\cal R}(\d V_c/\d z)(1+z)\Delta f/\dot f_r\,\d z\,\d\vartheta$；式中 $\Delta f_r=(1+z)\Delta f$，速率按源系时间计。若它远小于一，连续背景的平均不再代表一般的单次天空；某些频率箱完全空，另一些含一个特别重且近的源。高质量二元组还会在较低的终止频率离开积分，进一步压低高频谱。估计率时可把星系质量函数、具有近伴的比例、并合时间标度及黑洞质量与核球关系相乘积分；从星系配对到黑洞并合还有延迟和可能停滞。将一次可见近邻源的质量简单乘所有星系的数目，会同时漏掉质量权重、距离分布及频率箱的离散性。


可算出微小偏心率对圆轨道幂律的第一项偏离。用 \S\ref{sec:gw-a-sources} 的 Bessel 谐波展开，在给定轨道频率处有
\[
 g(1,e)=\tfrac{29}{192}e^2,\qquad
 g(2,e)=1-5e^2,\qquad
 g(3,e)=\tfrac{729}{64}e^2
 \qquad \text{（略去 }O(e^4)\text{）}.
\]
其和为 $1+157e^2/24$，恰与总功率 $F(e)$ 一致。注意这个核对是在同一轨道时刻作的。设 $e_2$ 是系统在 $f_{\rm orb}=f_r/2$ 的偏心率，辐射驱动的小偏心解 $e\propto f_{\rm orb}^{-19/18}$ 给 $e_n=e_2(n/2)^{19/18}$，各 $n$ 的贡献还须乘 $(2/n)^{2/3}$。于是谱相对圆轨道的修正因子为
\[
 1+e_2^2\left[
 \frac{29}{192}\left(\frac12\right)^{13/9}
 +\frac{729}{64}\left(\frac32\right)^{13/9}
 -5-\frac{157}{24}\right]+O(e_2^4).
\]
正的两项是别的轨道时刻搬来的功率。故在同一 $f_r$、三个相关谐波都已进入辐射驱动阶段且偏心率都很小时，可比单一圆轨道略增强；这不与低频被削弱矛盾。若系统在 $f_{{\rm orb},0}$ 才从环境驱动转为辐射驱动，则 $f_r/n<f_{{\rm orb},0}$ 的那些初期谐波不能仍沿这个小偏心外推。实际处理可给每个源抽取 $(m_1,m_2,e_0,a_0)$，由 $a=(GM)^{1/3}(2\pi f_r/n)^{-2/3}$ 和偏心演化积分求 $e_n$，执行谐波求和，再放入红移及质量积分。$M$ 在这里通过从初始间距到初始频率的转换进入，即使知道整群源的啁啾质量分布，仍不足以固定背景转折。

\subsection{星系对的计数怎样进入应变积分？}

为把 \eqref{eq:gw-cpopulation} 化成可估的数，考虑质量为 $M_g$ 的星系及与它配对的成员。记 $\Phi_g(M_g,z)=\d n_g/\d\ln M_g$ 为共动数密度；$\d f_{\rm pair}/\d\ln q_g$ 为在指定投影分离下找到质量比 $q_g$ 伴星的比例。若这类配对在源系保持可见构型的时间为 $\tau_{\rm pair}$，近稳态时，每共动体积、每源系时间的星系并合率密度可近似取
\[
 \frac{\d{\cal R}_g}{\d\ln M_g\,\d\ln q_g}
 \simeq \frac{\Phi_g(M_g,z)}{\tau_{\rm pair}(M_g,q_g,z)}
             \frac{\d f_{\rm pair}}{\d\ln q_g}.
\]
配对定义若任一成员都可入主目录，应避免一对数两次；有些投影近的星系不束缚，$\tau_{\rm pair}$ 的标定也随选择半径而变。这些量通过成对星系样本和动力学模型获得，和引力波探测器的噪声是独立的误差来源。若只知道一次快照中有多少对而不知道它们保持此状态多久，就无法从数密度唯一推到合并率。

假定第一个星系的中心黑洞质量满足 $\ln m=\ln m_0+\alpha\ln(M_g/M_{g,0})+\delta$，$\delta$ 为标准差 $\sigma_b$ 的零均值 Gaussian 散布。若双黑洞的质量比已给定，而两者质量共同按这个随机比例变化，则能谱的质量权重含 $\langle m^{5/3}\rangle=\bar m_{\rm med}^{\,5/3}\exp(25\sigma_b^2/18)$，$\bar m_{\rm med}$ 为中位质量。这来自正态分布的指数矩，并非将平均质量直接作 $5/3$ 次方。例 $\sigma_b=0.3\ln10$ 时，这个因子约为 $1.94$，应变平方受影响比星系数显著。若两颗星的黑洞质量散布独立，质量比也随机，必须对 $\mathcal M^{5/3}=m_1m_2/(m_1+m_2)^{1/3}$ 联合平均；不能仍乘同一个校正因子。

从星系相接到黑洞并合还存在延迟。令 $p(t_d|\vartheta)$ 为给定源参数的延迟分布，积分可小于一以容许一部分系统长期停滞，则黑洞并合率是 $\int_0^t{\cal R}_g(t-t_d)p(t_d|\vartheta)\,\d t_d$ 再对质量映射求平均。若在一个时段里 ${\cal R}_g$ 近常数且所有系统都会在短时间内合并，延迟对当前总率影响很小；在恒星形成或星系并合率变化剧烈的红移区间，同一延迟会明显移动 \eqref{eq:gw-cpopulation} 的积分权重。对于在最低频停留时间已可与宇宙演化相比的系统，连“按最终并合率乘静态停留时间”都须重审。可改为演化共动分布 $N(f_r,t)$，令每单位源系频率的数密度满足 $\partial_t N+\partial_{f_r}(\dot f_r N)=S_{\rm form}-S_{\rm rem}$。在给定频带没有源汇且近稳态时 $\dot f_r N$ 为常数，才还原前文的停留时间估计。若偏心率同时演化，还须给频率--偏心率平面加 $\partial_e(\dot e N)$。这清楚区分“以稳态简化积分”与“并合率确已直接测过”。

最后问许多源是否真的贡献近乎 Gaussian 的噪声。对一箱内独立均匀的相位 $\phi_i$，取单一正交分量 $X=\sum_i a_i\cos\phi_i$，则 $\langle X^2\rangle=\sum_i a_i^2/2$。四阶累积量为各源累积量之和，单源有 $\langle\cos^4\phi\rangle-3\langle\cos^2\phi\rangle^2=-3/8$，从而
\[
 \frac{\langle X^4\rangle-3\langle X^2\rangle^2}{\langle X^2\rangle^2}
 =-\frac{3}{2N_{\rm eff}},\qquad
 N_{\rm eff}=\frac{(\sum_i a_i^2)^2}{\sum_i a_i^4}.
\]
在给定一组振幅、仅对相位平均时，大量等振幅源给 $N_{\rm eff}=N$，一个极强近邻源主导时却接近一。若还对每箱源数和距离总体取平均，稀有大振幅源可以产生另一种尾部，此式描述的只是相位混合。它解释何以很平滑的理论平均谱也可对应有显著各向异性、少数源可单独拟合的真实时序。要把这样的源分布化为观者的实测残差，还需计时阵列的角相关与地球、脉冲星两端的响应，见 \S\ref{sec:gw-e-pta}。

\section{膨胀背景中的长度与时间}\label{sec:gw-dbg}

一列波离开双星后，沿途穿过的宇宙可以比它发射时大许多。久期以源的钟为准，频率却由接收者的钟测定；若要从今天的波谱追溯极早期的物理，还得问各频率有多久处在 Hubble 半径以内。本节取 $c=\hbar=k_{\rm B}=1$，保留 $G$，记约化 Planck 质量 $M_{\rm P}=(8\pi G)^{-1/2}$。下标 $0$ 指今天，并选 $a_0=1$。空间平直时度规为 $\mathrm ds^2=a^2(\eta)(-\mathrm d\eta^2+\mathrm d\bm x^2)$，$\mathrm dt=a\,\mathrm d\eta$。点与撇分别表示对共动观者的固有时 $t$、共形时 $\eta$ 求导。特别要区分 $H=\dot a/a$ 和 $\mathcal H=a'/a=aH$。

\subsection{一把共动的尺与一只共动的钟}

给共动格点 $\bm x$ 刻上记号，同一时刻测得的小段长度是 $a|\mathrm d\bm x|$。有质量的粒子从一个格点跑到另一个，格点上所见速度是 $\bm v=a\,\mathrm d\bm x/\mathrm dt$。由于平直空间平移对称，粒子的协变动量 $p_i=g_{i\mu}p^\mu$ 沿测地线不变。共动观者的空间标架是 $a^{-1}\partial_i$，故实际测得的动量是 $p_{\hat i}=p_i/a$。光子由此给出 $E=|\bm p|\propto a^{-1}$，冷物质则给出 $v\propto a^{-1}$。把光信号换成高频引力波，其波矢沿零测地线平行移动，亦有 $f_{\rm o}=f_{\rm e}/(1+z)$，其中 $1+z=1/a_{\rm e}$。

不用粒子语言也能看出红移。一个波峰由 $\eta_{\rm e}$ 出发，在 $\eta_0$ 到达，其行过的径向共动距离 $\chi=\eta_0-\eta_{\rm e}$。下一个波峰在同两条共动世界线之间传播，故 $\delta\eta_0=\delta\eta_{\rm e}$，两端各乘以本地 $a$ 即得 $\delta t_0/\delta t_{\rm e}=1+z$。角距离与光度距离也可由同一光锥读出。径向坐标为 $\chi$ 的球，面积为 $4\pi a^2 S_K^2(\chi)$，其中 $S_K(\chi)$ 在 $K>0,\,K=0,\,K<0$ 时分别为 $\sin(\sqrt K\chi)/\sqrt K,\,\chi,\,\sinh(\sqrt{-K}\chi)/\sqrt{-K}$。源上横向固有长度 $l$ 张角 $l/[a_{\rm e}S_K(\chi)]$。若各向同性源的光度为 $L$，光子每个能量少一因子 $1+z$，单位接收时内到达的个数又少同一因子。从而
\begin{equation}\label{eq:gw-d-dist}
 \chi(z)=\int_0^z\frac{\mathrm dz'}{H(z')},\qquad
 d_A=\frac{S_K(\chi)}{1+z},\qquad
 F=\frac{L}{4\pi d_L^2},\qquad d_L=(1+z)^2d_A .
\end{equation}
最后一步使用光子数守恒和几何光学。引力波的几何传播也有面积稀释，但把波幅解释成距离时还要处理源本身的频率标度。

\subsection{哪一项主宰膨胀}

共动体积 $a^3$ 中的流体满足 $\mathrm d(\rho a^3)=-p\,\mathrm d(a^3)$。若各组分不交换能量，状态方程 $p_i=w_i\rho_i$ 给出 $\mathrm d\ln\rho_i/\mathrm d\ln a=-3(1+w_i)$。尘埃只有数密度的 $a^{-3}$，辐射还多出每个量子红移的 $a^{-1}$，宇宙常数的负压则恰好补偿体积增大。由 Einstein 方程的 $00$ 分量及其时间导数得
\begin{align}
 H^2&=\frac{\sum_i\rho_i}{3M_{\rm P}^2}-\frac{K}{a^2},&
 \dot H&=-\frac{\sum_i(\rho_i+p_i)}{2M_{\rm P}^2}+\frac{K}{a^2},\nonumber\\
 \frac{\ddot a}{a}&=-\frac{\sum_i(\rho_i+3p_i)}{6M_{\rm P}^2}.
 \label{eq:gw-d-fried}
\end{align}
不能以密度随 $a$ 衰减的快慢判断加速，只能以 $\rho+3p$ 的符号判断。取 $\rho_{c0}=3M_{\rm P}^2 H_0^2$，设 $\Omega_{i0}=\rho_{i0}/\rho_{c0}$，$\Omega_{K0}=-K/H_0^2$。若有辐射、冷物质及恒定真空能，今天的约束是各 $\Omega_{i0}$ 与 $\Omega_{K0}$ 之和等于一；任意时刻则有
\[
 E^2(a)\equiv H^2(a)/H_0^2
 =\Omega_{r0}a^{-4}+\Omega_{m0}a^{-3}+\Omega_{K0}a^{-2}+\Omega_{\Lambda0}.
\]
若暗能量状态方程随红移变，将最后一项乘以 $\exp[3\int_0^z(1+w(z'))\mathrm dz'/(1+z')]$ 即可。这个因子进入式 \eqref{eq:gw-d-dist} 的积分，所以用距离反推 $w(z)$ 总会混入较低红移的历史。

平坦宇宙由一种常数 $w$ 流体主导时，$H\propto a^{-3(1+w)/2}$。积分一次得 $a\propto t^{2/[3(1+w)]}$，再用 $\mathrm d\eta=\mathrm dt/a$ 得 $a\propto |\eta-\eta_c|^{2/(1+3w)}$。这里排除了使分母为零的情形，$w=-1$ 是 \eqref{eq:dS-flat-slicing} 的指数膨胀，$w=-1/3$ 则给 $a\propto t$、$\eta\propto\ln t$。当 $w=1/3$ 时 $a\propto\eta$，$w=0$ 时 $a\propto\eta^2$。后一对幂恰好决定了张量波方程能否化为平直时空的振子。

在辐射与尘埃交接处不必拼接两个幂律。忽略曲率及真空能，式 \eqref{eq:gw-d-fried} 化为 $(a')^2=H_0^2(\Omega_{r0}+\Omega_{m0}a)$。取大爆炸处 $\eta=0$，两边除以平方根再积分，得到
\begin{equation}\label{eq:gw-d-equality}
 a(\eta)=H_0\sqrt{\Omega_{r0}}\,\eta+\frac14 H_0^2\Omega_{m0}\eta^2
 =a_{\rm eq}\left[\left(1+\frac{\eta}{\eta_1}\right)^2-1\right],
 \qquad \eta_1=\frac{2\sqrt{\Omega_{r0}}}{H_0\Omega_{m0}} .
\end{equation}
物质与辐射等密度要求 $a_{\rm eq}=\Omega_{r0}/\Omega_{m0}$，故 $\eta_{\rm eq}=(\sqrt2-1)\eta_1$。此时给扰动留下的共动尺度是 $k_{\rm eq}=a_{\rm eq}H_{\rm eq}=H_0\sqrt{2\Omega_{m0}/a_{\rm eq}}$。等密度是引力势演化的转折；原子形成及辐射解耦另由电子数与散射率决定，二者的时刻无需相同。

\subsection{热史留下的频标}

早期等温热浴的能量密度和熵密度分别写作 $\rho_r=(\pi^2/30)g_\rho T^4$、$s=(2\pi^2/45)g_sT^3$。这里 $g_\rho$ 对每个玻色态计一、对同温费米态计 $7/8$；若某类粒子早已脱耦，还要按它与光子的温度比的四次方计权。熵的相应权重是三次方。绝热膨胀保持 $sa^3$，故 $aTg_s^{1/3}$ 为常数。随温度下降，某种有质量的粒子湮灭会加热仍与它相耦合的热浴；早已脱耦的引力波保持 $f\propto a^{-1}$，不会分到这些熵。

以中微子为例。在瞬时脱耦近似下，电子、正电子湮灭前电磁热浴有 $2+(7/8)4=11/2$ 个熵自由度，湮灭后只剩光子的两个。中微子温度继续按 $a^{-1}$ 降低，故 $T_\nu/T_\gamma=(4/11)^{1/3}$。把任何额外的相对论性能量都用同一单位记帐，就得到
\begin{equation}\label{eq:gw-d-neff}
 \rho_r=\rho_\gamma\left[1+\frac78\left(\frac4{11}\right)^{4/3}N_{\rm eff}\right],
 \qquad
 \frac{\rho_{{\rm gw},0}}{\rho_{\gamma0}}
 =\frac78\left(\frac4{11}\right)^{4/3}\Delta N_{\rm eff}
\end{equation}
的等效关系。第二式适用于在所考察时代已在 Hubble 半径之内、平均后像自由辐射那样膨胀的波。尚被冻结的长波不能直接填入热浴密度。实际的中微子脱耦有有限宽度，故 $N_{\rm eff}$ 是能量密度的度量，无需恰为整数。

温度 $T_*$ 的辐射时代有 $H_*=(\pi^2g_{\rho *}/90)^{1/2}T_*^2/M_{\rm P}$。若一个源的典型角频率为 $\omega_*=\lambda H_*$，今天测得的频率为
\begin{equation}\label{eq:gw-d-fstar}
 f_0=\frac{\lambda}{2\pi}\frac{T_0}{T_*}
       \left(\frac{g_{s0}}{g_{s*}}\right)^{1/3}
       \left(\frac{\pi^2g_{\rho *}}{90}\right)^{1/2}
       \frac{T_*^2}{M_{\rm P}} .
\end{equation}
$T_*$ 越高，红移越大；但同一温度下 $H_*$ 增长得更快，净效应反而是 $f_0\propto T_*$。例如取 $g_{s*}\simeq g_{\rho *}\simeq100$、$T_*=100\,{\rm GeV}$、$T_0=2.725\,{\rm K}$、$g_{s0}\simeq3.91$，得 $a_*H_*\simeq1.65\times10^{-5}\,{\rm s}^{-1}$。若尺度定义为 $k_*/a_*=H_*$，相应循环频率还要除以 $2\pi$；若直接给出的是 $f_*/H_*$，则乘 $a_*H_*$ 即可。气泡相变的尺寸可远小于 $H_*^{-1}$，所以同一热史并不唯一指定峰频。

类似地，产生时占总能量的一小部分 $\omega_*^{(r)}=\rho_{{\rm gw},*}/\rho_{r,*}$ 到今天变为
\begin{equation}\label{eq:gw-d-gstar}
 \Omega_{{\rm gw},0}=\Omega_{r0}
 \frac{g_{\rho *}}{g_{\rho0}}\left(\frac{g_{s0}}{g_{s*}}\right)^{4/3}
 \omega_*^{(r)} .
\end{equation}
这是将波的 $a^{-4}$ 与热浴的 $g_\rho T^4$ 分别传播所得。对于一个频谱，应将 $\omega_*^{(r)}$ 换为每单位 $\ln f_*$ 的能量分数，在 $f_*=f_0/a_*$ 处取值；$\mathrm d\ln f_*=\mathrm d\ln f_0$，谱密度上没有额外的频率 Jacobian。


\subsection{径向尺度和横向尺度的一次换算}

给定 $H(z)$，对光锥在观者附近作 Taylor 展开可以检查距离拟合。令 $q_0=-\ddot a_0/(a_0H_0^2)$，由 $\dot z=-(1+z)H$（沿背景历史求导）得 $(\mathrm dH/\mathrm dz)_0=H_0(1+q_0)$。代 \eqref{eq:gw-d-dist} 逐阶积分，曲率项 $S_K(\chi)-\chi$ 直到 $\chi^3$ 才出现，故 $d_L=\frac1{H_0}\left[z+\frac{1-q_0}{2}z^2+O(z^3)\right]$。这里求的是沿观者过去光锥的距离，固定共动星系在观者\emph{不同时刻}的红移漂移却须对 $1+z=a_0/a_{\rm e}$ 连发两个光子，$\mathrm dt_{\rm e}/\mathrm dt_0=(1+z)^{-1}$，于是 $\mathrm dz/\mathrm dt_0=(1+z)H_0-H(z)$。两者在给宇宙学双星计时或转换频率导数时不能混为一谈。

在仅有尘埃、空间平直的模型，$H=H_0(1+z)^{3/2}$，故 $\chi=2[1-(1+z)^{-1/2}]/H_0$。角距离 $d_A=\chi/(1+z)$ 先增后减，导数为零的条件是 $(1+z)/H(z)=\chi$，代入得 $z=5/4$。一支同样固有尺寸的标准尺即使放得更早，视角也会重新变大；光度距离却继续增加。若把角尺寸用于标定位于宿主中的引力波源，直接用 $d_L$ 作横向尺度会差两个红移因子。对壳层计数，单位立体角、单位红移的共动体积为 $S_K^2(\chi)/H(z)$，物理体积还要乘 $(1+z)^{-3}$。一个给在每共动体积、每\emph{源时间}的事件率 $\mathcal R_s(z)$ 则转换为
\[
 \frac{\mathrm d\dot N_{\rm obs}}{\mathrm dz}
 =\frac{4\pi S_K^2(\chi)}{(1+z)H(z)}\,\mathcal R_s(z).
\]
体积若已经是共动的，只需一个时间膨胀因子。按物理密度计算弦圈时则还要三个体积因子；不同演化积分里常见的不同红移幂多源于此处的定义。

\subsection{解耦的两把尺}

对一个组分从相对论性到非相对论性的演化，可以从同一个相空间分布直接计算。将守恒的共动动量记为 $q$、$E_q=\sqrt{q^2+a^2m^2}$，$g$ 为内部态数。无碰撞组分有 $f(q)$ 不随 $\eta$ 变，在共动观者的标架中
\[
 \rho=\frac{g}{2\pi^2a^4}\int q^2 E_q f(q)\,\mathrm dq,\qquad
 p=\frac{g}{6\pi^2a^4}\int\frac{q^4}{E_q} f(q)\,\mathrm dq .
\]
微分 $a^4\rho$ 时只有 $E_q$ 贡献；用 $a\partial_a E_q=a^2m^2/E_q=E_q-q^2/E_q$，便有 $a\partial_a\rho=-3(\rho+p)$。因此一个脱耦组分从相对论性过渡到非相对论性时，$w(a)$ 一般不是常数。对中微子自由流阻尼，只应把当时能以相对论速度带走剪切的组分放进无质量核；今天从静质量得到的物质份额不能原封不动代进去。

光子最后散射附近，决定散射率的是自由电子，不是宇宙中所有电子。以纯氢为例，在化学平衡且气体不简并时，$\mu_p+\mu_e=\mu_{\rm H}$，非相对论理想气体有 $n_i=g_i(m_iT/2\pi)^{3/2}\exp[(\mu_i-m_i)/T]$。用电荷中性 $n_e=n_p$，令 $x_e=n_e/n_b$，忽略激发态及约为一的简并比，便得 Saha 式
\[
 \frac{x_e^2}{1-x_e}=\frac{1}{n_b}\left(\frac{m_eT}{2\pi}\right)^{3/2}
              e^{-B_{\rm H}/T},
 \qquad n_b=\eta_b\,\frac{2\zeta_{\rm R}(3)}{\pi^2}T^3 ,
\]
其中 $B_{\rm H}$ 是束缚能，$\zeta_{\rm R}$ 表示 Riemann zeta 函数。由于每个重子对应极多光子，$T$ 必须显著低于 $B_{\rm H}$ 才会大幅减少自由电子。Saha 式给平衡电离分数；真正的再组合需考虑 Lyman 光子的再吸收、两光子跃迁以及膨胀，$x_e$ 最终偏离平衡值。$\Gamma_T=n_e\sigma_T$ 与 $H$ 的比值给每 Hubble 时间的散射次数，光学深度 $\tau=\int\Gamma_T\,\mathrm dt$ 才给从某点到今天一路无散射的概率。后一量定义的可见度峰与任一局域反应率等于 $H$ 的时刻并非严格相同。

原初声学标尺可在尘埃加辐射的背景精确积出。将 $R_b=R_{\rm eq}a/a_{\rm eq}$ 与 $c_s=[3(1+R_b)]^{-1/2}$ 放入 $r_s(a)=\int_0^a c_s\,\mathrm da/(a^2H)$，用 $k_{\rm eq}$、$y=a/a_{\rm eq}$ 可以化为
\[
 r_s(a)=\frac{\sqrt{2/3}}{k_{\rm eq}}\int_0^y
             \frac{\mathrm dy'}{\sqrt{(1+y')(1+R_{\rm eq}y')}}
 =\frac{2}{k_{\rm eq}}\sqrt{\frac{2}{3R_{\rm eq}}}
   \ln\frac{\sqrt{1+R_b}+\sqrt{R_b+R_{\rm eq}}}{1+\sqrt{R_{\rm eq}}}.
\]
当 $R_{\rm eq}\to0$ 时展开对数，得 $r_s=\eta/\sqrt3$；重子数增加，流体同一时间内的声学距离变短。把这个长度投影到末次散射面的共动距离上，可以一边量出角声学尺度，一边检验 \eqref{eq:gw-d-dist} 的晚期距离。这是联合微波背景与标准汽笛判断膨胀史时，二者所测几何量的差别。

\subsection{从膨胀背景中剥出 Newton 动力学}

Hubble 半径以内、速度非相对论、势弱时，可取一团物质的固有位置 $\bm r=a\bm x$。均匀背景在该点的相对加速度为 $(\ddot a/a)\bm r$。定义特异加速度 $\bm g=-\nabla_{\bm x}\Psi/a$，设 $\ddot{\bm r}=(\ddot a/a)\bm r+\bm g$；将 $\bm r=a\bm x$ 求两次导数，背景项精确抵消，得 $\ddot{\bm x}+2H\dot{\bm x}=-\nabla_{\bm x}\Psi/a^2$。Poisson 方程的源因而是 $\rho-\bar\rho$，不会有对无限均匀密度直接积分产生的歧义。写 $\bm v=a\dot{\bm x}$，保留无压流体的非线性项，有
\begin{align*}
 \dot\delta+\frac1a\nabla\cdot[(1+\delta)\bm v]&=0,\\
 \dot{\bm v}+H\bm v+\frac1a(\bm v\cdot\nabla)\bm v&=-\frac1a\nabla\Psi,\qquad
 \nabla^2\Psi=4\pi G a^2\bar\rho\,\delta .
\end{align*}
这些是短波方程，不能用其齐次解代替大尺度 相对论初值。相反，小于 Hubble 半径并不自动意味着可忽略流体压强，须再与 Jeans 长度比较。

可用一团已束缚物质作检验。若 $\Psi$ 来自近似孤立的质量 $M$，固有半径满足 $\ddot r\simeq-GM/r^2+(\ddot a/a)r$（此处 $M$ 指相对均匀背景的质量过量）。两项之比的量级为 $H^2r^3/(GM)$，在内部动力学时间 $\sqrt{r^3/(GM)}\ll H^{-1}$ 的系统里膨胀只作微小潮汐修正。宇宙学的 Hubble 流主要描述共动观者之间的总体运动，不能把 $\dot a/a$ 当作本地每支杆都应随之变长的运动学定理。

另一可算的例子是尘埃宇宙内球对称的小扰动。若物质在共动半径 $x$ 内平均过密 $\bar\delta(<x)$，连续性方程积分给 $\dot{\bar\delta}=-3v_r/(ax)$（线性阶）。成长模 $\bar\delta\propto D(a)$，记 $f_g=\mathrm d\ln D/\mathrm d\ln a$，则 $v_r=-aH f_g\,x\bar\delta/3$，外层质量流向过密中心。对一傅里叶模同一式变成 $\nabla\cdot\bm v=-aH f_g\delta$。星系的视向本动速度与密度的相关测到的是 $f_g$ 乘聚集振幅；透镜却测沿路两度规势之和。第 3 章用它们比较引力理论，而下一节求这些可测量的密度和势怎样继承原初扰动。
\section{如何分开膨胀时空的扰动}\label{sec:gw-dpert}

在各向同性背景上，一阶场方程保持绕 $\bm k$ 的旋转对称性。一个球标量密度波不可能线性地驱动螺旋度为二的波。这个简单的选择定则使我们既能在同一天空图上辨认张量的印记，也能精确指出何时标量会成为它的源。以下 $k=|\bm k|>0$，采用 $\partial_i\to\mathrm i k_i$ 的 Fourier 约定；严格均匀的 $k=0$ 模式应另外检查它是否只改变背景解。

\subsection{空间分解与规范}

将三个空间坐标上的位移分成梯度与无散部分。一个对称空间张量则分为迹、无迹的二阶标量梯度、横矢量的对称梯度、横向无迹张量。对给定 $\bm k$，两次收缩 $k^i$ 给出两个标量，一次收缩后作横向投影给出两个矢量分量，余下两个分量便是张量。收集这些项，把度规写作
\begin{align}
 g_{00}/a^2&=-(1+2A),&g_{0i}/a^2&=\partial_i B+S_i,\nonumber\\
 g_{ij}/a^2&=(1-2\psi)\delta_{ij}+2\partial_i\partial_j E
 +\partial_i F_j+\partial_j F_i+h_{ij}. \label{eq:gw-d-svt}
\end{align}
这里 $\partial_iS_i=\partial_i F_i=\partial_i h_{ij}=h_{ii}=0$，在此分解中用 $\delta_{ij}$ 缩并空间指标。四个标量、两个横矢量、一个张量共给出 $4+4+2=10$ 个度规分量。

作被动坐标变换 $\widetilde\eta=\eta+T,\ \widetilde x^i=x^i+\partial^iL+L^i$，其中 $\partial_iL^i=0$。把背景度规的 Lie 导数从微扰度规扣去便得变换律。比如 $g_{00}$ 的 Lie 导数含 $T(-a^2)'+2(-a^2)T'$，故 $\widetilde A=A-\mathcal H T-T'$；$g_{0i}$ 的梯度项给 $\widetilde B=B+T-L'$。同理 $\widetilde E=E-L$、$\widetilde\psi=\psi+\mathcal H T$，而 $\widetilde S_i=S_i-L_i'$、$\widetilde F_i=F_i-L_i$。由此可见 $h_{ij}$ 本身已规范不变。定义 $\sigma=B-E'$，则 $\widetilde\sigma=\sigma+T$，组合
\begin{equation}\label{eq:gw-d-bardeen}
 \Psi=A+\mathcal H\sigma+\sigma',\qquad
 \Phi=\psi-\mathcal H\sigma,\qquad V_i=S_i-F_i'
\end{equation}
同样不变。前两者称 Bardeen 势。算题时可取 $B=E=F_i=0$，这样 $A,\psi,S_i$ 直接给出这三个物理量。规范不变仍不等于自由传播。广义相对论的矢量约束把 $k^2 V_i$ 同横向动量密度联系，无矢量应力时演化方程给 $(a^2V_i)'=0$。物质涡旋可以拖曳空间标架，自由的长程矢量引力波却不会由此产生。

物质的变量也随时间切片而变。背景密度 $\bar\rho(\eta)$ 给出 $\widetilde{\delta\rho}=\delta\rho-\bar\rho' T$；在上面的纵向规范中读到的密度恰是任意规范下的 $\delta\rho+\bar\rho'\sigma$。把 $\rho$ 换为 $p$ 得到压力的对应组合。另可问，在各处密度刚好相同的同时面上，空间本身弯曲了多少？令 $T=\delta\rho/\bar\rho'$ 消去密度扰动，得到 $\zeta=-\psi-\mathcal H\frac{\delta\rho}{\bar\rho'}$。这一定义在任意规范中成立；负号使空间度规在均匀密度面上带因子 $1+2\zeta$。选不同符号定义三曲率时，随后的势与 $\zeta$ 的关系也会反号，但功率谱不变。

\subsection{约束、声速与熵}

以下纵向规范度规的标量部分为 $a^2[-(1+2\Psi)\mathrm d\eta^2+(1-2\Phi)\mathrm d\bm x^2]$。写 $\delta T^0{}_0=-\delta\rho$、$\delta T^0{}_i=(\bar\rho+\bar p)v_i$，设速度散度 $\theta=\mathrm i k_iv_i$。定义标量剪切 $\Sigma$ 的号使得 $k^2(\Phi-\Psi)=12\pi G a^2(\bar\rho+\bar p)\Sigma$。线性 $00$、$0i$ 方程依次为
\begin{align}
 k^2\Phi+3\mathcal H(\Phi'+\mathcal H\Psi)&=-4\pi G a^2\delta\rho,\nonumber\\
 k^2(\Phi'+\mathcal H\Psi)&=4\pi G a^2(\bar\rho+\bar p)\theta .
 \label{eq:gw-d-constraints}
\end{align}
消去两式共有的时间导数，得 $-k^2\Phi=4\pi G a^2\bar\rho\Delta$，其中 $\Delta=\delta+3\mathcal H(1+w)\theta/k^2$、$\delta=\delta\rho/\bar\rho$。形式虽如 Poisson 方程，源却是共动切片的密度差。深处 Hubble 半径以内，$\Delta\simeq\delta$ 往往很好；长波极限，纵向规范密度差可非零而 $\Delta$ 却为 $O(k^2/\mathcal H^2)$，故不能拿 Newton 式从 $\delta$ 直接反求 $\Phi$。

每个独立守恒的流体满足
\begin{align}
 \delta'&=-(1+w)(\theta-3\Phi')
       -3\mathcal H\left(\frac{\delta p}{\bar\rho}-w\delta\right),\nonumber\\
 \theta'&=-\mathcal H(1-3w)\theta-\frac{w'}{1+w}\theta
       +k^2\left(\Psi+\frac{\delta p}{\bar\rho+\bar p}-\Sigma\right).
 \label{eq:gw-d-fluid}
\end{align}
第一式中 $3\Phi'$ 是共动小体积所受的度规压缩；第二式是梯度力与膨胀阻尼。对真空能 $\bar\rho+\bar p=0$ 时不可用 $\theta$ 作独立未知量。非相对论物质的 $w$、声速、$\Sigma$ 皆可忽略。将 $\theta$ 消去并忽略短波区势的时间导数，便得到 $\delta_m''+\mathcal H\delta_m'-4\pi G a^2\bar\rho_m\delta_m=0$。这也可由固有速度的 Euler 方程 $\dot{\bm v}+H\bm v=-a^{-1}\nabla\Psi$ 与连续性方程 $\dot\delta=-a^{-1}\nabla\cdot\bm v$ 推出；对 $a^{-1}$ 本身求导才会在固有时方程出现 $2H\dot\delta$。若有压强，则左边还多 $c_s^2k^2\delta$，相应的 Jeans 波数满足 $c_s^2k_J^2=4\pi G a^2\bar\rho$。

区分 $c_a^2=\bar p'/\bar\rho'$ 和实际压强扰动非常重要。一般地 $\delta p=c_a^2\delta\rho+\delta p_{\rm nad}$；最后一项在改变时间切片时不变。所谓\textbf{绝热初值}，就是所有组分在各处仅沿同一背景历史作微小时间平移，$\delta\rho_i/\bar\rho_i'$ 皆相等。因 $\bar\rho_i'=-3\mathcal H(1+w_i)\bar\rho_i$，光子和冷物质的绝热初值必须满足 $\delta_m=3\delta_\gamma/4$。定义 $S_{ij}=3(\zeta_i-\zeta_j)$，便有 $S_{m\gamma}=\delta_m-3\delta_\gamma/4$。可以取总密度扰动为零而 $S_{m\gamma}\ne0$；随后辐射与物质按不同速率稀释，原先的补偿会失效，这种初值称\textbf{等曲率扰动}。

把总能量连续性方程代入 $\zeta$ 的时间导数，$3\Phi'$ 项恰好抵消来自切片曲率的 $-\Phi'$，剩下
\begin{equation}\label{eq:gw-d-zeta}
 \zeta'=-\frac{\mathcal H}{\bar\rho+\bar p}\,\delta p_{\rm nad}-\frac{\theta}{3}
 \qquad\text{（纵向规范）}.
\end{equation}
对规则的超 Hubble 模式，$\theta$ 有两个空间梯度，故绝热时 $\zeta$ 守恒。若只保留一个绝热流体并令 $\Sigma=0$，由 \eqref{eq:gw-d-constraints} 的长波 $00$ 方程又得 $\delta=-2(\Phi+\Phi'/\mathcal H)$。常 $w$ 的常势解因而满足 $\zeta=-(5+3w)\Phi/[3(1+w)]$，在辐射主导到物质主导的转变中 $\Phi_{\rm m}/\Phi_{\rm r}=(3/5)/(2/3)=9/10$。这个转变可以只用守恒量求得，无需把波逐个数值推进穿过等密度时刻。

\subsection{从曲率初值到物质谱}

简单情形可将流体变量完全消去。当宇宙由一个无剪切、绝热组分主导时，$\Phi=\Psi$。空间 Einstein 方程的迹给出 $\Phi''+3\mathcal H\Phi'+(2\mathcal H'+\mathcal H^2)\Phi=4\pi G a^2\delta p$。用 $00$ 方程将右边的 $c_a^2\delta\rho$ 换成势，得
\begin{equation}\label{eq:gw-d-pot}
 \Phi''+3(1+c_a^2)\mathcal H\Phi'
 +[c_a^2 k^2+2\mathcal H'+(1+3c_a^2)\mathcal H^2]\Phi=0 .
\end{equation}
特别地，若平坦且 $c_a^2=w$ 恒定，方括号最后两项抵消。物质时代 $w=0$、$\mathcal H=2/\eta$，两个解为常数及 $\eta^{-5}$。辐射时代令 $y=k\eta/\sqrt3$，方程化为 $\partial_y^2\Phi+(4/y)\partial_y\Phi+\Phi=0$。规则解和由约束求出的密度差是
\begin{align}
 \Phi&=3\Phi_{\rm p}\frac{\sin y-y\cos y}{y^3},\nonumber\\
 \delta_r&=-2\big[(1+y^2)\Phi+y\partial_y\Phi\big].
 \label{eq:gw-d-phirad}
\end{align}
$y\to0$ 时分子为 $y^3/3$，故 $\Phi\to\Phi_{\rm p}$。$y\gg1$ 时势按 $y^{-2}$ 衰减，密度却留下 $6\Phi_{\rm p}\cos y$ 的声振荡。声速决定的是声学距离 $\eta/\sqrt3$，引力波的振荡自变量则是 $k\eta$。

一小份冷物质处在这个辐射背景中，会怎样？深处声学尺度以内，辐射势衰减，其短时驱动过去后式 \eqref{eq:gw-d-fluid} 给 $\delta_m''+\mathcal H\delta_m'\simeq0$。因 $\mathcal H=1/\eta$，增长只呈 $\ln\eta$。若继续追踪到等密度以后，令 $y_m=a/a_{\rm eq}$，近似把辐射当作平滑的膨胀组分，只留下冷物质自身的密度作短波引力源。代 $\mathrm d/\mathrm dt=y_m H\,\mathrm d/\mathrm dy_m$ 及 $4\pi G\bar\rho_m=(3/2)H^2 y_m/(1+y_m)$，得 Mészáros 方程
\[
 y_m(1+y_m)\frac{\mathrm d^2\delta_m}{\mathrm dy_m^2}
 +\left(1+\frac{3y_m}{2}\right)\frac{\mathrm d\delta_m}{\mathrm dy_m}
 -\frac32\delta_m=0 .
\]
一个解直接试得 $D_1=1+3y_m/2$。另取 $D_2=D_1 u$，代回得 $u'\propto[y_m\sqrt{1+y_m}\,D_1^2]^{-1}$。积分后可选
\[
 D_2=\left(1+\frac{3y_m}{2}\right)
 \ln\frac{\sqrt{1+y_m}+1}{\sqrt{1+y_m}-1}-3\sqrt{1+y_m}.
\]
在 $y_m\ll1$ 时它含 $\ln(4/y_m)$，在 $y_m\gg1$ 时前两个正幂抵消，剩下 $y_m^{-3/2}$ 的衰减模。故实际增长解的早期系数要同辐射势驱动阶段匹配。若只把 $D_1$ 从零时刻外推，就会漏掉短波在等密度前已积累的对数。

随机初值的振幅由二点函数定义。设 $\langle X(\bm k)X^*(\bm q)\rangle=(2\pi)^3\delta^3(\bm k-\bm q)P_X(k)$，引入 $\mathcal P_X=k^3P_X/(2\pi^2)$ 后，$\langle X(\bm x)^2\rangle=\int\mathcal P_X\,\mathrm d\ln k$。平移对称给出动量的 $\delta$ 函数，各向同性才使 $P_X$ 只依赖 $k$。在枢轴波数 $k_p$ 附近，以 $\mathcal P_\zeta=A_s(k/k_p)^{n_s-1}$ 定义标量倾斜；尺度不变是每一相同的对数间隔贡献相同，而非 $P_\zeta$ 为常数。

将增长因子 $D(a)$ 归一在物质时代 $D/a\to1$。把等密度转变对每一波数的影响写成 $T_m(k)$，取长波端 $T_m(0)=1$。深处 Hubble 半径以内，物质时代的势为 $-(3/5)\zeta T_m D/a$，用约束式得到
\begin{equation}\label{eq:gw-d-matterpower}
 \Delta_m(k,a)=\frac{2k^2}{5\Omega_{m0}H_0^2}\,T_m(k)D(a)\zeta(k),
 \qquad
 \mathcal P_{\Delta_m}
 =\left[\frac{2k^2T_m D}{5\Omega_{m0}H_0^2}\right]^2\mathcal P_\zeta .
\end{equation}
先在辐射时代进入 Hubble 半径的模式只作对数增长，到等密度时相对未受压低的 $k^2\zeta$ 已差一因子约 $(k_{\rm eq}/k)^2\ln(k/k_{\rm eq})$。物质谱的转折由此有动力学意义。若在晚期还有平滑的宇宙常数，增长方程有衰减解 $H(a)$，降阶法给增长解 $D(a)=(5\Omega_{m0}/2)E(a)\int_0^a\mathrm da'/[a'^3 E(a')^3]$，此式忽略辐射。在纯尘埃极限 $E=a^{-3/2}$，积分正比 $a^{5/2}$，确实还原 $D=a$。真空能主导后 $H$ 趋常数，积分收敛，密度差近于冻结而势按 $D/a$ 降低。穿越势阱的光子因而留下净红移或蓝移。


\subsection{一个冷物质模式怎样穿过辐射声波}

前面由量纲找到短波 $T_m$ 的对数。还可在辐射主导时从规则的绝热初值精确算出这个对数的系数。暂略去中微子剪切，将冷物质视为随辐射势运动的试验流体。取 $y=k\eta/\sqrt3$ 与式 \eqref{eq:gw-d-phirad} 的 $\Phi$；冷物质的 Euler 方程是 $(\eta\theta_c)'=\eta k^2\Phi$。用 $(\sin y/y)'=-(\sin y-y\cos y)/y^2$ 积分，舍去不规则的齐次速度解，得 $\eta\theta_c=9\Phi_{\rm p}\left(1-\frac{\sin y}{y}\right)$。在 $y\ll1$ 时有 $\theta_c\simeq k^2\eta\Phi_{\rm p}/2$，与初始引力加速相符。接着积分 $\delta_c'=-\theta_c+3\Phi'$ 并以 $\delta_c(0)=-3\Phi_{\rm p}/2$ 归一。定义余弦积分 $\operatorname{Ci}(y)=\gamma_{\rm E}+\ln y+\int_0^y(\cos u-1)\mathrm du/u$；由于
$\int_0^y[1-\sin u/u]\mathrm du/u=\gamma_{\rm E}+\ln y-\operatorname{Ci}(y)+\sin y/y-1$，得
\begin{equation}\label{eq:gw-d-coldrad}
 \delta_c=-\frac32\Phi_{\rm p}+3(\Phi-\Phi_{\rm p})
 -9\Phi_{\rm p}\left[\gamma_{\rm E}+\ln y-\operatorname{Ci}(y)
                       +\frac{\sin y}{y}-1\right].
\end{equation}
长波端方括号趋零，短波端则有 $\delta_c\simeq-9\Phi_{\rm p}(\ln y+\gamma_{\rm E}-1/2)$。辐射声波离开之后，冷物质仍保有它经过时赋予的速度；连续性方程对这个按 $\eta^{-1}$ 衰减的速度积分，恰积成对数。若从一个“无驱动、零速度”的冷物质初值出发，就会错过这一项。

现在取 $k\gg k_{\rm eq}$，在 $k^{-1}\ll\eta\ll\eta_{\rm eq}$ 的重叠区，将式 \eqref{eq:gw-d-coldrad} 与 Mészáros 的两支解相配。用 $y_m=a/a_{\rm eq}\simeq2\eta/\eta_1$，由式 \eqref{eq:gw-d-equality} 有 $\eta_1=2\sqrt2/k_{\rm eq}$。故 $y\simeq\sqrt{2/3}(k/k_{\rm eq})y_m$。此时 $D_1\simeq1$，$D_2\simeq\ln(4/y_m)-3$。写 $\delta_c=A D_1+B D_2$，比较 $\ln y_m$ 系数及常数，得
\[
 B=9\Phi_{\rm p},\qquad
 A=-9\Phi_{\rm p}
   \left[\ln\left(4\sqrt{\frac23}\frac{k}{k_{\rm eq}}\right)
   +\gamma_{\rm E}-\frac72\right].
\]
到 $y_m\gg1$，只有 $A$ 会乘上线性增长的大因子。这个渐近式只适用于对数已足够大的波数，并忽略重子、中微子及辐射--尘埃过渡中的残余驱动；不能让其在过渡区的外推零点充当真实的转移函数零点。它说明短波的物质谱并非仅凭一段 $D\sim a$ 的晚期增长就能重建。

\subsection{两种初始切片与两种熵}

再把超 Hubble 半径的关系写得完整些。在纵向规范中 $v_i=-\mathrm i k_i\theta/k^2$。为令总动量密度为零，取时间平移 $T=\theta/k^2$；共动曲率与均匀密度曲率分别为
\[
 \mathcal R=-\Phi-\mathcal H\frac{\theta}{k^2},\qquad
 \zeta=-\Phi+\frac{\delta}{3(1+w)},\qquad
 \zeta-\mathcal R=\frac{\Delta}{3(1+w)}
 =-\frac{k^2\Phi}{12\pi G a^2(\bar\rho+\bar p)} .
\]
这使常说的 $\mathcal R=\zeta$ 有了精确条件：须选规则势并可略两阶空间梯度。若 $\bar\rho+\bar p$ 极小，不能未经估计便省略右端。由 \eqref{eq:gw-d-constraints} 消去 $\theta$，得 $\mathcal R=-\Phi-2(\Psi+\Phi'/\mathcal H)/[3(1+w)]$（平坦背景）；这个形式便于从度规解取初值。

若两组分各自绝热且无相互物质交换，长波的 $\zeta_i$ 各守恒，因 $\delta\rho=\sum_i\delta\rho_i$，总 $\zeta$ 为以 $\bar\rho_i'$ 加权的平均。冷物质与辐射这一个例子可完全代数化。写 $S=3(\zeta_m-\zeta_r)$、$y_m=\bar\rho_m/\bar\rho_r$，便有
\[
 \zeta(y_m)=\frac{4\zeta_r+3y_m\zeta_m}{4+3y_m}
 =\zeta_r+\frac{y_m S}{4+3y_m},\qquad
 \frac{\mathrm d\zeta}{\mathrm d\ln a}=\frac{4y_m S}{(4+3y_m)^2}.
\]
若初始是纯冷物质等曲率模，辐射时代的总曲率近零，物质时代却趋于 $S/3$。这种转换发生在组分比重变化最快的阶段，与个别流体产生内禀熵是两回事。验算可从压力作起：$\delta p=\delta\rho_r/3$、$c_a^2=4/[3(4+3y_m)]$，于是 $\delta p_{\rm nad}=-(4\bar\rho_m/3)S/(4+3y_m)$。代入式 \eqref{eq:gw-d-zeta} 立即重得上式。

因而微波背景的大角度端点温度也会记住初值。末次散射时近于物质主导，有 $\Phi=\Psi\simeq-3\zeta_m/5$，而光子本征温度 $\delta_\gamma/4=\zeta_\gamma+\Phi$。故 $\Theta_*+\Psi_*=\zeta_\gamma-6\zeta_m/5$。绝热初值 $\zeta_m=\zeta_\gamma$ 给 $-\zeta/5$；纯冷物质等曲率初值 $\zeta_\gamma=0$ 则给 $-2S/5$。这是相同冷物质密度谱未必给相同大角度温度谱的一个直接反例。两种初值若相关，$C_\ell$ 的理论式应另留交叉谱，只有初始互不相关时才可把两条自谱直接相加。

\subsection{多组分初值的规则性}

自由传播的相对论组分在 $k\eta\ll1$ 也会影响初始势的比。以光子紧耦合、无质量中微子已脱耦的辐射背景为例，令 $R_\nu=\bar\rho_\nu/\bar\rho_r$。取规则的绝热增长模，$\Phi,\Psi$ 的领头项为常数，$00$ 约束即给 $\delta_\gamma=\delta_\nu=-2\Psi$。中微子 Euler 方程 $\theta_\nu'=k^2(\delta_\nu/4+\Psi-\Sigma_\nu)$ 在最低阶给 $\theta_\nu=k^2\eta\Psi/2$。其四极方程最低阶为 $\Sigma_\nu'=(4/15)\theta_\nu$，八极及更高矩还多一个 $k\eta$，故 $\Sigma_\nu=(k\eta)^2\Psi/15$。用 $k^2(\Phi-\Psi)=12\pi G a^2(\bar\rho_r+\bar p_r)R_\nu\Sigma_\nu=6\mathcal H^2R_\nu\Sigma_\nu$，得到
\[
 \Phi=\left(1+\frac{2R_\nu}{5}\right)\Psi,\qquad
 \zeta=-\Phi-\frac{\Psi}{2}
      =-\left(\frac32+\frac{2R_\nu}{5}\right)\Psi .
\]
剪切本身虽有 $(k\eta)^2$ 的抑制，约束中却同时除以 $k^2$，故两势在规则长波上的差别并不必须同阶消失。若取 $R_\nu=0$，便还原理想辐射流体先前使用的 $\Psi=-2\zeta/3$。同一批中微子对张量模造成的是时间卷积；标量初值的这个常系数不能拿来当作张量的阻尼因子。

初始化一组耦合方程时，可以先选规范不变的 $\zeta_i$、相对速度及各向异性应力的规则支，再用约束求势、用连续性方程求导数。反过来任意给每个组分的密度、速度，再独立指定 $\Phi,\Phi'$，会过定约束并激起发散的非物理初解。在标量、矢量和张量分开后，统计问题还须指定它们的原初相关；统计各向同性禁止一个标量与单个螺旋度为二的振幅有非零二点函数，因为绕 $\bm k$ 旋转可连续改变后者的相位。
\section{宇宙学张量波的传播}\label{sec:gw-dtensor}

\subsection{外部常模与内部行波}

将空间 Einstein 方程作横向无迹投影。用 $h_{ij}=\sum_{A=+,\times}e^A_{ij}h_A$，$e^A_{ij}e^{Bij}=2\delta^{AB}$，设 $\Pi_A=e^{Aij}\delta T^i{}_j/2$（空间指标在背景正交基中缩并）。单个模式满足
\begin{equation}\label{eq:gw-d-wave}
 h_A''+2\mathcal Hh_A'+k^2h_A=16\pi G a^2\Pi_A,\qquad
 u_A''+\left(k^2-\frac{a''}{a}\right)u_A=16\pi G a^3\Pi_A,\quad u_A=ah_A.
\end{equation}
这里的 $2\mathcal H$ 来自对 $a^2 h_A'$ 求导。源项是空间应力的剪切部分，理想流体在一阶只有各向同性压强；速度乘积 $(\rho+p)v_iv_j$ 要到二阶才出现。

考察无源的长波。略去 $k^2$ 后 $(a^2h_A')'=0$，故 $h_A=C_A+D_A\int^\eta\mathrm d\eta'/a^2(\eta')$。在标准辐射、物质膨胀阶段选合适积分常数，后项随时间衰减，留下所谓冻结的常模。波长很短时 $k^2\gg|a''/a|$，$u_A\sim\exp(\pm\mathrm i k\eta)$，故 $h_A$ 的包络按 $a^{-1}$ 下降。此时局部频率为 $k/(2\pi a)$，度规扰动的平方与频率的平方各给 $a^{-2}$，能量恰按 $a^{-4}$ 稀释。Hubble 半径不是信号不能越过的界面；这里的“进入”指膨胀时间尺度开始长于波的振荡时间，故约定 $k=aH$ 用来标记它。

令 $a\propto\eta^p$、$x=k\eta$，齐次方程写为 $h_{xx}+2p\,h_x/x+h=0$。置 $h=x^{1/2-p}b$ 得阶数 $\nu=p-1/2$ 的 Bessel 方程。对辐射、物质时代，规则解 $T_h=h/h_{\rm p}$ 分别为
\begin{equation}\label{eq:gw-d-trans}
 T_h^{(r)}(x)=j_0(x)=\frac{\sin x}{x},\qquad
 T_h^{(m)}(x)=\frac{3j_1(x)}{x}
   =3\frac{\sin x-x\cos x}{x^3}.
\end{equation}
两者都在 $x=0$ 归一为一。短波包络分别按 $\eta^{-1},\eta^{-2}$，与各自的 $a^{-1}$ 相符。第二个解在物质时代只适用于该模在物质时代仍有常模初值；不能把辐射时代已经振荡过的波重新从 $T_h(0)=1$ 开始计起。

例如在 $\eta_{\rm e}$ 作突然的等密度近似，要求 $a,a'$ 连续。辐射段 $a=A\eta$，物质段选 $a=A(\eta+\eta_{\rm e})^2/(4\eta_{\rm e})$。从式 \eqref{eq:gw-d-wave} 积分经过交接点，$h$ 和 $a^2 h'$ 连续。设 $F(x)=3j_1(x)/x$、$Y(x)=3y_1(x)/x$，这里 $y_1$ 是第二类球 Bessel 函数，物质段 $x=k(\eta+\eta_{\rm e})$。令 $q=k\eta_{\rm e}$，初态振幅取一，则匹配给出
\[
 \begin{pmatrix}F(2q)&Y(2q)\\ \partial_xF(2q)&\partial_xY(2q)\end{pmatrix}
 \binom{C}{D}=\binom{j_0(q)}{\partial_qj_0(q)} .
\]
行列式由球 Bessel 的 Wronskian 得 $9/(2q)^4$，故两个初始条件唯一指定 $C,D$。$q\ll1$ 时几乎只留下常模；$q\gg1$ 时两列基解都要保留以连续地传递当时的相位。使用连续背景 \eqref{eq:gw-d-equality} 数值积分时，只需在 $k\eta\ll1$ 设 $T_h=1+O(k^2\eta^2),\ T'_h=O(k^2\eta)$，交接常数便自动确定。


\subsection{用 Green 函数保持模式初值}

若热史中有多个短暂阶段，给每段各乘一个独立的振幅因子可能丢失相位。直接求解 $u=ah$ 的二阶初值问题较稳妥。取任意两支齐次解 $v_1,v_2$，令 $W=v_1v_2'-v_1'v_2$；$u$ 方程无一阶导，故 $W$ 与时间无关。对任意源 $J(\eta)=16\pi G a^3\Pi_A$，
\[
 G_k(\eta,\eta')=\Theta(\eta-\eta')
 \frac{v_1(\eta')v_2(\eta)-v_2(\eta')v_1(\eta)}{W},
 \qquad u(\eta)=u_{\rm hom}(\eta)+\int G_k(\eta,\eta')J(\eta')\,\mathrm d\eta' .
\]
因 $G_k(\eta',\eta')=0$、$\partial_\eta G_k|_{\eta'+0}=1$，积分经过源脉冲时 $u$ 连续、$u'$ 正好跳跃所需大小。物质时代 $a\propto\eta^2$，可取 $v_1=xj_1(x)$、$v_2=xy_1(x)$，$x=k\eta$，则 $W=k$。将 $xj_1=\sin x/x-\cos x$、$xy_1=-\cos x/x-\sin x$ 代入得
\[
 G_k^{(m)}=\frac{\Theta(\eta-\eta')}{k}
 \left[\left(1+\frac1{xx'}\right)\sin(x-x')
       +\left(\frac1x-\frac1{x'}\right)\cos(x-x')\right].
\]
$xx'\gg1$ 时还原 $\sin[k(\eta-\eta')]/k$。取 $k\to0$ 时两项各有表面发散，须一并展开；独立的 $k=0$ 齐次解为 $\eta^2,\eta^{-1}$，检验可得 $G_0^{(m)}=(\eta^2/\eta'-\eta'^2/\eta)/3$。这也是为何直接在较早的物质背景使用平直时空 Green 函数去算接近 Hubble 半径的源会出错。

无源时，任何连续热史都还容许一种从长波端启动的迭代。不用抵消很大的 $1/k$ 项，只需将 $(a^2h')'=-k^2a^2h$ 积两次。以 $\eta_i$ 时 $h=h_{\rm p},\,h'=0$，得
\[
 \frac{h(\eta)}{h_{\rm p}}=1-k^2\int_{\eta_i}^\eta
 \frac{\mathrm d\eta_1}{a^2(\eta_1)}
 \int_{\eta_i}^{\eta_1}a^2(\eta_2)\frac{h(\eta_2)}{h_{\rm p}}\,\mathrm d\eta_2 .
\]
右端先代一，再逐步代回可求 $k^2$ 级数。若 $a\propto\eta^p$ 且可取 $\eta_i=0$，内积给 $\eta_1^{2p+1}/(2p+1)$，故首项 $1-(k\eta)^2/[2(2p+1)]$。在 $p=1,2$ 时分别给 $1-x^2/6,\ 1-x^2/10$，与式 \eqref{eq:gw-d-trans} 一致。用已知背景 $a(\eta)$ 对双积分求值，就能在中间时刻给数值积分器指定更准确且无需猜测衰减模的初值。
\subsection{无碰撞中微子如何带走应力}

光子在热浴中频繁散射时剪切近乎为零。已脱耦的相对论性中微子却能沿固定方向飞过一段波长，其分布会记得过去哪些方向曾被度规拉长。取 $\bm n$ 为粒子传播方向、$q_\nu$ 为其共动动量大小。测地线给出张量造成的 $\mathrm d\ln q_\nu/\mathrm d\eta=-h_{ij}'n^in^j/2$。将无碰撞 Boltzmann 方程线性化，并对 $q_\nu^3\,\mathrm dq_\nu$ 积分。对 $\mathrm d\bar f/\mathrm d\ln q_\nu$ 分部积分产生 $-4$ 倍背景能量密度，故中微子的相对方向亮度 $I_\nu$ 满足
\[
 (\partial_\eta+\mathrm i k\mu)I_\nu=-2n^i n^j h_{ij}',\qquad
 \mu=\hat{\bm k}\cdot\bm n.
\]
若脱耦时没有独立张量亮度，沿特征积分所得为过去的 $h_{ij}'$ 与 $\exp[-\mathrm i k\mu(\eta-\eta')]$ 的卷积。再乘 $n_i n_j$ 对方向积分，即得 TT 应力。取 $\bm k$ 沿 $z$ 轴，有 $e^+_{ij}n^i n^j=(1-\mu^2)\cos2\varphi$，方位角平均使交叉偏振互不耦合。定义
\[
 K(s)=\frac{1}{16}\int_{-1}^{1}(1-\mu^2)^2e^{-\mathrm i s\mu}\,\mathrm d\mu
 =\frac{j_2(s)}{s^2},\qquad K(0)=\frac1{15}.
\]
用 $\Pi_A=-4\bar\rho_\nu\int_{\eta_{\rm dec}}^\eta K[k(\eta-\eta')]h_A'(\eta')\,\mathrm d\eta'$ 代入式 \eqref{eq:gw-d-wave}，得
\begin{equation}\label{eq:gw-d-nudamp}
 h_A''+2\mathcal H h_A'+k^2 h_A
 =-24 f_\nu \mathcal H^2\int_{\eta_{\rm dec}}^\eta
 K[k(\eta-\eta')]h_A'(\eta')\,\mathrm d\eta',
 \qquad f_\nu=\frac{\bar\rho_\nu}{\bar\rho_{\rm tot}} .
\end{equation}
早期常模的时间导数为零，故不受阻尼。效应最大在中微子已脱耦且模式进入 Hubble 半径时仍由辐射主导的波段。在标准三种轻中微子的辐射时代，$f_\nu\simeq0.405$；积分解的晚时振幅约为无中微子时的 $0.80$ 倍，能量则约为 $0.64$ 倍。不能把这个因子乘到任意频率上。若模在中微子脱耦前已深处半径内，或在物质主导后才进入，前者届时引力耦合已弱，后者 $f_\nu$ 已小。


式 \eqref{eq:gw-d-nudamp} 还可给数值求解一组含中微子的规则初值。取脱耦足够早、完整辐射时代中 $x=k\eta$、$f_\nu$ 常数，令 $h/h_{\rm p}=1+c_2x^2+c_4x^4+\cdots$。核的前两项为 $K(s)=1/15-s^2/210+\cdots$，故
\[
 \int_0^x K(x-u)\frac{\mathrm d}{\mathrm du}\frac{h(u)}{h_{\rm p}}\,\mathrm du
 =\frac{c_2x^2}{15}
   +\left(\frac{c_4}{15}-\frac{c_2}{1260}\right)x^4+\cdots .
\]
代入无量纲波方程，比较常数与 $x^2$ 项，得到
\[
 c_2=-\frac1{6+8f_\nu/5},\qquad
 c_4=-\frac{c_2(1-2f_\nu/105)}{20+8f_\nu/5}.
\]
$f_\nu=0$ 恢复 $\sin x/x$ 的 $-1/6$ 和 $1/120$；$f_\nu>0$ 时第一弯曲变小。若某模式在 $x_{\rm dec}\gg1$ 时才遇到自由流，积分下限应移到 $x_{\rm dec}$，同时以此前无碰撞应力可以忽略的波解给 $h,h'$。补作这个匹配之后，强行把 $x=0$ 的规则序列用到高频模上所产生的假阻尼便会消失。

还可从方向流的结构看出能量去向。惯性极限下 $K(s)=(1/16)\int(1-\mu^2)^2e^{-\mathrm is\mu}\mathrm d\mu$。若所有粒子沿波矢传播，$(1-\mu^2)^2$ 使应力源为零；沿不同方向飞行的粒子在过去经历的张量时相不同，对得出的局部四极作相位混合。无碰撞粒子保持沿各自相空间轨线的分布信息，但只要测量局域的张量振幅，流走的高阶角矩就像一种耗散。若出现新相互作用令粒子重新紧耦合，必须将有碰撞亮度方程的核重算，单用其均匀能量份额 $f_\nu$ 已不足以定出阻尼。
\subsection{从振幅传递到今天的能谱}

规定张量无量纲谱包含两偏振及其张量缩并，即 $\langle h_{ij}h_{ij}\rangle=\int\mathcal P_t(k,\eta)\,\mathrm d\ln k$。若初始只有规则常模且其后无源，则 $\mathcal P_t(k,\eta)=\mathcal P_t^{\rm p}(k)|T_h(k,\eta)|^2$。式 \eqref{eq:gw-effective-stress} 在共动正交标架中给 $\rho_{\rm gw}=\langle h_{ij}'h_{ij}'\rangle/(32\pi G a^2)$。对短波作若干振荡平均便有
\begin{equation}\label{eq:gw-d-omegat}
 \Omega_{\rm gw}(k,\eta)
 =\frac{\mathcal P_t^{\rm p}(k)}{12a^2H^2}\overline{|T_h'|^2}
 \simeq \frac1{12}\left(\frac{k}{aH}\right)^2
       \mathcal P_t^{\rm p}(k)\overline{|T_h|^2}.
\end{equation}
分母指该时刻的临界密度。第二等号对变化缓慢的短波包络成立；不能在一个恰好 $T_h=0$ 的固定相位用它替换实际导数。今天的 $k=2\pi f$，对照 \eqref{eq:gw-b-omega} 可知 $h_c^2(f)=\mathcal P_t^{\rm p}\overline{|T_h|^2}/2$。没有普适意义的“超 Hubble 波能量”可用同一短波应力直接定义。

对于在辐射时代入内的短波，式 \eqref{eq:gw-d-trans} 给 $\overline{|T_h|^2}\simeq1/(2k^2\eta^2)$，又因 $\mathcal H=1/\eta$，当时 $\Omega_{\rm gw}\simeq\mathcal P_t^{\rm p}/24$。把它按式 \eqref{eq:gw-d-gstar} 红移到今天，近尺度不变的原初谱就产生近乎平坦的高频平台。若波在物质时代入内，$\overline{|T_h|^2}$ 的包络按 $k^{-4}$，故同一时刻的 $\Omega_{\rm gw}\propto k^{-2}\mathcal P_t^{\rm p}$。等密度附近的斜率改变首先反映传播史，不能把它全部归给发射过程。

更一般地，在常 $w>-1/3$ 的时代，进入条件 $k=a_kH_k$ 给 $a_k\propto k^{-2/(1+3w)}$。之后包络正比 $a_k/a_0$，故若 $\mathcal P_t^{\rm p}\propto k^{n_t}$，有
\[
 \frac{\mathrm d\ln\Omega_{{\rm gw},0}}{\mathrm d\ln k}
 \simeq n_t+\frac{2(3w-1)}{3w+1}.
\]
物质主导、辐射主导、动能主导的 $w=0,1/3,1$ 分别给传播附加倾角 $-2,0,1$。式中忽略了自由度变化和自由流阻尼，分段使用时须另匹配峰值；极早期的再热持续多久，因而可能在频谱上留下转折。

\subsection{标准汽笛的距离为何会改变}

源的相位随着红移质量演化，绝对振幅却可定标其光度距离，这便是\textbf{标准汽笛}(standard siren)。要用它画 $d_L(z)$，红移仍须有独立信息，如电磁对应体所在的宿主星系。一组宿主候选可用各自的红移概率加权；逐个给没有对应体的源硬配最近的星系会偏置结果。低红移时 $d_L\simeq z/H_0$，本动速度的等效红移与 Hubble 红移之比近似 $v_{\rm pec}/(H_0d_L)$。距离较大时，此误差下降，弱透镜放大对单个事件的影响却上升。给定倾角的振幅方能换成距离；双偏振测量或电磁喷流方向有助于解除这一简并。

改变背景方程会经式 \eqref{eq:gw-d-dist} 改变光与引力波共同经历的几何距离。另有一类修改只改变张量的振幅输运。取光速张量二次作用量的动能系数为 $a^2M_*^2(\eta)$，自由方程立即成为 $h_A''+\left[2\mathcal H+(\ln M_*^2)'\right]h_A'+k^2h_A=0$。正则振子为 $aM_*h_A$。故短波平面振幅除 $a^{-1}$ 外还乘 $M_*^{-1}$；球面波再除以共动半径。把源处的局域发射幅度固定，观测端相对广义相对论的变化即 $M_*(z)/M_*(0)$。用同一个源模型拟合距离得到
\begin{equation}\label{eq:gw-d-gwdl}
 \frac{d_L^{\rm gw}(z)}{d_L^{\rm em}(z)}
 =\frac{M_*(0)}{M_*(z)}
 =\exp\left[\frac12\int_0^z\frac{\alpha_M(z')}{1+z'}\,\mathrm dz'\right],
 \qquad \alpha_M=\frac{\mathrm d\ln M_*^2}{\mathrm d\ln a}.
\end{equation}
若源区耦合也改变，必须把波形产生与传播分别计算，才能作这项距离比较。该式检验的量与群速度不同，故光和波同时到达仍容许振幅缓慢偏离。

例如含红外有效项 $-m^2R\Box^{-2}R/6$ 的非局域模型可引入辅助量 $U=-\Box^{-1}R,\ S=-\Box^{-1}U$。逆算符的边界条件须由原有效方程给定；不能把任意齐次解都当作新的可量子化粒子。其张量动能有效系数正比于 $1-m^2\bar S/3$，故 \eqref{eq:gw-d-gwdl} 的比值为 $\sqrt{[1-m^2\bar S(0)/3]/[1-m^2\bar S(z)/3]}$。由距离去拟合此类模型，应当在各自的背景解中同时拟合 $H_0,\Omega_{m0}$ 等参数；固定它们作图只给出某一组参数下的几何差异。


\subsection{由一组距离估计传播阻尼}

可将式 \eqref{eq:gw-d-gwdl} 的自由函数与背景几何区别得更明确。取平坦、有尘埃和常 $w$ 暗能量的背景并保持 $H_0,\Omega_{m0}$ 不变，$E^2(z)=\Omega_{m0}(1+z)^3+(1-\Omega_{m0})(1+z)^{3(1+w)}$。对积分求导得
\[
 \frac{\partial\ln d_L^{\rm em}}{\partial w}
 =-\frac{3(1-\Omega_{m0})}{2}
 \frac{\displaystyle\int_0^z
 (1+z')^{3(1+w)}\ln(1+z')\,\mathrm dz'/E^3(z')}
 {\displaystyle\int_0^z \mathrm dz'/E(z')}.
\]
增加 $w$ 使过去的暗能量密度更高，在相同红移处距离变短。若 $\alpha_M$ 近似为常数，另有 $\partial\ln d_L^{\rm gw}/\partial\alpha_M=\ln(1+z)/2$。二者的红移权重不同；一条高红移汽笛不能同时唯一确定 $w,\alpha_M$，而跨红移的几何距离与传播距离组合有机会分开。把每个模型的 $H_0$ 和物质份额也放开时，还会多出背景简并。

在本地作交叉检查。设 $q_0=-\ddot a_0/(a_0H_0^2)$、$\alpha_{M0}$ 为小红移处的阻尼改动。两项距离合并给
\[
 d_L^{\rm gw}=\frac{z}{H_0}
 \left[1+\frac{1-q_0+\alpha_{M0}}{2}z+O(z^2)\right].
\]
首阶的 Hubble 斜率不依赖该传播项；下一阶只量到 $-q_0+\alpha_{M0}$ 的组合。若独立的电磁距离测出 $q_0$，才可单列局域阻尼斜率。这个例子也说明为何某一组固定参数的理论距离相差百分之几，不等于资料就以同样精度排除了其中一个引力模型。

以有宿主星系的少量汽笛为例，若观测到距离 $\widehat d$、宿主红移 $\widehat z$，在误差小且以对数距离拟合时，单个残差可取 $r_d=\ln\widehat d-\ln d_L^{\rm gw}(\widehat z)$。 它的方差在独立近似下为
\[
 \sigma_{r_d}^2\simeq\sigma_{\ln d,{\rm wave}}^2
 +\left(\frac{\mathrm d\ln d_L^{\rm gw}}{\mathrm dz}\right)^2
    \sigma_z^2+\sigma_{\ln d,{\rm lens}}^2 .
\]
$\sigma_{\ln d,{\rm wave}}$ 含距离与倾角等波形参数的边缘化；小 $z$ 时第二项约 $\sigma_z^2/z^2$。弱透镜放大率 $\mu_{\rm L}$ 给应变多乘 $\sqrt{\mu_{\rm L}}$，所以单事件推断的对数距离偏移是 $-\ln\mu_{\rm L}/2$。只收超过阈值的事件会倾向于选大放大率或有利倾角；用事件母体作精密推断，还须在似然中以检测概率归一化。若只有候选宿主，便将上式的单一红移换成各宿主的红移、位置权重的概率混合，而不是把混合的平均红移当成确知的宿主。这样波形的统计量和宇宙学模型对应的距离曲线才是在比较同一种观测。
\section{微波背景上的引力波}\label{sec:gw-dcmb}

有一类张量波的波长同可见宇宙相比也不算短。局域仪器很难跟踪它的振荡，但从最后散射面飞来的光子已经走过极长的基线。同一片天空上，温度给出沿射线积累的频移，线偏振还记录了散射当时的方向分布。以下只保留一次扰动，另行指出二次透镜效应何时重要。

\subsection{天空的多极矩与样本方差}

去掉平均温度 $T_0$，在观者指向天球的单位矢量 $\bm n$ 上展开 $\Theta(\bm n)=\delta T/T_0=\sum_{\ell m}a_{\ell m}Y_{\ell m}(\bm n)$。实温度保证 $a_{\ell,-m}=(-1)^ma_{\ell m}^*$。统计各向同性要求 $\langle a_{\ell m}a^*_{\ell'm'}\rangle=C_\ell\delta_{\ell\ell'}\delta_{mm'}$。用加法定理可把这组数还原成角相关
\[
 \langle\Theta(\bm n)\Theta(\bm n')\rangle
 =\sum_{\ell=0}^{\infty}\frac{2\ell+1}{4\pi}C_\ell P_\ell(\bm n\cdot\bm n').
\]
单极已包含在平均温度中，偶极主要来自观者运动，通常从 $\ell=2$ 开始比较原初涨落。$C_\ell$ 是无量纲量，若图上的单位为温度平方，还要乘以 $T_0^2$。常画 $\ell(\ell+1)C_\ell/(2\pi)$，因为在大 $\ell$ 时，这近似给每单位 $\ln\ell$ 对温度方差的贡献。

为什么整片天空的理想测量也还会有误差？固定 $\ell$ 只有 $2\ell+1$ 个独立的实高斯数可用。估计量 $\widehat C_\ell=\sum_m|a_{\ell m}|^2/(2\ell+1)$ 的期望是 $C_\ell$，由四点函数的两种连通于这对估计量的 Wick 配对，得 $\operatorname{Var}(\widehat C_\ell)=2C_\ell^2/(2\ell+1)$。低多极仅有少数样本，故天空在该角尺度上恰好出现大涨落并不能唯一归因于张量。遮去部分天空还会混合相邻多极，偏振的 $E,B$ 分离也会受边界影响。

\subsection{光子穿过不断变化的势}

先忽略碰撞，两次散射之间黑体分布沿测地线保持其占有数。给光子动量一个方向依赖的小红移，等价于给黑体温度同样的相对扰动。令 $\bm s$ 为光子的传播方向，故在未扰动光路上 $\bm s=-\bm n$。正交标架中能量为 $p$，共动能量 $q=ap$ 的一阶测地线方程为
\[
 (\partial_\eta+\bm s\cdot\nabla)\ln q
 =\Phi'-\bm s\cdot\nabla\Psi-\tfrac12 s^is^j h'_{ij}.
\]
这里左端表示沿射线求变化率。用 $\bm s\cdot\nabla\Psi=\mathrm d\Psi/\mathrm d\eta-\Psi'$ 把引力频移分解为端点与沿途两项。若在 $\eta_*$ 瞬时最后散射，则在观测点 $\bm x_0=0$ 收到的温度为
\begin{align}
 \Theta(\bm n)
 &=\left[\frac{\delta_\gamma}{4}+\Psi-\bm n\cdot\bm v_b\right]_{*}
 -\Psi_0+\bm n\cdot\bm v_0 \nonumber\\
 &\quad+\int_{\eta_*}^{\eta_0}\left[\Phi'+\Psi'
      -\tfrac12 n^in^j h'_{ij}\right]
       \big(\eta,(\eta_0-\eta)\bm n\big)\,\mathrm d\eta .\label{eq:gw-d-los}
\end{align}
发射处用 $\rho_\gamma\propto T^4$ 得本征温度 $\delta_\gamma/4$，$\bm v_b$ 是最后一次散射的重子速度，朝向观者的运动产生蓝移。观测端标量势只移单极，速度只移偶极。余下端点势加本征温度为普通 Sachs--Wolfe 效应，沿途的 $\Phi'+\Psi'$ 称积分 Sachs--Wolfe 效应。若宇宙长期由纯尘埃主导，常势解令后者消失。晚期势开始衰减时光子穿入势阱所得与穿出所失不再相等，遂使温度与前景密度产生相关。

在最后散射时尚未进入 Hubble 半径的绝热物质时代模式上，$\delta_m=-2\Psi$，$\delta_\gamma/4=\delta_m/3$，故端点温度恰为 $\Psi/3=-\zeta/5$。若只留这项，把平面波 $\exp[\mathrm i k\chi_* \hat{\bm k}\cdot\bm n]$ 展开为球波，$\chi_*=\eta_0-\eta_*$，可得
\begin{align}
 a_{\ell m}^{(s)}
 &=-\frac{4\pi\mathrm i^\ell}{5}
 \int\frac{\mathrm d^3 k}{(2\pi)^3}
 \zeta(\bm k)j_\ell(k\chi_*)Y_{\ell m}^*(\hat{\bm k}),\nonumber\\
 C_\ell^{(s)}&=\frac{4\pi}{25}\int_0^\infty
 \mathcal P_\zeta(k)j_\ell^2(k\chi_*)\,\mathrm d\ln k.
 \label{eq:gw-d-clsw}
\end{align}
若 $\mathcal P_\zeta=A_s$ 恒定，换元 $x=k\chi_*$ 后，积分不含源面距离。利用 $\int_0^\infty j_\ell^2(x)\mathrm dx/x=1/[2\ell(\ell+1)]$（$\ell\ge1$），得 $\ell(\ell+1)C_\ell^{(s)}/(2\pi)=A_s/25$。这是一条大角尺度平台。对较短波，最后散射前的光子压力已造成声振荡，本征温度和速度须由 \eqref{eq:gw-d-fluid} 连同散射项求解；不能把平台延伸到声学峰。

\subsection{张量如何投影到温度图}

长波张量在沿途若几乎恒定，则式 \eqref{eq:gw-d-los} 的张量项也几乎为零；随模式入内演化而变的潮汐形变才积累温度差。为求张量对 $C_\ell$ 的贡献，先取某一模式的 $\hat{\bm k}$ 沿 $z$。在源方向 $(\vartheta,\varphi)$ 上，加、叉偏振的角因子依次为 $(1-\mu^2)\cos2\varphi$ 与 $(1-\mu^2)\sin2\varphi$，$\mu=\cos\vartheta$。因此这一个平面波只激发 $m=\pm2$。转动 $\hat{\bm k}$ 会把同一 $\ell$ 的各 $m$ 混合，却不改 $\sum_m|a_{\ell m}|^2$。

令 $x=k(\eta_0-\eta)$。投影到 $Y_{\ell 2}$ 后需算 $\int(1-\mu^2)P_\ell^2(\mu)e^{\mathrm i x\mu}\mathrm d\mu$，其中 $P_\ell^2$ 是缔合 Legendre 函数。用 $P_\ell^2=(1-\mu^2)P_\ell''$ 及 Legendre 方程，或两次分部积分后用球 Bessel 方程，得
\[
 \int_{-1}^1(1-\mu^2)P_\ell^2(\mu)e^{\mathrm i x\mu}\,\mathrm d\mu
 =-2\mathrm i^\ell\frac{(\ell+2)!}{(\ell-2)!}\frac{j_\ell(x)}{x^2}.
\]
例如 $\ell=2$ 时左边在 $x=0$ 是 $16/5$，右边由 $j_2(x)\sim x^2/15$ 得到同一值。方位角积分给因子 $\pi$，$Y_{\ell2}$ 的归一化再给 $\sqrt{(2\ell+1)(\ell-2)!/[4\pi(\ell+2)!]}$。将 $m=\pm2$ 一并求和，采用式 \eqref{eq:gw-d-omegat} 的两偏振谱约定，得到瞬时解耦近似下的张量温度谱
\begin{equation}\label{eq:gw-d-clt}
 C_\ell^{(t)}=\frac{\pi}{4}\frac{(\ell+2)!}{(\ell-2)!}
 \int_0^\infty \mathcal P_t^{\rm p}(k)
 \left|\int_{\eta_*}^{\eta_0}T_h'(k,\eta)
 \frac{j_\ell[k(\eta_0-\eta)]}{[k(\eta_0-\eta)]^2}\,\mathrm d\eta\right|^2\mathrm d\ln k .
\end{equation}
这个结果把时间演化与天球投影分开。$j_\ell$ 在 $x\ll\ell$ 时很小，所以某个 $\ell$ 对特定距离的波数选择约为 $k\chi\gtrsim\ell$。大 $k$ 的波在最后散射前已阻尼且其后经过许多振荡，正负的 $T_h'$ 又会沿视线抵消。张量温度谱在小角尺度上因此迅速下降。在温度图上，标量与张量若起初不相关，测到的是两个 $C_\ell$ 之和；一个天区的冷斑本身不会告诉我们它来自哪一项。

\subsection{末次散射有多厚}

设自由电子密度为 $n_e$，Thomson 截面为 $\sigma_T$，从 $\eta$ 到今天的光学深度是 $\tau(\eta)=\int_\eta^{\eta_0}a n_e\sigma_T\,\mathrm d\eta'$。途中无散射的概率为 $e^{-\tau}$，故最后散射在 $\eta$ 到 $\eta+\mathrm d\eta$ 的概率密度为 $g(\eta)=-\tau'e^{-\tau}$。若早期完全不透明，就有 $\int g\,\mathrm d\eta=1$。对端点源把 $\eta_*$ 换为对 $g$ 加权的积分，对沿途频移则乘 $e^{-\tau}$，因为在那之后再次散射会擦去方向记忆。

用一个宽度 $\sigma_\eta$ 的归一高斯近似 $g$ 的主峰，若源在峰内缓变，积分 $e^{\mathrm i k\mu(\eta_0-\eta)}g(\eta)$ 给出额外 $e^{-(k\mu\sigma_\eta)^2/2}$。故快速的视向波被抹平。这个因子源自沿最后散射层平均；另一个阻尼来自解耦前的光子随机行走，两个物理长度不必相等。在紧耦合的光子--重子流体中，定义重子负载 $R_b=3\rho_b/(4\rho_\gamma)$。共同惯性增大而压强仍来自光子，所以声速 $c_s^2=1/[3(1+R_b)]$，声学相位为 $k r_s$，$r_s=\int^{\eta_*}c_s\,\mathrm d\eta$。由 $k\chi_*\sim\ell$ 可知相邻声学峰粗略间隔 $\pi\chi_*/r_s$。这类标量振荡是隔离张量信号时需同时拟合的温度结构。

更一般地，将光子分布写为 $\bar f(q)-(\partial\bar f/\partial\ln q)\Theta(\eta,\bm k,\mu)$。在标量无碰撞区，其亮度方程为 $\Theta'+\mathrm i k\mu\Theta=\Phi'-\mathrm i k\mu\Psi$，这里 $\mu$ 改指传播方向与 $\bm k$ 的夹角。展开 $\Theta=\sum_\ell(-\mathrm i)^\ell(2\ell+1)\Theta_\ell P_\ell(\mu)$，利用 $(2\ell+1)\mu P_\ell=(\ell+1)P_{\ell+1}+\ell P_{\ell-1}$ 得
\[
 \Theta_\ell'=\frac{k}{2\ell+1}[\ell\Theta_{\ell-1}-(\ell+1)\Theta_{\ell+1}]
 \quad(\ell\ge2,\ \text{无碰撞}).
\]
引力势对低阶矩的驱动另由方程右端读出。散射率 $-\tau'$ 很大时碰撞会抑制四极及更高矩，偶极被拉向重子速度；解耦后自由传播把低阶矩逐级送往高阶。数值求解可先演化碰撞与势构成的局域源，再用径向的 $j_\ell$ 沿视线积分，从而无需把所有大 $\ell$ 从早期一直推进到今天。

\subsection{一个四极怎样产生线偏振}

取垂直于出射光的一对正交偏振基 $\bm e_1,\bm e_2$。若场的复分量为 $E_1,E_2$，则 Stokes 参量 $I=\langle|E_1|^2+|E_2|^2\rangle$、$Q=\langle|E_1|^2-|E_2|^2\rangle$、$U=2\operatorname{Re}\langle E_1 E_2^*\rangle$；$V$ 给互相关的虚部，描述圆偏振。若将基转过 $\gamma$，$Q\pm\mathrm i U$ 乘以 $e^{\mp2\mathrm i\gamma}$，同一个物理偏振方向在转过半周后已经回到原状。

Thomson 散射振幅正比于入射与出射偏振矢量的点积。设沿 $z$ 轴观测散射光，入射方向为 $(\vartheta,\varphi)$。对入射非偏振光先把两个偏振的投影平方求平均，得到沿 $x,y$ 的强度之差正比于 $-(n_x^2-n_y^2)$。包括交叉项后，每个电子散射到单位立体角的偏振强度为
\[
 (Q+\mathrm i U)_{\rm sc}
 =-\frac{3\sigma_T}{16\pi}
 \int \mathrm d\Omega'\,I_{\rm in}(\vartheta,\varphi)
       (1-\cos^2\vartheta)e^{2\mathrm i\varphi}.
\]
例如一束沿 $x$ 入射的光散射到 $z$ 后只留下 $y$ 偏振，上式给 $Q<0$。核正比于 $Y_{22}$，各向同性光或仅有偶极的入射光均不给净偏振。这也说明把最后散射面取到严格无厚度且之前严格紧耦合的极限，为何会错失偏振。散射之前，光子须先飞过有限的平均自由程才积累得起本地四极。再电离时自由电子重新出现，来自很远处的辐射已有大尺度四极，故另在低 $\ell$ 留下偏振。

$Q,U$ 随所选天球基而变，不能直接拿不同方向的 $Q$ 相乘。定义旋量球谐展开 $(Q\pm\mathrm iU)(\bm n)=-\sum_{\ell m}(a^E_{\ell m}\pm\mathrm i a^B_{\ell m})\,{}_{\pm2}Y_{\ell m}(\bm n)$。求和从 $\ell=2$ 开始。等价地，在天球度规 $\gamma_{ab}$ 上把偏振的对称无迹张量分解为两类。偶宇称部分由 $(D_aD_b-\gamma_{ab}D^2/2)E$ 产生，奇宇称部分由 $\epsilon_{c(a}D_{b)}D^c B$ 产生；$D$ 是天球协变导数。前者在反演下乘 $(-1)^\ell$，后者还多一个负号。小块平面天空中若横波矢的方向角为 $\phi_\ell$，可直接看到
\[
 E=Q\cos2\phi_\ell+U\sin2\phi_\ell,\qquad
 B=-Q\sin2\phi_\ell+U\cos2\phi_\ell
\]
（允许同时翻转 $E,B$ 的整体符号约定）。给一条轴对称的标量平面波，取径向基沿波矢与出射方向所张的平面。关于该平面作反射不改变波，却改变 $45^\circ$ 交叉分量的符号，故 $U=0$，线性标量只能产生 $E$。张量的两偏振打破此轴对称，一般 $E,B$ 皆有。

全天偏振数据遂给出 $TT,TE,EE,BB$ 等角谱。在统计宇称对称时，$TB$、$EB$ 变号而期望为零。$BB$ 中还含前景和引力透镜。后一项可从偏折 $\bm\alpha=\nabla\phi_{\rm L}$ 看出，将原偏振图在 $\bm n+\bm\alpha$ 取值，一阶增加 $\nabla\phi_{\rm L}\cdot\nabla(Q\pm\mathrm iU)$。在平面 Fourier 空间，这是两个不同方向波矢的卷积，偏振相对于总波矢不再纯为 $E$，遂产生 $B$。以重建的偏折与实测 $E$ 估计并扣掉这部分称为去透镜。复合时代与再电离时代的原初张量印记、较小角尺度的透镜及多频前景各有自己的频率或角尺度依赖，只有分开它们才能把 $BB$ 换成原初张量振幅。


\subsection{把碰撞项写在光子的屏面上}

如何同时求温度和偏振的碰撞项而不混淆两个方向的偏振基？用三维笛卡尔指标，将相干矩阵 $J_{ij}(\bm s)$ 定义在垂直于传播方向 $\bm s$ 的屏面上，故 $s_iJ_{ij}=s_jJ_{ij}=0$。取 $P_{ij}(\bm s)=\delta_{ij}-s_is_j$。屏面上 $\operatorname{tr}J=I$，不偏振的分布为 $J=I P/2$。在电子静系中弹性 Thomson 散射的增益和损失可以合写成
\begin{equation}\label{eq:gw-d-thomson}
 \left.\frac{\mathrm dJ_{ij}(\bm s)}{\mathrm d\eta}\right|_{\rm coll}
 =a n_e\sigma_T\left[
  -J_{ij}(\bm s)+\frac{3}{8\pi}P_i{}^k(\bm s)P_j{}^l(\bm s)
  \int J_{kl}(\bm s')\,\mathrm d\Omega'\right].
\end{equation}
两次投影来自散射振幅及其复共轭。对各向同性、不偏振入射，$\int P_{kl}(\bm s')\,\mathrm d\Omega'=8\pi\delta_{kl}/3$，正好抵消损失，证明归一化正确。把式 \eqref{eq:gw-d-thomson} 的迹对 $\bm s$ 再积分，两项同样抵消，反映散射不增减光子数。矩阵若最初实对称，将始终为实对称，故 Thomson 散射不从零产生 $V$。电子有小体速度时要把入射分布先 Doppler 变到电子静系再算碰撞；一阶上这给亮度增加偶极源，不会凭空制造四极。

令 $K_{ij}=\int J_{ij}(\bm s')\mathrm d\Omega'$。它的迹给各向同性增益，无迹部分给温度四极与已有线偏振的共同反馈。若输入纯为未偏振亮度 $I=I_0[1+b\,P_2(\hat{\bm k}\cdot\bm s')]$，在 $\hat{\bm k}=\hat{\bm z}$ 系中，$K_{zz}-K_{xx}=-(I_0 b/2)\int P_2(\mu')(\mu'^2-n_x'^2)\mathrm d\Omega'=-2\pi I_0 b/5$，这里用了 $\langle n_x'^2\rangle_\varphi=(1-\mu'^2)/2$。取出射方向的极基 $\bm e_\vartheta,\bm e_\varphi$，有 $(e_\vartheta^ie_\vartheta^j-e_\varphi^ie_\varphi^j)K_{ij}=(K_{zz}-K_{xx})(1-\mu^2)$，$\mu=\hat{\bm z}\cdot\bm s$。所以单次散射每单位光学深度的 $Q=-3I_0 b(1-\mu^2)/20,\ U=0$。沿对称轴看不到偏振，垂直轴最强；若把先前产生的 $Q$ 放回 $K_{ij}$，便得到偏振自身的再散射，完整 Boltzmann 测算需要这项反馈。

也能直接由无碰撞传播算出一次散射前如何获得四极。若上一次散射时只有标量单极 $M(\bm k)$，过一段共形时 $\Delta\eta$ 后，入射分布近于 $M e^{-\mathrm i k\mu\Delta\eta}$。展开至二阶，$\mu^2=(1+2P_2)/3$，故四极为 $-M(k\Delta\eta)^2 P_2/3$。若最初有速度偶极，则 $\mu$ 乘传播展式的一次项已可给四极，领头量级为 $k\Delta\eta$ 乘偶极。有限的平均自由程在充分紧耦合时很小，因此偏振通常远弱于温度；到最后散射层内既要有足够自由传播积出四极，也要还剩一些电子把它转换成偏振。

\subsection{从一次散射到极化的径向核}

对标量平面波，末次散射产生的 $Q$ 有 $(1-\mu^2)$ 角形，$U=0$。把所有局域四极、再散射和可见度收进实转移源 $S_P(k,\eta)$，传播到观者后形如
\[
 (Q\pm\mathrm i U)(\bm n,\bm k)
 =\int S_P(k,\eta)(1-\mu^2)e^{\mathrm i x\mu}\,\mathrm d\eta,
 \qquad x=k(\eta_0-\eta).
\]
在 $\hat{\bm k}$ 轴系中，只需 $m=0$ 的旋量球谐，其极角因子正比于 $\sqrt{(\ell-2)!/(\ell+2)!}\,P_\ell^2(\mu)$。前面算温度张量时的角积分可原封不动使用。略去取决于 $S_P$ 定义的共同常数，标量 $E$ 的径向函数因而为 $\sqrt{(\ell+2)!/(\ell-2)!}\,j_\ell(x)/x^2$，标量 $B$ 恒零。若 $g(\eta)$ 在再电离时代还另有一峰，此处是两段 $\eta$ 积分的\emph{振幅}先相加再平方，不能一律把两个独立的角功率谱相加。因为两个时代的四极可来自同一个原初 Fourier 模式。

对于单一张量波，局域平均 $K_{ij}$ 的张量部分在 $\hat{\bm k}=\hat{\bm z}$ 系中沿 $e^+_{ij}$ 或 $e^\times_{ij}$。例如入射产生 $K_{xx}=-K_{yy}$ 的加偏振，屏面给出的 $Q$ 正比 $(1+\mu^2)\cos2\varphi$，$U$ 正比 $-2\mu\sin2\varphi$。这是由 $\bm e_\vartheta=(\mu\cos\varphi,\mu\sin\varphi,-\sin\vartheta)$ 及 $\bm e_\varphi=(-\sin\varphi,\cos\varphi,0)$ 两次缩并所得；叉偏振把前式的 $\cos2\varphi$ 换成 $\sin2\varphi$，后式变 $2\mu\cos2\varphi$。因而 $Q\pm\mathrm i U$ 对每个 $m=\pm2$ 含 $(1\mp\mu)^2$ 或 $(1\pm\mu)^2$。这个角依赖能产生奇、偶两种旋量调和分量，和只伴 $1-\mu^2$ 的标量有实质区别。

以一个共同源归一化张量 $E,B$ 的径向投影，常用两函数
\begin{align*}
 \epsilon_\ell^{(t)}(x)&=\frac14
 \left[-j_\ell(x)+j_\ell''(x)+\frac{2j_\ell(x)}{x^2}
                         +\frac{4j_\ell'(x)}{x}\right],\\
 \beta_\ell^{(t)}(x)&=\frac12
 \left[j_\ell'(x)+\frac{2j_\ell(x)}{x}\right].
\end{align*}
它们可由 $(1\pm\mu)^2e^{\mathrm ix\mu}$ 的偶、奇组合投影得到：每一因子 $\mu$ 换作 $(1/\mathrm i)\partial_x$，然后在旋量谐函数中对 $\mu$ 分部积分。整体倍数取决于把局部 Thomson 源如何定标，两个函数的相对形状却固定。用球 Bessel 方程还可将 $\epsilon$ 化成 $[-2j_\ell+2j_\ell'/x+(\ell(\ell+1)+2)j_\ell/x^2]/4$。近 $x=0$，$j_\ell\sim x^\ell/(2\ell+1)!!$，故 $\epsilon\sim(\ell+1)(\ell+2)x^{\ell-2}/[4(2\ell+1)!!]$ 而 $\beta\sim(\ell+2)x^{\ell-1}/[2(2\ell+1)!!]$。一个正在本地散射、无空间相位变化的均匀张量四极可给 $\ell=2$ 的 $E$，$B$ 则还需跨屏面的相位变化。数值算小 $x$ 时用这些幂级数比直接相减两个大数稳定。

给定一组标量、张量的原初协方差 $\mathcal P_{IJ}$，这些径向核最终只进同一种投影式。取按球谐约定归一的转移函数 $\Delta^X_{\ell I}(k)$，有
\[
 C_\ell^{XY}=4\pi\int \mathrm d\ln k\sum_{IJ}
 \mathcal P_{IJ}(k)\,\Delta^X_{\ell I}(k)\Delta^{Y*}_{\ell J}(k),
 \qquad X,Y=T,E,B.
\]
例如 $TE$ 的正负由温度源与偏振源的相对相位决定；定义 $E$ 时若翻整体符号，$EE$ 不变而 $TE$ 也须随约定翻号。$BB$ 的振幅虽只涉及平方，估计它的协方差却会混入其它场及天空截断。可见度峰位置敏感于电离史，$T_h'$ 敏感于张量传播史，这两类输入应分别改变以跟踪近似的误差。

\subsection{温度的各声学源如何相加}

在紧耦合极限还可显式求出声峰初相。取纵向规范，$\Theta_0=\delta_\gamma/4$，设 $\theta_\gamma\simeq\theta_b=\theta$。略去重子气压与光子剪切，将两个组分的动量方程按惯性密度 $\rho_b,\ 4\rho_\gamma/3$ 加权相加，内部 Thomson 力正好抵消。于是
\[
 \Theta_0'=-\theta/3+\Phi',\qquad
 (1+R_b)\theta'+R_b\mathcal H\theta
 =k^2[\Theta_0+(1+R_b)\Psi].
\]
由于 $R_b'=\mathcal H R_b$，消去 $\theta$ 有 $[(1+R_b)\Theta_0']'+k^2\Theta_0/3=[(1+R_b)\Phi']'-(1+R_b)k^2\Psi/3$。若几个声周期内可视 $R_b$ 和势为常数，振子平衡点是 $\Theta_0=-(1+R_b)\Psi$。从零速度的初值演化有
\[
 (\Theta_0+\Psi)(\eta)
 =[\Theta_0(0)+(1+R_b)\Psi]\cos(kc_s\eta)-R_b\Psi .
\]
可见重子负载使压缩、稀疏两个极值相对零点不再对称。速度偶极由 $\theta=-3\Theta_0'$ 给出，在密度峰处为零，约相差四分之一周期。声峰的对比度因此能测重子惯性，温度与偏振交叉相关还保有初相。缓变的 $R_b$ 情形可用 $k\eta c_s$ 改为 $k\int c_s\mathrm d\eta$ 并在振幅中保留阻尼，辐射时代快速变化的势则须作为外驱力积分。

把这个解放回视线积分，还要区分单极的 $j_\ell(x)$ 和速度偶极的 $j_\ell'(x)$。因为 $-\bm n\cdot\bm v_b=\mathrm i\mu\,\theta_b/k$，对 $e^{\mathrm ix\mu}$ 作用一阶 $\partial_x$ 即得速度投影。因此单模标量温度的转移振幅可写（此处仅列单极、速度和势）
\[
 \int \mathrm d\eta\,\left\{
 g(\eta)\left[(\Theta_0+\Psi)j_\ell(x)+\frac{\theta_b}{k}j_\ell'(x)\right]
 +e^{-\tau}(\Phi'+\Psi')j_\ell(x)\right\}.
\]
温度四极与偏振共同反馈入碰撞项时还须加 $P_2(\mu)$ 的投影。求 $C_\ell$ 前先将这些振幅相加；例如两个源间的反相正可消掉一个声学极值的一部分。若只从峰对应 $k\chi_*\simeq\ell$ 作一一投影，会漏掉 $j_\ell'$ 与有限宽度产生的展宽。

\subsection{角谱的透镜混合及误差}

在一小块天空上，用二维 Fourier 矢量 $\bm\ell$。一次偏折给 $\delta P=\nabla\phi_{\rm L}\cdot\nabla P$，其中 $P=Q+\mathrm iU$。将梯度分别变为对应波矢，若未偏折的原初 $B=0$，有
\[
 \delta B(\bm\ell)=-\int\frac{\mathrm d^2L}{(2\pi)^2}
 [\bm L\cdot(\bm\ell-\bm L)]
 \sin[2(\varphi_{\ell-L}-\varphi_\ell)]\,
 \phi_{\rm L}(\bm L)E(\bm\ell-\bm L).
\]
两波矢平行或夹角直角时旋角因子为零。若 $E$ 与偏折势可近似为独立高斯场，平方平均后透镜 $C_\ell^{BB}$ 成为 $C_L^{\phi\phi}C_{|\ell-L|}^{EE}$ 对以上角核平方的卷积。这比用常数噪声来表示透镜更精确，因为测到偏折势和 $E$ 的相位后，可预报本片天空某些 $B$ 模式的具体取值。若重建噪声很大，实际能消去的只是其中相关的部分。

对理想全天且统计高斯的一组 $T,E,B$ 图，取均匀噪声并把角波束吸进噪声谱 $N_\ell^{XY}$，记 $\mathcal C_\ell^{XY}=C_\ell^{XY}+N_\ell^{XY}$。由 $2\ell+1$ 个样本的四点函数直接给
\[
 \operatorname{Cov}(\widehat C_\ell^{XY},\widehat C_{\ell'}^{WZ})
 =\frac{\delta_{\ell\ell'}}{2\ell+1}
 \left(\mathcal C_\ell^{XW}\mathcal C_\ell^{YZ}
     +\mathcal C_\ell^{XZ}\mathcal C_\ell^{YW}\right).
\]
例如即使真实的 $TB$ 期望为零，$\widehat C_\ell^{TB}$ 也有温度和 $B$ 自谱乘积造成的散布。若其它参数暂固定，小张量比下可写 $C_\ell^{BB}=r\,C_{\ell,r=1}^{BB,{\rm prim}}+C_\ell^{BB,{\rm lens}}$。由各多极的独立高斯近似，加权估计 $r$ 的逆方差约为
$\sum_{\ell}(2\ell+1)[C_{\ell,r=1}^{BB,{\rm prim}}/\mathcal C_\ell^{BB}]^2/2$。未能去净的透镜、多频前景余量及仪器噪声都进入分母。部分天空常粗略乘有效覆盖率，但若遮罩边界强烈混合多极或透镜非高斯性显著，还要计算真实的非对角协方差。这解释了为何不同角范围给出的张量限制不能只按所含多极数机械缩放。
\section{从真空涨落到原初张量谱}\label{sec:gw-dinfl}

\subsection{哪些波模曾经相互接触}

第 \ref{chpt:BH} 章已区分粒子视界与事件视界。这里需要的是共动 Hubble 半径 $1/(aH)=1/\dot a$。它随时间缩小的充要条件是 $\ddot a>0$。在普通辐射或物质时代它增大，同一个固定的共动波长最终会进入其中。反向外推，今天早已能比较温度的区域，在最后散射时却可能相隔许多当时的粒子视界。若在热大爆炸之前经过一段加速膨胀，这些区域可来自同一个较小的因果域，曾经的短波会横跨越来越多的 Hubble 半径。曲率困难也与同一个量有关，因为 $|\Omega_{\rm tot}-1|=|K|/(aH)^2$。$aH$ 在暴涨中增长得越多，原来有界的曲率份额就被压得越低。

设暴涨末时刻为 $t_{\rm end}$，一条波数 $k$ 在 $a_kH_k=k$ 时离开 Hubble 半径。其后的膨胀倍数记为 $e^{N_k}=a_{\rm end}/a_k$。把今天测得的 $k/(a_0H_0)$ 分步拆开，有
\[
 \frac{k}{a_0H_0}=e^{-N_k}
 \frac{a_{\rm end}}{a_{\rm reh}}\frac{a_{\rm reh}}{a_0}\frac{H_k}{H_0}.
\]
若再热前平均状态方程为 $\bar w$，第一距离比为 $(\rho_{\rm reh}/\rho_{\rm end})^{1/[3(1+\bar w)]}$；再热后第二距离比由熵守恒确定。故常用的“枢轴尺度在末期前五六十个 $e$ 折离开”隐含能标和热史，计算模型的 $n_s,r$ 时应说明 $N_k$ 的取值。平直问题和视界问题给的是足够长暴涨的条件，本身并未指定驱动加速的物质。

\subsection{单个滚动标量}

一个最小耦合、正则标量场的均匀能量与压强分别是 $\dot\phi^2/2+V$ 和 $\dot\phi^2/2-V$，能量守恒给 $\ddot\phi+3H\dot\phi+V_{,\phi}=0$。用式 \eqref{eq:gw-d-fried} 消去能量，得到
\begin{equation}\label{eq:gw-d-slow}
 H^2=\frac{\dot\phi^2/2+V}{3M_{\rm P}^2},\qquad
 \dot H=-\frac{\dot\phi^2}{2M_{\rm P}^2},\qquad
 \epsilon_1=-\frac{\dot H}{H^2}
          =\frac{\dot\phi^2}{2M_{\rm P}^2H^2}.
\end{equation}
因 $\ddot a/a=H^2(1-\epsilon_1)$，单要膨胀加速只需 $\epsilon_1<1$。若还要求存在缓慢变化的吸引解，则 $\dot\phi^2/2\ll V$ 且 $|\ddot\phi|\ll3H|\dot\phi|$，场每单位时间下降的距离近于 $-V_{,\phi}/(3H)$。势的无量纲斜率和曲率定义为 $\epsilon_V=(M_{\rm P}^2/2)(V_{,\phi}/V)^2$、$\eta_V=M_{\rm P}^2V_{,\phi\phi}/V$。用 $3H^2M_{\rm P}^2\simeq V$ 及慢滚速度可化为
\[
 \epsilon_1\simeq\epsilon_V,\qquad
 \epsilon_2\equiv\frac{\dot\epsilon_1}{H\epsilon_1}
 \simeq4\epsilon_V-2\eta_V,\qquad
 N_k\simeq\frac{1}{M_{\rm P}^2}\int_{\phi_{\rm end}}^{\phi_k}
           \frac{V}{V_{,\phi}}\,\mathrm d\phi .
\]
$\epsilon_2$ 的正号按向未来增加的 $\ln a$ 定义；$N_k$ 则计末时刻以前剩余的膨胀量，随时间减少，若混用两者会把谱倾角的符号反过来。若场单调，还可把 $H$ 看作 $\phi$ 的函数，由 $\dot H=H_{,\phi}\dot\phi$ 精确得到 $\dot\phi=-2M_{\rm P}^2 H_{,\phi}$ 和 $V=3M_{\rm P}^2 H^2-2M_{\rm P}^4 H_{,\phi}^2$；这给检查慢滚近似提供了不必预先丢掉动能的办法。

\subsection{在膨胀空间选正频}

先考察一个质量 $m$、曲率耦合 $\xi$ 的自由标量 $\varphi$。将二次作用量用共形坐标展开，动能和梯度都带 $a^2$。置 $v=a\varphi$ 并对交叉项 $-\mathcal H(v^2)'/2$ 分部积分，可写成
\[
 S=\frac12\int\mathrm d\eta\,\mathrm d^3x
 \left[(v')^2-(\nabla v)^2-
 \left(a^2m^2+(6\xi-1)\frac{a''}{a}\right)v^2\right].
\]
用到了平坦背景的 $R=6a''/a^3$。若 $m=0,\xi=1/6$，共形重标度后仍是恒频振子，无由膨胀本身造成的正负频混合。最小耦合 $\xi=0$ 则留下有效势 $a''/a$。对实标量作算符展开
\[
 \widehat v(\eta,\bm x)=\int\frac{\mathrm d^3 k}{(2\pi)^3}
 \left[v_k(\eta)\widehat a_{\bm k} e^{\mathrm i\bm k\cdot\bm x}
 +v_k^*(\eta)\widehat a_{\bm k}^\dagger e^{-\mathrm i\bm k\cdot\bm x}\right],
 \qquad [\widehat a_{\bm k},\widehat a_{\bm q}^\dagger]
 =(2\pi)^3\delta^3(\bm k-\bm q).
\]
正则动量是 $v'$。等时对易关系恰要求 $v_k v_k^{*\prime}-v_k^*v_k'=\mathrm i$，模式方程保证此 Wronskian 守恒。另一组同样归一的解可写为 $\widetilde v_k=\alpha_k v_k+\beta_k v_k^*$，于是 $|\alpha_k|^2-|\beta_k|^2=1$。将同一 $\widehat v$ 用两套模式展开，得到 $\widetilde a_{\bm k}=\alpha_k^*a_{\bm k}-\beta_k^*a^\dagger_{-\bm k}$。若初态为 $a$ 的真空，以新模式计数的平均占有数是 $|\beta_k|^2$。正负动量成对出现源于空间平移不变。时空若没有全局类时 Killing 场，不可能仅从正范数条件中挑出唯一的真空；还需给出准备初态的物理边界条件。

对精确 de Sitter 膨胀 $a=-1/(H\eta)$、$\eta<0$。取最小耦合无质量场，当 $-k\eta\gg1$ 时空间在一次振荡内几乎不变，选正频模式 $e^{-\mathrm i k\eta}/\sqrt{2k}$，这给出 Bunch--Davies 初态。此后 $v_k''+(k^2-2/\eta^2)v_k=0$，代 $v_k=e^{-\mathrm i k\eta}(1+b/\eta)/\sqrt{2k}$，由 $\eta^{-2}$ 项解得 $b=-\mathrm i/k$，$\eta^{-3}$ 项自动消去。故
\begin{equation}\label{eq:gw-d-bd}
 v_k=\frac{e^{-\mathrm i k\eta}}{\sqrt{2k}}
       \left(1-\frac{\mathrm i}{k\eta}\right),\qquad
 |\varphi_k|^2=\frac{H^2}{2k^3}(1+k^2\eta^2).
\end{equation}
在 $-k\eta\ll1$ 时，每对数波数的方差 $\mathcal P_{\delta\varphi}=k^3|\varphi_k|^2/(2\pi^2)$ 渐近 $(H/2\pi)^2$，不依赖 $k$。用“在 $k=aH$ 处取值”估算谱，是用该时刻的缓变 $H$ 去标记\emph{冻结后的}振幅；若真把式 \eqref{eq:gw-d-bd} 在 $|k\eta|=1$ 的瞬时值当作最后常模，就会多出因子二。

\subsection{时空与场共有一个标量自由度}

$\delta\phi$ 在挪动时间切片时变化 $\widetilde{\delta\phi}=\delta\phi-\bar\phi'T$。切到 $\delta\phi=0$ 的同时面，空间扰动为 $\psi_{\rm com}=\psi+\mathcal H\delta\phi/\bar\phi'$。这里记 $\mathcal R=-\psi_{\rm com}$，则在守恒的绝热长波上 $\mathcal R\simeq\zeta$。反过来，在空间平坦切片 $\psi=0$ 上，$\mathcal R=-H\,\delta\phi_{\rm flat}/\dot\phi$。这可从时钟来理解：若相同势面在相邻位置到达的时间相差 $\delta t=-\delta\phi_{\rm flat}/\dot\phi$，它们在空间膨胀中相差 $H\delta t$ 个 $e$ 折。

解去 lapse 与 shift 的线性约束，标量二次作用量只剩一个变量，
\begin{equation}\label{eq:gw-d-ms}
 S_s^{(2)}=\frac12\int z_s^2[(\mathcal R')^2-(\nabla\mathcal R)^2]
          \,\mathrm d\eta\,\mathrm d^3x,\qquad
 z_s=\frac{a\dot\phi}{H},\qquad
 v_k''+\left(k^2-\frac{z_s''}{z_s}\right)v_k=0,\quad v=z_s\mathcal R .
\end{equation}
最后一步直接来自展开 $(v'-z_s'v/z_s)^2$ 并分部积分，这便是 Mukhanov--Sasaki 方程。$z_s^2=2a^2\epsilon_1M_{\rm P}^2$；若 $\dot\phi=0$，用单调标量给时空定时的步骤本身退化，不能把精确 de Sitter 的 $\epsilon_1=0$ 代入慢滚标量谱。令 $k\to0$ 还原 $(z_s^2\mathcal R')'=0$，说明慢滚吸引解除常模外另有一个衰减模。若 $z_s$ 本身迅速减小，这项也可能增长，故常模主导不是任意单标量背景的定理。

为求缓慢演化留下的倾斜，先算 $(\ln z_s)'=\mathcal H(1+\epsilon_2/2)$，又有 $\mathcal H'=\mathcal H^2(1-\epsilon_1)$。忽略慢滚量乘积及其高阶导数，得
\begin{align}
 \frac{z_s''}{z_s}&=\mathcal H^2\left(2-\epsilon_1+\tfrac32\epsilon_2\right)
 \simeq\frac{2+3\epsilon_1+\tfrac32\epsilon_2}{\eta^2},\nonumber\\
 \frac{a''}{a}&=\mathcal H^2(2-\epsilon_1)
 \simeq\frac{2+3\epsilon_1}{\eta^2}.\label{eq:gw-d-indices}
\end{align}
第二个等号用了 $\mathcal H\simeq-(1+\epsilon_1)/\eta$。把有效势写作 $(\nu^2-1/4)/\eta^2$，标量 $\nu_s\simeq3/2+\epsilon_1+\epsilon_2/2$，张量 $\nu_t\simeq3/2+\epsilon_1$。标准归一的正频解为 $\sqrt{-\pi\eta}\exp[\mathrm i(\nu+1/2)\pi/2]H_\nu^{(1)}(-k\eta)/2$。当 $-k\eta$ 很小时，$H_\nu^{(1)}(x)\simeq-\mathrm i\Gamma(\nu)(2/x)^\nu/\pi$，从而固定晚时刻的无量纲谱正比 $k^{3-2\nu}$。这直接给 $n_s-1=-2\epsilon_1-\epsilon_2$，解释了“近尺度不变”偏离零的量级。

以最低阶归一化，$\mathcal R=-H\delta\phi_{\rm flat}/\dot\phi$ 及 \eqref{eq:gw-d-bd} 给 $\mathcal P_{\mathcal R}^{\rm p}=H_k^2/(8\pi^2\epsilon_{1k}M_{\rm P}^2)$。另一种求斜率的办法是沿 $k=a_kH_k$ 求导：$\mathrm d\ln k=(1-\epsilon_1)\mathrm d\ln a$，故 $\mathrm d\ln(H_k^2/\epsilon_{1k})/\mathrm d\ln k=(-2\epsilon_1-\epsilon_2)/(1-\epsilon_1)$。保留一阶即与 Hankel 指数一致；分母不能单独视为完整的二阶修正，振幅的下一阶也会贡献。

\subsection{张量归一化及慢滚一致性关系}

Einstein 作用量在 TT 部分的二阶项为
\[
 S_t^{(2)}=\frac{M_{\rm P}^2}{8}\int a^2
  [(h_{ij}')^2-(\partial_lh_{ij})^2]\,\mathrm d\eta\,\mathrm d^3x .
\]
由于 $e^A_{ij}e^{Bij}=2\delta^{AB}$，每个偏振的正则振子是 $v_A=aM_{\rm P}h_A/\sqrt2$，其有效势正是 \eqref{eq:gw-d-indices} 的 $a''/a$。用式 \eqref{eq:gw-d-bd} 的冻结解，每偏振 $|h_A|^2\simeq H_k^2/(M_{\rm P}^2 k^3)$；$\mathcal P_t$ 还需对两个偏振各乘张量缩并因子二。故
\begin{equation}\label{eq:gw-d-r}
 \mathcal P_t^{\rm p}=\frac{2H_k^2}{\pi^2M_{\rm P}^2},\qquad
 r=\frac{\mathcal P_t^{\rm p}}{\mathcal P_{\mathcal R}^{\rm p}}=16\epsilon_{1k},
 \qquad n_t=\frac{\mathrm d\ln\mathcal P_t^{\rm p}}{\mathrm d\ln k}
 \simeq -2\epsilon_{1k}=-\frac r8 .
\end{equation}
这个一致性关系用到了单场、正则动能、慢滚吸引解及绝热真空。额外轻场可在超 Hubble 半径上继续产生曲率扰动，非正则动能则可改变标量声速，因而即使张量产生仍遵从同一机制，比值 $r$ 也会变。给定枢轴处的 $A_s$ 和 $r$，有 $H_k=\pi M_{\rm P}\sqrt{rA_s/2}$，势能标度 $V_k^{1/4}\simeq[(3\pi^2/2)r A_s]^{1/4}M_{\rm P}$。张量谱的绝对振幅遂有单独测量暴涨能标的意义。

再看场位移。由式 \eqref{eq:gw-d-slow} 有 $|\mathrm d\phi/\mathrm d\ln a|=M_{\rm P}\sqrt{2\epsilon_1}$，沿慢滚解便得 $|\Delta\phi|/M_{\rm P}=\int\sqrt{r(N)/8}\,\mathrm dN$。若 $r$ 在数十个 $e$ 折间始终可观，场位移可能超过 $M_{\rm P}$。这个结论关乎势函数在一长段场空间上的形状，却不要求能量密度接近 $M_{\rm P}^4$；后者由 $A_sr$ 定标。只测到单一枢轴的 $r$，也不能在未知它如何演化时把积分直接写成 $N_k\sqrt{r_k/8}$。


\subsection{由标量约束求曲率作用量}

式 \eqref{eq:gw-d-ms} 把四个标量度规扰动和物质场化成了一只振子。下面给出消元过程。在 \S\ref{sec:Hamilton} 的空间分解中，用固有时写 $\mathrm ds^2=-N^2\mathrm dt^2+g^{(3)}_{ij}(\mathrm dx^i+N^i\mathrm dt)(\mathrm dx^j+N^j\mathrm dt)$。取均匀 $\phi$ 切片 $\delta\phi=0$，消去标量空间剪切，令 $g^{(3)}_{ij}=a^2e^{2\mathcal R}\delta_{ij}$、$N=1+\alpha$、$N_i=\partial_i b$。这时 $\mathcal R$ 是对数空间尺度的扰动。定义 $E_{ij}=(\dot g^{(3)}_{ij}-D_iN_j-D_jN_i)/2$，则舍去边界项后，作用量的标量部分可由
\[
 S=\frac12\int\mathrm dt\,\mathrm d^3x\,\sqrt{g^{(3)}}\!
 \left\{M_{\rm P}^2\left[N R^{(3)}+\frac{E_{ij}E^{ij}-E^2}{N}\right]
       +\frac{\dot{\bar\phi}^{\,2}}{N}-2NV(\bar\phi)\right\}
\]
展开。$N,N_i$ 无时间导数，故先对它们求约束。动量约束 $D_j[(E^j{}_i-\delta^j_i E)/N]=0$ 的梯度部分给 $\partial_i(H\alpha-\dot{\mathcal R})=0$，对 $k\ne0$ 即 $\alpha=\dot{\mathcal R}/H$。能量约束在此代换后化为
\[
 \frac{\nabla^2 b}{a^2}=-\frac{\nabla^2\mathcal R}{a^2H}
                       +\epsilon_1\dot{\mathcal R}.
\]
有些写法用 $N^i$ 的势，其 Poisson 方程会少一因子 $a^2$，因为此处定义的是 $N_i$ 的势。若将 $b=-\mathcal R/H+\chi$，只须解 $\nabla^2\chi=a^2\epsilon_1\dot{\mathcal R}$。空间边界决定可加的齐次解，对衰减扰动或周期盒且 $k\ne0$ 时已固定。

代回作用量时无须先求 $N,N_i$ 的二阶解；二阶的 时延或位移矢量 只乘背景约束，后者为零。空间曲率由三维共形变换写成 $R^{(3)}=a^{-2}e^{-2\mathcal R}[-4\nabla^2\mathcal R-2(\nabla\mathcal R)^2]$。对 $a e^{\mathcal R}\nabla^2\mathcal R$ 分部积分留下 $(\nabla\mathcal R)^2$，不能直接把看似 Laplace 算符的项整体丢掉。将混合的 $\dot{\mathcal R}\nabla^2\mathcal R$ 先对空间再对时间分部积分，并用 $\dot H=-\dot{\bar\phi}^{\,2}/(2M_{\rm P}^2)$，结果是
\[
 S^{(2)}=M_{\rm P}^2\int\mathrm dt\,\mathrm d^3x\,a^3\epsilon_1
 \left[\dot{\mathcal R}^2-\frac{(\nabla\mathcal R)^2}{a^2}\right].
\]
换回共形时即为式 \eqref{eq:gw-d-ms}。时间、空间项带同一正系数，所以正则单场的传播声速是 $1$；它通常不等于背景绝热导数 $c_a^2=\dot p/\dot\rho$。慢滚时 $\dot\rho<0$ 而动能变化小，$c_a^2$ 甚至可近于 $-1$；将这数当作标量模振荡的声速平方会误判出梯度不稳定。把场视作有内禀非绝热压强的流体，或保留这里的曲率作用量，两种处理才相容。

\subsection{一个可以精确检查倾斜的背景}

指数势 $V=V_0\exp[-\lambda\phi/M_{\rm P}]$ 有精确的幂律解。试取 $a\propto t^p$、$\phi=(2M_{\rm P}/\lambda)\ln(t/t_c)$，则动能及势均按 $t^{-2}$，式 \eqref{eq:gw-d-slow} 的 $\dot H$ 方程给 $p=2/\lambda^2$，其余方程固定 $V_0t_c^2=M_{\rm P}^2p(3p-1)$。若 $p>1$，有 $\epsilon_1=1/p<1$、$\epsilon_2=0$，却不必假定 $\epsilon_1$ 极小。取未来的共形时趋零，$a\propto(-\eta)^{-1/(1-\epsilon_1)}$，有 $z_s\propto a$，故标量与张量的正则振子具有同一精确 Hankel 指数
\[
 \nu=\frac{3-\epsilon_1}{2(1-\epsilon_1)},\qquad
 n_s-1=n_t=3-2\nu=-\frac{2\epsilon_1}{1-\epsilon_1}.
\]
两类振子从同一正频真空开始，又因 $z_s/a$ 恒定，谱之比仍为 $r=16\epsilon_1$。故严格得到 $r=-8n_t/(1-n_t/2)$，在首阶还原 $r\simeq-8n_t$。无限延长的指数势模型本身没有结束暴涨的过程；只是考察某段演化中的局域谱，便无需另设 $\phi_{\rm end}$。

精确解还可核对振幅的首阶慢滚项。对常 $\epsilon_1$ 定义 $x=-k\eta$，在离开时 $x_k=1/(1-\epsilon_1)$。由归一的 Hankel 解取 $x\ll1$ 极限，得到每类功率相对按 $H_k$ 写的最低阶结果都乘
\[
 F(\epsilon_1)=
 \left[2^{\nu-3/2}\frac{\Gamma(\nu)}{\Gamma(3/2)}
             (1-\epsilon_1)^{\nu-1/2}\right]^2 .
\]
$\epsilon_1\to0$ 时 $F\to1$。令 $C=\gamma_{\rm E}+\ln2-2$，由 $\psi_{\Gamma}(3/2)=2-\gamma_{\rm E}-2\ln2$ 展开 $\Gamma$ 函数（$\psi_{\Gamma}$ 是它的对数导数），得 $F=1-2(C+1)\epsilon_1+O(\epsilon_1^2)$。一般慢滚时，$z_s$ 比 $a$ 多出随时间改变的 $\sqrt{\epsilon_1}$，相同求解给
\begin{align*}
 \mathcal P_{\mathcal R}^{\rm p}
 &=\frac{H_k^2}{8\pi^2M_{\rm P}^2\epsilon_{1k}}
    [1-2(C+1)\epsilon_{1k}-C\epsilon_{2k}+O(\epsilon_i^2)],\\
 \mathcal P_t^{\rm p}
 &=\frac{2H_k^2}{\pi^2M_{\rm P}^2}
    [1-2(C+1)\epsilon_{1k}+O(\epsilon_i^2)] .
\end{align*}
首阶振幅和首阶倾角不是同一展开精度：对方括号作 $\ln k$ 导数时又多一个慢滚阶次。若只需首阶斜率，先前在 $k=aH$ 给参数加下标已经足够。

谱倾斜自身也可缓慢改变。定义 $\epsilon_3=\dot\epsilon_2/(H\epsilon_2)$（若 $\epsilon_2=0$ 则直接用其导数），有 $\alpha_s=\mathrm dn_s/\mathrm d\ln k\simeq-2\epsilon_1\epsilon_2-\epsilon_2\epsilon_3$、$\alpha_t\simeq-2\epsilon_1\epsilon_2$。以 $L=\ln(k/k_p)$ 展开 $\ln\mathcal P_t=\ln\mathcal P_t(k_p)+n_t(k_p)L+\alpha_t(k_p)L^2/2+\cdots$。微波背景与干涉仪对应的 $L$ 可达几十，枢轴附近很小的二阶项也会被 $L^2$ 放大。势若在这段场程内拐弯，截断的 Taylor 展开甚至会给无效的大振幅，应逐模用其真实离开历史求谱。
\subsection{三种势的逐项计算}

取单项式势 $V=\lambda\phi^p$，$\phi>0$。有 $\epsilon_V=p^2M_{\rm P}^2/(2\phi^2)$、$\eta_V=p(p-1)M_{\rm P}^2/\phi^2$。以 $\epsilon_V(\phi_{\rm end})=1$ 估计暴涨结束，$\phi_{\rm end}^2=p^2M_{\rm P}^2/2$。积分慢滚方程得 $\phi_k^2/M_{\rm P}^2=2pN_k+p^2/2$。在同一近似下
\[
 r=\frac{16p}{4N_k+p},\qquad
 n_s-1=-6\epsilon_V+2\eta_V=-\frac{2(p+2)}{4N_k+p}.
\]
例如 $p=2$ 时 $n_s\simeq1-2/N_k$ 而 $r\simeq8/N_k$。调节 $\lambda$ 会改变 $A_s$，却不会在给定 $N_k$ 时改变以上两个比值。另一例是小场在山顶附近的 $V=V_0[1-(\phi/\mu)^p+\cdots]$、$p>2$。在 $\phi/\mu\ll1$ 时 $|\eta_V|\gg\epsilon_V$，$N_k\simeq\mu^p/[p(p-2)M_{\rm P}^2\phi_k^{p-2}]$（假定上端对积分可略）。因而 $n_s-1\simeq-2(p-1)/[(p-2)N_k]$，而 $r\simeq8p^2M_{\rm P}^2\phi_k^{2p-2}/\mu^{2p}$ 可以很小。若 $p=2$ 上式积分应改为对数，不能把 $p=2$ 硬代进去。

还有一例无需一开始便设基本标量。考虑 $f(R)=R+R^2/(6M^2)$，引入辅助量 $s$，以 $f'(s)(R-s)+f(s)$ 代 $f(R)$；对 $s$ 变分得 $f''(s)(R-s)=0$，故在此模型中二者等价。设 $F=f'(s)=1+s/(3M^2)>0$，作 Einstein 标架变换 $\widetilde g_{\mu\nu}=F g_{\mu\nu}$，参看 \S\ref{sec:scalar} 的标量--张量变换。规范标量 $\phi=\sqrt{3/2}\,M_{\rm P}\ln F$ 的势为
\[
 V=\frac{M_{\rm P}^2}{2}\frac{sF-f(s)}{F^2}
   =\frac34 M_{\rm P}^2M^2
   \left(1-e^{-\sqrt{2/3}\,\phi/M_{\rm P}}\right)^2 .
\]
令 $y=\exp(\sqrt{2/3}\,\phi/M_{\rm P})$。直接求导得 $\epsilon_V=4/[3(y-1)^2]$、$\eta_V=4(2-y)/[3(y-1)^2]$，而 $N_k=(3/4)[y_k-y_{\rm end}-\ln(y_k/y_{\rm end})]$。平台上 $y_k\gg1$，取首项即 $y_k\simeq4N_k/3$。代回得 Starobinsky 模型的 $n_s\simeq1-2/N_k$、$r\simeq12/N_k^2$。它同二次单项式有相同的标量倾角首项，张量比却差一个 $N_k$ 的阶；仅比较温度平台上的斜率不容易分开这两种势。

\subsection{一次模式匹配的量子计算}

也可全程用量子振子计算遗迹引力波。为看清物理，取暴涨末 $\eta_1<0$，之前 $a=-1/(H\eta)$，之后 $a=(\eta-2\eta_1)/(H\eta_1^2)$，使 $a,a'$ 在 $\eta_1$ 连续。辐射阶段 $a''=0$，选模式 $e^{-\mathrm i k(\eta-\eta_1)}/\sqrt{2k}$。令 $x=k\eta_1$，把入射的 \eqref{eq:gw-d-bd} 在后段展开为 $[\alpha e^{-\mathrm i k(\eta-\eta_1)}+\beta e^{\mathrm i k(\eta-\eta_1)}]/\sqrt{2k}$。匹配 $v,v'$ 给
\begin{align*}
 \alpha+\beta&=e^{-\mathrm i x}(1-\mathrm i/x),\\
 -\mathrm i(\alpha-\beta)&=e^{-\mathrm i x}(-\mathrm i-1/x+\mathrm i/x^2),\\
 \alpha&=e^{-\mathrm i x}(1-\mathrm i/x-1/(2x^2)),\qquad
 \beta=e^{-\mathrm i x}/(2x^2).
\end{align*}
直接代回有 $|\alpha|^2-|\beta|^2=1$。从初态真空产生的每偏振平均粒子数 $n_k=1/(4x^4)$。在辐射时代晚期，$k/a$ 已足够大，可赋予每粒子能量 $k/a$。扣去零点能，两个偏振合计的物理能量谱为
\[
 \frac{\mathrm d\rho_{\rm gw}}{\mathrm d\ln k}
 =2\frac{4\pi k^3}{(2\pi)^3a^3}\frac{k}{a}|\beta|^2
 =\frac{1}{4\pi^2a^4\eta_1^4}.
\]
用 $\eta_1^{-1}=-a_1H$ 及当时的背景密度红移，份额为 $H^2/(12\pi^2M_{\rm P}^2)=\mathcal P_t^{\rm p}/24$，与经典传递所得相同。若初态已有对称的占有数 $n$， 末态为 $n+(2n+1)|\beta|^2$，常数项即自发产生，比例于 $n$ 的项是受激部分。在 $|x|\gg1$ 时，突然拼接的高频尾取决于交接的理想化；真实再热时变过程若对该模式绝热，粒子产生会更弱。故不宜把以上常能谱积分到无限频率。

两种计算对测量的侧重点不同。角谱分析用冻结的原初方差作为初值，配 \eqref{eq:gw-d-trans} 得不同时代的随机场。粒子算法在入内后的正频基底上给出能量和占有数。留下的长波是成对模式的相干量子态，具有固定的增长与衰减分量关系；若误作两个任意独立初相的古典行波，就会失去穿越 Hubble 半径时选定的驻波相位。

\section{早期宇宙中运动的物质如何发波}\label{sec:gw-dsource}

暴涨时无物质剪切的量子振子也能产生张量涨落。热宇宙中还有另一种生产机制：相变、凝聚场破裂、流体及线缺陷的运动都提供经典的横向无迹应力。对这类源，猜一个应力的典型频率还不够。相同空间图样若仅保持一瞬，与保持许多振荡周期所产生的谱可以很不一样。

\subsection{非等时应力相关}

保留式 \eqref{eq:gw-d-wave} 的源，在辐射时代 $a''=0$，并设 $\eta_i$ 之前无这类波。用迟滞 Green 函数得
\begin{align}
 u_A(k,\eta)&=\frac{16\pi G}{k}\int_{\eta_i}^{\min(\eta,\eta_f)}
 a^3(\eta_1)\sin[k(\eta-\eta_1)]\Pi_A(k,\eta_1)\,\mathrm d\eta_1,\nonumber\\
 h_A'(k,\eta)&=\frac{16\pi G}{a(\eta)}
 \int_{\eta_i}^{\eta_f} a^3(\eta_1)\cos[k(\eta-\eta_1)]
 \Pi_A(k,\eta_1)\,\mathrm d\eta_1-\mathcal H h_A(k,\eta)
 \qquad (\eta>\eta_f). \label{eq:gw-d-sourcegreen}
\end{align}
短波时略去末项。定义 $\langle\Pi_A(\bm k,\eta_1)\Pi_B^*(\bm q,\eta_2)\rangle=(2\pi)^3\delta^3(\bm k-\bm q)\mathcal U_{AB}(k;\eta_1,\eta_2)$。这是\textbf{非等时相关函数}。对晚时刻若干周期平均，两个余弦的乘积只剩 $\cos[k(\eta_1-\eta_2)]/2$。将 $e^A:e^B=2\delta^{AB}$、相空间测度 $k^3/(2\pi^2)$ 以及本地时间导数的 $a^{-1}$ 都计入，得
\begin{equation}\label{eq:gw-d-uetc}
 \frac{\mathrm d\rho_{\rm gw}}{\mathrm d\ln k}
 =\frac{4G k^3}{\pi a^4}
 \int_{\eta_i}^{\eta_f}\!\mathrm d\eta_1\,\mathrm d\eta_2\,
 a_1^3a_2^3\cos[k(\eta_1-\eta_2)]
 \operatorname{Re}\sum_A\mathcal U_{AA}(k;\eta_1,\eta_2).
\end{equation}
右边为实，可由最终振幅的绝对值平方看出总结果非负。独立源的非等时相关相加，所以能谱相加；若源具有确定的相互相位，还会有交叉项。$\mathcal U$ 的量纲含一个共动体积，不能直接用局部应力平方代替而不指定空间相关长度。

若在短暂源期间可固定 $a$，并写 $\Pi_A(k,\eta)=p_A(k)b(\eta)$，则双重积分成为 $|\int b(\eta)e^{\mathrm i k\eta}\mathrm d\eta|^2\sum_A P_{p_A}(k)$。矩形持续时间 $\Delta\eta$ 给 $4\sin^2(k\Delta\eta/2)/k^2$：在 $k\Delta\eta\ll1$ 时强度正比 $\Delta\eta^2$，高频却发生端点相消。若源每经过相关时间 $\tau_c$ 就丢失相位，且总持续时间远大于 $\tau_c$，取 $\mathcal U=P_\Pi e^{-|\eta_1-\eta_2|/\tau_c}$ 便得双积分近于 $2P_\Pi\Delta\eta\,\tau_c/(1+k^2\tau_c^2)$。此时强度只与持续时间成正比。湍流、声波、快速气泡壁的结果相异，一部分原因正是这些相干时标不同。

对因果生成且具有有限空间和时间相关长度的源，在低到无法分辨此相关尺度的 $k$，相应应力谱的模量通常趋于有限常数，而 $\cos[k(\eta_1-\eta_2)]\to1$。因而式 \eqref{eq:gw-d-uetc} 的红外端随 $k^3$，三个幂全来自三维相空间。取此极限仍须在观测时代波已入内时定义能量；长寿命的标度源或具有超 Hubble 初始关联的源不满足有限时间或有限相关尺度的前提，低频斜率需重算。

\subsection{标量场的梯度源}

一个实标量的应力中，$\delta^i{}_j$ 项全被 TT 投影消掉，留下 $a^{-2}\partial_i\phi\partial_j\phi$。这说明即使均匀背景按极大幅度振荡，没有空间梯度也不会产生经典张量波。令 $P_{ij}=\delta_{ij}-\hat k_i\hat k_j$，则对称张量的投影核为 $\Lambda_{ij,lm}=P_{i(l}P_{m)j}-P_{ij}P_{lm}/2$。在动量空间中先作卷积，然后才以\emph{总}波矢 $\bm k$ 投影，
\[
 \Pi_A(\bm k,\eta)
 =\frac{1}{2a^2}\int\frac{\mathrm d^3p}{(2\pi)^3}
       e^A_{ij}(\hat{\bm k})p_i p_j
       \phi(\bm p,\eta)\phi(\bm k-\bm p,\eta).
\]
原式有 $-p_i(k-p)_j$，其中 $k_j$ 一项被横向性消去，余下正号。若 $\bm p$ 恰平行 $\bm k$，便不能贡献剪切。这一点在用数值格点估算谱时很重要，离散梯度和 TT 投影所用的波矢必须相容。

设 $\phi$ 为统计均匀的零均值高斯场，$\langle\phi(\bm p,\eta_1)\phi^*(\bm q,\eta_2)\rangle=(2\pi)^3\delta^3(\bm p-\bm q)C_\phi(p;1,2)$。四点函数有两组交叉配对，其结果相同；各自内部的配对只给齐次 $k=0$ 分量。取 $\bm k$ 沿 $z$，对 $\varphi$ 有 $e^+_{ij}p_ip_j=p^2(1-\mu^2)\cos2\varphi$、$e^\times_{ij}p_ip_j=p^2(1-\mu^2)\sin2\varphi$。平方和后得
\[
 \sum_A\mathcal U_{AA}
 =\frac{1}{2a_1^2a_2^2}\int\frac{\mathrm d^3p}{(2\pi)^3}
 p^4(1-\mu^2)^2 C_\phi(p;1,2)
 C_\phi\big(\sqrt{k^2+p^2-2kp\mu};1,2\big).
\]
各向同性只消掉 $\bm k$ 方向和方位角，不能省去 $\mu$ 的权重。若场在相变中已成为高度非高斯的团块，四点函数另有连通项；只从标量的等时能谱猜测波谱，就可能同时漏掉时相关与连通四点函数。

\subsection{预热与参数共振}

慢滚结束后，若势在谷底近于 $m^2\phi^2/2$，均匀 $\phi\simeq\Phi(t)\sin mt$ 作阻尼振荡。在很多振荡上的平均有 $\langle\dot\phi^2\rangle\simeq m^2\langle\phi^2\rangle$，故平均压强近零、$\Phi\propto a^{-3/2}$。能量转入其它场不一定是逐粒子的缓慢衰变。例如势含 $g^2\phi^2\chi^2/2$ 时，$\chi$ 的各模有周期的有效质量。取 $X_k=a^{3/2}\chi_k$，由 $\ddot\chi_k+3H\dot\chi_k+[k^2/a^2+g^2\phi^2]\chi_k=0$ 得
\[
 \ddot X_k+
 \left[\frac{k^2}{a^2}+g^2\Phi^2\sin^2 mt
 -\frac94H^2-\frac32\dot H\right]X_k=0.
\]
物质主导的平均背景中最后两项抵消。若一周期内 $a,\Phi$ 变化缓慢，令 $z=mt$，可化为 Mathieu 方程
\[
 \frac{\mathrm d^2X_k}{\mathrm dz^2}
 +[A_k-2q\cos(2z)]X_k=0,\qquad
 q=\frac{g^2\Phi^2}{4m^2},\qquad A_k=\frac{k^2}{a^2m^2}+2q .
\]
周期系数容许 $X_k=e^{\mu_k z}P_k(z)$ 的 Floquet 解，其中 $P_k$ 是周期或反周期函数。若 $\operatorname{Re}\mu_k>0$，场振幅指数增长，占有数在增长段近于 $e^{2\operatorname{Re}\mu_k mt}$。窄带极限 $q\ll1$ 时，设 $X=C(z)e^{\mathrm iz}+D(z)e^{-\mathrm iz}$，只留可与基频共振的项，得 $2\mathrm i C'+(A_k-1)C-qD=0$ 与 $-2\mathrm iD'+(A_k-1)D-qC=0$。其特征值满足 $\mu_k^2=[q^2-(A_k-1)^2]/4$，故第一不稳定带为 $|A_k-1|<q$。膨胀会让同一个共动 $k$ 扫过不稳定带，不能只看瞬时 $\mu_k$ 的极大值。

$q\gg1$ 时每次 $\phi$ 过零都会使 $\omega_k^2=k^2/a^2+g^2\phi^2$ 急变，宽带模式可被连续放大。在过零附近 $\phi\simeq\dot\phi_0(t-t_0)$，绝热性指标 $|\dot\omega_k|/\omega_k^2$ 在 $g|\dot\phi_0||t-t_0|\sim k/a$ 附近最大，于是非绝热范围的估计为 $(k/a)^2\lesssim g|\dot\phi_0|\sim gm\Phi$。激发的 $\chi$ 与 $\phi$ 再散射，扰动回馈到背景，直至线性 Floquet 解失效。此时同一物理尺度的场梯度依式 \eqref{eq:gw-d-uetc} 发波。峰频由 $k_*/a_*$ 决定，而暴涨能量何时转入热浴还控制之后的红移，不能一律将其当作 $H_*$。

另一机制是\textbf{快子预热}。例如混合势 $V=\lambda(\chi^2-v^2)^2/4+g^2\phi^2\chi^2/2+V_{\rm sr}(\phi)$，$\chi=0$ 附近的有效质量平方是 $g^2\phi^2-\lambda v^2$。当 $\phi$ 滚过临界值 $\sqrt\lambda v/g$，足够长波的 $\chi_k$ 即有 $k^2/a^2<|m_{\rm eff}^2|$；以冻结背景估计，其振幅按 $\exp[\int\sqrt{|m_{\rm eff}^2|-k^2/a^2}\,\mathrm dt]$ 增长。邻近区域选择不同真空方向，随后场畴碰撞、梯度急增，便能发出引力波。这个不稳定性无需等待背景许多次过零；它与窄带参数共振的初始时间尺度和峰波数都不同。


\subsection{过零的非绝热产生与反作用}

对宽共振，可把单次过零独立解开。令 $q_{\rm phys}=k/a(t_0)$、$k_c^2=g|\dot\phi_0|$，在 $|t-t_0|\ll H^{-1},m^{-1}$ 内 $\omega_k^2\simeq q_{\rm phys}^2+k_c^4(t-t_0)^2$。置 $\tau=k_c(t-t_0)$、$\kappa=q_{\rm phys}/k_c$，有 $X_{\tau\tau}+(\tau^2+\kappa^2)X=0$，这是一只频率经过最小值的 Weber 振子。它的复转折点在 $\tau=\pm\mathrm i\kappa$。从 $0$ 到 $\mathrm i\kappa$ 的相位积分为
\[
 \int_0^{\mathrm i\kappa}\sqrt{\tau^2+\kappa^2}\,\mathrm d\tau
 =\mathrm i\kappa^2\int_0^1\sqrt{1-u^2}\,\mathrm du
 =\frac{\mathrm i\pi\kappa^2}{4}.
\]
解的正负频连接给从真空产生的占有数 $|\beta_\kappa|^2=e^{-\pi\kappa^2}$。因而实际激发并非在 $k_c$ 处突然截断，而是随 $q_{\rm phys}^2/k_c^2$ 指数压低。一次过零后每物理体积的 $\chi$ 个数是一个三维高斯积分，
\[
 n_\chi=\int\frac{\mathrm d^3q_{\rm phys}}{(2\pi)^3}
           e^{-\pi q_{\rm phys}^2/k_c^2}=\frac{k_c^3}{8\pi^3}.
\]
在 $\phi$ 离开零点后粒子有效质量 $g|\phi|$ 增大，背景场必须为它们作功，即使在很短的产生时间内其质量曾接近零。

每次过零的 Bogoliubov 矩阵都保持 $|\alpha|^2-|\beta|^2=1$。若在此前已有相反动量成对的占有数 $n$，将上次产生的正、负频振幅传播到下一次过零后，新占有数可写成
\[
 n_{\rm new}=p+(1+2p)n
 +2\sqrt{p(1+p)n(1+n)}\cos\vartheta,\qquad p=e^{-\pi\kappa^2}.
\]
$\vartheta$ 含两次过零之间 $\int\omega_k\mathrm dt$ 的相位以及过零处的散射相位。稳定的周期运动使得某些模反复相长，某些模反复相消；膨胀会令不同周期的相位漂移。若近似随机相位，交叉项平均趋小，却仍留下 $1+2p$ 乘旧占有数的受激增长。非绝热加热由此不同于给所有 $k$ 各乘一个恒定衰变率。

当 $g^2\langle\chi^2\rangle$ 与背景的 $m^2$ 可比，平均背景本身的振荡频率已改变，使用既定 Mathieu 参数的步骤失效。粗估 $\langle\chi^2\rangle\simeq\int n_k\,\mathrm d^3q_{\rm phys}/[(2\pi)^3\omega_k]$ 时，应扣去未激发真空的半量子，数值结果还应检验网格紫外截断。即使此时有大的平均 $\chi^2$，波仍只由带方向的梯度应力产生。数值演化常在每个时间步先由标量场求未经投影的 $T_{ij}$，用与离散梯度一致的 $P_{ij}$ 投影，再以驱动方程更新两支张量模；若先对场做径向平均才计算 $T_{ij}$，各向异性会在平方之前丢失。

一个格点谱的独立核查是能量预算。若生成峰物理波数 $k_{\rm phys}$，梯度份额 $\rho_\nabla/\rho_{\rm tot}=\epsilon_\nabla$，且持续的相干时间约 $k_{\rm phys}^{-1}$，则 $h\sim16\pi G\rho_\nabla/k_{\rm phys}^2$。忽略几何因子，峰的能量份额量级为 $\epsilon_\nabla^2(H_*/k_{\rm phys})^2$。共振尺度越短，虽然空间梯度更强，但在固定梯度能量份额下波的产生效率越小。若积分许多周期或有长寿命的非球对称团块，需改用式 \eqref{eq:gw-d-uetc} 的时相关。将绝对格点能量谱换为今天的 $\Omega_{\rm gw}$ 还有一个经常漏掉的间隔：若数值演化停在源关闭 $a_f$，再热到 $a_{\rm reh}$ 才形成辐射浴，两时刻间背景具有平均 $w$，则相对份额须先乘 $(a_{\rm reh}/a_f)^{3w-1}$；只有之后才能用热浴熵守恒作红移。

\subsection{源的谱峰可以怎样由相关函数算出}

给一可解析的有限源，以显示长度与寿命分别在哪个积分中出现。设在某一很短的宇宙阶段可视 $a=a_*$，$\sum_A\mathcal U_{AA}$ 的空间形状为 $\Pi_0^2 R_c^3\exp[-k^2R_c^2/2]$，其中 $R_c$ 为共动相关长度。若相位在持续时段内完全相关，且振幅包络是 $b(\eta)=\exp[-(\eta-\eta_*)^2/(2\Delta_c^2)]$，代入式 \eqref{eq:gw-d-uetc} 的非等时部分 $b(\eta_1)b(\eta_2)$ 得 $|\widetilde b(k)|^2=2\pi\Delta_c^2e^{-k^2\Delta_c^2}$。故去掉频率无关因子后，
\[
 \frac{\mathrm d\rho_{\rm gw}}{\mathrm d\ln k}
 \propto k^3 \exp[-k^2(R_c^2/2+\Delta_c^2)],\qquad
 k_{\rm peak}=\sqrt{\frac{3}{R_c^2+2\Delta_c^2}} .
\]
$\Delta_c\gg R_c$ 时同一应力缓慢改变，源的高 $k$ 虽仍有功率，却很少在波的本征频率上作净功，峰主要由时间傅里叶宽度决定。$\Delta_c\ll R_c$ 时是空间尺度限制了传播波数。这一算例采用高斯相关以便积分，未假定其高频端是所有因果源的通式。

若改为统计定常但每隔 $\tau_c$ 失相的源，可以保留有限总历时 $D_c$ 精确积分。令 $z_c=\tau_c^{-1}-\mathrm i k$，得
\[
 \int_0^{D_c}\!\mathrm d\eta_1\!\int_0^{D_c}\!\mathrm d\eta_2\,
 e^{-|\eta_1-\eta_2|/\tau_c}\cos[k(\eta_1-\eta_2)]
 =2\operatorname{Re}\!\left[\frac{D_c}{z_c}
             -\frac{1-e^{-z_cD_c}}{z_c^2}\right].
\]
当 $\tau_c\to\infty$ 时它还原 $2(1-\cos kD_c)/k^2$；当 $D_c/\tau_c\gg1$ 时则给 $2D_c\tau_c/(1+k^2\tau_c^2)$ 的首项。前者可有精确的相消频率，后者的功率对持续时间作近线性积累。因而仅凭一幅相变末期的等时应力图样，无论空间分辨率多高，都不足以在这两类谱之间作选择。
\section{一阶相变留下的波谱}\label{sec:gw-dphase}

\subsection{自由能的两个极小值}

以温度 $T$ 的有效势 $V(\phi,T)$ 表示固定序参量时的自由能密度。一条平滑交叉只使平衡序参量连续改变；若自由能同时有两个相隔势垒的局部极小，降温经过二者等深的临界温度 $T_c$，系统可继续滞留在较高的假真空，随后才通过成核转入稳定相。后者才会产生由大尺度气泡驱动的剪切应力。不能仅因温度到了电弱或 QCD 尺度便假设必有气泡；需要先检验所用粒子模型的有限温度势与相图。

例如考虑近临界展开 $V=D(T^2-T_b^2)\phi^2-E T\phi^3+\lambda\phi^4/4$，取 $D,E,\lambda>0$ 且 $E^2<D\lambda$。在 $T_c$ 同时满足 $V(\phi_c,T_c)=V(0,T_c)$ 和 $V_{,\phi}(\phi_c,T_c)=0$。先把前一式除以 $\phi_c^2$，后一式除以 $\phi_c$，从后一式减去两倍前一式，得 $-ET_c\phi_c+\lambda\phi_c^2/2=0$，于是
\[
 \frac{\phi_c}{T_c}=\frac{2E}{\lambda},\qquad
 T_c^2=\frac{T_b^2}{1-E^2/(D\lambda)},\qquad
 V(\phi,T_c)=\frac{\lambda}{4}\phi^2(\phi-\phi_c)^2 .
\]
两相虽等自由能，能量密度 $\rho=V-T\partial_TV$ 仍可不等。有 $\Delta V=V_{\rm false}-V_{\rm true}$、$\Delta\rho=\Delta V-T\partial_T\Delta V$，在 $T_c$ 有 $\Delta V=0$ 而 $\Delta\rho=-T_c\partial_T\Delta V$。求导时 $\phi_{\rm min}(T)$ 的隐含变化乘 $V_{,\phi}=0$，故只须偏导温度项。这是压力驱动力 $\Delta p=\Delta V$ 与热浴可释放的能量不同的一个例子。通常以 $\alpha=\Delta\rho/\rho_{r,*}$ 定义转变强度；若采用应力迹定义的 $\alpha$，驱动效率的拟合也须使用同一约定。


\subsection{热自由能中势垒从何而来}

上面的 $T\phi^3$ 有时来自玻色子的热运动。设一类有 $g$ 个内部态的热平衡玻色子，其能量 $E_p=(p^2+m^2(\phi))^{1/2}$ 随缓变序参量而改变。把每个谐振子的热配分函数相乘，零点项另行重整化，单位体积的热自由能为 $\Delta V_B=gT\int\frac{\mathrm d^3p}{(2\pi)^3} \ln(1-e^{-E_p/T})$。当 $m=0$，以 $\ln(1-e^{-x})=-\sum_{n\ge1}e^{-nx}/n$ 展开，积分 $\int_0^\infty x^2e^{-nx}\mathrm dx=2/n^3$，故 $\Delta V_B(0)=-g\pi^2T^4/90$。为求质量依赖，先对 $m^2$ 微分便少一个待处理的对数：
\[
 \frac{\partial\Delta V_B}{\partial m^2}
 =\frac g2\int\frac{\mathrm d^3p}{(2\pi)^3}
                \frac{1}{E_p(e^{E_p/T}-1)}.
\]
直接在 $m=0$ 处积分，得 $gT^2/24$。若企图再作关于 $m^2$ 的普通 Taylor 级数，下一导数却会在 $p=0$ 不可积；红外玻色子决定下一项是非解析的。为取其系数，在软区 $E_p\ll T$ 用 $1/(e^{E_p/T}-1)\simeq T/E_p$，并减去 $m=0$ 时已计的软贡献，得
\[
 \frac{gT}{4\pi^2}\int_0^\infty\left(\frac{p^2}{p^2+m^2}-1\right)\mathrm dp
 =-\frac{gTm}{8\pi}.
\]
这次差值积分由 $p\sim m$ 主宰，上限延到无穷不改所求的线性 $m$ 项。对 $m^2$ 积回去便有
\[
 \Delta V_B=-\frac{g\pi^2T^4}{90}+\frac{gT^2m^2}{24}
             -\frac{gT m^3}{12\pi}
             +O\!\left(m^4\ln\frac{m}{T}\right),\qquad m\ll T .
\]
若质量来自正的 $m^2=c_b\phi^2$，对 $\phi\ge0$ 这第三项就是 $-g c_b^{3/2}T\phi^3/(12\pi)$。费米子的热势为 $-gT\int\ln(1+e^{-E_p/T})\mathrm d^3p/(2\pi)^3$，零质量项为 $-(7/8)g\pi^2T^4/90$，二次项为 $gT^2m^2/48$。费米分布在 $E_p\to0$ 时有界，没有相同的经典软模增强，因而无此阶 $Tm^3$ 项。$\phi$ 无关的 $T^4$ 对比较同一温度下两相的自由能位置可相消，对总压强、熵、流体焓的导数却未必可删。

热浴还会对长波玻色子产生屏蔽。若软自由度的自能为 $\Pi_b(T)$，在热零模中用 $m^2+\Pi_b$ 代 $m^2$，须扣除按同一方法已计算过的贡献。典型重求和修正是 $-gT[(m^2+\Pi_b)^{3/2}-m^3]/(12\pi)$，其中与 $\phi$ 无关的项可按能量零点处理。当 $m^2\ll\Pi_b$ 时，$(m^2+\Pi_b)^{3/2}=\Pi_b^{3/2}+(3/2)\sqrt{\Pi_b}\,m^2+\cdots$，三次势垒会减弱。哪些热模有何种屏蔽取决于场内容与相互作用。近临界区若涨落强，微扰有效势不够，须以更可靠的有限温度计算确定转变阶数。这就是从有形式上的玻色零模到预言一个可发波的强一阶转变，中间还需作的物理判断。

对前述多项式势，可以不靠作图就跟踪两个极小值。除 $\phi=0$ 外，极值满足 $\lambda\phi^2-3ET\phi+2D(T^2-T_b^2)=0$。若 $9E^2<8D\lambda$，温度从高处降至
\[
 T_1^2=\frac{T_b^2}{1-9E^2/(8D\lambda)},\qquad
 \phi_1=\frac{3ET_1}{2\lambda}
\]
时首次出现重根；再降低温度，较小根是势垒顶部，较大根是真空候选。$T_1>T_c>T_b$，在 $T_1$ 至 $T_c$ 间新极小仍较高，在 $T_c$ 至 $T_b$ 间 $\phi=0$ 已成亚稳局部极小，在 $T_b$ 则 $V_{,\phi\phi}(0,T)$ 变号、原点失去局部稳定性。若宇宙一直未成核就进入不稳定区，长波直接滚落的动力学应与新相气泡在势垒下成核分开。以 $T_c$ 处 $\phi_c=2ET_c/\lambda$ 求两相能量差，
\[
 \Delta\rho_c=T_c[2DT_c\phi_c^2-E\phi_c^3]
 =\frac{8DE^2T_c^4}{\lambda^2}
 \left(1-\frac{E^2}{D\lambda}\right)>0.
\]
用这个量求靠近 $T_c$ 的驱动力，才可与 $\sigma$ 一起给出薄壁半径的温度依赖。
\subsection{鞍点为何给出成核率}

先回看一维亚稳态。Euclid 时间 $\tau=\mathrm i t$ 将权重变为 $e^{-S_E}$，其中 $S_E=\int[\dot q^2/2+V(q)-V(q_f)]\,\mathrm d\tau$。非平凡的经典轨道从亚稳点 $q_f$ 出发，在反转势 $-V$ 中走到转折点再返回，称\textbf{反弹解}(bounce)。一次积分给 $\dot q^2/2=V(q)-V(q_f)$，往返作用量遂为 $2\int_{q_f}^{q_{\rm turn}}\sqrt{2[V(q)-V(q_f)]}\,\mathrm dq$。这与势垒的 WKB 指数相同。反弹发生在何时不影响作用量，故沿 $\dot q_{\rm b}$ 的涨落是零模，把它换成反弹中心坐标的积分就得到观测时间因子。另有一个负模，沿适当陡降路径延拓后给亚稳自由能虚部，故结果刻画逃逸率。简单地用平衡分配函数在两个极小点附近各做一次正定高斯积分，并不能算出成核。

场论有限温度时 Euclid 时间周期为 $1/T$。零温极限常由四维旋转对称的反弹主导，而热激发充分时主鞍点几乎不随 Euclid 时间变，故指数由 $S_3/T$ 控制。取球对称构形，
\begin{align}
 S_3(T)&=4\pi\int_0^\infty r^2
 \left[\tfrac12(\partial_r\phi)^2+V(\phi,T)-V(\phi_f,T)\right]\mathrm dr,\nonumber\\
 \partial_r^2\phi+\frac2r\partial_r\phi&=V_{,\phi},\qquad
 \partial_r\phi(0)=0,\quad\phi(\infty)=\phi_f .\label{eq:gw-d-bounce}
\end{align}
这是含摩擦的反转势力学问题，半径充当“时间”。在中心设 $\phi(0)=\phi_0$，规则展开给 $\phi(r)=\phi_0+V_{,\phi}(\phi_0)r^2/6+\cdots$。若从某个 $\phi_0$ 释放后越过 $\phi_f$，改选更接近不越过一端的初值；由过冲、欠冲的分界求出反弹。边界值不能任意取为精确真极小，否则满足初值的常解只会静止在那里。计算中还可用标度恒等式检验：设 $K_\phi=\int(\nabla\phi)^2/2$，$U_\phi=\int[V(\phi)-V_f]$，代缩放 $\phi(r/\lambda)$，$S_3(\lambda)=\lambda K_\phi+\lambda^3U_\phi$，故驻值时 $K_\phi+3U_\phi=0$，$S_3=2K_\phi/3>0$。沿这条缩放方向二阶导为 $6U_\phi<0$，正给出临界气泡的不稳定方向。

成核率每单位物理时间、单位物理体积可写作 $\Gamma(T)=A(T)e^{-S_3(T)/T}$。三个平移零模转成泡心位置后留下 $(S_3/2\pi T)^{3/2}$ 等前因子；量纲上 $A$ 常近于 $T^4$ 乘缓变因子。若温度降得很低，Euclid 时间依赖不能略去，须比较四维反弹 $S_4$ 及三维热反弹的指数。以下对热率的公式只用于其确实占主导的温度区间。


\subsection{一条可解反弹的零模与负模}

用一个具体亚稳势能看清楚为什么只留\emph{一个}负模。令质量为一、$V(q)=\omega^2q^2(1-q/q_e)/2$，$q=0$ 为局部极小，$q_e$ 是能量相同的转折点。零温反弹的一次积分 $\dot q^2=\omega^2q^2(1-q/q_e)$ 可分离变量，给出
\[
 q_b(\tau)=q_e\,\operatorname{sech}^2[\omega(\tau-\tau_0)/2],\qquad
 B=2\int_0^{q_e}\omega q\sqrt{1-q/q_e}\,\mathrm dq
   =\frac{8}{15}\omega q_e^2.
\]
虚时间中心 $\tau_0$ 任意，$B$ 与之无关。令 $x=\omega(\tau-\tau_0)/2$。二次涨落算符 $\mathcal M=-\partial_\tau^2+V''(q_b)$ 化为
$\mathcal M=(\omega^2/4)[-\partial_x^2+4-12\operatorname{sech}^2x]$。直接微分 $q_b$ 的方程得 $\mathcal M\dot q_b=0$。$\dot q_b$ 关于 $x=0$ 为奇函数，恰有一个节点；按一维束缚态的节点定理，零本征值之下恰好还有一个没有节点的束缚态。试取 $u_-(x)=\operatorname{sech}^3x$，由
$u_-''/u_-=9-12\operatorname{sech}^2x$ 立即得 $\mathcal M u_-=-(5\omega^2/4)u_-$。这一个负模并非计算错误，它刻画了反弹解处在越过势垒的分界。

归一化零模 $u_0=\dot q_b/\sqrt B$，因为一次积分给 $\int\dot q_b^2\mathrm d\tau=B$。若把其系数写 $c_0$，整体平移反弹有 $\delta q_b=-\dot q_b\,\delta\tau_0$，故 $|\mathrm dc_0|=\sqrt B\,|\mathrm d\tau_0|$。原先发散的高斯零模积分应换成 $\sqrt{B/(2\pi)}\int\mathrm d\tau_0$，不是在行列式里任意丢掉一个小本征值。对其余正模积分得行列式比，负模沿陡降线延拓。围绕亚稳支作延拓时取逃逸方向的一半高斯，产生虚因子 $\mathrm i/2$。以 Euclid 总时间 $\mathcal T$ 观察，一枚反弹相对假真空振幅的因子可写 $\mathrm i\mathcal T K e^{-B}/2$，其中 $K$ 为实正行列式前因子。若有 $n$ 枚相距足够远的反弹，每个中心可独立积分，再除以不可区分排列 $n!$；稀薄极限求和得
\[
 \sum_{n=0}^\infty\frac{(\mathrm i\mathcal T K e^{-B}/2)^n}{n!}
 =\exp(\mathrm i\mathcal T K e^{-B}/2).
\]
与 $\exp[-(E_f-\mathrm i\Gamma/2)\mathcal T]$ 比较，得 $\Gamma=Ke^{-B}$。于是一次穿越的虚振幅变为单位实时间的概率衰减率。密集反弹相互影响时独立中心的指数求和不再有效。

热场论中空间平移给三个零模，而静态三维反弹在虚时间方向并没有第四个独立中心。将 $\phi_b(\bm x-\bm x_0)$ 对 $x_0^i$ 求导，得到 $u_i\propto\partial_i\phi_b$。球对称性及 $K_\phi=3S_3/2$ 给 $\int(\partial_i\phi_b)^2\mathrm d^3x=S_3$（固定一个 $i$）。于是三个积分合成 $\mathcal V(S_3/(2\pi T))^{3/2}$，除以物理体积 $\mathcal V$ 后正是热成核前因子的平移部分。反弹涨落的静态算符是 $-\nabla^2+V''(\phi_b)$，展开在球谐上，径向有效势另有 $\ell(\ell+1)/r^2$。平移零模在 $\ell=1$，缩放的负方向在 $\ell=0$。若其负本征值为 $-\omega_-^2$，虚时间 Fourier 模标 $n$ 给本征值 $(2\pi nT)^2-\omega_-^2$。足够高温时只有 $n=0$ 保留这个负模；当 $2\pi T$ 降过 $\omega_-$ 时又出现时变的不稳定方向，说明选静态三维鞍点作为主反弹的近似需要重审。这比仅比较温度与势垒高度更贴近反弹的时间尺度。
\subsection{薄壁极限的临界气泡}

在 $T$ 稍低于 $T_c$ 时，两相自由能差 $\Delta V>0$ 小，气泡半径可比壁厚大很多。用近乎平面且静止的界面计算表面张力 $\sigma$。在等深势 $V_c$ 上，$\partial_z^2\phi=V_{c,\phi}$ 的一次积分给 $(\partial_z\phi)^2/2=V_c-V_{\rm min}$，故
\[
 \sigma=\int_{-\infty}^\infty(\partial_z\phi)^2\mathrm dz
 =\int_{\phi_f}^{\phi_t}\sqrt{2(V_c-V_{\rm min})}\,|\mathrm d\phi|.
\]
对上一小节的四次势，这个积分等于 $\sqrt{\lambda/2}\,\phi_c^3/6$。反弹自由能现可只对半径变分：$S_3(R)=4\pi R^2\sigma-4\pi R^3\Delta V/3$。极小气泡主要损失表面积能，不能增长；半径足够大时体积能获益占优。两者的分界及指数为
\[
 R_c=\frac{2\sigma}{\Delta V},\qquad
 S_3(R_c)=\frac{16\pi\sigma^3}{3\Delta V^2},\qquad
 \left.\frac{\mathrm d^2S_3}{\mathrm dR^2}\right|_{R_c}=-8\pi\sigma .
\]
若 $\Delta V\simeq\Delta\rho_c(T_c-T)/T_c$，则在微小过冷处 $S_3/T\propto(T_c-T)^{-2}$，说明成核几乎不可能恰在 $T_c$ 发生。薄壁条件随过冷加深而变差，须回式 \eqref{eq:gw-d-bounce}。四维零温类比有 $S_4(R)=2\pi^2R^3\sigma-\pi^2R^4\Delta V/2$，给 $R_c^{(4)}=3\sigma/\Delta V$、$S_4=27\pi^2\sigma^4/(2\Delta V^3)$，不能把三维解直接除以一个温度外推到零温。

\subsection{多久能填满空间}

一个 Hubble 体积中平均出现一枚气泡给出成核温度的估计。在辐射主导且热自由度近乎不变时，$\dot T=-HT$，故条件写成 $\int_{T_n}^{\infty}\Gamma(T)\,\mathrm dT/[T H^4(T)]\sim1$。仅就指数精度，令 $\Gamma\sim H^4$ 得 $S_3/T\sim4\ln(T/H)$；电弱量级时这可达一百余，故用指数相当大的鞍点计算仍有每 Hubble 体积一次的机会。成核温度、泡相互接触的温度和释放热量后的温度应分开。同样把成核率在碰撞时期附近展开为 $\Gamma(t)\simeq\Gamma_*e^{\beta(t-t_*)}$，忽略前因子的缓变，得
\begin{equation}\label{eq:gw-d-beta}
 \frac{\beta}{H_*}\simeq
 \left.T\,\frac{\mathrm d}{\mathrm dT}\frac{S_3(T)}{T}\right|_* .
\end{equation}
若强过冷使真空能主导或再加热显著，$\dot T=-HT$ 的假定失效，应该直接求 $\mathrm d\ln\Gamma/\mathrm dt$。

令壁相对等离子体背景以速度 $v_w(t)$ 推进，在 $t'$ 成核到 $t$ 所及的共动半径 $r(t,t')=\int_{t'}^t v_w(t'')\mathrm dt''/a(t'')$。取一点，它在一给定时刻仍未被任何泡覆盖的概率是 $P_f(t)=e^{-I(t)}$，其中 Poisson 过程的期望覆盖次数为 $I(t)=\frac{4\pi}{3}\int^t\Gamma(t')a^3(t')r^3(t,t')\,\mathrm dt'$。$\Gamma$ 取可延伸定义在全空间上的理想成核率；对实际已转变区抑制成核的处理须同所用泡增长模型一致。在 $\beta\gg H_*$、$v_w$ 恒定的快变极限，$a^3(t')r^3\simeq v_w^3(t-t')^3$，于是 $I=8\pi v_w^3\Gamma(t)/\beta^4$。当 $I\sim1$ 时才有相当空间份额完成转变。把实际形成的泡密度估计为 $n_b=\int\Gamma(t)e^{-I(t)}\mathrm dt$，代 $\mathrm dI=\beta I\mathrm dt$ 得 $n_b\simeq\beta^3/(8\pi v_w^3)$，平均间距 $R_*=n_b^{-1/3}\simeq(8\pi)^{1/3}v_w/\beta$。这给波长一把可计算的尺。若膨胀不能忽略，还需检查 $a^3P_f$ 是否真正减少；仅有某些位置成核并不保证宇宙由新相占据。

\subsection{壁、声波与湍流分别作什么功}

以壁静系比较紧贴壁两侧的完美流体，令 $+$、$-$ 分别为壁前、壁后，相对壁的流速大小 $v_\pm$ 均取正，焓密度 $w_{\!h}=\rho+p$。能量、法向动量流连续要求
\[
 w_{\!h+}\gamma_+^2v_+=w_{\!h-}\gamma_-^2v_-=J,\qquad
 Jv_++p_+=Jv_-+p_- .
\]
由第二式 $J(v_+-v_-)=p_--p_+$，再以第一式写 $w_{\!h\pm}=J(1/v_\pm-v_\pm)$，消去 $J$ 得
\[
 v_+v_-=\frac{p_+-p_-}{\rho_+-\rho_-},\qquad
 \frac{v_+}{v_-}=\frac{\rho_-+p_+}{\rho_++p_-}.
\]
若给两相的状态方程及一侧温度，这两个关系选出可行的燃烧解。壁超声速推进而前方流体静止可以给爆轰支；亚声速推进可伴随提前加热流体的冲击前沿。局域跃迁条件本身不能给出壁速，还须用壁的微观摩擦和流体远边界条件。各向同性地膨胀的孤立球泡只有球对称应力，对外无张量辐射；泡碰撞破坏对称性，流体形成的非球对称声壳在碰撞后还可持续做功。

设参与某一源的能量份额 $K_X=\kappa_X\alpha/(1+\alpha)$，$\kappa_X$ 为释放能转入壁或流体的效率。若应力以尺度 $R_*$、幅度 $\Pi\sim K_X\rho_*$ 变化，由波方程得 $h\sim G\Pi R_*^2$、$\dot h\sim G\Pi R_*$。持续一个过境时间后 $\rho_{\rm gw}/\rho_*\sim K_X^2(H_*R_*)^2$。这一估计适用于短时的壁碰撞源；只保留尚未碰撞的壁面、已碰撞部分立即失活，称包络近似。若湍流或声壳可维持 $\tau_s\gg R_*$，按非等时相关分成若干有效独立时间段，粗估改为 $K_X^2(H_*R_*)(H_*\tau_s)$。声源寿命受膨胀或激波形成的限制，常估为 $\tau_s\lesssim\min(H_*^{-1},R_*/\overline U_f)$，$\overline U_f$ 为流体均方根速度。直接把声波一直算到一个 Hubble 时间会在短寿命情况下高估信号。

区分这三项还影响峰形。壁碰撞主要反映过境时间与碰撞几何；流体的纵向声波有声壳厚度与泡间距两种长度，足够高频时可分辨其厚度；涡旋与磁场则按湍流谱在一段惯性区内分配剪切能。常用一个归一在峰上的形状函数表示有限带宽的模拟结果，例如声源经验式 $S_{\rm sw}(x)=x^3[7/(4+3x^2)]^{7/2}$、$x=f/f_{\rm sw}$。对数导数为 $3-21x^2/(4+3x^2)$，故 $x=1$ 是极值，低频 $x^3$、高频 $x^{-4}$。此式的 $x^{-4}$ 是一类模型的近似；不同壁速、声壳及源寿命可能改变高频结构，不是从因果性推出的指数。

作一数量级估算。取 $T_*=100\,{\rm GeV}$、$\beta/H_*=100$、$v_w$ 近于一，若源处循环峰频约 $\beta$，则观测峰约 $10^{-3}\,{\rm Hz}$；若给出峰的角频率约 $\beta$，则要再除以 $2\pi$。取 $\alpha=0.1,\ \kappa_X=0.1$，得到 $K_X\simeq9\times10^{-3}$。若 $R_*\sim10^{-2}H_*^{-1}$，短时估计给 $\Omega_{{\rm gw},*}\sim10^{-8}$，声源维持 $0.1H_*^{-1}$ 则可多一个约十的因子。将这些产生时的份额乘式 \eqref{eq:gw-d-gstar} 的辐射份额及熵因子之后，才是今天临界密度归一的谱；直接将 $\Omega_{{\rm gw},*}$ 同干涉仪的 $\Omega_{{\rm gw},0}$ 曲线比较会差若干数量级。


\subsection{流过移动界面的粒子}

壁速度决定的既是谱峰频率，又是能量在标量壁和等离子体之间的分配。沿壁法线建坐标 $z$，取 $z\to-\infty$ 为真相、$z\to+\infty$ 为假相。在壁静系中，未扰动流体以 $u^z=-\gamma v$ 流向壁。若温度横跨壁近于不变，给标量方程加一个线性耗散项 $\Box\phi-V_{,\phi}=\eta_w u^\mu\partial_\mu\phi$。静态平面壁方程为 $\phi''-V_{,\phi}=-\eta_w\gamma v\phi'$；这里两撇是法向空间导数。乘以 $\phi'$ 沿 $z$ 积分，动能边界项消失，得
\[
 p_{\rm drive}=V(\phi_f,T)-V(\phi_t,T)
 =\int_{-\infty}^{\infty}\eta_w\gamma v\,(\partial_z\phi)^2\,\mathrm dz .
\]
右端在低速时近似正比壁速，$\eta_w$ 由热浴微观相互作用决定。如果壁两侧的温度显著不同，$\int V_{,\phi}\phi'\mathrm dz=[V]-\int V_{,T}T'\mathrm dz$ 中最后一项亦要保留；把两侧不同温度的势直接相减当作机械推力会漏掉流体的热力学功。上式只是将热浴过程浓缩成一个速度阻力的低速模型，不能把它的 $\gamma v$ 外推到任意高速。

较微观地，场依赖质量 $m_i(\phi)$ 的粒子对壁方程贡献为 $g_i(\partial_\phi m_i^2)\int f_i\,\mathrm d^3p/[(2\pi)^3\,2E_i]$。局部平衡的 $f_i^{(0)}$ 已可写成 $\partial_\phi\Delta V_i$，故在使用热有效势的方程中，只应再放 $\delta f_i=f_i-f_i^{(0)}$ 的偏离。用 Boltzmann 方程的弛豫时间近似可看到阻力的符号：流体经历不断变化的 $m_i(\phi(z))$，若碰撞未能瞬时调整占有数，粒子分布对壁的作用会滞后于平衡值。其灵敏度含 $(\partial_\phi m_i^2)^2$、壁的厚度、散射时间及热权重；较强耦合可能同时改变势垒和摩擦，故从 $V(\phi,T)$ 本身不能唯一读出 $v_w$。

对极快而薄的平面壁，一颗粒子穿越时可近似保持能量及切向动量。若从假相到真相其质量平方增 $\Delta m_i^2$，则初末法向动量大小之差为 $|p_{z,f}|-|p_{z,t}|\simeq\Delta m_i^2/(2|p_{z,f}|)$。入射通量还乘 $|p_z|/E$，两个 $|p_z|$ 抵消，故该过程的阻力贡献近于
\[
 p_{{\rm part},i}\simeq g_i\Delta m_i^2
      \int\frac{\mathrm d^3p}{(2\pi)^3}\frac{f_i^{\rm in}}{2E_i}.
\]
超相对论壁前的热粒子几乎都从同一侧入射，$\mathrm d^3p/E$ 与分布函数组合可在上游流体静系求值。若上游粒子近于无质量，此积分是玻色子的 $T^2/24$、费米子的 $T^2/48$，同热自由能的二次项一致。于是只计这种穿越动量转移时的高速阻力趋于有限值。用裸势差减去这里的粒子动量转移，或用热势加非平衡修正，是两套记帐方法，不可把平衡项扣两次。若单粒子高速阻力已足以抵消驱动，则壁可达到终端速度；若不够，是否持续加速还取决于高能辐射等其它动量转移渠道。因而“壁以近光速运动”本身不足以断定能量主要留在壁中。

例如忽略膨胀，单个球泡很大后壁能量为 $E_w=4\pi R^2\sigma\gamma_w$，释放的自由能约为 $4\pi R^3\Delta p/3$。若 $\gamma_w$ 已稳定，两者之比 $3\sigma\gamma_w/(R\Delta p)$ 随 $R^{-1}$ 降低，转变越晚，标量壁份额越小。若有近似恒定的正净压 $p_{\rm net}$ 加速壁，则 $\mathrm dE_w/\mathrm dR\simeq4\pi R^2p_{\rm net}$，大 $R$ 时 $\gamma_w\simeq R p_{\rm net}/(3\sigma)$，能量比例可保持有限。将壁碰撞功率估的是 $E_w$ 中的可辐射部分；把全部释放热量都代入壁源会破坏这个随 $R$ 的标度关系。

\subsection{袋模型中的壁流速}

可用两相的辐射袋模型把通量匹配做到底。令壁前后能量、压强为
\[
 \rho_+=a_+T_+^4+\epsilon,\quad p_+=a_+T_+^4/3-\epsilon,\qquad
 \rho_-=a_-T_-^4,\quad p_-=a_-T_-^4/3 .
\]
两侧各自的热浴声速皆 $1/\sqrt3$，$\epsilon$ 是常真空能差。此处定义 $\alpha_+=\epsilon/(a_+T_+^4)$、$r_b=a_-T_-^4/(a_+T_+^4)$，其中 $\alpha_+$ 不等于任意定义下的 $\Delta\rho/\rho_r$。上面的跃迁关系化为
\[
 v_+v_-=\frac{(1-r_b)/3-\alpha_+}{1-r_b+\alpha_+},
 \qquad
 \frac{v_+}{v_-}=\frac{r_b+1/3-\alpha_+}{1+\alpha_++r_b/3}.
\]
从第一式求 $1-r_b=\alpha_+(1+v_+v_-)/(1/3-v_+v_-)$（分母为零须回原式），代第二式并收集 $v_+$ 的各次幂，得
\[
 (1+\alpha_+)v_+^2
 -\left(v_-+\frac1{3v_-}\right)v_+
 +\frac13-\alpha_+=0 .
\]
故若记 $A=v_-/2+1/(6v_-)$，有 $v_+=[A\pm\sqrt{A^2+\alpha_+^2+2\alpha_+/3-1/3}]/(1+\alpha_+)$。需选满足各方温度、熵流和远端边界的物理根。以 $v_-=1/\sqrt3$ 作声速离壁的例子，两根成 $(1/\sqrt3\pm\sqrt{\alpha_+^2+2\alpha_+/3})/(1+\alpha_+)$，上根给 Jouguet 爆轰。$\alpha_+\to0$ 时它趋声速；$\alpha_+=0.1$ 时约为 $0.78$。此结果是该边界条件下的一支解。若壁前流体已被冲击波加热或运动，$T_+$ 和 $v_+$ 均非远处背景量；尤其不能把算出的 $v_+$ 自动认成普遍的壁速。

流体动能的效率也可按剖面积分。以球对称自相似变量 $\xi=r/t$ 表示单泡诱导的径向流 $v(\xi)$，用未扰动流体静系定义体积。$\rho_{\rm lab}=w_{\!h}\gamma^2-p$，扣去本地静能量 $\rho=w_{\!h}-p$ 得 $w_{\!h}\gamma^2v^2$。若壁半径 $v_w t$，流入集体运动的能量相对同体积释放能量的比为
\[
 \kappa_v=\frac{3}{v_w^3\,\epsilon}
  \int_0^\infty \xi^2\,w_{\!h}(\xi)\gamma^2(\xi)v^2(\xi)\,\mathrm d\xi .
\]
激波或声壳的支撑可延至壁以外，所以上限不能一律截在 $v_w$。$\gamma^2v^2$ 的量级及非零区间的宽度同时影响效率。在小速度声壳中，若近似 $w_{\!h}$ 和 $v$ 在厚度 $\Delta R$ 的壳内恒定，则动能正比 $R^2\Delta R\,v^2$；与同一个 $R$ 但更薄的壳相比，前者发声功率可以更大。给定 $v_w,\alpha_+$ 后，还要提供流体剖面或经校准的 $\kappa_v$，才能将微观势接到引力波振幅上。

\subsection{两列纵向声波也能产生横向引力波}

还可把温度势、壁摩擦和波源放进同一守恒记帐。为看清力的方向，先把温度无关的部分 $V_0(\phi)$ 计入标量场的能动张量，把热浴的自由能和压强计入流体。局域近 Minkowski 度规上有 $\partial_\mu T_\phi^{\mu\nu}=(\Box\phi-V_{0,\phi})\partial^\nu\phi$。若将摩擦对场的作用记 $J^\nu=\eta_w(u\cdot\partial\phi)\partial^\nu\phi$，同一项须以负号出现在流体方程 $\partial_\mu T_{\rm fl}^{\mu\nu}=-J^\nu$ 中（由各相自由能分配产生的可逆项另以相反号成对计入）。与 $u_\nu$ 缩并，摩擦给流体静系的热化功率 $u_\nu J^\nu=\eta_w(u\cdot\partial\phi)^2\ge0$。故一方损失的机械能在另一方中有明确去向。数值演化场与热浴时应使用同一能量交换项；若只给场加阻尼、再让独立流体无源自由膨胀，就会丢失本应激起声波的热量与动量。

总应力在局部正交标架中的 TT 源只需
\[
 \Pi_{ij}^{\rm TT}=\Lambda_{ij,lm}
 \big[\partial_l\phi\,\partial_m\phi
       +w_{\!h}\gamma^2 v_l v_m\big].
\]
压强项随局域温度大幅变化仍正比 $\delta_{lm}$，投影后为零；理想流体的速度却以二次项耦合。若只在一阶流体方程中寻找一个与 $h_{ij}$ 同螺旋度的声波本征模，自然会漏掉这一源。为说明生成条件，在均匀介质中取两列小振幅声波，物理波矢为 $\bm p,\bm q$，速度分别沿两波矢，角频率为 $\omega_p=c_s p$、$\omega_q=c_s q$。速度乘积的和频项有 $\bm k=\bm p+\bm q,\qquad \omega=\omega_p+\omega_q=c_s(p+q)$。一只自由引力波的色散关系是 $\omega=k$。若两列声波长度相同 $p=q$，设方向夹角为 $\vartheta$，则 $k=2p\cos(\vartheta/2)$，相干驱动要求 $\cos(\vartheta/2)=c_s$。辐射流体取 $c_s=1/\sqrt3$，给 $\cos\vartheta=-1/3$。此时两速度相对 $\bm k$ 皆有横向分量，剪切投影可非零。若仅一列平面声波与自身相乘，$\bm k$ 平行速度，TT 投影就彻底消失。差频支 $\omega=c_s|p-q|$ 却总小于等于 $c_s|\bm p-\bm q|$，对 $0<c_s<1$ 不可能在非零 $\bm k$ 精确满足自由波色散。

实际声壳在一段波数上有功率、持续有限时间，故上述共振条件有宽度。设驱动含 $e^{-\mathrm i\omega t}$、在 $0<t<D$ 内持续，式 \eqref{eq:gw-d-sourcegreen} 中相位匹配的积分平方为 $4\sin^2[(k-\omega)D/2]/(k-\omega)^2$。$D$ 很大时单模在共振点的峰高按 $D^2$，宽度按 $D^{-1}$，对连续模式积分后的总面积按 $D$ 增长，因 $\int_{-\infty}^\infty4\sin^2(x D/2)\mathrm dx/x^2=2\pi D$。这说明声波寿命延长可提高总波功率，却不应给一个连续、随机源的整个频带统一乘 $D^2$。冲击波使相干声能转成热和涡运动后，纵波共振这一累积机制随之结束。

\subsection{涡流和磁场的长度分布}

相变产生的流动若具有很大的 Reynolds 数 $\mathrm{Re}=U R_*/\nu_{\rm kin}$，黏性不能在一个大尺度涡旋的翻转时间内将其耗尽。能量可从驱动长度 $R_*$ 向小尺度传递，直到 $\mathrm{Re}$ 降到一附近。为了估量各尺度的影响，设单位质量、单位时间流过惯性区的动能为 $\varepsilon$，长度 $l$ 上的速度差为 $v_l$。若局域的传递时间 $\tau_l\sim l/v_l$，通量 $v_l^2/\tau_l$ 近似与 $l$ 无关，便有 $v_l\sim(\varepsilon l)^{1/3}$、$\tau_l\sim l^{2/3}\varepsilon^{-1/3}$。定义一维速度谱 $\int E_v(k)\mathrm dk=\langle v^2\rangle/2$，则 $kE_v(k)\sim v_l^2$ 给 $E_v(k)\propto\varepsilon^{2/3}k^{-5/3}$。若改用三维 Fourier 协方差标记 $P_v(k)$，其相空间还带 $k^2$，同一结果写成 $P_v\propto k^{-11/3}$。两种指数不可混用。

在高导电等离子体中，磁场也会给出剪切。以有理化单位、流体静系中的电场略去，磁应力为 $B^2\delta_{ij}/2-B_iB_j$，故波源中还有 $-\Lambda_{ij,lm}B_lB_m$。动能、磁能即使按同一惯性区幂律分配，波谱也不等于 $E_v(k)$ 或 $E_B(k)$：需先算二次应力的卷积，对每个尺度放入它自己的相关时间，再代入式 \eqref{eq:gw-d-uetc}。大涡往往存活较久，却对应较小 $k$；小涡速度较小且很快去相。$R_*$ 附近的峰来自这些量与波色散的共同选择。若两种场相关，还要保留 $\langle(v_iv_j)(B_lB_m)\rangle$ 的交叉项。只知道总能量中某一固定比例转成湍流，还不足以从第一原理预言峰的宽度和高频斜率。
\section{宇宙弦与背景的积分限制}\label{sec:gw-dstrings}

\subsection{线缺陷的能量集中在哪里}

若复标量在相变后满足 $|\phi|=v$ 而相位任意，真空流形是一圆周。环绕一条轴方位角从 $0$ 变到 $2\pi$，相位可绕 $2\pi n$，单值性要求 $n$ 为整数。只要大圈上始终有 $|\phi|=v$，连续变形不能改变绕数。故大圈张成的面内必有一点 $\phi=0$，沿空间延伸即成线缺陷。绕直线轴可取 $\phi=vF(r)e^{\mathrm i n\varphi}$，$F(0)=0,\ F(\infty)=1$。小半径区由势能和径向梯度决定芯宽，远处仍有角梯度 $|\nabla\phi|^2\simeq n^2v^2/r^2$。于是对\textbf{整体弦}，张力含 $\int_\delta^Rr\,\mathrm dr/r^2=\ln(R/\delta)$，依赖红外截断 $R$（如相邻弦距离）。它还耦合到传播的 Goldstone 模式；把全部失能归给引力辐射常常不适用。

若对称性为局域 $U(1)$，$D_\mu\phi=(\partial_\mu-\mathrm i e A_\mu)\phi$。在芯外 $A_{\hat\varphi}\simeq n/(er)$ 可把相位梯度抵消，但总磁通 $\oint\bm A\cdot\mathrm d\bm l=2\pi n/e$ 留在芯中。标量和矢量的质量给两种芯宽，远处场强指数衰减，每单位长度能量 $\mu$ 有限。若观测波长及弦的曲率半径皆远大于芯宽，才可将它近于零厚、有恒定张力的相对论弦。此时能量与张力量纲相同，引力强度由无量纲 $G\mu$ 表示。稳直、静止的局域弦有线能密度等于张力，Newton 意义的横向吸引可相互抵消，外部几何却有与 $G\mu$ 正比的锥缺角；弯曲且运动的弦圈的时变剪切则可辐射。


\subsection{局域弦芯的一个可积极限}

薄弦的 $\mu$ 可以由场论芯结构定标。取 Abelian Higgs 模型的横截面能量为
\[
 \mu=\int\mathrm d^2x\left[|D_x\phi|^2+|D_y\phi|^2
 +\frac{B^2}{2}+\frac{\lambda}{2}(|\phi|^2-v^2)^2\right].
\]
在临界耦合 $\lambda=e^2$ 且 $n>0$ 时，可利用 $[D_x,D_y]\phi=-\mathrm i eB\phi$，分部积分后配方为
\[
 \mu=\int\mathrm d^2x\left\{|(D_x+\mathrm iD_y)\phi|^2+
 \frac12[B-e(v^2-|\phi|^2)]^2\right\}
 +ev^2\int B\,\mathrm d^2x .
\]
有限能量的边界条件令横截面无穷远的全导数不贡献。最后一项由磁通恰等于 $2\pi n v^2$，故 $\mu\ge2\pi n v^2$。若两个平方均为零，就达到这个下界；负绕数选相反的符号。这个结果固定的是给定规范化及临界耦合下的张力，不能无条件套在所有拓扑线缺陷上。

取极坐标 $\phi=vF(r)e^{\mathrm in\varphi}$、$A_{\hat\varphi}=n A(r)/(er)$，便有 $B=nA'/(er)$。上述平方为零等价于 $F'=n(1-A)F/r$、$A'=e^2v^2r(1-F^2)/n$。近芯处 $F\sim r^n$、$A\sim e^2v^2r^2/(2n)$；芯外令 $F=1-f$、$A=1-b$，线性化给 $f''+f'/r-2e^2v^2 f=0$，衰减支为 $K_0(\sqrt2 e v r)$。可见局域弦的张力收敛来自规范场对角梯度的屏蔽。若改为整体弦，类似的 $b$ 不存在，远端长程相位场正是早先 $\ln(R/\delta)$ 的来源。观测频率对应的弦上波长接近芯宽时，Nambu--Goto 的频谱将受到这些具体场自由度的修正。
\subsection{世界面上的左右两列波}

以两个参数 $\sigma^a$ 描述弦在时空中的世界面 $X^\mu(\sigma^a)$，诱导度规为 $\gamma_{ab}=g_{\mu\nu}\partial_aX^\mu\partial_bX^\nu$。最低阶有效作用量只有世界面面积，$S_{\rm NG}=-\mu\int\mathrm d^2\sigma\sqrt{-\det\gamma}$。它对世界面重新标记不变。变分 $X^\mu$，得到
\[
 \frac{1}{\sqrt{-\gamma}}\partial_a
 (\sqrt{-\gamma}\gamma^{ab}\partial_bX^\mu)
 +\Gamma^\mu{}_{\nu\rho}\gamma^{ab}\partial_a X^\nu\partial_bX^\rho=0.
\]
世界面张力试图使面积驻定。一个远小于 Hubble 半径的弦圈在一周期内可用 Minkowski 背景。取时间规范 $X^0=t$，并在世界面选正交、等共形长度标记，使
\[
 \dot{\bm X}\cdot\bm X'=0,\qquad
 \dot{\bm X}^{\,2}+\bm X'^{\,2}=1,\qquad
 \ddot{\bm X}-\bm X''=0 .
\]
这里点、撇暂指对 $t,\sigma$ 求导。能量元为 $\mu\,\mathrm d\sigma$，故 $\sigma$ 从 $0$ 到不变长度 $L$ 给圈能量 $\mu L$，而某一时刻的几何弧长为 $\int|\bm X'|\mathrm d\sigma$，二者并不相等。一般解可写
\[
 \bm X=\tfrac12[\bm a(\sigma-t)+\bm b(\sigma+t)],
 \qquad |\bm a'|=|\bm b'|=1 .
\]
由每个切矢都在单位球上，约束在左右传播中自动保持。闭合圈在质心静系中可取 $\bm a,\bm b$ 皆 $L$ 周期，平均导数各为零；$\bm X(t+L/2,\sigma+L/2)=\bm X(t,\sigma)$，因而一般物理图形以 $L/2$ 为周期，谐频是 $f_{{\rm e},n}=2n/L$。

在两条单位切线曲线 $\bm a'$ 与 $-\bm b'$ 的交点，$\bm X'=0$、$|\dot{\bm X}|=1$，弦瞬间形成\textbf{尖点}(cusp)。实际芯宽和辐射回馈会使最高 Lorentz 因子有限。若某一行波的切矢有跳变，则折痕以单位速度沿弦运动，称\textbf{折点}(kink)。两段弦相遇并重新连结时易产生折点。并非每个圈每个周期必有同样数目的尖点，左右单位曲线的形状决定它们是否相交；高频谐波的权重因而依赖圈的历史。

\subsection{高次谐波来自何处}

变分背景度规得薄弦能动张量，在上述规范中为
\[
 T^{\mu\nu}(t,\bm x)
 =\mu\int_0^L
 [\dot X^\mu\dot X^\nu-X'^\mu X'^\nu]\,
 \delta^3(\bm x-\bm X(t,\sigma))\,\mathrm d\sigma .
\]
用 $\dot X^0=1,X'^0=0$ 即得 $\int T^{00}\mathrm d^3x=\mu L$。把 $\bm X$ 的左右行波代入 $\int T^{ij}\exp[\mathrm i\omega(t-\hat{\bm n}\cdot\bm X)]$，以 $\sigma\pm t$ 为自变量，积分核及相位可分成两个一维积分的对称乘积。空间方向的横向投影须在最后做，沿观测方向的领头切矢会被它消掉。对一束指向尖点速度的辐射，左右两个相位都具有驻点；一般方向却无共同驻点，高谐波因快速振荡而相消。

具体看一边。设沿 $\hat{\bm n}$ 的单位切线为 $\bm b'(s_0)$，令 $s=s_0+u$。相位含 $s-\hat{\bm n}\cdot\bm b(s)$，一次项因 $\bm b'=\hat{\bm n}$ 为零；由 $\bm b'^2=1$ 得 $\bm b'\cdot\bm b''=0$，二次项亦零。再求一次导数得 $\bm b'\cdot\bm b'''=-|\bm b''|^2$，故首个非零项正比于 $u^3|\bm b''|^2$。若曲率半径约 $L$，有效积分区为 $\Delta u\sim(L^2/\omega)^{1/3}$。切矢的常数部分纯纵向，首个横向项 $\bm b''u$，故每个积分的横向量级为 $(\Delta u)^2/L\sim L^{1/3}\omega^{-2/3}$。左右两边相乘，远处普通 Fourier 应变的高频包络即
\[
 |\widetilde h_{\rm cusp}(f)|\sim
 \frac{G\mu L^{2/3}}{D\, f^{4/3}}
 \qquad(D\ll H_0^{-1},\quad fL\gg1).
\]
这里 Fourier 应变量纲是时间；有时画的是 $f|\widetilde h|$，频率指数便为 $-1/3$。若与尖点方向夹角 $\vartheta$，相位另有 $\omega\vartheta^2 u$ 的一次项。在 $\Delta u$ 内小于一的条件给锥半角 $\vartheta_m\sim(fL)^{-1/3}$。折点的某一边改为切矢的不连续，分部积分给端点幅度 $\sim f^{-1}$，另一边仍可有驻点，故单个折点爆发常有 $|\widetilde h|\propto f^{-5/3}$，辐射较呈扇形。以上只描述可分辨弦结构而尚看不到芯宽的区间。

圈的总功率须用相对论源的波公式求。量纲上只有 $G\mu^2$，故写第 $n$ 个谐波的功率为 $P_n=\Gamma_nG\mu^2$、$\Gamma=\sum_n\Gamma_n$。光滑生尖点的理想圈高次谐波常有 $\Gamma_n\propto n^{-4/3}$；这可由场的 $f^{-4/3}$、能流的 $f^2$ 及束开角平方 $f^{-2/3}$ 相乘验证。高频回馈会截断幂律，总和有限。由 $\mathrm d(\mu L)/\mathrm dt=-\Gamma G\mu^2$ 得 $L(t)=L_i-\Gamma G\mu(t-t_i),\qquad \tau_{\rm loop}=\frac{L_i}{\Gamma G\mu}$。若 $\Gamma G\mu\ll1$，即使宏观上圈会消失，它仍可在长度显著变化前重复振荡很多次。某些特殊对称圈的某些 $\Gamma_n$ 可为零；已知总 $\Gamma$ 不能唯一决定频谱。


\subsection{圆圈与一般圈的区别}

取一个特别的平面闭圈以检验世界面约束。设 $R_0=L/(2\pi)$，以 $\bm e_r(\varphi)=(\cos\varphi,\sin\varphi,0)$ 表示平面内的径向单位矢。令 $\bm X=R_0\cos(t/R_0)\bm e_r(\sigma/R_0)$。$\dot{\bm X}$ 径向、$\bm X'$ 切向，点积为零；它们模平方的和是 $\sin^2(t/R_0)+\cos^2(t/R_0)=1$。固有半径 $R(t)=R_0|\cos(t/R_0)|$，弦芯长度收缩时动能增大，而总能量 $\mu L$ 不变。对一小段，瞬时弧长 $\mathrm dl=|\bm X'|\mathrm d\sigma$，$\gamma=(1-\dot{\bm X}^2)^{-1/2}=1/|\bm X'|$，故 $\mu\gamma\,\mathrm dl=\mu\,\mathrm d\sigma$。这个恒等式说明从观测时的几何弧长估计圈能量，必须同时计运动的 Lorentz 因子。

若暂用最低辐射矩估量，圆圈的质量二阶矩是 $\int T^{00}x_ix_j\mathrm d^3x=(\mu L R^2/2)\operatorname{diag}(1,1,0)$，去迹后给 $\mu L R^2\operatorname{diag}(1,1,-2)/6$。因 $R^2=R_0^2[1+\cos(2t/R_0)]/2$，其时间变化带最低谐频 $f=1/(\pi R_0)=2/L$，与用左右行波得到的周期相同。三阶导量级 $\mu LR_0^2/R_0^3\sim\mu$，故 $G\dddot Q^2$ 的量纲给 $G\mu^2$。但当 $R$ 接近零时整个圈都接近光速，它是高度非一般的同步收缩解；不能以非相对论四极式求其精确 $\Gamma$，更不能把整圈的同步事件当成一般圈的孤立尖点。在有限厚度、自引力或自交显著的时段，薄弦平直背景解自身也会失效。
\subsection{从弦网计数到今天的频谱}

若宇宙膨胀时长弦只随坐标网格拉长且不相交，能量密度约按 $a^{-2}$ 衰减，迟早会干扰背景。实际长弦互换端点、切出小圈并辐射，可能趋于\textbf{标度解}：每个视界尺度内保持近似一定数目的长弦，特征长度 $\xi\sim t$，故 $\rho_\infty\sim\mu/\xi^2\sim\mu/t^2$，与 $\rho_{\rm tot}\sim1/(Gt^2)$ 的比为 $O(G\mu)$。这只约束整体能量流，圈在什么尺寸生成还须有分布模型。以 $\mathcal F(L_i,t_i)$ 记每物理体积、时间、初始长度的圈出生率，忽略圈间再作用，某时刻每单位长度的物理数密度是
\begin{equation}\label{eq:gw-d-loopnumber}
 n(L,t)=\int^t
 \left(\frac{a_i}{a}\right)^3
 \mathcal F\big(L+\Gamma G\mu(t-t_i),t_i\big)\,\mathrm dt_i .
\end{equation}
积分里用到 $\partial L_i/\partial L=1$，只计还活着的 $L>0$。若作单尺度产生的例子 $\mathcal F=C\,t_i^{-4}\delta(L_i-\alpha_\ell t_i)$，有 $t_i=(L+\Gamma G\mu t)/(\alpha_\ell+\Gamma G\mu)$，故
\[
 n(L,t)=\frac{C}{(\alpha_\ell+\Gamma G\mu)t_i^4}
 \left(\frac{a_i}{a}\right)^3,\qquad
 0<L<\alpha_\ell t ,
\]
还须 $t_i$ 晚于成网时刻。$\alpha_\ell\gg\Gamma G\mu$ 给长寿命圈，同一观测时刻的圈既可能刚生于物质时代，也可能是辐射时代遗留，不能以单一时代的标度统统替代。

在时刻 $t$、物理长度 $L$ 的圈以 $2n/L$ 发射，今天频率是 $2na(t)/L$。直接将每单位体积注入的功率乘 $\mathrm dt$ 并按 $a^4$ 红移，转到每对数频率有
\begin{align}
 \rho_{c0}\Omega_{{\rm gw},0}(f)
 &=\int \mathrm dt\,a^4(t)
   \sum_{n\ge1} P_n\int_0^\infty\mathrm dL\,n(L,t)
   \delta\left(\ln f-\ln\frac{2na(t)}{L}\right)\nonumber\\
 &=\frac{2G\mu^2}{f}\sum_{n\ge1}n\Gamma_n
   \int\mathrm dt\,a^5(t)\,
    n\left(\frac{2na(t)}{f},t\right).
 \label{eq:gw-d-loopomega}
\end{align}
第二式的一个 $a$ 来自 $\delta(\ln f-\ln f_n)=L_n\delta(L-L_n)$，$L_n=2na/f$。若忽略这一 Jacobian 会同时算错辐射时代权重及 $f$ 的幂。图的高频端可收到早期小圈和大圈的高谐波，低频端通常对晚期的大圈敏感。重联概率较低会改变长弦密度、圈形成率与小尺度结构，不能仅在单圈的张力中吸收。

还可检查为何辐射时代容易给平台。若 $a=A_rt^{1/2}$ 且 $n(L,t)=t^{-4}N(L/t)$，式 \eqref{eq:gw-d-loopomega} 中 $a^5 n\propto t^{-3/2}N[2nA_r/(f\sqrt t)]$。令 $x=2nA_r/(f\sqrt t)$，则 $t^{-3/2}\mathrm dt=-f\,\mathrm dx/(nA_r)$，显式 $f/n$ 恰与前面的 $n/f$ 抵消。若有效积分区覆盖 $N(x)$ 的主要支撑，每个谐波都是近乎平坦的谱。转入物质主导、热自由度变化、圈尚未形成或谐波截止时，边界或权重不再标度，平台便会弯曲。这个结论来自源的连续产生和红移平衡，同暴涨谱平坦的物理原因不同。

对独立尖点爆发，设每周期平均有 $N_c$ 个，则时刻 $t$ 每物理体积每固有时的爆发率是 $\int n(L,t)(2N_c/L)\,\mathrm dL$。要变成观者每单位时间的率，乘红移壳的物理体积 $4\pi S_K^2(\chi)\,\mathrm dz/[H(z)(1+z)^3]$、时间膨胀 $(1+z)^{-1}$ 以及指向观者的概率 $(1-\cos\vartheta_m)/2\simeq\vartheta_m^2/4$；在束宽内以 $(1+z)f$ 算源频。非常强而罕见的爆发对形式上的二阶矩贡献大，却未必在一次有限观测中出现。若搜寻平稳背景，应检查某强度以上爆发的预期次数或占空比，并给可单独分辨的事件单独计数。

\subsection{另一种泵浦场与蓝谱}

标准慢滚张量谱稍向低频倾斜。一种能使谱走向相反的机制是张量的动能系数随时间迅速改变。设每偏振二次作用量正比于 $\int z_t^2[(h')^2-(\nabla h)^2]$，则 $v=z_th$ 满足 $v_k''+[k^2-z_t''/z_t]v_k=0$。若 $z_t\propto|\eta|^p$，有效势为 $p(p-1)/\eta^2$，Hankel 阶数 $\nu=|p-1/2|$。在这一阶段起自绝热真空的波，于某一固定时刻处在长波极限时有 $\mathcal P_t\propto k^{3-2\nu}$；$\nu=0$ 时 Hankel 函数还带对数。在随后的匹配不额外引入这段频带内的急剧尺度依赖时，辐射时代入内的能谱保留这个斜率。

例如在前大爆炸的弦宇宙学模型中，低能作用量的曲率项带 $e^{-\varphi}$，故在弦标架中 $z_t\propto a e^{-\varphi/2}$。设低曲率时 $\varphi$（dilaton）随弦耦合增长，三维各向同性膨胀支的幂解（忽略原点平移）可写作 $a\propto(-t)^{-1/\sqrt3}$、$\varphi=-(1+\sqrt3)\ln(-t)+{\rm const.}$，$t<0$。用 $\mathrm d\eta=\mathrm dt/a$ 得 $-\eta\propto(-t)^{1+1/\sqrt3}$，代 $z_t$ 的幂，恰有 $z_t\propto(-\eta)^{1/2}$。故 $\nu=0$，低频谱约随 $f^3$，可伴对数修正。在一种参数模型中，高曲率阶段取 $a\propto(-\eta)^{-1}$、$\varphi=-2b\ln(-\eta)+{\rm const.}$，则 $p=b-1$，相应较高频段随 $f^{3-2|b-3/2|}$。曲折处与终止频率分别反映进、出高曲率阶段的 $aH$，斜率与幅度依赖这段未知的高曲率动力学。这个算例说明不能将从微波背景测得的慢滚红倾斜当成所有极早期背景的先验。


\subsection{从非瞬时冻出求轻元素的响应}

第 3 章用冻结温度估算了 $n/p$ 对引力变化的敏感度。对一个随时间开启或截止的引力波背景，单取 $T_f$ 可能丢掉反应持续一段时间的影响。设 $X_n=n_n/(n_n+n_p)$，将相反转换的总速率写为 $\lambda=\lambda_{np}+\lambda_{pn}$，相应平衡值为 $X_{\rm eq}=\lambda_{pn}/\lambda$。在原子核尚未形成、忽略重子数变化时，动力学方程是 $\dot X_n=-\lambda(X_n-X_{\rm eq})$。写 $x=\ln a$、$L(x)=\lambda/H$。若早期完全平衡，其后精确的滞后量为
\[
 X_n(x)-X_{\rm eq}(x)=-\int_{x_i}^{x}
 X_{\rm eq}'(u)\exp\!\left[-\int_u^x L(v)\,\mathrm dv\right]\mathrm du ,
\]
此处撇指对 $x$。当 $\lambda/H\gg1$，早出平衡的偏差立刻被后续反应抹掉；$\lambda/H\ll1$ 后指数趋近一，剩余组分大致记住此前的丰度。自由中子衰变可作为 $\lambda_{np}$ 的一项计入。若将它放在冻结后的另一个阶段单独近似，便应从此处的晚时 $\lambda$ 中拿掉，避免重复计时。

现在仅以额外辐射改动 $H$，保持在给定 $a$ 的温度、热浴反应率与 $X_{\rm eq}$ 不变。设 $H\to H(1+\delta_H)$，一阶有 $\delta L=-L\delta_H$。在未扰动轨迹 $X_n^{(0)}$ 周围变分，得
\[
 \delta X_n(x)=\int_{x_i}^x\!\mathrm du\,
 L(u)[X_n^{(0)}(u)-X_{\rm eq}(u)]\,\delta_H(u)
 \exp\!\left[-\int_u^xL(v)\,\mathrm dv\right].
\]
这给出了一个明确的时间权重。平衡前括号近零，即使当时已多了辐射，反应仍会抹去短暂的改变；完全冻出后 $L$ 很小，弱过程也几乎不再改变比例。通常 $X_{\rm eq}$ 随冷却下降，实际 $X_n^{(0)}$ 比它大，故正的额外膨胀使末态中子增加。能量份额若随某些张量模陆续入内而改变，便应把相应 $\delta_H(u)$ 放进积分；恒定 $\Delta N_{\rm eff}$ 的粗估相当于它在敏感区变化缓慢。

在电子、正电子尚相对论性且各轻粒子近于同温的时间，标准瞬时模型给 $g_\rho=2+(7/8)(4+6)=43/4$。以一个中微子种类及其反粒子的能量为单位，$\Delta g_\rho=(7/4)\Delta N_{\rm eff}$。故在固定 $T$ 处，
\[
 \frac{H^2}{H_{\rm std}^2}=1+\frac7{43}\Delta N_{\rm eff},
 \qquad
 \delta_H\simeq\frac7{86}\Delta N_{\rm eff}.
\]
电子湮灭后则要用式 \eqref{eq:gw-d-neff} 的较低中微子温度，不能在整段核合成直接使用 $7/43$。作为检验，以尖锐冻结近似和 $\lambda\propto T^5$ 代入 $\lambda=H$，在自由度缓变时 $\delta\ln T_f\simeq\delta_H/3$。$X_{\rm eq}\simeq[1+\exp(Q/T)]^{-1}$、$Q=m_n-m_p$ 给 $\delta X_{\rm eq}=X_{\rm eq}(1-X_{\rm eq})(Q/T_f)\delta\ln T_f$，与上面宽权重的正号一致。后续若以 $t_d$ 为能稳定烧成氦的时刻，尚有 $X_n(t_d)\simeq X_n(t_f)e^{-(t_d-t_f)/\tau_n}$。额外辐射还使固定温度区间所花时间缩短，这部分由 $\delta(t_d-t_f)=-\int_{x_f}^{x_d}\delta_H\,\mathrm dx/H$ 计算（这项表达式先固定两端，端温移动另计）。因而氦的响应同时含转换率时间权重与中子寿命时钟；组分数界是对这两种热史效应的联合约束。
\subsection{能量限制究竟限制哪段频率}

远短于 Hubble 半径、随机各向同性、自由传播的波会像额外辐射那样加速膨胀。第 3 章已讨论轻元素丰度怎样测量膨胀率。对遗迹张量谱，应用式 \eqref{eq:gw-d-neff} 还要把实测的额外辐射单位换成对数频率积分。令 $h_0=H_0/(100\,{\rm km\,s^{-1}\,Mpc^{-1}})$，今天的黑体光子给 $\Omega_{\gamma0}h_0^2\simeq2.47\times10^{-5}(T_0/2.725\,{\rm K})^4$。故可用的额外辐射若为 $\Delta N_{\rm eff}^{\max}$，有
\begin{equation}\label{eq:gw-d-integral}
 \int_{f_{\rm min}}^{f_{\rm max}}\mathrm d\ln f\,h_0^2\Omega_{{\rm gw},0}(f)
 \ \lesssim\
 5.6\times10^{-6}\left(\frac{T_0}{2.725\,{\rm K}}\right)^4
 \Delta N_{\rm eff}^{\max}.
\end{equation}
右边只有单位换算系数固定，允许的 $\Delta N_{\rm eff}$ 来自所选资料与模型。下限 $f_{\rm min}$ 应按检验时刻取约 $aH/(2\pi)$；在核合成以后才产生的波不会改变核合成。若使用解耦时的额外辐射界，下限比核合成时低，但相应约束来自声学尺度、扩散长度及扰动对引力势的反馈，不能无条件把每种初值都当作热中微子。极低频的原初张量还可通过 \eqref{eq:gw-d-clt} 直接影响大角度温度与偏振，其限制途径又不同。

例如 $h_0^2\Omega_{{\rm gw},0}=A$ 的平台跨 $N$ 个十倍频程，式 \eqref{eq:gw-d-integral} 左边为 $A N\ln10$。若在 $f_{\rm p}$ 附近只有窄峰 $A\exp[-(\ln(f/f_{\rm p}))^2/(2s^2)]$，并远离积分边界，面积改为 $\sqrt{2\pi}sA$。同样高的峰可因宽度很窄而满足积分界。对蓝倾斜 $\Omega\propto f^\alpha$，固定 $f_{\rm max}$ 的顶值 $\Omega_{\rm max}$ 后，积分为 $\Omega_{\rm max}[1-(f_{\rm min}/f_{\rm max})^\alpha]/\alpha$；$\alpha\to0$ 连续还原对数平台。任何向高频增长的模型都须给出真实源尺度或热史导出的上限，不能以测量频段上端充当整个谱的截止。干涉仪或脉冲计时的频段约束则有各自的噪声权重，通常是在预设谱形上估一个幅度。与积分界联用时，应把同一个模型的全谱同时放进这两种权重。

\section{用脉冲星计时阵列测波}\label{sec:gw-e-pta}
\subsection{一对相隔千年的波形}
双星脉冲星的轨道衰减给出一对天体在源处损失的能量。另一类计时观测使用许多视线，把每颗脉冲星当作天然射电钟，测别处发出的引力波在光路上留下的时差。沿用第三章\hyperlink{subsection.3.12.9}{「从六种潮汐形变到干涉仪读数」}的约定，$\bm p_a$ 从地球指向距离 $L_a$ 的脉冲星，$\bm N$ 沿波的传播方向，$z_a$ 为周期的相对变化，残差为其时间积分。令 $\mu_a=\bm N\cdot\bm p_a$，$\tau_a=L_a(1+\mu_a)/c$。则第三章的两端项在频域合成一个响应
\begin{equation}\label{eq:gw-e-pta}
 \tilde r_a(f)=\sum_A \frac{F_a^A(\bm N)}{2\pi\i f}
 (1-\e^{-2\pi\i f\tau_a})\,\tilde h_A(f),\qquad
 F_a^A=\frac{p_a^ip_a^j e_{ij}^A}{2(1+\mu_a)} .
\end{equation}
这里先写单一来向，频率为相对于质心计时的引力波频率。$F_a^A$ 的表观极点来自沿射电路径改变积分变量。若 $\mu_a\to-1$，分子也随 $1-\mu_a^2$ 缩小，两端项之差又含一个 $1+\mu_a$，整条响应有有限极限。不能在这一小块天空先令脉冲星项平均消失，再取共线极限。

取频率演化可略的单一线偏振 $h=h_0\cos(\omega t+\phi)$，取中点相位 $v=\omega(t-\tau_a/2)+\phi$，积分给出 $r_a=(2F_a h_0/\omega)\sin(\omega\tau_a/2)\cos v$。积分常数由计时模型吸收。距离未知使得相位 $\omega\tau_a$ 几乎未知，两端的差既可增强也可降低单颗星的信号。形式上的 $1/f$ 使低频更有利；若波周期远长于观测时长，光滑残差却与脉冲星自转频率及其导数高度退化，实际响应很快降下。在另一极限 $\omega\tau_a\ll1$，两端差展开后 $r_a\simeq F_a\tau_a h(t)$，分母中的 $1+\mu_a$ 消去，源、收两端尝到近乎相同的场。

对远方准圆旋近源，地球端与脉冲星端即使在一个十年时段内各自近似单色，也未必有相同频率。设观测系满足 $\dot f=K f^{11/3}$。将方程从较早的脉冲星波相位积分到地球端，得
\[
 f_p^{-8/3}-f_e^{-8/3}=\tfrac83K\tau_a,\qquad
 f_p=f_e\left(1+\tfrac83 K\tau_a f_e^{8/3}\right)^{-3/8}.
\]
若 $\tau_a\dot f_e/f_e\ll1$，两线相差约 $\tau_a\dot f_e$。要在数据中分开它们还需 $|f_p-f_e|T\gtrsim1$。由此一颗脉冲星实际提供了更早的源时刻，若距离与来向足够准确，可将两端相位连起来约束旋近速率。相反，忽略源在数千年基线上的演化而只拟合两个自由相位，会漏掉一段频谱结构。此处 $\tau_a$ 为观者系时差，$K$ 含红移后的质量。

再考虑一次记忆跃迁 $h=\Delta h\,\Theta(t-t_B)$，式 \eqref{eq:gw-e-pta} 中频移先有台阶，到 $t_B+\tau_a$ 又撤去。令 $u_+=\max(0,u)$，则残差为 $F_a\Delta h[(t-t_B)_+-(t-t_B-\tau_a)_+]$。地球项在所有脉冲星同时出现，斜率的符号随方向而变。远端台阶常在历史数据之外。一条直线本会被频率拟合吸收，故事件在数据中段留下的折角比刚好出现在记录起点的事件容易辨认。这一例也表明单看一颗星的自转突变，常不能确定是否发生了传播效应。

\subsection{角相关的积分}
第三章已给 Hellings--Downs 曲线。它的对数项如何从两个光路的几何中出现？取各向同性、无偏振张量背景，先考虑两颗不同的脉冲星，$2\pi fL_a/c,2\pi fL_b/c\gg1$，两远端的位置也相隔许多引力波波长。快速变化的远端相位在角或频率平均后消失，余下地球端角函数 $\Gamma_{ab}=\frac{3}{2}\int\frac{\d^2\bm N}{4\pi}\sum_A F_a^A F_b^A$。系数这样选，是为了把同一脉冲星两端都计入的自相关归一为一。若两颗星在空间上极相近或观测频率很低，应保留式 \eqref{eq:gw-e-pta} 中的两个复因子再积分。以下推导可与 \href{https://arxiv.org/abs/1412.1142}{Jenet 与 Romano} 所讨论的天线相关作比较。

设 $\kappa=\bm p_a\cdot\bm p_b$，$u=\bm p_a\cdot\bm N$，$v=\bm p_b\cdot\bm N$。以 $\bm p_a$ 为极轴，绕它的方位角为 $\varphi$，则 $v=\kappa u+\sqrt{(1-u^2)(1-\kappa^2)}\cos\varphi$。用前面偏振完备式把两个响应缩并，得
\[
  \sum_A F_a^A F_b^A=
 \frac{2(\kappa-uv)^2-(1-u^2)(1-v^2)}{4(1+u)(1+v)} .
\]
此式只含三个内积。方位平均写成横线，有 $\bar v=\kappa u$、$\overline{1/(1+v)}=1/|u+\kappa|$，后一个等式由 $\int\d\varphi/(a+b\cos\varphi)=2\pi/\sqrt{a^2-b^2}$ 得到。这里看似还有极点；先做除法，
\[
 \frac{(\kappa-uv)^2}{1+v}
 =u^2v-u^2-2\kappa u+\frac{(\kappa+u)^2}{1+v},
\]
方位平均后的最后一项实为 $|u+\kappa|$。因此
\[
 \Gamma_{ab}=\frac{3}{16}\int_{-1}^{1}
 \frac{\kappa u^3-u^2-3\kappa u-1+2|u+\kappa|}{1+u}\,\d u .
\]
将被积函数拆为 $\kappa u^2-(1+\kappa)u+1-2\kappa$ 加余项。在 $u<-\kappa$ 上余项为 $-2$，在 $u>-\kappa$ 上则为 $2+4(\kappa-1)/(1+u)$。后半段正好给一个对数。令 $x=(1-\kappa)/2$，积分结果为 $\Gamma_{ab}=1/2-x/4+(3/2)x\ln x$。例如两条不同视线趋于平行时得 $1/2$，反向时得 $1/4$，直角时为 $3/8-(3/4)\ln2<0$。普通的相关系数可以为负，任意加权和的自相关仍必须非负。

\begin{figure}[htbp]
 \centering
 \begin{tikzpicture}
  \begin{axis}[width=.76\textwidth,height=4.8cm,axis lines=left,
    xmin=0,xmax=180,ymin=-.22,ymax=.56,
    xtick={0,45,90,135,180},ytick={-.15,0,.25,.5},
    tick label style={font=\small},xlabel={两条视线的夹角（度）},ylabel={$\Gamma_{ab}$}]
   \addplot[thin,gray,domain=0:180] {0};
   \addplot[very thick,domain=0:180,samples=181]
    {0.5-(1-cos(x))/8+0.75*(1-cos(x))*ln(max(0.000001,(1-cos(x))/2))};
  \end{axis}
 \end{tikzpicture}
 \caption{各向同性张量背景给两颗不同脉冲星的地球端相关。谷底在 $x=\exp(-5/6)$ 处。重合视线的极限与同一颗星的自相关有别，后者还含脉冲星端的贡献。}
 \label{fig:gw-e-hd}
\end{figure}

同一颗星还须计算 $(1-\e^{\i\omega\tau_a})(1-\e^{-\i\omega\tau_a})$，其平均为二。因此常把上式写成对所有脉冲星都适用的 $\Gamma_{ab}=(1+\delta_{ab})/2-x/4+(3/2)x\ln x$。再将式 \eqref{eq:gw-e-pta} 的时间积分与式 \eqref{eq:gw-b-omega} 的谱约定配对，两颗星的单边残差互谱为
\begin{equation}\label{eq:gw-e-res}
 S_{ab}^r(f)=\Gamma_{ab}\,\frac{h_c^2(f)}{12\pi^2 f^3}.
\end{equation}
检验其系数。对一颗星，$\sum_A(F_a^A)^2=(1-u)^2/4$，天区平均为 $1/3$。两端平均增倍，每偏振谱为 $S_g/2$，故频移自谱为 $S_g/3$，再除 $(2\pi f)^2$ 正合上式。当双黑洞群给 $h_c\propto f^{-2/3}$ 时，残差功率谱正比于 $f^{-13/3}$。在低频端计时拟合造成的压低必须连同这一红谱一并计算。

可进一步问，这个角相关有无恰好同钟差或历表误差相同的部分？对不同星只取地球项，定义 $C_\ell=2\pi\int_{-1}^1\Gamma(\kappa)P_\ell(\kappa)\d\kappa$。常数和线性项仅能贡献 $\ell=0,1$，但加上对数项后直接积分给 $C_0=C_1=0$。对 $\ell\ge2$，代入 Rodrigues 公式，分部积分 $\ell$ 次，边界的 $(1-\kappa^2)^\ell$ 使端点项消失；而 $[(1-\kappa)\ln(1-\kappa)]^{(\ell)}=(\ell-2)!/(1-\kappa)^{\ell-1}$。于是只需算 $\int_{-1}^1(1-\kappa)(1+\kappa)^\ell\d\kappa=2^{\ell+2}/[(\ell+1)(\ell+2)]$，得
\[
 C_\ell=\frac{6\pi}{(\ell-1)\ell(\ell+1)(\ell+2)},\qquad
 \Gamma(\kappa)=\sum_{\ell\ge2}\frac{2\ell+1}{4\pi}C_\ell P_\ell(\kappa).
\]
可见曲线以四极为主，同时含所有更高阶。若以同一地面钟的漂移污染各星，角相关是单极；质心位置误差的各向同性矢量协方差给偶极。只有在全天连续均匀覆盖时，这三类模板才严格正交。真实阵列的视线、观测时段和噪声均不均匀，须在实际脉冲星对上求模板间的重叠。

\subsection{时间拟合怎样改变红谱}
实测的时刻既不等间距，也常由多台射电望远镜在不同射频上取得。令一颗星的初始残差为列矢量 $\bm y$，在初始参数附近，自转、天体测量、双星轨道及色散参数的改变量给 $\bm y=M\delta\bm\lambda+\bm s+\bm n$。$M$ 的每列为定时模型对一个参数的偏导。若暂设随机项的协方差为 $C$，加权最小二乘给
\[
 R=I-M(M^{\mathsf T}C^{-1}M)^{-1}M^{\mathsf T}C^{-1},\qquad
 \bm r=R\bm y,\qquad C_{\rm post}=R C R^{\mathsf T}.
\]
因 $RM=0$、$R^2=R$，拟合后的协方差在原空间必然奇异。可改取矩阵 $G$，其列为 $M^{\mathsf T}$ 的零空间的一组正交基，则 $\bm y'=G^{\mathsf T}\bm y$ 不含 $\delta\bm\lambda$，协方差 $G^{\mathsf T}C G$ 在可测空间内可逆。这个投影解释了为何先扣掉最佳轨道，再把剩下的数当独立白噪声，会高估红背景的显著性。

例如居中的均匀记录拟合 $1,t,t^2$，分别对应相位、频率及频率导数。对 $\omega T\ll1$，入射复模板 $\exp(\i\omega t)$ 的前三项全落在拟合空间，正交剩余的振幅至多为 $O[(\omega T)^3]$。以 $v=2t/T$ 为变量，展开 $\exp(\i\omega Tv/2)=\sum_{\ell\ge0}(2\ell+1)\i^\ell j_\ell(\omega T/2)P_\ell(v)$。去掉 $\ell=0,1,2$ 后，残余功率的比例为 $1-\sum_{\ell=0}^2(2\ell+1)j_\ell^2(\omega T/2)$。利用 $j_3(z)\simeq z^3/105$，首项为 $7(\omega T/2)^6/105^2$。这比简单规定低频截止 $1/T$ 更准确。拟合位置、视差及自行，还会在年周频率及其谐波附近挖去灵敏度。实际不规则采样时可直接以矩阵计算同一投影。

射电传播的色散时延随观测射频 $\nu_{\rm rad}^{-2}$ 改变，而引力波改变光行时所生的残差对射频基本无色。两个频段若真在同时测量，写 $r_i=s+d/\nu_i^2+n_i$，消去色散可得 $s_{\rm est}=(\nu_1^2r_1-\nu_2^2r_2)/(\nu_1^2-\nu_2^2)$。两个频率过于接近，热噪声会被差分放大。若相隔日久才取得两个频段，$d(t)$ 也会改变，须将传播介质的时间协方差纳入模拟。脉冲轮廓的单次涨落、星自身的红色计时噪声、设备白噪声等大多只出现在自相关中。所有星都含相近的红谱尚不能单独确定地球端的角相关。

\subsection{阵列的最优二次统计量}
将各星投影后的时间序列叠成 $\bm y$，无背景时的协方差记为 $P$。由式 \eqref{eq:gw-e-res} 把给定谱形变换到各对观测时刻，再作计时投影，得到单位谱幅的协方差 $S$。例如平稳的互谱对时差 $\Delta t$ 给出 $\int_0^\infty S^r_{ab}(f)\cos(2\pi f\Delta t)\d f$，此式与有限观测和低频投影一同解释。有背景时 $C=P+A S$，$A$ 乘功率而非应变。Gaussian 对数似然为 $-\tfrac12[\ln\det C+\bm y^{\mathsf T}C^{-1}\bm y]$ 加常数。利用 $\partial_A C^{-1}=-C^{-1}SC^{-1}$，在 $A=0$ 展开，得
\[
 \widehat A=
 \frac{\bm y^{\mathsf T}P^{-1}S P^{-1}\bm y-
             \operatorname{tr}(P^{-1}S)}{\operatorname{tr}(P^{-1}S P^{-1}S)},\qquad
 \sigma_A^2\simeq \frac{2}{\operatorname{tr}(P^{-1}S P^{-1}S)} .
\]
这相当于将似然的一阶导数除以零假设下的信息，适用于弱信号并假定噪声参数已知。若只取不同星的交叉块作统计量，且 $P$ 按星分块对角，则分子中迹项为零。此时即使每颗星的自谱模型有小误差也不易生出虚假的固定角相关，但权重和误差仍会受影响。信号较强时，多个星对共享的同一地球项本身涨落，须用完整协方差或联合似然；不能把天球上同一随机实现的所有星对当成独立重抽样。

若采样近乎规则，先忽略拟合与不等噪声以看尺度。令 $S_0=h_c^2/(12\pi^2 f^3)$，白噪声残差单边谱为 $P_a$。由互谱方差可得弱信号信噪比 $\rho^2\simeq2T\sum_{a<b}\Gamma_{ab}^2\int S_0^2/(P_aP_b)\d f$。对白色独立时刻误差 $\sigma_a$ 和间隔 $\Delta t$，有 $P_a\simeq2\sigma_a^2\Delta t$。若 $S_0\propto f^{-\gamma}$ 且 $\gamma>1/2$，积分受低频截止 $\sim1/T$ 支配，故 $\rho\propto T^\gamma$。这陡峭增益只在新增的最低频格依然由仪器或脉冲星噪声主导时成立。在该格已由共同红背景主导后，分母也随信号增大，不能继续使用弱信号的时间幂律。

作一个量纲估计。频率每年一周、特征应变 $10^{-15}$ 的背景，每对数频带在一颗星上给约 $h_c/(\sqrt{12}\pi f)\simeq3\,\mathrm{ns}$ 的均方根残差。若来自准圆双星群，频率减为原来的十分之一，该量以 $f^{-5/3}$ 增大约 $46$ 倍。但波周期此时已有十年，拟合损失、可用的独立频格数及星体自噪声均不可略。正因如此，增加观测年数、增加有精确定时的天空方向及改善单次计时精度，是三种不同的实验改进。
